%\documentclass[11pt]{hmcpset}
\documentclass[12pt]{article}
\usepackage[margin=1in,top=1in]{geometry}
\usepackage{graphicx}  % Required for inserting images
\usepackage{tcolorbox}  % Required for problem statement boxes
\usepackage[calc,sets,vectors,phys,misc,LinAlg,quantum,geo]{shortmath}
%\usepackage[symbol]{footmisc}
\usepackage[utf8]{inputenc}
\usepackage{multicol}
\usepackage{hyperref}
\usepackage{mathtools}
\usepackage{multirow}
\usepackage{ulem}
\usepackage{array}
\newcommand{\PreserveBackslash}[1]{\let\temp=\\#1\let\\=\temp}
\newcolumntype{C}[1]{>{\PreserveBackslash\centering}p{#1}}
\usepackage{pdflscape}
\usepackage{longtable}
\usepackage{algorithm}
\usepackage{algpseudocode}
\usepackage{listings}
\lstset{
basicstyle=\small\ttfamily,
columns=flexible,
breaklines=true
}

\makeatletter
\@addtoreset{section}{part}
\makeatother

\renewcommand{\thefootnote}{\fnsymbol{footnote}}
\newcommand{\ang}{\mathring{\mathrm{A}}}

\newtcolorbox[]{problem}[1][]{title=#1}
\newenvironment{solution}[1][]{\def\myVar{#1}}{\hfill{\myVar}\bigskip}
\newenvironment{solution*}{}{}
\newtcolorbox[]{codefile}[1][]{title=#1}

\usepackage{enumitem}
\setlist[enumerate,1]{label=(\alph*)}

\renewcommand{\S}{\mathcal{S}}
\renewcommand{\R}{\mathcal{R}}
\renewcommand{\C}{\mathcal{C}}
\newcommand{\G}{\mathcal{G}}

\title{Optimal Degree Path Planning}
\date{November 19, 2025}
\author{HMCOR (Nathan Nguyen, Cole Plepel, and Joshua Zhong)}

\begin{document}

\maketitle

\iffalse
\begin{enumerate}
    \item \textbf{Pre/Co-Requisite structures}
    \item Which courses were offered in which of the last 4-8 semesters (Can be obtained from Hyperschedule)
    \item Number of credits (Can be obtained from Hyperschedule)
    \item Maximum number of repetitions (???)
    \item HMC Departmental HSA status (Can be obtained from Hyperschedule)
    \item 5C HSA status (Can be obtained from Hyperschedule)
    \item Core/Graduation requirement status (Can be obtained from the Student Handbook)
    \item Major requirement status (Can be obtained from the Student Handbook)
\end{enumerate}

\begin{verbatim}
set Courses;  # A subset of requirements
set Semesters;
set Requirements;
set GradRequirements;  # A subset of Requirements

param prereqCourse{Courses, Requirements} binary;
param coreqCourse{Courses, Requirements} binary;
param coreqGroup{Requirements, Requirements} binary;
param reqCount{Requirements} >= 0;
param interest{Courses};
param credits{Courses};
param minSemesterCredits >= 0;
param maxSemesterCredits >= 0;
param highSemesterCredits >= 0; # number of non-overloaded credits considered "high" (15?)
param minGradCredits >= 0;
param isOffered{Courses, Semesters} binary;  # only allow core courses in the proper semesters
param creditPenalty >= 0; # amount of objective function penalty incurred for each "high" credit (over 15)
param fragilityPenalty >= 0;

var takeCourse{Courses, Semesters};
var satisfied{Requirements, Semesters};
var isHighCredits{Semesters};

maximize Total_Interest:
    sum {s in Semesters, c in Courses} interest[c] * takeCourse[c, s]
    - sum{s in Semesters} (isHighCredits[s] * creditPenalty)
    + fragilityPenalty * sum{?} ?;  # to be filled in

subject to Satisfaction{r in Requirements, s in Semesters}: # maybe there's a better name for this...
    reqCount[r] * satisfied[r, s] <=
        sum{rr in Requirements, rr != r} (coreqGroup[r,rr] * satisfied[rr, s]) +  # Boolean logic
        sum{c in Courses} (prereqCourse[r, c] * credits[c] * sum{ss in Semesters, ss < s} takeCourse[c, ss]) +  # Prerequisite courses
        sum{c in Courses} (coreqCourse[r, c] * credits[c] * sum{ss in Semesters, ss <= s} takeCourse[c, ss]) +  # Corequisite courses
        reqCount[r] * satisfied[r, s-1];  # Satisfied requirements stay satisfied

subject to Graduation{r in GradRequirements}:
    satisfied[r, postgrad] = 1;

subject to TotalCredits:
    sum{c in Courses, s in Semesters} credits[c] * takeCourse[c, s] >= minGradCredits;

subject to NoOverload{s in Semesters}:
    sum{c in Courses} credits[c] * takeCourse[c, s] <= maxSemesterCredits;

subject to HighCredits{s in Semesters}:
    sum{c in Courses} credits[c] * takeCourse[c, s] - isHighCredits[s] <= highSemesterCredits;

subject to FullTime{s in Semesters}:
    sum{c in Courses} credits[c] * takeCourse[c, s] >= minSemesterCredits;

subject to Prerequisites{c in Courses, s in Semesters}:
    takeCourse[c, s] <= isOffered[c, s] * satisfied[c, s]

subject to NoRepeats{c in Courses}:
    sum{s in Semesters} takeCourse[c, s] <= 1;

\end{verbatim}

\newpage 
\fi

\begin{abstract}
    Four-year course planning to maximize student interest while balancing graduation requirements, prerequisites, and semesterly workloads is a challenging optimization problem faced by many students at Harvey Mudd College.
    This problem is especially relevant for first-year students seeking to understand how well different majors align with their curricular interests, yet it is also especially challenging, as these students are the least familiar with the graduation requirements and course prerequisites.
    To address this, we develop a model for four-year course planning, seeking to optimize a student's total cumulative interest in their planned coursework, subject to constraints corresponding to graduation requirements, course requisites, and semesterly workload limitations.
    After gathering a student's preferences, we express the optimization problem as a mixed integer program and solve it with a standard solver.
    We then analyze the sensitivity of our model to the choice of penalty parameters controlling the tradeoffs between the total course desirability, robustness to conflicts, and stressfulness of the plans produced by our model.
\end{abstract}

\part{Executive Summary}

\section{Intoduction}

Course scheduling can be hard, especially for first-year students seeking to understand possible paths through majors while navigating unfamiliar graduation and major requirements.
Even sophomores and upperclassmen could benefit from optimized scheduling tools designed to plan curricular paths tailored to students' interests, which minimize stressful, high-workload semesters while meeting graduation and major requirements.
To address this challenge, we at HMCOR have created a four-year course scheduler for Harvey Mudd College students by modeling the problem as a mixed integer linear program.

We used data from the Hyperschedule.io API, where we obtained information on course offerings within a two-year span, which included information about course names, credit amounts, discipline, and semesters offered.
We also collected course prerequisite and corequisite data from the departmental website of each STEM department at Harvey Mudd College, which we augmented through discussions with current students majoring in various STEM disciplines.
Further, we manually collected data from the Harvey Mudd Course Catalog for information about graduation requirements, such as required core classes, major classes, and completed credits.
Finally, we let the user of the model input their own data to express interest in potential classes and majors.
Given all of this data, we used a linear program to construct an optimal schedule that maximizes the students' interest in their classes while still ensuring the students graduate within 4 years without overloading.

\section{Model}

We modeled the problem of 4-year course scheduling as a mixed-integer program, seeking a schedule that maximizes the student's interest in the courses they are taking, minimizes the courses that are taken at the last possible time, and minimizes the number of per-semester credits over a certain threshold. The constraints of the model enforce major requirements, prerequisites/corequisites, and minimum and maximum allowable credit counts for each semester and for graduation, and they prohibit retaking most courses multiple times.

Since students hold different degrees of interest in different courses and must meet different major requirements for graduation, the model also relies on user input. So, we ran the model on several instances of sample input data in order to gather information on the general tendencies of the program. 

\section{Results}

The output produced when we ran the model for various majors and preferences does indeed resemble possible schedules that students of those majors could actually pursue. Core, major, and HSA requirements are met, prerequisites and corequisites are respected, and credit counts are within a reasonable range. The model has a notable tendency to de-prioritize major-required classes, particularly those that the student is not interested in. This reflects the belief that the student will always be allowed to take a required course if they have reached the final possible offering of that course and they need it to graduate.

In our analysis of the model, we see that the model first prioritizes including as many of the desired courses in the schedule as possible. Next, it prioritizes limiting the number of high-credit semesters, which we define as a semester with over 15 credits taken. Lastly, the model prioritizes taking desired classes as early as possible to the extent that this goal does not interfere with the previous two priorities. By changing some of the parameters, the model's priorities can be altered in accordance with the users' desires; for instance, avoiding high-credit semesters could be given higher or lower priority.

While the model presented in this report provides a working baseline, there are numerous ways to improve its quality and usability. Since credit counts are not always in direct correspondence with a semester's actual workload, the model is prone to suggesting schedules that would be exceedingly difficult in practice, despite maintaining low credit loads in each semester. Factoring in more subtle scheduling best-practices with the help of stakeholders, along with a more user-friendly interface, would help greatly in creating a model that can actually assist Mudders in their course-scheduling endeavors.

\clearpage

\part{Technical Report}

\section{Introduction}

Planning four years of classes given major requirements, idiosyncratic interests, and complex prerequisite structures can be a challenge for any Mudder.
Furthermore, first-year students often want to have an idea of their future coursework prospects in different majors to help decide their area of study.
To address this, we worked to solve the problem of four-year course scheduling at Harvey Mudd College, developing a mixed integer linear programming model of the problem.
We provide current students, especially freshmen, with the ability to plan out their four-year course schedules given their desired classes and major.
Moreover, we want to help students plan their schedules to reduce overscheduling or risk of being in danger of missing requirements.
This project also helps freshmen get an idea of what their schedules would look like in different majors, thus helping them decide what they want to study.

The primary stakeholders of this problem are students who want to plan out their future schedule.
First-year and major advisors are also stakeholders, as integral parts of the schedule planning process.
Lastly, the various academic departments at Harvey Mudd also play a role, as they expect students entering their majors to understand the coursework they are intended to take and plan accordingly.
We expect the students to value their interest in the courses they plan to take and the preparation they will receive for their post-graduation plans.
The former concern is directly reflected in the objective function we maximize.
The latter concern is also partially addressed by allowing students to indicate their degree of interest in each course.
Since major requirements are also designed to provide students with the necessary preparation for their post-graduation plans, constraining schedules to fulfill major and graduation requirements also serves to satisfy this objective, as well as the students' more basic desire to graduate on time.
Similarly, advisors' objectives are to ensure that students create a schedule that is not unnecessarily stressful and ensures flexibility for students' changing plans.
By allowing students and advisors to efficiently generate well-optimized schedules designed to align with both student interests and graduation requirements, we aim to help reduce the scheduling burden placed on advisors.
Improving the overall level of interest on the part of the students in terms of the classes they take may also improve the experience for instructors, as engaged students make for a better time teaching a class.
Further, by penalizing schedules that have little flexibility or require especially busy semesters, we hope to help students and advisors plot less stressful college trajectories.

To develop our model, we collected course data, graduation requirements, and relevant college policies from Hyperschedule.io, the Harvey Mudd College Student Handbook, and the Harvey Mudd College website.
We then customized the model with preference data provided by the student users of our model via an in-person survey.
We used this data to specify the parameters of a mixed integer linear programming model we developed for the course planning problem, which we solved using standard software.
% After solving this model, we also reran the model with each of the planned offerings of each course taken removed individually from the list of predicted course offerings to observe the robustness of the plan to uncertainties in future course offerings and conflicts.

The rest of this report is organized as follows.
In Section \ref{sec:model}, we specify our model formulation as a mixed integer linear program.
Then, in Section \ref{sec:data}, we describe our methods and sources of data collection.
We proceed to analyze our model in Sections \ref{sec:results} and \ref{sec:sensitivity}.
Specifically, we interrogate the model's sensitivity to the penalty parameters, evaluate the model's robustness to scheduling conflicts, and identify general trends in the observed scheduling strategy of the model.
Finally, we conclude our report and present directions for further improvement of our model in Section \ref{sec:conclusion}.

\section{Model Formulation}\label{sec:model}

Our course scheduler is a mixed integer linear program, maximizing the student's interest in a schedule's courses, with penalties for undesirable schedule traits and subject to scheduling constraints of prerequisites, graduation requirements, and other general considerations within a schedule, such as a lack of overloading.
In creating our model, we made some simplifying assumptions.
First, we assume no classes ever conflict with other classes, since we can’t predict which times which courses will be offered.
We also assume that the roster of course offerings repeats every four semesters, since most classes in the catalog are offered at least once every two years. We do not factor in ``topics classes'' since predicting when these classes will be offered is impractical.
This assumption also implies that we are assuming professors never go on sabbatical and that no professors will leave in the four years scheduled, since some electives are only taught by a single professor.

For our baseline model, we assume that the user has taken no relevant electives and no core classes, since we are prioritizing the problem of freshmen trying to create a four-year plan.
For simplicity, we assume in particular that the scheduling period begins in the fall of an odd-numbered year and terminates after eight semesters, in the spring of another odd-numbered year.
We also assume that the user can get into any class they want, since it is impractical to accurately predict when a class will fill and whether a student’s PERM request will be accepted.
A direction for future improvement could involve predicting when courses fill and how often PERMs are approved to account for the likelihood of the user being able to register for a given class in a given semester.

\subsection{Sets}
\begin{table}[H]
    \centering
    \begin{tabular}{c|l}
        \hline 
        \textbf{Symbol} &  \textbf{Description} \\
        \hline 
        $\S$ & Semesters\\ 
        $\R$ & Requirement Groups \\ 
        $\C$ & Courses\\
        $\G$ & Graduation Requirements
    \end{tabular}
    \caption{Sets}
    \label{tab:sets}
\end{table}

To represent our model, we define an ordered set $\S=\{1,2,3,4,5,6,7,8\}$ of semesters and a set $\R$ of ``requirement groups,'' which we use to encode Boolean logic conditions such as course prerequisites and graduation requirements.
We additionally define a set $\C\subseteq\R$ of courses and a set $\G\subseteq\R$ of graduation requirements.
The set of courses $\C$ contains the available courses to be scheduled based on historical course offerings, as well a number of generic courses which indicate opportunities for the user to fill in any course meeting certain requirements.
In particular, we add 12 generic STEM courses satisfying no requirements, 5 generic electives in each STEM department (Biology, Chemistry, Physics, Math, Computer Science, and Engineering), 6 generic 5C HSA courses, 4 generic Harvey Mudd departmental HSA courses, and one generic Harvey Mudd departmental writing-intensive HSA course.
Note that these generic course offerings are never chosen by the model if the student expresses any interest in a specific class that satisfies the corresponding requirement.

The sets $\C$ and $\G$ are defined as subsets of the set of requirement groups so that the requirement group associated with a course represents the prerequisites for taking the course, and the graduation requirements are simply requirements that must be met to graduate.
Each major is also represented by a requirement group in $\R$, although we do not define explicit subsets for these special elements.
In the nomenclature of our model implementation presented in Appendix \ref{app:data}, the graduation requirements group $\G$ consists of the requirement group identifiers \verb|Core|, \verb|HSAs|, and a group corresponding to the user's major.

\subsection{Decision Variables}
\begin{table}[H]
    \centering
    \begin{tabular}{c|l}
        \hline 
        \textbf{Symbol} &  \textbf{Description} \\
        \hline 
         $x_{c,s}$&  Binary variable indicating whether course $c$ is taken during semester $s$\\ 
         $y_{r,s}$& Binary variable, true if requirement group $r$ is satisfied upon starting semester $s$ \\ 
         $h_s$ & Nonnegative variable indicating the number of penalized credits in semester $s$
    \end{tabular}
    \caption{Variables}
    \label{tab:variables}
\end{table}

In our model, we use the binary decision variables $x_{c,s}$ for $c\in\C$ and $s\in\S$ to encode whether a course, $c$, is taken in semester, $s$. The decision variable is 1 if the course is taken, and 0 otherwise.
We additionally define the binary decision variables $y_{r,s}$ for $r\in\R$ and $s\in\S$ to encode whether a requirement $r$ is satisfied at the beginning of semester $s$.
Finally, we define $h_s$ for $s\in\S$ to represent the number of penalized credits taken in semester $s$ above the soft credit limit.

\subsection{Parameters}

\begin{table}[H]
    \centering
    \begin{tabular}{c|l}
        \hline 
        \textbf{Symbol} & \textbf{Description} \\
        \hline 
         $p_{c,r}$&  Binary parameter determining whether $c$ is a prerequisite in group $r$\\ 
         $c_{c,r}$ & Binary parameter determining whether $c$ is a corequisite in group $r$\\
          $g_{r',r}$& Binary parameter determining whether $r'$ is included in $r$ \\ 
          $n_r$&  Number of requirements and credits of requisites that satisfy $r$\\ 
         $I_c$ & Interest in $c$\\ 
          $H$& Threshold for ``high-credit'' semester\\
          $O$ & Penalty per credit over threshold\\ 
          $C_c$& Credit amount of $c$\\
          $G$ & Graduation credit requirement\\ 
          $m$ & Minimum credit amount in a semester\\ 
          $M$ & Maximum credit amount in a semester\\ 
          $\gamma$ & Anti-fragility weight\\
          $o_{c,s}$ & Binary parameter determining whether $c$ is offered in $s$.
    \end{tabular}

    \caption{Parameters}
    \label{tab:parameters}
\end{table}

Our model involves several parameters.
First, we define binary parameters $p_{c,r}$ and $c_{c,r}$ for $c\in\C$ and $r\in\R$, which are 1 if course $c$ is a prerequisite or corequisite for requirement group $r$, respectively.
Similarly, $g_{r',r}$ is 1 if requirement $r'\in\R$ is included in requirement group $r\in\R$, or 0 otherwise.
The positive parameters $n_r$ for $r\in\R$ then define the number of requirements and credits of requisites needed to satisfy the requirement group $r$.
These parameters together define the Boolean logic system of our model, which we use to ensure the generated schedule meets all graduation and major requirements and obeys course prerequisites and corequisites.

The parameters $I_c$ for $c\in\C$ represent the user's degree of interest in taking course $c$.
We note that $I_c$ will differ from a negative default value only for courses for which the user has chosen to specify a preference, and for the added generic electives, which we assign 0 interest.
We then define the positive parameters $H$ and $O$ representing the cutoff for a ``high-credit'' semester and the penalty per credit over this threshold, respectively.

The parameter $C_c$ for $c\in\C$ represents the credit amount of course $c$, and $G$ is the number of credits that must be taken by graduation.
Additionally, parameters $m$ and $M$ represent the minimum and maximum number of credits that can be taken in a specific semester, respectively.
The positive parameter $\gamma$ defines the weight of the ``anti-fragility'' term in the objective function, with higher values of $\gamma$ biasing the model more strongly towards taking more interesting classes earlier.
Finally, $o_{c,s}$ for $c\in\C$ and $s\in\S$ is a binary parameter that is 1 if course $c$ is offered in semester $s$, or 0 otherwise.

\subsection{Objective Function}

Our objective function, which we aim to maximize, is the total interest the user has in the courses included in the schedule, less penalties. This is given by
\begin{equation*}
    \sum_{s\in\S}\sum_{c\in\C} I_c x_{c,s} - \sum_{s}Oh_s - \gamma\sum_{s\in\S}\sum_{c\in\C} I_co_{c,s}\left(1-\sum_{s'>s\in\S}x_{c,s'}\right).
\end{equation*}
The first summation accounts for the student's total interest in the courses they would take under the plan, and the second summation penalizes semesters with credit counts exceeding the semesterly ``high-credit'' threshold.
The third summation discourages planning to take courses at the last possible time by penalizing each missed opportunity to take each desired course.

\subsection{Constraints}

Our first and most complex constraint implements the defining features of the requirement groups to build our Boolean logic system of prerequisites and corequisites.
\begin{equation*}
    n_r y_{r,s} \le \sum_{r'\in\R\setminus\{r\}} g_{r',r} y_{r',s}
    + \sum_{c\in\C} p_{c,r} C_{c} \sum_{s'<s\in\S} x_{c,s'}
    + \sum_{c\in\C} c_{c,r} C_{c} \sum_{s'\le s\in\S} x_{c,s'}
    \qquad\forall r\in\R, s\in\S.
\end{equation*}
This ensures $y_{r,s}$ can only be 1 if at least $n_r$ units of requirements in requirement group $r$ are met.
The first summation on the right-hand side accounts for dependencies between requirement groups, while the second and third summations account for the requirement's prerequisite and corequisite courses.

Next, we enforce the prerequisites of each course taken with the constraint
\begin{equation*}
    x_{c,s} \le o_{c,s} y_{c,s} \qquad\forall c\in\C, s\in\S,
\end{equation*}
and we enforce graduation requirements with the constraint
\begin{equation*}
    y_{r,\max S} = 1\qquad\forall r\in\G.
\end{equation*}
The total number of credits required to graduate is enforced by the constraint
\begin{equation*}
    \sum_{c\in\C} \sum_{s\in\S} C_c x_{c,s} \ge G.
\end{equation*}
Overloads and underloads are prohibited by the following constraints
\begin{gather*}
    m \le \sum_{c\in\C} C_c x_{c,s} \le M \qquad\forall s\in\S.
\end{gather*}
Meanwhile, the constraint that
\begin{equation*}
    \sum_{c\in\C} C_c x_{c,s} - h_s \le H
\end{equation*}
ensures credit loads over the ``high-credit'' threshold are reflected in the decision variables $h_s$ for $s\in\S$.

Finally, we prohibit taking the same course more than once with the requirement that
\begin{equation*}
    \sum_{s\in\S} x_{c,s} \le 1 \qquad\forall c\in\C.
\end{equation*}
Courses that may be repeated are accounted for by duplicating those courses in the instantiation data and assigning each duplicate the same interest level.

\section{Data Collection}\label{sec:data}

We primarily used the \href{https://hyperschedule.io}{Hyperschedule.io} API to collect the static model parameters, and for the Fall 2025 and Spring 2025 course data, we opened Hyperschedule on a local machine to extract the Hyperschedule test data.
All data from Hyperschedule was collected in JSON format, and parsed by a custom parser written in Python.
This data supplied the elements of the set of courses, $\C$, as well as the values of the parameters $C_c$ and $o_{c,s}$, under the assumption that the roster of course offerings exactly repeats every eight semesters.
%Music performance classes and physical education classes were duplicated to produce eight and three copies of each, respectively, to allow these courses to appear multiple times in the optimized schedules.
Classes which may only be taken in specific years, such as Core classes and senior capstones, have this restriction enforced by ensuring $o_{c,s}$ is zero for these classes in semesters they may not be taken, even if the course is otherwise offered.
The specific values of $C_c$ and $o_{c,s}$ we use can be found in Tables \ref{tab:offerings} and \ref{tab:credits} in Appendix \ref{app:data}.

We also collect data from the users of our model.
In particular, we ask each user to indicate their interest or disinterest in each course on which they have a preference on a scale of 0 (disinterested) to 5 (extremely interested), providing values for the parameters $I_c$.
For any course $c$ for which the user does not indicate interest or disinterest, $I_c$ is assigned the default interest value of $-10$, except for the set of added generic courses, which are each assigned an interest value of 0.
This encourages the model to suggest these generic courses instead of randomly selecting unrated courses where possible.
Example values of $I_c$ can be found in Tables \ref{tab:math-in}, \ref{tab:engr-in}, and \ref{tab:cs-in} in Appendix \ref{app:sample}.

The user is also asked to select a major, reflected in the inclusion of a corresponding major-specific requirement group in the set of graduation requirements $\G$.
The user may also elect to change the ``high-credit'' threshold $H$ and the ``high-credit'' penalty $O$, which default to 15 and 1.5, respectively.
Similarly, the user could elect to modify the value of the ``anti-fragility'' weight $\gamma$ from its default value of 0.05.

The minimum and maximum credit load per semester, $m$ and $M$, are set to 12 and 18, respectively, while the minimum number of credits to graduate, $G$, is set to 120, all in accordance with the student handbook.
Major requirements were gathered from the majors' pages on the Harvey Mudd College website, which are updated to reflect the requirements of the most recent student handbook.
These requirements were collected manually and recorded in a Python program, which then generated the appropriate AMPL data file.
These requirements constitute elements of the set of requirements, $\R$, as well as corresponding values of the parameters $p_{c,r}$, $c_{c,r}$, and $g_{r',r}$.
Graduation requirements were similarly collected and encoded from the student handbook.
The particular values of $g_{r',r}$, $p_{c,r}$, and $c_{c,r}$ we used are detailed in Tables \ref{tab:reqs-groups}, \ref{tab:reqs-prereqs}, and \ref{tab:reqs-coreqs} in Appendix \ref{app:data}.

In order to acquire course prerequisite and corequisite data, we collected the data by parsing the content of departmental websites listing course offerings. We used a Python script, modified from a Google Gemini output, that parsed the websites and wrote the prerequisite information into text files. The prompt and code output are included in Appendix \ref{app:ai}. This output was slightly modified to allow the url of each departmental website and the name of the output file to be specified. The generated text files were then parsed by a Python program to generate a corresponding AMPL data file specifying the course prerequisite and corequisite relationships.

\section{Results}\label{sec:results}

We ran the model on sample data sets for several of the majors at Harvey Mudd, including math, engineering, and computer science. The specific user inputs and model outputs are listed in Appendix \ref{app:sample}. One interesting feature of the model output is the way HSA courses are distributed across the student's time at Mudd. For example, in the sample mathematics schedule, the model suggests the student take one HSA course in their sophomore year, 6 HSA courses in the junior year, and the final 3 in the senior year. Meanwhile, the sophomore fall semester consists of four mathematics courses along with the required Engineering Systems course. In practice, people often prefer to have some balance of technical classes and HSAs, especially when the technical classes are ones that are considered difficult or high-workload. However, the model cannot currently factor in the difficulty of courses. Instead, courses are generally sequenced by interest, beginning with the courses with user-specified positive interest, then 0 interest generic courses to fulfill generic HSA requirements, and finally required major-specific courses with the default negative interest.

We can also see that in the engineering program, which is somewhat reputed for pushing students to delay HSAs to be taken later in their time at Mudd, no HSAs are scheduled until the junior spring semester. This indicates that even with no knowledge of how people ``typically'' plan their engineering program, the model seems to agree that the engineering requirements encourage students to front-load their technical coursework.

We can also examine when major-required courses are scheduled compared to technical electives that the student is interested in. In the sample math schedule, in which the user input included preferences for Operations Research and Statistical Theory, those courses were both scheduled in the sophomore fall. Meanwhile, the key major-required classes in Mathematical Analysis I and Abstract Algebra I are delayed until the following semester. In general, unless the student expresses interest in required courses, they tend to get delayed in the model for as long as possible. This is because in the objective function, the anti-fragility penalty for taking a course with fewer ``fallback'' times to take the course later is also proportional to the student's interest in the course. This speaks to an underlying principle of the model, that ``uninteresting'' courses that are required to graduate need not be prioritized, as the student will have no trouble securing seats in these classes. As such, the model pursues a strategy in which the student prioritizes the optional courses that they find interesting. In practice, these are indeed the courses the student is less certain to be able to register for, making this a useful strategy. %However, one downside of the current formulation is that the student's level of interest in major-required courses, even those that have no alternatives in the major program, can affect the output of the model. This is potentially something that could be changed about the model in the future, since if a student is required to take a course, and they will always get into it, then their interest in that course should not be a factor in making decisions.
Similarly, the model encourages students to take more interesting required courses earlier, which may have the side effect of producing additional opportunities for students to gain depth in these areas through further independent study or research after taking these courses.

\section{Sensitivity Analysis}\label{sec:sensitivity}

For our sensitivity analysis, we ran the model with altered parameters. Holding all other parameters at their usual values detailed in Appendix \ref{app:data}, we first ran the model witth $\gamma = 0$, then $\gamma = 0.5$, then $O = 0$, then $O = 10$, and finally set the default interest $I_c = 0$.

When we tested $\gamma = 0$, the courses taken were the same, with the exception of placeholders being exchanged for others, but the order was different, as detailed in Table \ref{tab:fragility-0-out} in Appendix \ref{app:sensitivity}. The model also took PHYS050 HM at the same time as ENGR079 HM, creating a 17-credit semester sophomore fall. We are unsure why this is happening, as we verified that this is being penalized in the model by the ``high-credit'' penalty. Otherwise, the suggested schedules were effectively unchanged from the baseline model. Moreover, the model reached a 0\% optimality gap and ran much faster with fewer simplex iterations. 

When we tested $\gamma = 0.5,$ the schedule was the same, with the exception of placeholders being exchanged for other placeholders, as detailed in Table \ref{tab:fragility-05-out} in Appendix \ref{app:sensitivity}. This makes sense, as we are already taking all of our desired courses as early as possible in our original model with $\gamma = 0.05,$ so decreasing $\gamma$ would not cause any differences. We believe that there may only be a difference in the output if we increase $\gamma$ such that reduction in fragility obtained by taking a class earlier overpowers the effect of the high credit penalty. For example, if we desired taking PHYS50 HM, the output likely would take it at the same time as ENGR79 HM, creating a 17 credit semester, overpowering the penalty from taking a high credit semester. 

When we tested $O = 0$, the courses taken were the same, but the order was different, as detailed in Table \ref{tab:credits-0-out} in Appendix \ref{app:sensitivity}. Generally, desired classes were taken earlier. For example, a high-priority class like MATH152 PO, which has no prerequisites in our model, was taken in the first semester. This makes sense, as taking the class earlier would reduce the fragility penalty effect from this class. Moreover, some semesters have 18 credits for similar reasoning, as the extra space freed up future classes to be taken earlier. Moreover, the model ran much faster with fewer simplex iterations. 

When we tested $O = 10$, CLES049A HM was taken in the first semester, as detailed in Table \ref{tab:credits-10-out} in Appendix \ref{app:sensitivity}. We assume this is because an increased ``high-credit'' penalty significantly inflates the objective function value such that the effect of taking the undesired class CLES049A HM is not pushed out when the optimality gap threshold is reached. Moreover, the sophomore fall semester has 14 credits instead of 16 credits, and the sophomore spring semester has 15 credits instead of 16 credits. This makes sense, as we are more heavily punishing high credit semesters. However, the tradeoff is that we take one less class when $O = 10$ compared to the baseline. 

When we tested the default interest $I_c = 0$, the outputted schedule takes many more courses for each semester, as many semesters are comprised of several 0.1, 0.75 and 1.5-credit courses. Since these classes have low credit counts, taking these classes does not violate the overloading constraint. Moreover, since these classes have zero interest, taking these classes does not directly affect the objective function value. We hypothesize that this is being done by the model so that semesters with PHYS050 HM, MATH198 HM, and ENGR079 HM can be taken with 15 credits while still not violating total graduation requirements. Under normal parameter values, we are forced to take 16 credit semesters, as taking 13 credit semesters would make it impossible to graduate, but now we are able to take several 15 credit semesters while still ensuring graduation. Moreover, the model ran much faster with fewer simplex iterations. The full model output from this experiment is detailed in Table \ref{tab:interest-0-out} in Appendix \ref{app:sensitivity}.

Overall, our analysis implies the model primarily prioritizes taking the most desired courses. Next, it prioritizes limiting the number of high-credit semesters, which we define as a semester with over 15 credits taken. Lastly, it prioritizes taking classes as early as possible the least. We can see that these priorities are dependent, as expected, upon the penalization and weight parameters, allowing the user to reflect their own priorities in the model's objective function. 

\section{Conclusion and Recommendations}\label{sec:conclusion}

We have seen that mixed integer programming provides a viable method of creating hypothetical 4-year course schedules, while factoring in requirements for the core, majors, and HSAs, pre/corequisites, student interest levels, desirable credit counts, and schedule fragility. Simply by knowing a student's major and course preferences, the model can incorporate their desired courses into a schedule that the student could hypothetically follow. 

There are still many improvements that could be made to the model. A primary concern is the workload required by the model in certain semesters. To avoid particularly difficult course schedules, students and faculty often recommend and adhere to certain principles for multi-semester planning. For example, math majors will advise each other to take Math 131 and Math 171 in separate semesters. Physics majors often follow a similar guideline for Physics 111 and Physics 117 if possible. The engineering department strongly recommends that students take Engineering 101 and 102 in the junior year rather than as a senior to mitigate the risk of not passing these classes. More generally, people tend to avoid taking too much technical coursework in a single semester. None of these agreed-upon ``best practices'' are factored into our current model. So, it would be interesting to attempt to quantify some of them and factor them into the objective function or constraints. For instance, a model could place limits on the ``effective workload'' of each semester if such ``effective workloads'' could be estimated for each class.

Given more time, the model would also benefit from having a more user-friendly front-end that is available to all Mudd students. This would require substantially more discussions with stakeholders, especially students and faculty from as many departments as possible to ensure consistent quality across majors. It would also be beneficial to implement a more flexible version of the model that can take as input the courses a student has already taken, and return an optimized plan for the remaining semesters. In conclusion, with a good deal of additional work, the baseline model introduced in this report has significant potential to become a widely applicable and useful tool.

\part{Appendices}

\section{AI Usage}\label{app:ai}

\subsection{AI Prompt Used} 
\begin{lstlisting}
Write a python script that is able to parse this website in order to extract the prerequisites for each course and writes the information to a txt file. The prerequisite data is only loaded once you click on the aria for each course. The website is: https://www.hmc.edu/cs/academic-programs/course-descriptions/. Make the code functional for other course description pages, such as https://www.hmc.edu/engineering/curriculum/courses/. 

The aria for one specific course in the html looks something like this: <a class="wp-block-mudd-acalog-widget-courses-course-link collapsed" href="https://catalog.hmc.edu/preview_course_nopop.php?catoid=26&amp;coid=8451" data-bs-target="#course-03ncg" data-bs-toggle="collapse" aria-expanded="false" aria-controls="course-03ncg" role="button">Introduction to Biology and Computer Science</a>

When opened, the aria looks as follows: <li class="wp-block-mudd-acalog-widget-courses-course"><div class="wp-block-mudd-acalog-widget-courses-course-title-wrapper"><span>CSCI121 HM</span><span><a class="wp-block-mudd-acalog-widget-courses-course-link" href="https://catalog.hmc.edu/preview_course_nopop.php?catoid=26&amp;coid=8247" data-bs-target="#course-5630m" data-bs-toggle="collapse" aria-expanded="true" aria-controls="course-5630m" role="button">Software Development</a></span></div><div class="collapse show" data-acalog-ajax="/catalog/28/course/6639/" data-acalog-ajax-type="course" id="course-5630m" data-acalog-data="courses" data-acalog-catalog-legacy-id="26" data-acalog-course-legacy-id="8247" style=""><div class="acalog-course-body"><p><strong>Credit(s):</strong> 3</p><p><strong>Instructor(s):</strong> Staff</p><p><strong>Description:</strong> Introduction to the discipline concerned with the design and implementation of software systems. The course presents a historical perspective on software development practice and explores modern, agile techniques for eliciting software requirements, designing and implementing software architecture and modules, robust testing practices, and project management. Student teams design, develop, and test a substantial software project. </p><p><strong>Prerequisite(s):</strong> <a class="acalog-permalink-link" href="https://catalog.hmc.edu/preview_course_nopop.php?catoid=26&amp;coid=8243">CSCI070 HM</a>&nbsp;<br></p></div></div></li>

Make sure that the content is loaded before trying to extract the prerequisite data. 

Make sure you use selenium in order to load the prerequisite data since it isn't loaded until the data is retrieved from the server.

A sample description will look like this:

Credit(s): 3
Instructor(s): Staff
Description: Introduction to the discipline concerned with the design and implementation of software systems. The course presents a historical perspective on software development practice and explores modern, agile techniques for eliciting software requirements, designing and implementing software architecture and modules, robust testing practices, and project management. Student teams design, develop, and test a substantial software project.
Prerequisite(s): CSCI070 HM 

\end{lstlisting}

\subsection{AI Model Code Output} 
\begin{lstlisting}
import time
import sys
from selenium import webdriver
from selenium.webdriver.chrome.service import Service
from selenium.webdriver.chrome.options import Options
from selenium.webdriver.common.by import By
from selenium.webdriver.support.ui import WebDriverWait
from selenium.webdriver.support import expected_conditions as EC
from selenium.common.exceptions import TimeoutException, NoSuchElementException, ElementClickInterceptedException
from webdriver_manager.chrome import ChromeDriverManager

def scrape_hmc_courses(url, output_file):
    """
    Scrapes HMC course description pages for prerequisites and details.
    """
    
    # Setup Chrome options
    chrome_options = Options()
    chrome_options.add_argument("--headless")  # Run in background
    chrome_options.add_argument("--no-sandbox")
    chrome_options.add_argument("--disable-dev-shm-usage")
    # Set a window size to ensure elements are "visible" for clicking
    chrome_options.add_argument("--window-size=1920,1080")

    print(f"Initializing Browser...")
    # specific service setup to handle driver installation automatically
    service = Service(ChromeDriverManager().install())
    driver = webdriver.Chrome(service=service, options=chrome_options)

    try:
        print(f"Navigating to: {url}")
        driver.get(url)

        # Wait for the main list to load
        wait = WebDriverWait(driver, 10)
        
        # Locate all the course toggle links (the aria buttons)
        # Based on the HTML provided: class="wp-block-mudd-acalog-widget-courses-course-link"
        course_links = wait.until(EC.presence_of_all_elements_located(
            (By.CSS_SELECTOR, "a.wp-block-mudd-acalog-widget-courses-course-link")
        ))
        
        print(f"Found {len(course_links)} courses. Starting extraction...")

        with open(output_file, "w", encoding="utf-8") as f:
            
            # We iterate by index because clicking elements might refresh parts of the DOM
            # causing "StaleElementReferenceException" if we hold onto the element objects
            for i in range(len(course_links)):
                try:
                    # Re-find the list to avoid stale elements
                    current_links = driver.find_elements(By.CSS_SELECTOR, "a.wp-block-mudd-acalog-widget-courses-course-link")
                    if i >= len(current_links):
                        break
                    
                    link = current_links[i]
                    
                    # Get the target ID (e.g., #course-03ncg) to know what to wait for
                    target_id = link.get_attribute("data-bs-target")
                    course_title = link.text.strip()
                    
                    # Attempt to find the course code (sibling span)
                    # Structure: <span>CODE</span><span><a ...>Title</a></span>
                    try:
                        parent_span = link.find_element(By.XPATH, "./..")
                        grandparent_div = parent_span.find_element(By.XPATH, "./..")
                        course_code = grandparent_div.find_element(By.TAG_NAME, "span").text.strip()
                    except:
                        course_code = "UNKNOWN"

                    print(f"Processing: {course_code} - {course_title}")

                    # Scroll to element to ensure it's clickable
                    driver.execute_script("arguments[0].scrollIntoView({block: 'center'});", link)
                    time.sleep(0.5) # Short buffer for scroll to settle

                    # Click to expand (Trigger AJAX)
                    # Using JS click is often safer for these types of overlay widgets
                    driver.execute_script("arguments[0].click();", link)

                    # Wait for the body to become visible and contain text
                    # The content loads into a div with class 'acalog-course-body' inside the target_id
                    body_selector = f"{target_id} .acalog-course-body"
                    
                    try:
                        wait.until(EC.visibility_of_element_located((By.CSS_SELECTOR, body_selector)))
                    except TimeoutException:
                        # Sometimes it's already loaded or takes longer
                        print(f"  Warning: Timeout waiting for details on {course_code}")
                        continue

                    # Extract the content div
                    content_div = driver.find_element(By.CSS_SELECTOR, body_selector)
                    
                    # Get all paragraphs to format them nicely
                    paragraphs = content_div.find_elements(By.TAG_NAME, "p")
                    
                    # Format output for this course
                    f.write(f"Course: {course_code} {course_title}\n")
                    f.write("-" * 40 + "\n")
                    
                    for p in paragraphs:
                        # Get text and clean it up
                        text = p.text.strip()
                        # If the text is empty, skip
                        if not text:
                            continue
                        f.write(f"{text}\n")
                    
                    f.write("\n" + "="*40 + "\n\n")
                    
                    # Optional: Collapse it back to keep DOM clean, though not strictly necessary
                    # driver.execute_script("arguments[0].click();", link)
                    
                except Exception as e:
                    print(f"Error processing course index {i}: {e}")
                    continue

        print(f"Done! Data saved to {output_file}")

    except Exception as main_e:
        print(f"Critical Error: {main_e}")
    finally:
        driver.quit()

if __name__ == "__main__":
    # Default URL if none provided
    target_url = "https://www.hmc.edu/cs/academic-programs/course-descriptions/"
    
    # Allow command line argument for different pages
    if len(sys.argv) > 1:
        target_url = sys.argv[1]
        
    output_filename = "course_prerequisites.txt"
    scrape_hmc_courses(target_url, output_filename)
\end{lstlisting}

\section{Data}\label{app:data}

\begin{table}[H]
    \centering
    \begin{tabular}{|c|c|l|}
        \hline
        Parameter & Value & Explanation \\\hline
        $H$ & 15 & Credit threshold for a ``high-credit'' semester \\
        $O$ & 1 & Penalty per credit over threshold \\
        $m$ & 12 & Minimum credit amount per semester \\
        $M$ & 18 & Maximum credit amount per semester \\
        $G$ & 120 & Minimum total credits required to graduation \\
        $\gamma$ & 0.05 & Weight of the ``anti-fragility'' penalty term \\\hline
    \end{tabular}
    \caption{Scalar parameter values}
    \label{tab:parameters}
\end{table}

The set $\G$ is defined as $\G=\{\verb|Core|,\verb|HSAs|,\mathcal{M}\}$, where $\mathcal{M}$ is selected from Table \ref{tab:majors} and indicates the student's major.

\begin{table}[H]
    \centering
    \begin{tabular}{|ccccc|}
        \hline
        \verb|Maj_BioChem| & \verb|Maj_BioClim| & \verb|Maj_Bio| & \verb|Maj_CS| & \verb|Maj_CSClimate| \\
        \verb|Maj_ChemClimate| & \verb|Maj_Chem| & \verb|Maj_CsPhys| & \verb|Maj_Engr| & \verb|Maj_Math| \\
        \verb|Maj_MathCS| & \verb|Maj_MathCompBio| & \verb|Maj_MathPhys| & \verb|MolBio| & \verb|Maj_Phys| \\\hline
    \end{tabular}
    \caption{Major identifiers.}
    \label{tab:majors}
\end{table}
%\in\{, \ , \ , \ , \ , \ , \ , \ , \ , \ , \ , \ , \ , \ , \ \}$ indicates the student's major.

{\footnotesize
\begin{longtable}{|l|l||l|l||l|l|}
    \hline
    Course ID & Credits & Course ID & Credits & Course ID & Credits\\\hline
    \verb|AFRI010ASAF| & 3 & \verb|AFRI010A_AF| & 3 & \verb|AFRI010BSAF| & 3\\\hline
    \verb|AFRI010B_AF| & 3 & \verb|AFRI014__AF| & 3 & \verb|AFRI020__PZ| & 3\\\hline
    \verb|AFRI110__PZ| & 3 & \verb|AFRI114__AF| & 3 & \verb|AFRI116__AF| & 3\\\hline
    \verb|AFRI120__PZ| & 3 & \verb|AFRI121__AF| & 3 & \verb|AFRI124__AF| & 3\\\hline
    \verb|AFRI125__AF| & 3 & \verb|AFRI128__AF| & 3 & \verb|AFRI134__PO| & 3\\\hline
    \verb|AFRI138__PO| & 3 & \verb|AFRI149__AF| & 3 & \verb|AFRI150__PZ| & 3\\\hline
    \verb|AFRI159__AF| & 3 & \verb|AFRI180__AF| & 3 & \verb|AFRI190__AF| & 3\\\hline
    \verb|AFRI190__AF2| & 3 & \verb|AFRI191C_AF| & 3 & \verb|AFRI191C_AF2| & 3\\\hline
    \verb|AFRI191__AF| & 3 & \verb|AFRI191__AF2| & 3 & \verb|AFRI195D_AF| & 3\\\hline
    \verb|AFRI195D_AF2| & 3 & \verb|AMST103__JT| & 3 & \verb|AMST103__PO| & 3\\\hline
    \verb|AMST103__PZ| & 3 & \verb|AMST113__SC| & 3 & \verb|AMST118__PO| & 3\\\hline
    \verb|AMST140__SC| & 3 & \verb|AMST179A_HM| & 3 & \verb|AMST180__PZ| & 3\\\hline
    \verb|AMST180__SC| & 3 & \verb|AMST190__PO| & 3 & \verb|AMST190__PO2| & 3\\\hline
    \verb|AMST191__PO| & 3 & \verb|AMST191__PO2| & 3 & \verb|AMST191__SC| & 3\\\hline
    \verb|AMST191__SC2| & 3 & \verb|ANTH001__PZ| & 3 & \verb|ANTH002__PO| & 3\\\hline
    \verb|ANTH002__PZ| & 3 & \verb|ANTH002__SC| & 3 & \verb|ANTH003__PZ| & 3\\\hline
    \verb|ANTH005__PZ| & 3 & \verb|ANTH009__PZ| & 3 & \verb|ANTH012__PZ| & 3\\\hline
    \verb|ANTH015__PO| & 3 & \verb|ANTH022__SC| & 3 & \verb|ANTH025__SC| & 3\\\hline
    \verb|ANTH027__SC| & 3 & \verb|ANTH053__PO| & 3 & \verb|ANTH055__PZ| & 3\\\hline
    \verb|ANTH070__PZ| & 3 & \verb|ANTH075__PZ| & 3 & \verb|ANTH077__PO| & 3\\\hline
    \verb|ANTH080__PO| & 3 & \verb|ANTH087__SC| & 3 & \verb|ANTH099__PZ| & 3\\\hline
    \verb|ANTH100__PO| & 3 & \verb|ANTH101__PZ| & 3 & \verb|ANTH104__PO| & 3\\\hline
    \verb|ANTH105__PO| & 3 & \verb|ANTH105__PZ| & 3 & \verb|ANTH106__PZ| & 3\\\hline
    \verb|ANTH107__PO| & 3 & \verb|ANTH107__PZ| & 3 & \verb|ANTH107__SC| & 3\\\hline
    \verb|ANTH109__SC| & 3 & \verb|ANTH111__PZ| & 3 & \verb|ANTH112__PZ| & 3\\\hline
    \verb|ANTH113__SC| & 3 & \verb|ANTH114__PZ| & 3 & \verb|ANTH114__SC| & 3\\\hline
    \verb|ANTH115__SC| & 3 & \verb|ANTH116__PO| & 3 & \verb|ANTH123__PO| & 3\\\hline
    \verb|ANTH124__SC| & 3 & \verb|ANTH126__SC| & 3 & \verb|ANTH137__PZ| & 3\\\hline
    \verb|ANTH143__SC| & 3 & \verb|ANTH153__SC| & 3 & \verb|ANTH184__PO| & 3\\\hline
    \verb|ANTH185P_SC| & 3 & \verb|ANTH185__SC| & 3 & \verb|ANTH189P_PO| & 3\\\hline
    \verb|ANTH190__PO| & 3 & \verb|ANTH190__PO2| & 3 & \verb|ANTH190__SC| & 3\\\hline
    \verb|ANTH190__SC2| & 3 & \verb|ANTH191__SC| & 3 & \verb|ANTH191__SC2| & 3\\\hline
    \verb|ARBC001__CM| & 3 & \verb|ARBC002__CM| & 3 & \verb|ARBC033__CM| & 3\\\hline
    \verb|ARBC044__CM| & 3 & \verb|ARBC130__CM| & 3 & \verb|ARBC141__CM| & 3\\\hline
    \verb|ARBC166__CM| & 3 & \verb|ARCN101__SC| & 3 & \verb|ARCN115__SC| & 3\\\hline
    \verb|ARCN120__SC| & 3 & \verb|ARCN144__SC| & 3 & \verb|ARCN191__SC| & 3\\\hline
    \verb|ARCN191__SC2| & 3 & \verb|ARCT179A_HM| & 3 & \verb|ARCT179B_HM| & 3\\\hline
    \verb|ARCT179C_HM| & 3 & \verb|ARHI001A_PO| & 3 & \verb|ARHI001B_PO| & 3\\\hline
    \verb|ARHI114__PO| & 3 & \verb|ARHI120__PO| & 3 & \verb|ARHI123__PO| & 3\\\hline
    \verb|ARHI124__PO| & 3 & \verb|ARHI127__PO| & 3 & \verb|ARHI130__PO| & 3\\\hline
    \verb|ARHI131__PO| & 3 & \verb|ARHI135__PZ| & 3 & \verb|ARHI138__PO| & 3\\\hline
    \verb|ARHI139__PO| & 3 & \verb|ARHI140__PO| & 3 & \verb|ARHI141A_PO| & 3\\\hline
    \verb|ARHI144B_PO| & 3 & \verb|ARHI150__SC| & 3 & \verb|ARHI151__SC| & 3\\\hline
    \verb|ARHI154__SC| & 3 & \verb|ARHI158__SC| & 3 & \verb|ARHI164__SC| & 3\\\hline
    \verb|ARHI174__PO| & 3 & \verb|ARHI178__PO| & 3 & \verb|ARHI179E_HM| & 3\\\hline
    \verb|ARHI179__PO| & 3 & \verb|ARHI183__SC| & 3 & \verb|ARHI185__SC| & 3\\\hline
    \verb|ARHI186A_PZ| & 3 & \verb|ARHI186B_PZ| & 3 & \verb|ARHI186C_SC| & 3\\\hline
    \verb|ARHI186D_PZ| & 3 & \verb|ARHI186G_PO| & 3 & \verb|ARHI186L_PO| & 3\\\hline
    \verb|ARHI186M_SC| & 3 & \verb|ARHI186N_SC| & 3 & \verb|ARHI186W_PO| & 3\\\hline
    \verb|ARHI186Y_PO| & 3 & \verb|ARHI187__SC| & 3 & \verb|ARHI189__SC| & 3\\\hline
    \verb|ARHI190__PO| & 3 & \verb|ARHI190__PO2| & 3 & \verb|ARHI191__PO| & 3\\\hline
    \verb|ARHI191__PO2| & 3 & \verb|ARHI191__SC| & 3 & \verb|ARHI191__SC2| & 3\\\hline
    \verb|ART_005__PO| & 3 & \verb|ART_010__PO| & 3 & \verb|ART_011__PZ| & 3\\\hline
    \verb|ART_012__PZ| & 3 & \verb|ART_015__PZ| & 3 & \verb|ART_016__PZ| & 3\\\hline
    \verb|ART_017X_PZ| & 3 & \verb|ART_020__PO| & 3 & \verb|ART_020__PZ| & 3\\\hline
    \verb|ART_021__PO| & 3 & \verb|ART_022__PO| & 3 & \verb|ART_024__PO| & 3\\\hline
    \verb|ART_025A_PO| & 3 & \verb|ART_026__PO| & 3 & \verb|ART_027A_PO| & 3\\\hline
    \verb|ART_027__PO| & 3 & \verb|ART_028__PO| & 3 & \verb|ART_030__PZ| & 3\\\hline
    \verb|ART_033__PO| & 3 & \verb|ART_060__HM| & 1.5 & \verb|ART_075__PZ| & 3\\\hline
    \verb|ART_101__SC| & 3 & \verb|ART_102__SC| & 3 & \verb|ART_105A_PO| & 3\\\hline
    \verb|ART_105__SC| & 3 & \verb|ART_106__SC| & 3 & \verb|ART_111__PO| & 3\\\hline
    \verb|ART_112__PZ| & 3 & \verb|ART_113__PZ| & 3 & \verb|ART_114__PO| & 3\\\hline
    \verb|ART_115__PO| & 3 & \verb|ART_116__SC| & 3 & \verb|ART_119__PZ| & 3\\\hline
    \verb|ART_120__PZ| & 3 & \verb|ART_121__SC| & 3 & \verb|ART_122__PO| & 3\\\hline
    \verb|ART_123__PO| & 3 & \verb|ART_123__SC| & 3 & \verb|ART_125__PZ| & 3\\\hline
    \verb|ART_125__SC| & 3 & \verb|ART_126B_PO| & 3 & \verb|ART_126B_SC| & 3\\\hline
    \verb|ART_126__PZ| & 3 & \verb|ART_127__PZ| & 3 & \verb|ART_129__SC| & 3\\\hline
    \verb|ART_132__SC| & 3 & \verb|ART_133__PZ| & 3 & \verb|ART_134__SC| & 3\\\hline
    \verb|ART_135__SC| & 3 & \verb|ART_136__SC| & 3 & \verb|ART_139__PO| & 3\\\hline
    \verb|ART_140__PO| & 3 & \verb|ART_141__SC| & 3 & \verb|ART_142__SC| & 3\\\hline
    \verb|ART_143__SC| & 3 & \verb|ART_144__SC| & 3 & \verb|ART_145__SC| & 3\\\hline
    \verb|ART_146__SC| & 3 & \verb|ART_147__SC| & 3 & \verb|ART_148__SC| & 3\\\hline
    \verb|ART_149__SC| & 3 & \verb|ART_150__SC| & 3 & \verb|ART_151__SC| & 3\\\hline
    \verb|ART_179M_HM| & 3 & \verb|ART_179S_HM| & 3 & \verb|ART_179T_HM| & 3\\\hline
    \verb|ART_179V_HM| & 3 & \verb|ART_179W_HM| & 3 & \verb|ART_179X_HM| & 3\\\hline
    \verb|ART_179Y_HM| & 3 & \verb|ART_181D_SC| & 3 & \verb|ART_181G_SC| & 3\\\hline
    \verb|ART_181M_SC| & 3 & \verb|ART_181__SC| & 3 & \verb|ART_188__PO| & 3\\\hline
    \verb|ART_189__PZ| & 3 & \verb|ART_190__PO| & 3 & \verb|ART_190__PO2| & 3\\\hline
    \verb|ART_192__PO| & 3 & \verb|ART_192__PO2| & 3 & \verb|ART_192__SC| & 3\\\hline
    \verb|ART_192__SC2| & 3 & \verb|ART_193__SC| & 3 & \verb|ART_193__SC2| & 3\\\hline
    \verb|ART_199__PZ| & 3 & \verb|ART_199__PZ2| & 3 & \verb|ASAM055__AA| & 3\\\hline
    \verb|ASAM070__PO| & 3 & \verb|ASAM077B_PZ| & 3 & \verb|ASAM082__PZ| & 3\\\hline
    \verb|ASAM085__PZ| & 3 & \verb|ASAM086__HM| & 3 & \verb|ASAM088__PZ| & 3\\\hline
    \verb|ASAM089__PZ| & 3 & \verb|ASAM090__PZ| & 3 & \verb|ASAM094__PZ| & 3\\\hline
    \verb|ASAM105B_PZ| & 3 & \verb|ASAM115__CM| & 3 & \verb|ASAM115__PO| & 3\\\hline
    \verb|ASAM123__HM| & 3 & \verb|ASAM124__HM| & 3 & \verb|ASAM125__AA| & 9\\\hline
    \verb|ASAM126__HM| & 3 & \verb|ASAM128__PZ| & 3 & \verb|ASAM130__PZ| & 3\\\hline
    \verb|ASAM135B_PZ| & 3 & \verb|ASAM142__PO| & 3 & \verb|ASAM143__PO| & 3\\\hline
    \verb|ASAM144__PO| & 3 & \verb|ASAM151__AA| & 3 & \verb|ASAM160__AA| & 3\\\hline
    \verb|ASAM161__AA| & 3 & \verb|ASAM172__PZ| & 3 & \verb|ASAM175__PZ| & 3\\\hline
    \verb|ASAM179B_AA| & 3 & \verb|ASAM179H_AA| & 9 & \verb|ASAM179I_AA| & 3\\\hline
    \verb|ASAM179P_AA| & 9 & \verb|ASAM179Q_AA| & 9 & \verb|ASAM179R_HM| & 3\\\hline
    \verb|ASAM187__AA| & 3 & \verb|ASAM190A_PO| & 3 & \verb|ASAM190A_PO2| & 3\\\hline
    \verb|ASAM191__PO| & 3 & \verb|ASAM191__PO2| & 3 & \verb|ASIA081__PO| & 3\\\hline
    \verb|ASIA083__PO| & 3 & \verb|ASIA190__PO| & 3 & \verb|ASIA190__PO2| & 3\\\hline
    \verb|ASIA191__PO| & 3 & \verb|ASIA191__PO2| & 3 & \verb|ASIA192__PO| & 3\\\hline
    \verb|ASIA192__PO2| & 3 & \verb|ASTR001L_PO| & 0.1 & \verb|ASTR001__PO| & 3\\\hline
    \verb|ASTR002__PO| & 3 & \verb|ASTR051L_PO| & 0.1 & \verb|ASTR051__PO| & 3\\\hline
    \verb|ASTR062__HM| & 3 & \verb|ASTR066LXKS| & 0.1 & \verb|ASTR066L_KS| & 3\\\hline
    \verb|ASTR101L_PO| & 0.1 & \verb|ASTR101__PO| & 3 & \verb|ASTR121__PO| & 1.5\\\hline
    \verb|ASTR122__HM| & 2 & \verb|ASTR123__PO| & 1.5 & \verb|ASTR125__PO| & 1.5\\\hline
    \verb|ASTR127__PO| & 1.5 & \verb|ASTR128__HM| & 2 & \verb|BIOL001D_PO| & 3\\\hline
    \verb|BIOL007L_PO| & 0.1 & \verb|BIOL007__PO| & 3 & \verb|BIOL023__HM| & 1\\\hline
    \verb|BIOL040__PO| & 3 & \verb|BIOL041C_PO| & 3 & \verb|BIOL041E_PO| & 3\\\hline
    \verb|BIOL042L_KS| & 6 & \verb|BIOL043LXKS| & 0.1 & \verb|BIOL043L_KS| & 3\\\hline
    \verb|BIOL044LXKS| & 0.1 & \verb|BIOL044L_KS| & 3 & \verb|BIOL046R_HM| & 0.1\\\hline
    \verb|BIOL046__HM| & 3 & \verb|BIOL047__PO| & 3 & \verb|BIOL054__HM| & 1\\\hline
    \verb|BIOL057LXKS| & 0.1 & \verb|BIOL057L_KS| & 3 & \verb|BIOL059L_KS| & 3\\\hline
    \verb|BIOL099__KS| & 3 & \verb|BIOL101__HM| & 3 & \verb|BIOL103__HM| & 2\\\hline
    \verb|BIOL104A_PO| & 3 & \verb|BIOL107__PO| & 3 & \verb|BIOL108__HM| & 3\\\hline
    \verb|BIOL108__PO| & 3 & \verb|BIOL109__HM| & 3 & \verb|BIOL110__HM| & 3\\\hline
    \verb|BIOL111__HM| & 2 & \verb|BIOL111__KS| & 3 & \verb|BIOL112__KS| & 1.5\\\hline
    \verb|BIOL112__PO| & 3 & \verb|BIOL113L_KS| & 3 & \verb|BIOL113__HM| & 3\\\hline
    \verb|BIOL116__PO| & 3 & \verb|BIOL119__HM| & 3 & \verb|BIOL119__KS| & 0.75\\\hline
    \verb|BIOL120__KS| & 1.5 & \verb|BIOL121__HM| & 3 & \verb|BIOL121__PO| & 3\\\hline
    \verb|BIOL122__HM| & 3 & \verb|BIOL131L_KS| & 3 & \verb|BIOL131__KS| & 3\\\hline
    \verb|BIOL133__PO| & 3 & \verb|BIOL135L_KS| & 3 & \verb|BIOL138L_KS| & 0.1\\\hline
    \verb|BIOL138__KS| & 3 & \verb|BIOL140__PO| & 3 & \verb|BIOL141L_KS| & 3\\\hline
    \verb|BIOL142L_KS| & 3 & \verb|BIOL143__KS| & 3 & \verb|BIOL145__KS| & 3\\\hline
    \verb|BIOL146L_KS| & 3 & \verb|BIOL147__KS| & 3 & \verb|BIOL151L_KS| & 3\\\hline
    \verb|BIOL152__PO| & 3 & \verb|BIOL154__HM| & 3 & \verb|BIOL154__KS| & 3\\\hline
    \verb|BIOL156L_KS| & 3 & \verb|BIOL157L_KS| & 3 & \verb|BIOL160__HM| & 3\\\hline
    \verb|BIOL160__PO| & 3 & \verb|BIOL161__HM| & 1 & \verb|BIOL162__PO| & 3\\\hline
    \verb|BIOL163L_PO| & 0.1 & \verb|BIOL163__PO| & 3 & \verb|BIOL164__KS| & 3\\\hline
    \verb|BIOL168L_KS| & 3 & \verb|BIOL169L_PO| & 0.1 & \verb|BIOL169__PO| & 3\\\hline
    \verb|BIOL170L_KS| & 3 & \verb|BIOL170__PO| & 3 & \verb|BIOL173L_KS| & 1.5\\\hline
    \verb|BIOL173__PO| & 3 & \verb|BIOL174__PO| & 3 & \verb|BIOL175__KS| & 3\\\hline
    \verb|BIOL176__KS| & 3 & \verb|BIOL177__KS| & 3 & \verb|BIOL182L_KS| & 3\\\hline
    \verb|BIOL182__HM| & 3 & \verb|BIOL183__KS| & 3 & \verb|BIOL184L_KS| & 3\\\hline
    \verb|BIOL184__HM| & 1 & \verb|BIOL185G_HM| & 3 & \verb|BIOL187JLKS| & 3\\\hline
    \verb|BIOL187M_KS| & 3 & \verb|BIOL187P_KS| & 3 & \verb|BIOL188L_KS| & 3\\\hline
    \verb|BIOL188__HM| & 3 & \verb|BIOL189L_KS| & 0.1 & \verb|BIOL189__HM| & 3\\\hline
    \verb|BIOL190L_KS| & 3 & \verb|BIOL190L_KS2| & 3 & \verb|BIOL190__PO| & 1.5\\\hline
    \verb|BIOL190__PO2| & 1.5 & \verb|BIOL191F_PO| & 3 & \verb|BIOL191F_PO2| & 3\\\hline
    \verb|BIOL191H_PO| & 1.5 & \verb|BIOL191H_PO2| & 1.5 & \verb|BIOL191__HM| & 0.1\\\hline
    \verb|BIOL191__HM2| & 0.1 & \verb|BIOL191__HM3| & 0.1 & \verb|BIOL191__HM4| & 0.1\\\hline
    \verb|BIOL191__KS| & 3 & \verb|BIOL191__KS2| & 3 & \verb|BIOL193__HM| & 3\\\hline
    \verb|BIOL193__HM2| & 3 & \verb|BIOL194A_PO| & 3 & \verb|BIOL194A_PO2| & 3\\\hline
    \verb|BIOL194B_PO| & 3 & \verb|BIOL194B_PO2| & 3 & \verb|BIOL195__HM| & 6\\\hline
    \verb|BIOL195__HM2| & 6 & \verb|BIOL196__KS| & 0.75 & \verb|BIOL196__KS2| & 0.75\\\hline
    \verb|BIOL197__HM| & 1 & \verb|BIOL197__HM2| & 1 & \verb|BIOL197__KS| & 1.5\\\hline
    \verb|BIOL197__KS2| & 1.5 & \verb|BIOL900__HM| & 3 & \verb|BIOL901__HM| & 3\\\hline
    \verb|BIOL902__HM| & 3 & \verb|BIOL903__HM| & 3 & \verb|BIOL904__HM| & 3\\\hline
    \verb|CASA101__PZ| & 4.5 & \verb|CASA105__PZ| & 4.5 & \verb|CASA200__PZ| & 3\\\hline
    \verb|CGH_100__JT| & 3 & \verb|CGH_301X_CG| & 3 & \verb|CGS_010__PZ| & 3\\\hline
    \verb|CGS_050__PZ| & 3 & \verb|CGS_060__PZ| & 3 & \verb|CGS_075__PZ| & 3\\\hline
    \verb|CGS_080__PZ| & 3 & \verb|CGS_085__PZ| & 3 & \verb|CGS_101__PZ| & 1.5\\\hline
    \verb|CGS_120__PZ| & 3 & \verb|CGS_122__PZ| & 3 & \verb|CGS_127__PZ| & 3\\\hline
    \verb|CGS_146__PZ| & 3 & \verb|CGS_150__PZ| & 3 & \verb|CGS_151__PZ| & 3\\\hline
    \verb|CGS_190__PZ| & 3 & \verb|CGS_190__PZ2| & 3 & \verb|CGS_199__PZ| & 3\\\hline
    \verb|CGS_199__PZ2| & 3 & \verb|CHEM001ALPO| & 0.1 & \verb|CHEM001A_PO| & 3\\\hline
    \verb|CHEM001BLPO| & 0.1 & \verb|CHEM001B_PO| & 3 & \verb|CHEM014L_KS| & 3\\\hline
    \verb|CHEM015LXKS| & 0.1 & \verb|CHEM015L_KS| & 3 & \verb|CHEM023A_PO| & 3\\\hline
    \verb|CHEM023BLPO| & 0.1 & \verb|CHEM023B_PO| & 3 & \verb|CHEM024__HM| & 1\\\hline
    \verb|CHEM029L_KS| & 3 & \verb|CHEM042L_KS| & 3 & \verb|CHEM042R_HM| & 0.1\\\hline
    \verb|CHEM042__HM| & 4 & \verb|CHEM048__HM| & 3 & \verb|CHEM051L_PO| & 0.1\\\hline
    \verb|CHEM051__HM| & 3 & \verb|CHEM051__PO| & 3 & \verb|CHEM052__HM| & 3\\\hline
    \verb|CHEM053__HM| & 1 & \verb|CHEM056__HM| & 3 & \verb|CHEM058__HM| & 1\\\hline
    \verb|CHEM060__HM| & 3 & \verb|CHEM070L_KS| & 3 & \verb|CHEM080__HM| & 3\\\hline
    \verb|CHEM081LXJT| & 0.1 & \verb|CHEM081L_JT| & 3 & \verb|CHEM103__HM| & 3\\\hline
    \verb|CHEM104__HM| & 3 & \verb|CHEM105__HM| & 3 & \verb|CHEM106__PO| & 3\\\hline
    \verb|CHEM109__HM| & 1 & \verb|CHEM110ALPO| & 0.1 & \verb|CHEM110A_PO| & 3\\\hline
    \verb|CHEM110BLPO| & 0.1 & \verb|CHEM110B_PO| & 3 & \verb|CHEM110__HM| & 1\\\hline
    \verb|CHEM111__HM| & 1 & \verb|CHEM112__HM| & 1 & \verb|CHEM112__PO| & 3\\\hline
    \verb|CHEM114__HM| & 3 & \verb|CHEM115L_PO| & 0.1 & \verb|CHEM115__PO| & 3\\\hline
    \verb|CHEM116LXKS| & 0.1 & \verb|CHEM116L_KS| & 3 & \verb|CHEM116__HM| & 3\\\hline
    \verb|CHEM117LXKS| & 0.1 & \verb|CHEM117L_KS| & 3 & \verb|CHEM121__KS| & 3\\\hline
    \verb|CHEM122__HM| & 2 & \verb|CHEM122__KS| & 3 & \verb|CHEM123__KS| & 1.5\\\hline
    \verb|CHEM126L_KS| & 3 & \verb|CHEM127L_KS| & 3 & \verb|CHEM128__KS| & 3\\\hline
    \verb|CHEM140__KS| & 3 & \verb|CHEM147L_PO| & 0.1 & \verb|CHEM147__PO| & 3\\\hline
    \verb|CHEM150__HM| & 1 & \verb|CHEM151__HM| & 3 & \verb|CHEM151__PO| & 3\\\hline
    \verb|CHEM152__HM| & 3 & \verb|CHEM156__PO| & 3 & \verb|CHEM158A_PO| & 3\\\hline
    \verb|CHEM158BLPO| & 0.1 & \verb|CHEM158B_PO| & 3 & \verb|CHEM161L_PO| & 0.1\\\hline
    \verb|CHEM161__PO| & 3 & \verb|CHEM164__PO| & 1.5 & \verb|CHEM170__HM| & 3\\\hline
    \verb|CHEM171__HM| & 2 & \verb|CHEM175__PO| & 3 & \verb|CHEM177__KS| & 3\\\hline
    \verb|CHEM181__KS| & 1.5 & \verb|CHEM181__PO| & 3 & \verb|CHEM182__HM| & 3\\\hline
    \verb|CHEM184__HM| & 1 & \verb|CHEM187__PO| & 3 & \verb|CHEM188L_KS| & 3\\\hline
    \verb|CHEM189L_KS| & 0.1 & \verb|CHEM190L_KS| & 3 & \verb|CHEM190L_KS2| & 3\\\hline
    \verb|CHEM191__KS| & 3 & \verb|CHEM191__KS2| & 3 & \verb|CHEM191__PO| & 1.5\\\hline
    \verb|CHEM191__PO2| & 1.5 & \verb|CHEM193M_HM| & 2 & \verb|CHEM193M_HM2| & 2\\\hline
    \verb|CHEM193P_HM| & 2 & \verb|CHEM193P_HM2| & 2 & \verb|CHEM194__HM| & 2\\\hline
    \verb|CHEM194__HM2| & 2 & \verb|CHEM194__PO| & 1.5 & \verb|CHEM194__PO2| & 1.5\\\hline
    \verb|CHEM196__KS| & 0.75 & \verb|CHEM196__KS2| & 0.75 & \verb|CHEM197__HM| & 3\\\hline
    \verb|CHEM197__HM2| & 3 & \verb|CHEM197__KS| & 1.5 & \verb|CHEM197__KS2| & 1.5\\\hline
    \verb|CHEM198__HM| & 3 & \verb|CHEM198__HM2| & 3 & \verb|CHEM199__HM| & 0.1\\\hline
    \verb|CHEM199__HM2| & 0.1 & \verb|CHEM199__HM3| & 0.1 & \verb|CHEM199__HM4| & 0.1\\\hline
    \verb|CHEM900__HM| & 3 & \verb|CHEM901__HM| & 3 & \verb|CHEM902__HM| & 3\\\hline
    \verb|CHEM903__HM| & 3 & \verb|CHEM904__HM| & 3 & \verb|CHIN001A_PO| & 3\\\hline
    \verb|CHIN001B_PO| & 3 & \verb|CHIN002__PO| & 3 & \verb|CHIN011__PO| & 0.75\\\hline
    \verb|CHIN013__PO| & 0.75 & \verb|CHIN051A_PO| & 3 & \verb|CHIN051B_PO| & 3\\\hline
    \verb|CHIN051H_PO| & 3 & \verb|CHIN055__PO| & 3 & \verb|CHIN111A_PO| & 3\\\hline
    \verb|CHIN111B_PO| & 3 & \verb|CHIN123__PO| & 3 & \verb|CHIN124__PO| & 3\\\hline
    \verb|CHIN129__PO| & 3 & \verb|CHIN131__PO| & 3 & \verb|CHIN150__PO| & 3\\\hline
    \verb|CHIN153__PO| & 3 & \verb|CHIN192A_PO| & 1.5 & \verb|CHIN192A_PO2| & 1.5\\\hline
    \verb|CHLT013__CH| & 3 & \verb|CHLT060__CH| & 3 & \verb|CHLT072__CH| & 3\\\hline
    \verb|CHLT085__PZ| & 3 & \verb|CHLT103__CH| & 3 & \verb|CHLT138__CH| & 3\\\hline
    \verb|CHLT151__CH| & 3 & \verb|CHLT170__CH| & 3 & \verb|CHLT181__CH| & 3\\\hline
    \verb|CHNT168__PO| & 3 & \verb|CHNT178__PO| & 3 & \verb|CHNT180__PO| & 3\\\hline
    \verb|CHST015__CH| & 3 & \verb|CHST028__CH| & 3 & \verb|CHST064__CH| & 3\\\hline
    \verb|CHST066__CH| & 3 & \verb|CHST067__CH| & 3 & \verb|CHST074__CH| & 3\\\hline
    \verb|CHST081__PO| & 3 & \verb|CHST102__PO| & 3 & \verb|CHST120__CH| & 3\\\hline
    \verb|CHST124__PO| & 3 & \verb|CHST125__CH| & 3 & \verb|CHST128__CH| & 3\\\hline
    \verb|CHST130__CH| & 3 & \verb|CHST132__CH| & 3 & \verb|CHST136__CH| & 3\\\hline
    \verb|CHST182__CH| & 3 & \verb|CHST184D_CH| & 3 & \verb|CHST184__CH| & 3\\\hline
    \verb|CHST185C_CH| & 3 & \verb|CHST189B_PO| & 3 & \verb|CHST190__CH| & 3\\\hline
    \verb|CHST190__CH2| & 3 & \verb|CHST191__CH| & 3 & \verb|CHST191__CH2| & 3\\\hline
    \verb|CHST192__CH| & 3 & \verb|CHST192__CH2| & 3 & \verb|CLAS001__PO| & 3\\\hline
    \verb|CLAS010__SC| & 3 & \verb|CLAS012__SC| & 3 & \verb|CLAS019__SC| & 3\\\hline
    \verb|CLAS023__SC| & 3 & \verb|CLAS064__PO| & 3 & \verb|CLAS088__PZ| & 3\\\hline
    \verb|CLAS112__PO| & 3 & \verb|CLAS114__PO| & 3 & \verb|CLAS116A_PO| & 3\\\hline
    \verb|CLAS116__SC| & 3 & \verb|CLAS117__PO| & 3 & \verb|CLAS121__JT| & 3\\\hline
    \verb|CLAS133__SC| & 3 & \verb|CLAS150BEPZ| & 3 & \verb|CLAS161__PZ| & 3\\\hline
    \verb|CLAS162__PZ| & 3 & \verb|CLAS175__PZ| & 3 & \verb|CLAS190__PO| & 3\\\hline
    \verb|CLAS190__PO2| & 3 & \verb|CLAS191__PO| & 3 & \verb|CLAS191__PO2| & 3\\\hline
    \verb|CLAS191__SC| & 3 & \verb|CLAS191__SC2| & 3 & \verb|CLAS192__PO| & 3\\\hline
    \verb|CLAS192__PO2| & 3 & \verb|CLES049A_HM| & 0.5 & \verb|CLES101__HM| & 3\\\hline
    \verb|CLES120__HM| & 3 & \verb|CLES121__HM| & 1.5 & \verb|CLES122__HM| & 1.5\\\hline
    \verb|CLES130__HM| & 3 & \verb|CLES131__HM| & 3 & \verb|CLES179A_HM| & 3\\\hline
    \verb|CLES179B_HM| & 3 & \verb|CLES179C_HM| & 3 & \verb|CLES197__HM| & 1\\\hline
    \verb|CLES197__HM2| & 1 & \verb|CLES199__HM| & 0.1 & \verb|CLES199__HM2| & 0.1\\\hline
    \verb|COGS011__PZ| & 3 & \verb|COGS123L_JT| & 3 & \verb|COGS182__PZ| & 3\\\hline
    \verb|CORE001__SC| & 3 & \verb|CORE002__SC| & 3 & \verb|CORE003__SC| & 3\\\hline
    \verb|CORE010A_SC| & 3 & \verb|CORE079R_HM| & 0.1 & \verb|CORE079__HM| & 3\\\hline
    \verb|CORE099R_HM| & 0.1 & \verb|CORE099__HM| & 3 & \verb|CSCI004__PZ| & 3\\\hline
    \verb|CSCI005L_HM| & 0.1 & \verb|CSCI005__HM| & 3 & \verb|CSCI022__SC| & 3\\\hline
    \verb|CSCI035__HM| & 3 & \verb|CSCI036P_PZ| & 3 & \verb|CSCI036__CM| & 3\\\hline
    \verb|CSCI036__PZ| & 3 & \verb|CSCI040__CM| & 3 & \verb|CSCI046L_CM| & 0.1\\\hline
    \verb|CSCI046__CM| & 3 & \verb|CSCI050__PO| & 1.5 & \verb|CSCI051ALPO| & 0.1\\\hline
    \verb|CSCI051A_PO| & 3 & \verb|CSCI051L_PO| & 0.1 & \verb|CSCI051PLPO| & 0.1\\\hline
    \verb|CSCI051P_PO| & 3 & \verb|CSCI051__PO| & 3 & \verb|CSCI054__PO| & 3\\\hline
    \verb|CSCI060__HM| & 3 & \verb|CSCI062L_PO| & 0.1 & \verb|CSCI062__PO| & 3\\\hline
    \verb|CSCI070__HM| & 3 & \verb|CSCI081__HM| & 3 & \verb|CSCI101__PO| & 3\\\hline
    \verb|CSCI105L_PO| & 0.1 & \verb|CSCI105__HM| & 3 & \verb|CSCI105__PO| & 3\\\hline
    \verb|CSCI120__HM| & 3 & \verb|CSCI123__HM| & 3 & \verb|CSCI131__HM| & 3\\\hline
    \verb|CSCI133__HM| & 3 & \verb|CSCI134__HM| & 3 & \verb|CSCI140__HM| & 3\\\hline
    \verb|CSCI140__PO| & 3 & \verb|CSCI143__CM| & 3 & \verb|CSCI144__HM| & 3\\\hline
    \verb|CSCI145__CM| & 3 & \verb|CSCI148__CM| & 3 & \verb|CSCI151__HM| & 3\\\hline
    \verb|CSCI152__HM| & 3 & \verb|CSCI152__PO| & 3 & \verb|CSCI153__HM| & 3\\\hline
    \verb|CSCI158__HM| & 3 & \verb|CSCI158__PO| & 3 & \verb|CSCI158__PZ| & 3\\\hline
    \verb|CSCI159__HM| & 3 & \verb|CSCI159__PO| & 3 & \verb|CSCI181AAHM| & 3\\\hline
    \verb|CSCI181AAPO| & 3 & \verb|CSCI181ACHM| & 3 & \verb|CSCI181AFHM| & 3\\\hline
    \verb|CSCI181AGHM| & 3 & \verb|CSCI181ALHM| & 3 & \verb|CSCI181AOHM| & 3\\\hline
    \verb|CSCI181APHM| & 3 & \verb|CSCI181AQHM| & 3 & \verb|CSCI181ARHM| & 3\\\hline
    \verb|CSCI181ATHM| & 3 & \verb|CSCI181AUHM| & 3 & \verb|CSCI181CAPO| & 3\\\hline
    \verb|CSCI181DTPO| & 3 & \verb|CSCI181DVPO| & 3 & \verb|CSCI181G_PO| & 3\\\hline
    \verb|CSCI181NWPO| & 3 & \verb|CSCI181RTPO| & 3 & \verb|CSCI181R_PO| & 3\\\hline
    \verb|CSCI181S_PO| & 3 & \verb|CSCI181V_HM| & 3 & \verb|CSCI181V_PZ| & 3\\\hline
    \verb|CSCI181W_PO| & 3 & \verb|CSCI181__CM| & 3 & \verb|CSCI183__HM| & 3\\\hline
    \verb|CSCI184__HM| & 3 & \verb|CSCI186__HM| & 3 & \verb|CSCI188__PO| & 0.1\\\hline
    \verb|CSCI189__HM| & 1 & \verb|CSCI190__PO| & 0.1 & \verb|CSCI190__PO2| & 0.1\\\hline
    \verb|CSCI191__PO| & 3 & \verb|CSCI191__PO2| & 3 & \verb|CSCI192__PO| & 1.5\\\hline
    \verb|CSCI192__PO2| & 1.5 & \verb|CSCI195__HM| & 0.1 & \verb|CSCI195__HM2| & 0.1\\\hline
    \verb|CSCI195__HM3| & 0.1 & \verb|CSCI195__HM4| & 0.1 & \verb|CSCI198__SC| & 3\\\hline
    \verb|CSCI198__SC2| & 3 & \verb|CSCI900__HM| & 3 & \verb|CSCI901__HM| & 3\\\hline
    \verb|CSCI902__HM| & 3 & \verb|CSCI903__HM| & 3 & \verb|CSCI904__HM| & 3\\\hline
    \verb|CSMT183__HM| & 3 & \verb|CSMT184__HM| & 3 & \verb|DANC010__PO| & 3\\\hline
    \verb|DANC012P_PO| & 0.75 & \verb|DANC012__PO| & 1.5 & \verb|DANC050P_PO| & 0.75\\\hline
    \verb|DANC050__PO| & 1.5 & \verb|DANC051P_PO| & 0.75 & \verb|DANC051__PO| & 1.5\\\hline
    \verb|DANC068__SC| & 3 & \verb|DANC071A_SC| & 3 & \verb|DANC071B_SC| & 1.5\\\hline
    \verb|DANC081A_SC| & 3 & \verb|DANC081B_SC| & 1.5 & \verb|DANC083A_SC| & 3\\\hline
    \verb|DANC083B_SC| & 1.5 & \verb|DANC100A_SC| & 3 & \verb|DANC100B_SC| & 1.5\\\hline
    \verb|DANC102__SC| & 3 & \verb|DANC108A_SC| & 3 & \verb|DANC108B_SC| & 1.5\\\hline
    \verb|DANC111A_SC| & 3 & \verb|DANC111B_SC| & 1.5 & \verb|DANC114A_SC| & 3\\\hline
    \verb|DANC114B_SC| & 1.5 & \verb|DANC117B_SC| & 1.5 & \verb|DANC119__PO| & 3\\\hline
    \verb|DANC120P_PO| & 0.75 & \verb|DANC120__PO| & 1.5 & \verb|DANC122P_PO| & 0.75\\\hline
    \verb|DANC122__PO| & 1.5 & \verb|DANC123__PO| & 3 & \verb|DANC124P_PO| & 0.75\\\hline
    \verb|DANC124__PO| & 1.5 & \verb|DANC130__PO| & 3 & \verb|DANC131__SC| & 3\\\hline
    \verb|DANC135__PO| & 3 & \verb|DANC136B_SC| & 1.5 & \verb|DANC136__PO| & 3\\\hline
    \verb|DANC138__PO| & 3 & \verb|DANC139__PO| & 3 & \verb|DANC140__PO| & 3\\\hline
    \verb|DANC142P_PO| & 0.75 & \verb|DANC142__PO| & 1.5 & \verb|DANC150A_PO| & 3\\\hline
    \verb|DANC150C_PO| & 1.5 & \verb|DANC151P_PO| & 0.75 & \verb|DANC151__PO| & 1.5\\\hline
    \verb|DANC152IOSC| & 3 & \verb|DANC152P_PO| & 0.75 & \verb|DANC152__PO| & 1.5\\\hline
    \verb|DANC153P_PO| & 0.75 & \verb|DANC153__PO| & 1.5 & \verb|DANC159__SC| & 3\\\hline
    \verb|DANC160__PO| & 3 & \verb|DANC161__SC| & 3 & \verb|DANC162A_SC| & 3\\\hline
    \verb|DANC162B_SC| & 1.5 & \verb|DANC163__SC| & 3 & \verb|DANC166P_PO| & 0.75\\\hline
    \verb|DANC166__PO| & 1.5 & \verb|DANC170__PO| & 3 & \verb|DANC172__PO| & 0.75\\\hline
    \verb|DANC173__PO| & 1.5 & \verb|DANC174__PO| & 0.75 & \verb|DANC175__PO| & 3\\\hline
    \verb|DANC176__PO| & 1.5 & \verb|DANC180P_PO| & 0.75 & \verb|DANC180__PO| & 3\\\hline
    \verb|DANC180__SC| & 3 & \verb|DANC181P_PO| & 0.75 & \verb|DANC181__PO| & 1.5\\\hline
    \verb|DANC182__PO| & 0.75 & \verb|DANC190__SC| & 1.5 & \verb|DANC190__SC2| & 1.5\\\hline
    \verb|DANC191__SC| & 3 & \verb|DANC191__SC2| & 3 & \verb|DANC192__PO| & 1.5\\\hline
    \verb|DANC192__PO2| & 1.5 & \verb|DANC193__SC| & 0.1 & \verb|DANC193__SC2| & 0.1\\\hline
    \verb|DS__001__SC| & 3 & \verb|DS__002R_PO| & 3 & \verb|DS__002__SC| & 3\\\hline
    \verb|DS__108__PZ| & 3 & \verb|DS__180__CM| & 3 & \verb|DS__183__SC| & 3\\\hline
    \verb|DS__190__PO| & 1.5 & \verb|DS__190__PO2| & 1.5 & \verb|DS__191__SC| & 3\\\hline
    \verb|DS__191__SC2| & 3 & \verb|EA__010__PO| & 3 & \verb|EA__010__PZ| & 3\\\hline
    \verb|EA__010__SC| & 3 & \verb|EA__020__PO| & 3 & \verb|EA__030L_KS| & 3\\\hline
    \verb|EA__030L_PO| & 0.1 & \verb|EA__030__PO| & 3 & \verb|EA__031__PZ| & 3\\\hline
    \verb|EA__032__PZ| & 3 & \verb|EA__034__PZ| & 3 & \verb|EA__039__PZ| & 3\\\hline
    \verb|EA__055LCKS| & 3 & \verb|EA__086__PZ| & 3 & \verb|EA__086__SC| & 3\\\hline
    \verb|EA__097__PZ| & 3 & \verb|EA__100L_KS| & 3 & \verb|EA__100__KS| & 3\\\hline
    \verb|EA__101L_PO| & 0.1 & \verb|EA__101__PO| & 3 & \verb|EA__103L_KS| & 3\\\hline
    \verb|EA__103__KS| & 3 & \verb|EA__103__PZ| & 3 & \verb|EA__105__KS| & 3\\\hline
    \verb|EA__106__KS| & 3 & \verb|EA__106__PZ| & 3 & \verb|EA__112__PZ| & 3\\\hline
    \verb|EA__126__PZ| & 3 & \verb|EA__129__PZ| & 3 & \verb|EA__130__PZ| & 3\\\hline
    \verb|EA__133__PZ| & 3 & \verb|EA__134__PZ| & 3 & \verb|EA__135__PZ| & 3\\\hline
    \verb|EA__138__PZ| & 3 & \verb|EA__141__PZ| & 3 & \verb|EA__142__PZ| & 3\\\hline
    \verb|EA__143__PZ| & 3 & \verb|EA__144__PZ| & 3 & \verb|EA__150__PZ| & 3\\\hline
    \verb|EA__153__PZ| & 3 & \verb|EA__163__PZ| & 3 & \verb|EA__170__PO| & 3\\\hline
    \verb|EA__171__PO| & 3 & \verb|EA__180__PO| & 3 & \verb|EA__182__PO| & 0.75\\\hline
    \verb|EA__185__PO| & 3 & \verb|EA__188L_KS| & 3 & \verb|EA__189F_PO| & 3\\\hline
    \verb|EA__189G_PO| & 3 & \verb|EA__189L_KS| & 0.1 & \verb|EA__189M_PO| & 3\\\hline
    \verb|EA__190L_KS| & 3 & \verb|EA__190L_KS2| & 3 & \verb|EA__190__PO| & 3\\\hline
    \verb|EA__190__PO2| & 3 & \verb|EA__191H_PO| & 1.5 & \verb|EA__191H_PO2| & 1.5\\\hline
    \verb|EA__191__KS| & 3 & \verb|EA__191__KS2| & 3 & \verb|EA__191__PO| & 3\\\hline
    \verb|EA__191__PO2| & 3 & \verb|EA__196__KS| & 0.75 & \verb|EA__196__KS2| & 0.75\\\hline
    \verb|EA__197__KS| & 1.5 & \verb|EA__197__KS2| & 1.5 & \verb|EA__197__PZ| & 3\\\hline
    \verb|EA__197__PZ2| & 3 & \verb|ECON020__PZ| & 0.1 & \verb|ECON030__PZ| & 3\\\hline
    \verb|ECON050__CM| & 3 & \verb|ECON050__SC| & 3 & \verb|ECON051__PO| & 3\\\hline
    \verb|ECON051__PZ| & 3 & \verb|ECON051__SC| & 3 & \verb|ECON052__PO| & 3\\\hline
    \verb|ECON052__PZ| & 3 & \verb|ECON052__SC| & 3 & \verb|ECON053__HM| & 3\\\hline
    \verb|ECON054__HM| & 3 & \verb|ECON057B_PO| & 3 & \verb|ECON057__PO| & 3\\\hline
    \verb|ECON086__CM| & 3 & \verb|ECON089A_PO| & 3 & \verb|ECON091__PZ| & 3\\\hline
    \verb|ECON097__CM| & 3 & \verb|ECON100__CM| & 3 & \verb|ECON101__CM| & 3\\\hline
    \verb|ECON101__PO| & 3 & \verb|ECON101__SC| & 3 & \verb|ECON102__CM| & 3\\\hline
    \verb|ECON102__PO| & 3 & \verb|ECON102__SC| & 3 & \verb|ECON104__PZ| & 3\\\hline
    \verb|ECON105__PZ| & 3 & \verb|ECON107__CM| & 3 & \verb|ECON107__PO| & 3\\\hline
    \verb|ECON111__SC| & 3 & \verb|ECON114__SC| & 3 & \verb|ECON115__PO| & 3\\\hline
    \verb|ECON117__PO| & 3 & \verb|ECON118__CM| & 3 & \verb|ECON120__CM| & 3\\\hline
    \verb|ECON120__PO| & 3 & \verb|ECON120__SC| & 3 & \verb|ECON121__PO| & 3\\\hline
    \verb|ECON122__CM| & 3 & \verb|ECON122__PO| & 3 & \verb|ECON123__PO| & 3\\\hline
    \verb|ECON124__PO| & 3 & \verb|ECON125__CM| & 3 & \verb|ECON125__PO| & 3\\\hline
    \verb|ECON125__PZ| & 3 & \verb|ECON125__SC| & 3 & \verb|ECON126__CM| & 3\\\hline
    \verb|ECON126__PO| & 3 & \verb|ECON127__PO| & 3 & \verb|ECON128__PO| & 3\\\hline
    \verb|ECON129__CM| & 3 & \verb|ECON129__PO| & 3 & \verb|ECON130__PO| & 3\\\hline
    \verb|ECON131__PO| & 3 & \verb|ECON132__PO| & 3 & \verb|ECON133__SC| & 3\\\hline
    \verb|ECON134B_CM| & 3 & \verb|ECON134E_CM| & 3 & \verb|ECON134__CM| & 3\\\hline
    \verb|ECON134__SC| & 3 & \verb|ECON135__CM| & 3 & \verb|ECON135__SC| & 3\\\hline
    \verb|ECON136__CM| & 3 & \verb|ECON138__CM| & 3 & \verb|ECON139__CM| & 3\\\hline
    \verb|ECON140__CM| & 3 & \verb|ECON140__PZ| & 3 & \verb|ECON141__CM| & 3\\\hline
    \verb|ECON142__CM| & 3 & \verb|ECON142__PZ| & 3 & \verb|ECON143__PO| & 3\\\hline
    \verb|ECON145__CM| & 3 & \verb|ECON145__PZ| & 3 & \verb|ECON146__HM| & 3\\\hline
    \verb|ECON150__CM| & 3 & \verb|ECON150__PO| & 3 & \verb|ECON151__CM| & 3\\\hline
    \verb|ECON152__CM| & 3 & \verb|ECON152__PO| & 3 & \verb|ECON154__CM| & 3\\\hline
    \verb|ECON154__HM| & 3 & \verb|ECON154__PO| & 3 & \verb|ECON155__CM| & 3\\\hline
    \verb|ECON155__PO| & 3 & \verb|ECON156__PO| & 3 & \verb|ECON157__CM| & 3\\\hline
    \verb|ECON157__PO| & 3 & \verb|ECON158__CM| & 3 & \verb|ECON159__PO| & 3\\\hline
    \verb|ECON160__CM| & 3 & \verb|ECON161__PO| & 3 & \verb|ECON165__CM| & 3\\\hline
    \verb|ECON165__PO| & 3 & \verb|ECON165__PZ| & 3 & \verb|ECON166__PO| & 3\\\hline
    \verb|ECON167__PO| & 3 & \verb|ECON168__CM| & 3 & \verb|ECON170__SC| & 3\\\hline
    \verb|ECON171__CM| & 3 & \verb|ECON171__PZ| & 3 & \verb|ECON172__CM| & 3\\\hline
    \verb|ECON172__PZ| & 3 & \verb|ECON175__PZ| & 3 & \verb|ECON175__SC| & 3\\\hline
    \verb|ECON176__CM| & 3 & \verb|ECON177__CM| & 3 & \verb|ECON179F_HM| & 3\\\hline
    \verb|ECON179G_HM| & 3 & \verb|ECON179__CM| & 3 & \verb|ECON180__CM| & 1.5\\\hline
    \verb|ECON180__PZ| & 3 & \verb|ECON181__CM| & 3 & \verb|ECON184__SC| & 3\\\hline
    \verb|ECON185__CM| & 3 & \verb|ECON186__CM| & 3 & \verb|ECON190__PO| & 3\\\hline
    \verb|ECON190__PO2| & 3 & \verb|ECON191H_SC| & 3 & \verb|ECON191H_SC2| & 3\\\hline
    \verb|ECON191__CM| & 3 & \verb|ECON191__CM2| & 3 & \verb|ECON191__SC| & 3\\\hline
    \verb|ECON191__SC2| & 3 & \verb|ECON194A_CM| & 1.5 & \verb|ECON194A_CM2| & 1.5\\\hline
    \verb|ECON194B_CM| & 1.5 & \verb|ECON194B_CM2| & 1.5 & \verb|ECON195__PO| & 0.1\\\hline
    \verb|ECON195__PO2| & 0.1 & \verb|ECON196__CM| & 3 & \verb|ECON196__CM2| & 3\\\hline
    \verb|ECON197SBCM| & 1.5 & \verb|ECON197SBCM2| & 1.5 & \verb|ECON197S_CM| & 3\\\hline
    \verb|ECON197S_CM2| & 3 & \verb|ECON198__PZ| & 3 & \verb|ECON198__PZ2| & 3\\\hline
    \verb|ECON199__CM| & 1.5 & \verb|ECON199__CM2| & 1.5 & \verb|ECON199__PZ| & 0.1\\\hline
    \verb|ECON199__PZ2| & 0.1 & \verb|ECON316__CG| & 3 & \verb|EDUC301G_CG| & 3\\\hline
    \verb|EDUC424__CG| & 3 & \verb|EEP_191__SC| & 3 & \verb|EEP_191__SC2| & 3\\\hline
    \verb|ENGL001__PZ| & 3 & \verb|ENGL002__PZ| & 3 & \verb|ENGL010B_PZ| & 3\\\hline
    \verb|ENGL010__PO| & 3 & \verb|ENGL011A_PZ| & 3 & \verb|ENGL012B_AF| & 3\\\hline
    \verb|ENGL013__PZ| & 3 & \verb|ENGL015__PO| & 3 & \verb|ENGL016__PZ| & 3\\\hline
    \verb|ENGL017__PO| & 3 & \verb|ENGL018__PO| & 3 & \verb|ENGL019__PO| & 3\\\hline
    \verb|ENGL020__PO| & 3 & \verb|ENGL023__PO| & 3 & \verb|ENGL030__PZ| & 3\\\hline
    \verb|ENGL031__PZ| & 3 & \verb|ENGL036__PZ| & 3 & \verb|ENGL037__PZ| & 3\\\hline
    \verb|ENGL040__PO| & 3 & \verb|ENGL041__PO| & 3 & \verb|ENGL043__PO| & 3\\\hline
    \verb|ENGL044__PO| & 3 & \verb|ENGL045__PO| & 3 & \verb|ENGL046__PO| & 3\\\hline
    \verb|ENGL046__PZ| & 3 & \verb|ENGL047__PZ| & 3 & \verb|ENGL048__PO| & 3\\\hline
    \verb|ENGL048__PZ| & 3 & \verb|ENGL049__PZ| & 3 & \verb|ENGL055A_PO| & 3\\\hline
    \verb|ENGL060__PZ| & 3 & \verb|ENGL061__PZ| & 3 & \verb|ENGL064A_PO| & 3\\\hline
    \verb|ENGL064B_PO| & 3 & \verb|ENGL064C_PO| & 3 & \verb|ENGL072__PO| & 3\\\hline
    \verb|ENGL074__PO| & 3 & \verb|ENGL080__PO| & 3 & \verb|ENGL084__PO| & 3\\\hline
    \verb|ENGL086__PO| & 3 & \verb|ENGL087F_PO| & 3 & \verb|ENGL087H_PO| & 1.5\\\hline
    \verb|ENGL093__PO| & 3 & \verb|ENGL095__PO| & 3 & \verb|ENGL099__PO| & 3\\\hline
    \verb|ENGL101__PZ| & 3 & \verb|ENGL101__SC| & 3 & \verb|ENGL102__PZ| & 3\\\hline
    \verb|ENGL102__SC| & 3 & \verb|ENGL110__PZ| & 3 & \verb|ENGL111__SC| & 3\\\hline
    \verb|ENGL113S_SC| & 3 & \verb|ENGL115__SC| & 3 & \verb|ENGL116__SC| & 3\\\hline
    \verb|ENGL118S_SC| & 3 & \verb|ENGL120__PZ| & 3 & \verb|ENGL121__SC| & 3\\\hline
    \verb|ENGL124__SC| & 3 & \verb|ENGL125C_AF| & 3 & \verb|ENGL125__SC| & 3\\\hline
    \verb|ENGL128__PZ| & 3 & \verb|ENGL130__PZ| & 3 & \verb|ENGL131__SC| & 3\\\hline
    \verb|ENGL133S_SC| & 3 & \verb|ENGL141__PZ| & 3 & \verb|ENGL141__SC| & 3\\\hline
    \verb|ENGL145__SC| & 3 & \verb|ENGL146__PO| & 3 & \verb|ENGL148__SC| & 3\\\hline
    \verb|ENGL149S_SC| & 3 & \verb|ENGL151__PO| & 3 & \verb|ENGL151__PZ| & 3\\\hline
    \verb|ENGL153__PO| & 3 & \verb|ENGL154__SC| & 3 & \verb|ENGL157__PO| & 3\\\hline
    \verb|ENGL160S_SC| & 3 & \verb|ENGL161__PO| & 3 & \verb|ENGL162__PZ| & 3\\\hline
    \verb|ENGL163__PO| & 3 & \verb|ENGL164S_SC| & 3 & \verb|ENGL165__PO| & 3\\\hline
    \verb|ENGL165__SC| & 3 & \verb|ENGL168S_SC| & 3 & \verb|ENGL169__PO| & 3\\\hline
    \verb|ENGL170A_PO| & 3 & \verb|ENGL170H_PO| & 3 & \verb|ENGL170J_PO| & 3\\\hline
    \verb|ENGL170P_PO| & 3 & \verb|ENGL170R_PO| & 3 & \verb|ENGL170U_PO| & 3\\\hline
    \verb|ENGL170X_PO| & 3 & \verb|ENGL174__SC| & 3 & \verb|ENGL175__PZ| & 3\\\hline
    \verb|ENGL177__SC| & 3 & \verb|ENGL180__SC| & 3 & \verb|ENGL181__PZ| & 3\\\hline
    \verb|ENGL181__SC| & 3 & \verb|ENGL183A_PO| & 3 & \verb|ENGL183B_PO| & 3\\\hline
    \verb|ENGL188__PO| & 3 & \verb|ENGL190__SC| & 3 & \verb|ENGL190__SC2| & 3\\\hline
    \verb|ENGL191__PO| & 1.5 & \verb|ENGL191__PO2| & 1.5 & \verb|ENGL191__SC| & 3\\\hline
    \verb|ENGL191__SC2| & 3 & \verb|ENGL192__SC| & 3 & \verb|ENGL192__SC2| & 3\\\hline
    \verb|ENGL193__SC| & 3 & \verb|ENGL193__SC2| & 3 & \verb|ENGL194__PZ| & 3\\\hline
    \verb|ENGL194__PZ2| & 3 & \verb|ENGL194__SC| & 3 & \verb|ENGL194__SC2| & 3\\\hline
    \verb|ENGL195B_PO| & 1.5 & \verb|ENGL195B_PO2| & 1.5 & \verb|ENGL195__PO| & 3\\\hline
    \verb|ENGL195__PO2| & 3 & \verb|ENGL195__SC| & 3 & \verb|ENGL195__SC2| & 3\\\hline
    \verb|ENGL197N_SC| & 3 & \verb|ENGL197N_SC2| & 3 & \verb|ENGL199T_SC| & 3\\\hline
    \verb|ENGL199T_SC2| & 3 & \verb|ENGR004L_HM| & 0.1 & \verb|ENGR004__HM| & 4\\\hline
    \verb|ENGR015__HM| & 1 & \verb|ENGR025__HM| & 1 & \verb|ENGR026__HM| & 1\\\hline
    \verb|ENGR072__HM| & 3 & \verb|ENGR079P_HM| & 0.1 & \verb|ENGR079__HM| & 4\\\hline
    \verb|ENGR080__HM| & 3 & \verb|ENGR082__HM| & 3 & \verb|ENGR083__HM| & 3\\\hline
    \verb|ENGR084__HM| & 3 & \verb|ENGR085A_HM| & 1.5 & \verb|ENGR085__HM| & 3\\\hline
    \verb|ENGR086__HM| & 3 & \verb|ENGR091__HM| & 1 & \verb|ENGR101__HM| & 3\\\hline
    \verb|ENGR102__HM| & 3 & \verb|ENGR111__HM| & 3 & \verb|ENGR112__HM| & 3\\\hline
    \verb|ENGR113__HM| & 3 & \verb|ENGR114__HM| & 3 & \verb|ENGR122__HM| & 0.1\\\hline
    \verb|ENGR124__HM| & 0.1 & \verb|ENGR131__HM| & 3 & \verb|ENGR132__HM| & 3\\\hline
    \verb|ENGR133__HM| & 3 & \verb|ENGR134__HM| & 3 & \verb|ENGR138__HM| & 3\\\hline
    \verb|ENGR151__HM| & 4 & \verb|ENGR154__HM| & 3 & \verb|ENGR155__HM| & 4\\\hline
    \verb|ENGR157__HM| & 3 & \verb|ENGR164__HM| & 3 & \verb|ENGR171__HM| & 3\\\hline
    \verb|ENGR175__HM| & 3 & \verb|ENGR177__HM| & 3 & \verb|ENGR178__HM| & 3\\\hline
    \verb|ENGR180__HM| & 3 & \verb|ENGR181__HM| & 3 & \verb|ENGR183__HM| & 3\\\hline
    \verb|ENGR185A_HM| & 1.5 & \verb|ENGR185B_HM| & 1.5 & \verb|ENGR186__HM| & 3\\\hline
    \verb|ENGR187__HM| & 3 & \verb|ENGR190BAHM| & 3 & \verb|ENGR190BAHM2| & 3\\\hline
    \verb|ENGR190BFHM| & 1 & \verb|ENGR190BFHM2| & 1 & \verb|ENGR190BGHM| & 3\\\hline
    \verb|ENGR190BGHM2| & 3 & \verb|ENGR190BHHM| & 3 & \verb|ENGR190BHHM2| & 3\\\hline
    \verb|ENGR190BIHM| & 3 & \verb|ENGR190BIHM2| & 3 & \verb|ENGR190BJHM| & 3\\\hline
    \verb|ENGR190BJHM2| & 3 & \verb|ENGR191__HM| & 3 & \verb|ENGR191__HM2| & 3\\\hline
    \verb|ENGR205__HM| & 3 & \verb|ENGR207__HM| & 3 & \verb|ENGR208__HM| & 3\\\hline
    \verb|ENGR278__HM| & 3 & \verb|ENGR900__HM| & 3 & \verb|ENGR901__HM| & 3\\\hline
    \verb|ENGR902__HM| & 3 & \verb|ENGR903__HM| & 3 & \verb|ENGR904__HM| & 3\\\hline
    \verb|ENG_5300_KG| & 3 & \verb|ENTR179A_HM| & 1.5 & \verb|ENTR179B_HM| & 1.5\\\hline
    \verb|FGSS026__SC| & 3 & \verb|FGSS036__SC| & 3 & \verb|FGSS177__SC| & 3\\\hline
    \verb|FGSS183__SC| & 3 & \verb|FGSS184__SC| & 3 & \verb|FGSS186__SC| & 3\\\hline
    \verb|FGSS188E_SC| & 3 & \verb|FGSS188__SC| & 3 & \verb|FGSS191__SC| & 3\\\hline
    \verb|FGSS191__SC2| & 3 & \verb|FGSS192__SC| & 3 & \verb|FGSS192__SC2| & 3\\\hline
    \verb|FGSS198__SC| & 3 & \verb|FGSS198__SC2| & 3 & \verb|FHS_010__CM| & 3\\\hline
    \verb|FLAN191__SC| & 3 & \verb|FLAN191__SC2| & 3 & \verb|FREN001L_CM| & 0.1\\\hline
    \verb|FREN001L_SC| & 0.1 & \verb|FREN001__PO| & 3 & \verb|FREN001__PZ| & 3\\\hline
    \verb|FREN001__SC| & 3 & \verb|FREN002L_CM| & 0.1 & \verb|FREN002L_SC| & 0.1\\\hline
    \verb|FREN002__PO| & 3 & \verb|FREN002__PZ| & 3 & \verb|FREN002__SC| & 3\\\hline
    \verb|FREN011__PO| & 0.75 & \verb|FREN013__PO| & 0.75 & \verb|FREN022__PO| & 3\\\hline
    \verb|FREN033L_CM| & 0.1 & \verb|FREN033L_SC| & 0.1 & \verb|FREN033__CM| & 3\\\hline
    \verb|FREN033__PO| & 3 & \verb|FREN033__SC| & 3 & \verb|FREN044L_CM| & 0.1\\\hline
    \verb|FREN044L_SC| & 0.1 & \verb|FREN044__CM| & 3 & \verb|FREN044__PO| & 3\\\hline
    \verb|FREN044__SC| & 3 & \verb|FREN045L_SC| & 0.1 & \verb|FREN045__SC| & 3\\\hline
    \verb|FREN100__CM| & 3 & \verb|FREN100__PZ| & 3 & \verb|FREN101A_PO| & 3\\\hline
    \verb|FREN101B_PO| & 3 & \verb|FREN103__PO| & 3 & \verb|FREN104__PO| & 3\\\hline
    \verb|FREN105__PO| & 3 & \verb|FREN107__PO| & 3 & \verb|FREN110__PO| & 3\\\hline
    \verb|FREN112__PO| & 3 & \verb|FREN113__SC| & 3 & \verb|FREN115__CM| & 3\\\hline
    \verb|FREN125__SC| & 3 & \verb|FREN126__SC| & 3 & \verb|FREN127__SC| & 3\\\hline
    \verb|FREN129__PO| & 3 & \verb|FREN133__CM| & 3 & \verb|FREN135__CM| & 3\\\hline
    \verb|FREN150C_PO| & 3 & \verb|FREN153__SC| & 3 & \verb|FREN161__PZ| & 3\\\hline
    \verb|FREN179__SC| & 3 & \verb|FREN180__PO| & 3 & \verb|FREN181__SC| & 3\\\hline
    \verb|FREN185__PO| & 3 & \verb|FREN187A_SC| & 3 & \verb|FREN188__PO| & 3\\\hline
    \verb|FREN191__PO| & 1.5 & \verb|FREN191__PO2| & 1.5 & \verb|FREN191__SC| & 3\\\hline
    \verb|FREN191__SC2| & 3 & \verb|FREN192__PO| & 1.5 & \verb|FREN192__PO2| & 1.5\\\hline
    \verb|FREN193__PO| & 0.1 & \verb|FREN193__PO2| & 0.1 & \verb|FWS_010__CM| & 3\\\hline
    \verb|GEOG145__HM| & 3 & \verb|GEOG179H_HM| & 3 & \verb|GEOG179I_HM| & 3\\\hline
    \verb|GEOG179J_HM| & 3 & \verb|GEOG179K_HM| & 3 & \verb|GEOL015__PO| & 3\\\hline
    \verb|GEOL020A_PO| & 3 & \verb|GEOL020C_PO| & 3 & \verb|GEOL020E_PO| & 3\\\hline
    \verb|GEOL025__PO| & 3 & \verb|GEOL115__PO| & 3 & \verb|GEOL125__PO| & 3\\\hline
    \verb|GEOL127__PO| & 3 & \verb|GEOL129__PO| & 3 & \verb|GEOL131__PO| & 3\\\hline
    \verb|GEOL133__PO| & 3 & \verb|GEOL143__PO| & 3 & \verb|GEOL181__PO| & 3\\\hline
    \verb|GEOL183__PO| & 3 & \verb|GEOL185__PO| & 3 & \verb|GEOL189C_PO| & 3\\\hline
    \verb|GEOL189G_PO| & 3 & \verb|GEOL189I_PO| & 3 & \verb|GEOL192__PO| & 1.5\\\hline
    \verb|GEOL192__PO2| & 1.5 & \verb|GERM001__PO| & 3 & \verb|GERM001__SC| & 3\\\hline
    \verb|GERM002__PO| & 3 & \verb|GERM002__SC| & 3 & \verb|GERM011__PO| & 0.75\\\hline
    \verb|GERM012__SC| & 0.75 & \verb|GERM013__PO| & 0.75 & \verb|GERM033__PO| & 3\\\hline
    \verb|GERM033__SC| & 3 & \verb|GERM044__SC| & 3 & \verb|GERM102__PO| & 3\\\hline
    \verb|GERM107__SC| & 3 & \verb|GERM112__SC| & 3 & \verb|GERM152__PO| & 3\\\hline
    \verb|GERM154F_PO| & 3 & \verb|GERM191__PO| & 1.5 & \verb|GERM191__PO2| & 1.5\\\hline
    \verb|GERM191__SC| & 3 & \verb|GERM191__SC2| & 3 & \verb|GERM193__PO| & 1.5\\\hline
    \verb|GERM193__PO2| & 1.5 & \verb|GFS_036__PZ| & 3 & \verb|GFS_170__PZ| & 3\\\hline
    \verb|GLAS001__PZ| & 0.75 & \verb|GLAS180IOPZ| & 3 & \verb|GLAS194A_PZ| & 1.5\\\hline
    \verb|GLAS194A_PZ2| & 1.5 & \verb|GLAS194B_PZ| & 1.5 & \verb|GLAS194B_PZ2| & 1.5\\\hline
    \verb|GOVT020H_CM| & 3 & \verb|GOVT020__CM| & 3 & \verb|GOVT041B_CM| & 0.1\\\hline
    \verb|GOVT041__CM| & 1.5 & \verb|GOVT054__CM| & 3 & \verb|GOVT055__CM| & 3\\\hline
    \verb|GOVT060__CM| & 3 & \verb|GOVT070H_CM| & 3 & \verb|GOVT070__CM| & 3\\\hline
    \verb|GOVT071__CM| & 1.5 & \verb|GOVT080__CM| & 3 & \verb|GOVT094__CM| & 3\\\hline
    \verb|GOVT097__CM| & 3 & \verb|GOVT100__CM| & 3 & \verb|GOVT101__CM| & 3\\\hline
    \verb|GOVT102__CM| & 3 & \verb|GOVT103__CM| & 3 & \verb|GOVT107__CM| & 3\\\hline
    \verb|GOVT110__CM| & 3 & \verb|GOVT112A_CM| & 3 & \verb|GOVT112B_CM| & 3\\\hline
    \verb|GOVT115__CM| & 3 & \verb|GOVT116__CM| & 3 & \verb|GOVT117__CM| & 3\\\hline
    \verb|GOVT118__CM| & 3 & \verb|GOVT121__CM| & 3 & \verb|GOVT123__CM| & 3\\\hline
    \verb|GOVT128__JT| & 3 & \verb|GOVT129A_CM| & 3 & \verb|GOVT129E_CM| & 3\\\hline
    \verb|GOVT129__CM| & 3 & \verb|GOVT131__CM| & 3 & \verb|GOVT132E_CM| & 3\\\hline
    \verb|GOVT134__CM| & 3 & \verb|GOVT136C_CM| & 3 & \verb|GOVT137A_CM| & 3\\\hline
    \verb|GOVT137B_CM| & 3 & \verb|GOVT137__CM| & 3 & \verb|GOVT138C_CM| & 3\\\hline
    \verb|GOVT139__CM| & 3 & \verb|GOVT140__CM| & 3 & \verb|GOVT141__CM| & 3\\\hline
    \verb|GOVT142C_CM| & 3 & \verb|GOVT142E_CM| & 3 & \verb|GOVT142__CM| & 3\\\hline
    \verb|GOVT143C_CM| & 3 & \verb|GOVT144D_CM| & 3 & \verb|GOVT145E_CM| & 3\\\hline
    \verb|GOVT146A_CM| & 3 & \verb|GOVT146__CM| & 3 & \verb|GOVT147__CM| & 3\\\hline
    \verb|GOVT148__CM| & 3 & \verb|GOVT149__CM| & 3 & \verb|GOVT150C_CM| & 3\\\hline
    \verb|GOVT151__CM| & 3 & \verb|GOVT152C_CM| & 3 & \verb|GOVT152__CM| & 3\\\hline
    \verb|GOVT154E_CM| & 3 & \verb|GOVT156E_CM| & 3 & \verb|GOVT157S_CM| & 3\\\hline
    \verb|GOVT164__CM| & 3 & \verb|GOVT166__CM| & 3 & \verb|GOVT173C_CM| & 3\\\hline
    \verb|GOVT174C_CM| & 3 & \verb|GOVT177C_CM| & 3 & \verb|GOVT177__CM| & 3\\\hline
    \verb|GOVT178__CM| & 3 & \verb|GOVT179C_CM| & 3 & \verb|GOVT179__CM| & 3\\\hline
    \verb|GOVT182__CM| & 3 & \verb|GOVT191__CM| & 3 & \verb|GOVT191__CM2| & 3\\\hline
    \verb|GOVT192__CM| & 3 & \verb|GOVT192__CM2| & 3 & \verb|GOVT999__CM| & 0.1\\\hline
    \verb|GREK022__PO| & 3 & \verb|GREK022__SC| & 3 & \verb|GREK033__PO| & 3\\\hline
    \verb|GREK033__SC| & 3 & \verb|GREK044__PO| & 3 & \verb|GREK044__SC| & 3\\\hline
    \verb|GRMT014__PO| & 1.5 & \verb|GRMT114__SC| & 3 & \verb|GRMT115__SC| & 3\\\hline
    \verb|GRMT131__PO| & 3 & \verb|GRMT164__PO| & 3 & \verb|GWS_026__PO| & 3\\\hline
    \verb|GWS_036__PO| & 3 & \verb|GWS_070__PO| & 3 & \verb|GWS_081__PO| & 3\\\hline
    \verb|GWS_140__PO| & 3 & \verb|GWS_162__PO| & 3 & \verb|GWS_166__PO| & 3\\\hline
    \verb|GWS_170__PO| & 3 & \verb|GWS_172__PO| & 3 & \verb|GWS_173__PO| & 3\\\hline
    \verb|GWS_179A_HM| & 3 & \verb|GWS_179B_HM| & 3 & \verb|GWS_180__PO| & 3\\\hline
    \verb|GWS_183__PO| & 3 & \verb|GWS_184__PO| & 3 & \verb|GWS_190__PO| & 3\\\hline
    \verb|GWS_190__PO2| & 3 & \verb|GWS_191__PO| & 3 & \verb|GWS_191__PO2| & 3\\\hline
    \verb|HIST008__PO| & 3 & \verb|HIST009__SC| & 3 & \verb|HIST011__PO| & 3\\\hline
    \verb|HIST011__PZ| & 3 & \verb|HIST012__PZ| & 3 & \verb|HIST014__PO| & 3\\\hline
    \verb|HIST017__CH| & 3 & \verb|HIST020__PO| & 3 & \verb|HIST021__PO| & 3\\\hline
    \verb|HIST025__CH| & 3 & \verb|HIST025__PZ| & 3 & \verb|HIST031__CH| & 3\\\hline
    \verb|HIST032__CH| & 3 & \verb|HIST034__CH| & 3 & \verb|HIST036__PO| & 3\\\hline
    \verb|HIST037__PO| & 3 & \verb|HIST040__AF| & 3 & \verb|HIST041__AF| & 3\\\hline
    \verb|HIST043__PO| & 3 & \verb|HIST044__PZ| & 3 & \verb|HIST046__PZ| & 3\\\hline
    \verb|HIST047__PO| & 3 & \verb|HIST048__SC| & 3 & \verb|HIST054__CM| & 3\\\hline
    \verb|HIST055__CM| & 3 & \verb|HIST056__CM| & 3 & \verb|HIST059__CM| & 3\\\hline
    \verb|HIST060__PO| & 3 & \verb|HIST060__PZ| & 3 & \verb|HIST061__PO| & 3\\\hline
    \verb|HIST062__PO| & 3 & \verb|HIST064__PO| & 3 & \verb|HIST065__PZ| & 3\\\hline
    \verb|HIST066__PZ| & 3 & \verb|HIST069__PZ| & 3 & \verb|HIST070B_SC| & 3\\\hline
    \verb|HIST070__PO| & 3 & \verb|HIST071__PO| & 3 & \verb|HIST072__CM| & 3\\\hline
    \verb|HIST072__PO| & 3 & \verb|HIST081__HM| & 3 & \verb|HIST082__HM| & 3\\\hline
    \verb|HIST083__PZ| & 3 & \verb|HIST090__CM| & 3 & \verb|HIST096__CM| & 3\\\hline
    \verb|HIST097__CM| & 3 & \verb|HIST099__CM| & 3 & \verb|HIST100__CM| & 3\\\hline
    \verb|HIST101ADPO| & 3 & \verb|HIST101A_PO| & 3 & \verb|HIST101B_PO| & 3\\\hline
    \verb|HIST101D_PO| & 3 & \verb|HIST101F_PO| & 3 & \verb|HIST101HAPO| & 3\\\hline
    \verb|HIST101H_PO| & 3 & \verb|HIST101N_PO| & 3 & \verb|HIST101Q_PO| & 3\\\hline
    \verb|HIST101S_CH| & 3 & \verb|HIST101T_CH| & 3 & \verb|HIST101U_PO| & 3\\\hline
    \verb|HIST101Y_PO| & 3 & \verb|HIST101__CM| & 3 & \verb|HIST102__PO| & 3\\\hline
    \verb|HIST105__CM| & 3 & \verb|HIST105__PO| & 3 & \verb|HIST106__CM| & 3\\\hline
    \verb|HIST108__CM| & 3 & \verb|HIST108__SC| & 3 & \verb|HIST109__CM| & 3\\\hline
    \verb|HIST111__SC| & 3 & \verb|HIST112A_CM| & 3 & \verb|HIST112__CM| & 3\\\hline
    \verb|HIST113__CM| & 3 & \verb|HIST113__PO| & 3 & \verb|HIST113__SC| & 3\\\hline
    \verb|HIST114C_CM| & 3 & \verb|HIST114__CM| & 3 & \verb|HIST115__CM| & 3\\\hline
    \verb|HIST116IOCM| & 3 & \verb|HIST116__CM| & 3 & \verb|HIST117__SC| & 3\\\hline
    \verb|HIST118__PO| & 3 & \verb|HIST119__CM| & 3 & \verb|HIST120E_CM| & 3\\\hline
    \verb|HIST120__CM| & 3 & \verb|HIST121__CM| & 3 & \verb|HIST122__CM| & 3\\\hline
    \verb|HIST123__CM| & 3 & \verb|HIST123__PO| & 3 & \verb|HIST125__CM| & 3\\\hline
    \verb|HIST125__PO| & 3 & \verb|HIST126__CM| & 3 & \verb|HIST128__CM| & 3\\\hline
    \verb|HIST128__PO| & 3 & \verb|HIST129__CM| & 3 & \verb|HIST130__CH| & 3\\\hline
    \verb|HIST132__CM| & 3 & \verb|HIST133A_CM| & 3 & \verb|HIST133B_CM| & 3\\\hline
    \verb|HIST133__AF| & 3 & \verb|HIST134__CM| & 3 & \verb|HIST134__SC| & 3\\\hline
    \verb|HIST135__CM| & 3 & \verb|HIST136C_CM| & 3 & \verb|HIST139E_CM| & 3\\\hline
    \verb|HIST139__PO| & 3 & \verb|HIST140C_CM| & 3 & \verb|HIST140__AF| & 3\\\hline
    \verb|HIST141__SC| & 3 & \verb|HIST142__CM| & 3 & \verb|HIST142__PZ| & 3\\\hline
    \verb|HIST143__SC| & 3 & \verb|HIST144__SC| & 3 & \verb|HIST146__AF| & 3\\\hline
    \verb|HIST146__PO| & 3 & \verb|HIST146__SC| & 3 & \verb|HIST147__CM| & 3\\\hline
    \verb|HIST147__PO| & 3 & \verb|HIST148__PO| & 3 & \verb|HIST149A_SC| & 3\\\hline
    \verb|HIST149__SC| & 3 & \verb|HIST150C_CM| & 3 & \verb|HIST150D_CM| & 3\\\hline
    \verb|HIST150__CM| & 3 & \verb|HIST150__PZ| & 3 & \verb|HIST151__CM| & 3\\\hline
    \verb|HIST151__HM| & 3 & \verb|HIST151__SC| & 3 & \verb|HIST154__CM| & 3\\\hline
    \verb|HIST154__SC| & 3 & \verb|HIST155__SC| & 3 & \verb|HIST157__CM| & 3\\\hline
    \verb|HIST157__PZ| & 3 & \verb|HIST158__CM| & 3 & \verb|HIST158__PZ| & 3\\\hline
    \verb|HIST159__SC| & 3 & \verb|HIST160__PO| & 3 & \verb|HIST161__CM| & 3\\\hline
    \verb|HIST164__SC| & 3 & \verb|HIST165__PO| & 3 & \verb|HIST166__CM| & 3\\\hline
    \verb|HIST167__PO| & 3 & \verb|HIST168__CM| & 3 & \verb|HIST168__PO| & 3\\\hline
    \verb|HIST169__CM| & 3 & \verb|HIST169__PO| & 3 & \verb|HIST170__PO| & 3\\\hline
    \verb|HIST171__PO| & 3 & \verb|HIST173__CM| & 3 & \verb|HIST174__CM| & 3\\\hline
    \verb|HIST174__PO| & 3 & \verb|HIST175__CM| & 3 & \verb|HIST175__PO| & 3\\\hline
    \verb|HIST176__CM| & 3 & \verb|HIST178__CM| & 3 & \verb|HIST178__PZ| & 3\\\hline
    \verb|HIST179L_HM| & 3 & \verb|HIST179__CM| & 3 & \verb|HIST180__PO| & 3\\\hline
    \verb|HIST180__SC| & 3 & \verb|HIST181__PZ| & 3 & \verb|HIST182__CM| & 3\\\hline
    \verb|HIST184__PO| & 3 & \verb|HIST185__PO| & 3 & \verb|HIST186__PO| & 3\\\hline
    \verb|HIST188__CM| & 3 & \verb|HIST189F_PO| & 3 & \verb|HIST189G_PO| & 3\\\hline
    \verb|HIST190__CM| & 3 & \verb|HIST190__CM2| & 3 & \verb|HIST190__PO| & 3\\\hline
    \verb|HIST190__PO2| & 3 & \verb|HIST191__PO| & 3 & \verb|HIST191__PO2| & 3\\\hline
    \verb|HIST191__SC| & 3 & \verb|HIST191__SC2| & 3 & \verb|HIST192__PO| & 3\\\hline
    \verb|HIST192__PO2| & 3 & \verb|HIST192__SC| & 3 & \verb|HIST192__SC2| & 3\\\hline
    \verb|HIST196__CM| & 3 & \verb|HIST196__CM2| & 3 & \verb|HIST197__CM| & 3\\\hline
    \verb|HIST197__CM2| & 3 & \verb|HIST197__PZ| & 3 & \verb|HIST197__PZ2| & 3\\\hline
    \verb|HIST199__PZ| & 3 & \verb|HIST199__PZ2| & 3 & \verb|HMSC133__SC| & 3\\\hline
    \verb|HMSC138__SC| & 3 & \verb|HMSC191__SC| & 3 & \verb|HMSC191__SC2| & 3\\\hline
    \verb|HMSC192__SC| & 3 & \verb|HMSC192__SC2| & 3 & \verb|HSA_000__5C| & 3\\\hline
    \verb|HSA_000__HM| & 3 & \verb|HSA_001__5C| & 3 & \verb|HSA_001__HM| & 3\\\hline
    \verb|HSA_002__5C| & 3 & \verb|HSA_002__HM| & 3 & \verb|HSA_003__5C| & 3\\\hline
    \verb|HSA_003__HM| & 3 & \verb|HSA_004__5C| & 3 & \verb|HSA_005__5C| & 3\\\hline
    \verb|HSA_006__5C| & 3 & \verb|HSA_007__5C| & 3 & \verb|HSA_008__5C| & 3\\\hline
    \verb|HSA_009__5C| & 3 & \verb|HSA_010__HM| & 3 & \verb|HSA_100__HM| & 3\\\hline
    \verb|HUM_195J_SC| & 3 & \verb|HUM_195J_SC2| & 3 & \verb|HUM_196__PO| & 1.5\\\hline
    \verb|HUM_196__PO2| & 1.5 & \verb|ID__009__PO| & 0.75 & \verb|ID__076__JT| & 3\\\hline
    \verb|ID__077__JT| & 3 & \verb|ID__080__CM| & 1.5 & \verb|ID__081__CM| & 1.5\\\hline
    \verb|ID__191D_SC| & 3 & \verb|ID__191D_SC2| & 3 & \verb|ID__191H_SC| & 3\\\hline
    \verb|ID__191H_SC2| & 3 & \verb|ID__191S_SC| & 3 & \verb|ID__191S_SC2| & 3\\\hline
    \verb|ID__199P1PO| & 1.5 & \verb|ID__199P1PO2| & 1.5 & \verb|ID__199P2PO| & 1.5\\\hline
    \verb|ID__199P2PO2| & 1.5 & \verb|ID__199P3PO| & 1.5 & \verb|ID__199P3PO2| & 1.5\\\hline
    \verb|ID__199P4PO| & 1.5 & \verb|ID__199P4PO2| & 1.5 & \verb|ID__199S1PO| & 1.5\\\hline
    \verb|ID__199S1PO2| & 1.5 & \verb|ID__199S2PO| & 1.5 & \verb|ID__199S2PO2| & 1.5\\\hline
    \verb|ID__199S3PO| & 1.5 & \verb|ID__199S3PO2| & 1.5 & \verb|ID__199S4PO| & 1.5\\\hline
    \verb|ID__199S4PO2| & 1.5 & \verb|IR__100__PO| & 3 & \verb|IR__101__PO| & 3\\\hline
    \verb|IR__118__PO| & 3 & \verb|IR__190__PO| & 3 & \verb|IR__190__PO2| & 3\\\hline
    \verb|IR__191__PO| & 3 & \verb|IR__191__PO2| & 3 & \verb|IR__192__SC| & 3\\\hline
    \verb|IR__192__SC2| & 3 & \verb|IST_302__CG| & 3 & \verb|ITAL001__SC| & 3\\\hline
    \verb|ITAL002__SC| & 3 & \verb|ITAL022__SC| & 3 & \verb|ITAL033__SC| & 3\\\hline
    \verb|ITAL044__SC| & 3 & \verb|ITAL131__SC| & 3 & \verb|ITAL134__SC| & 3\\\hline
    \verb|ITAL140__SC| & 3 & \verb|ITAL141__SC| & 3 & \verb|ITAL150__SC| & 3\\\hline
    \verb|ITAL187C_SC| & 3 & \verb|ITAL191__SC| & 3 & \verb|ITAL191__SC2| & 3\\\hline
    \verb|JAPN001A_PO| & 3 & \verb|JAPN001B_PO| & 3 & \verb|JAPN011__PO| & 0.75\\\hline
    \verb|JAPN012A_PO| & 0.75 & \verb|JAPN012B_PO| & 0.75 & \verb|JAPN013__PO| & 0.75\\\hline
    \verb|JAPN014A_PO| & 0.75 & \verb|JAPN014B_PO| & 0.75 & \verb|JAPN051A_PO| & 3\\\hline
    \verb|JAPN051B_PO| & 3 & \verb|JAPN111A_PO| & 3 & \verb|JAPN111B_PO| & 3\\\hline
    \verb|JAPN125__PO| & 3 & \verb|JAPN127__PO| & 3 & \verb|JPNT171__PO| & 3\\\hline
    \verb|JPNT175__PO| & 3 & \verb|KORE001__CM| & 3 & \verb|KORE002__CM| & 3\\\hline
    \verb|KORE033__CM| & 3 & \verb|KORE044__CM| & 3 & \verb|KORE120__CM| & 3\\\hline
    \verb|KRNT130__CM| & 3 & \verb|LAMS191__PO| & 3 & \verb|LAMS191__PO2| & 3\\\hline
    \verb|LAST184__PO| & 3 & \verb|LAST190__PO| & 3 & \verb|LAST190__PO2| & 3\\\hline
    \verb|LAST191__PO| & 3 & \verb|LAST191__PO2| & 3 & \verb|LAST191__SC| & 3\\\hline
    \verb|LAST191__SC2| & 3 & \verb|LAS_180__PO| & 3 & \verb|LATN022__PO| & 3\\\hline
    \verb|LATN033__PO| & 3 & \verb|LATN033__SC| & 3 & \verb|LATN044__PO| & 3\\\hline
    \verb|LATN044__SC| & 3 & \verb|LATN103__PO| & 1.5 & \verb|LEAD010__CM| & 3\\\hline
    \verb|LEAD041__CM| & 1.5 & \verb|LEAD101__HM| & 3 & \verb|LEAD134__CM| & 3\\\hline
    \verb|LEAD143__CM| & 3 & \verb|LEAD144__CM| & 3 & \verb|LEAD149__CM| & 3\\\hline
    \verb|LEAD151__HM| & 3 & \verb|LGCS010__PO| & 3 & \verb|LGCS010__PZ| & 3\\\hline
    \verb|LGCS011__PO| & 3 & \verb|LGCS101__PO| & 3 & \verb|LGCS104__PO| & 3\\\hline
    \verb|LGCS105__PO| & 3 & \verb|LGCS106__PO| & 3 & \verb|LGCS108__PO| & 3\\\hline
    \verb|LGCS110__PZ| & 3 & \verb|LGCS112__PZ| & 3 & \verb|LGCS113__PO| & 3\\\hline
    \verb|LGCS114__PO| & 3 & \verb|LGCS115__PZ| & 3 & \verb|LGCS117__PO| & 3\\\hline
    \verb|LGCS118__PO| & 3 & \verb|LGCS119__PO| & 3 & \verb|LGCS120__PO| & 3\\\hline
    \verb|LGCS121__PO| & 3 & \verb|LGCS123__PO| & 3 & \verb|LGCS124__PO| & 3\\\hline
    \verb|LGCS125__PO| & 3 & \verb|LGCS128__PO| & 3 & \verb|LGCS132__PO| & 3\\\hline
    \verb|LGCS135__PO| & 3 & \verb|LGCS145__PO| & 3 & \verb|LGCS166__PZ| & 3\\\hline
    \verb|LGCS183__PO| & 3 & \verb|LGCS185__PO| & 3 & \verb|LGCS188__PO| & 3\\\hline
    \verb|LGCS190__PO| & 3 & \verb|LGCS190__PO2| & 3 & \verb|LGCS191__PO| & 1.5\\\hline
    \verb|LGCS191__PO2| & 1.5 & \verb|LIT_031__CM| & 3 & \verb|LIT_034__CM| & 3\\\hline
    \verb|LIT_035__HM| & 3 & \verb|LIT_037__CM| & 3 & \verb|LIT_038__CM| & 3\\\hline
    \verb|LIT_057__CM| & 3 & \verb|LIT_058__CM| & 3 & \verb|LIT_060__CM| & 3\\\hline
    \verb|LIT_061__CM| & 3 & \verb|LIT_062__CM| & 3 & \verb|LIT_066__CM| & 3\\\hline
    \verb|LIT_067__CM| & 3 & \verb|LIT_068__CM| & 3 & \verb|LIT_081__CM| & 3\\\hline
    \verb|LIT_090__CM| & 3 & \verb|LIT_093__CM| & 3 & \verb|LIT_097__CM| & 3\\\hline
    \verb|LIT_099A_CM| & 3 & \verb|LIT_099B_CM| & 3 & \verb|LIT_099C_CM| & 3\\\hline
    \verb|LIT_099__CM| & 3 & \verb|LIT_100__CM| & 3 & \verb|LIT_102__CM| & 3\\\hline
    \verb|LIT_107__HM| & 3 & \verb|LIT_110__HM| & 3 & \verb|LIT_111__CM| & 3\\\hline
    \verb|LIT_118__CM| & 3 & \verb|LIT_119__CM| & 3 & \verb|LIT_120__CM| & 3\\\hline
    \verb|LIT_125__CM| & 3 & \verb|LIT_129IOCM| & 3 & \verb|LIT_129__CM| & 3\\\hline
    \verb|LIT_130__CM| & 3 & \verb|LIT_131__CM| & 3 & \verb|LIT_132__CM| & 3\\\hline
    \verb|LIT_134C_CM| & 3 & \verb|LIT_134__CM| & 3 & \verb|LIT_135__CM| & 3\\\hline
    \verb|LIT_138__CM| & 3 & \verb|LIT_139__CM| & 3 & \verb|LIT_141__HM| & 3\\\hline
    \verb|LIT_143__CM| & 3 & \verb|LIT_144__HM| & 3 & \verb|LIT_148__CM| & 3\\\hline
    \verb|LIT_152__CM| & 3 & \verb|LIT_154__CM| & 3 & \verb|LIT_156__HM| & 3\\\hline
    \verb|LIT_161__AF| & 3 & \verb|LIT_162__AF| & 3 & \verb|LIT_163__AF| & 3\\\hline
    \verb|LIT_165__AF| & 3 & \verb|LIT_179ABHM| & 3 & \verb|LIT_179AHHM| & 3\\\hline
    \verb|LIT_179AIHM| & 3 & \verb|LIT_179AJHM| & 3 & \verb|LIT_179AKHM| & 3\\\hline
    \verb|LIT_179Y_HM| & 3 & \verb|LIT_182__CM| & 3 & \verb|LIT_183__CM| & 3\\\hline
    \verb|LIT_185__CM| & 3 & \verb|LIT_199__CM| & 1.5 & \verb|LIT_199__CM2| & 1.5\\\hline
    \verb|MATH001__PO| & 3 & \verb|MATH015B_PZ| & 3 & \verb|MATH019L_HM| & 0.1\\\hline
    \verb|MATH019__HM| & 4 & \verb|MATH022__SC| & 3 & \verb|MATH023__SC| & 3\\\hline
    \verb|MATH025__PZ| & 3 & \verb|MATH030__CM| & 3 & \verb|MATH030__PO| & 3\\\hline
    \verb|MATH030__PZ| & 3 & \verb|MATH030__SC| & 3 & \verb|MATH031A_CM| & 3\\\hline
    \verb|MATH031H_PO| & 3 & \verb|MATH031__CM| & 3 & \verb|MATH031__PO| & 3\\\hline
    \verb|MATH031__PZ| & 3 & \verb|MATH031__SC| & 3 & \verb|MATH032H_CM| & 3\\\hline
    \verb|MATH032S_PO| & 3 & \verb|MATH032__CM| & 3 & \verb|MATH032__PO| & 3\\\hline
    \verb|MATH032__PZ| & 3 & \verb|MATH032__SC| & 3 & \verb|MATH052__CM| & 3\\\hline
    \verb|MATH052__PZ| & 3 & \verb|MATH055__CM| & 3 & \verb|MATH055__HM| & 3\\\hline
    \verb|MATH055__PZ| & 3 & \verb|MATH055__SC| & 3 & \verb|MATH056__HM| & 3\\\hline
    \verb|MATH058B_PO| & 3 & \verb|MATH058__PO| & 3 & \verb|MATH060C_CM| & 3\\\hline
    \verb|MATH060__CM| & 3 & \verb|MATH060__PO| & 3 & \verb|MATH060__PZ| & 3\\\hline
    \verb|MATH060__SC| & 3 & \verb|MATH061__PO| & 3 & \verb|MATH062__HM| & 3\\\hline
    \verb|MATH067__PO| & 3 & \verb|MATH073__HM| & 3 & \verb|MATH082__HM| & 3\\\hline
    \verb|MATH093__HM| & 1 & \verb|MATH101__PO| & 3 & \verb|MATH102__PO| & 3\\\hline
    \verb|MATH102__PZ| & 3 & \verb|MATH102__SC| & 3 & \verb|MATH103__PO| & 3\\\hline
    \verb|MATH104__HM| & 3 & \verb|MATH106__HM| & 3 & \verb|MATH108B_PZ| & 3\\\hline
    \verb|MATH111__CM| & 3 & \verb|MATH112__PO| & 3 & \verb|MATH113__PO| & 3\\\hline
    \verb|MATH119__HM| & 3 & \verb|MATH131__CM| & 3 & \verb|MATH131__HM| & 3\\\hline
    \verb|MATH131__PO| & 3 & \verb|MATH131__SC| & 3 & \verb|MATH132__HM| & 3\\\hline
    \verb|MATH132__PO| & 3 & \verb|MATH135__CM| & 3 & \verb|MATH135__PO| & 3\\\hline
    \verb|MATH135__SC| & 3 & \verb|MATH136__HM| & 3 & \verb|MATH137__CM| & 3\\\hline
    \verb|MATH137__PO| & 3 & \verb|MATH138__PO| & 3 & \verb|MATH139__SC| & 3\\\hline
    \verb|MATH140__CM| & 3 & \verb|MATH142__PO| & 3 & \verb|MATH144__CM| & 3\\\hline
    \verb|MATH144__PZ| & 3 & \verb|MATH145A_PZ| & 3 & \verb|MATH145__PO| & 3\\\hline
    \verb|MATH147__HM| & 3 & \verb|MATH147__PO| & 3 & \verb|MATH148__CM| & 3\\\hline
    \verb|MATH150__PO| & 3 & \verb|MATH151__CM| & 3 & \verb|MATH151__PO| & 3\\\hline
    \verb|MATH152__CM| & 3 & \verb|MATH152__PO| & 3 & \verb|MATH153__PO| & 3\\\hline
    \verb|MATH154__PO| & 3 & \verb|MATH155__PO| & 3 & \verb|MATH156__CM| & 3\\\hline
    \verb|MATH156__HM| & 3 & \verb|MATH157__HM| & 1.5 & \verb|MATH158__HM| & 3\\\hline
    \verb|MATH158__PO| & 3 & \verb|MATH160__CM| & 3 & \verb|MATH162__PZ| & 3\\\hline
    \verb|MATH163__CM| & 3 & \verb|MATH164__HM| & 3 & \verb|MATH165__CM| & 3\\\hline
    \verb|MATH165__HM| & 3 & \verb|MATH171__CM| & 3 & \verb|MATH171__HM| & 3\\\hline
    \verb|MATH171__PO| & 3 & \verb|MATH172__CM| & 3 & \verb|MATH173__PO| & 3\\\hline
    \verb|MATH174__HM| & 3 & \verb|MATH175__CM| & 3 & \verb|MATH175__HM| & 3\\\hline
    \verb|MATH175__SC| & 3 & \verb|MATH176__HM| & 3 & \verb|MATH176__PO| & 3\\\hline
    \verb|MATH177__HM| & 3 & \verb|MATH177__PO| & 3 & \verb|MATH179__HM| & 3\\\hline
    \verb|MATH180__CM| & 3 & \verb|MATH180__HM| & 3 & \verb|MATH180__PO| & 3\\\hline
    \verb|MATH181__HM| & 3 & \verb|MATH181__PO| & 3 & \verb|MATH182__PZ| & 3\\\hline
    \verb|MATH183__PO| & 3 & \verb|MATH183__SC| & 3 & \verb|MATH185__PO| & 3\\\hline
    \verb|MATH185__SC| & 3 & \verb|MATH186__CM| & 3 & \verb|MATH187__HM| & 3\\\hline
    \verb|MATH187__PO| & 3 & \verb|MATH189AAHM| & 3 & \verb|MATH189AHHM| & 3\\\hline
    \verb|MATH189AIHM| & 3 & \verb|MATH189AJHM| & 3 & \verb|MATH189B_PO| & 3\\\hline
    \verb|MATH190__CM| & 1.5 & \verb|MATH190__CM2| & 1.5 & \verb|MATH190__PO| & 1.5\\\hline
    \verb|MATH190__PO2| & 1.5 & \verb|MATH191__CM| & 1.5 & \verb|MATH191__CM2| & 1.5\\\hline
    \verb|MATH191__PO| & 1.5 & \verb|MATH191__PO2| & 1.5 & \verb|MATH191__SC| & 3\\\hline
    \verb|MATH191__SC2| & 3 & \verb|MATH193__HM| & 3 & \verb|MATH193__HM2| & 3\\\hline
    \verb|MATH194__PZ| & 1.5 & \verb|MATH194__PZ2| & 1.5 & \verb|MATH195__CM| & 3\\\hline
    \verb|MATH195__CM2| & 3 & \verb|MATH196__HM| & 1 & \verb|MATH196__HM2| & 1\\\hline
    \verb|MATH197__HM| & 3 & \verb|MATH197__HM2| & 3 & \verb|MATH198__HM| & 1\\\hline
    \verb|MATH198__HM2| & 1 & \verb|MATH198__PZ| & 0.1 & \verb|MATH198__PZ2| & 0.1\\\hline
    \verb|MATH294__CG| & 3 & \verb|MATH389L_CG| & 3 & \verb|MATH398__CG| & 3\\\hline
    \verb|MATH454__CG| & 3 & \verb|MATH458__CG| & 3 & \verb|MATH900__HM| & 3\\\hline
    \verb|MATH901__HM| & 3 & \verb|MATH902__HM| & 3 & \verb|MATH903__HM| & 3\\\hline
    \verb|MATH904__HM| & 3 & \verb|MCBI117__HM| & 3 & \verb|MCBI118A_HM| & 1.5\\\hline
    \verb|MCBI118B_HM| & 1.5 & \verb|MCBI199__HM| & 0.1 & \verb|MCBI199__HM2| & 0.1\\\hline
    \verb|MCBI199__HM3| & 0.1 & \verb|MCBI199__HM4| & 0.1 & \verb|MCSI195__PZ| & 3\\\hline
    \verb|MCSI195__PZ2| & 3 & \verb|MENA191__SC| & 3 & \verb|MENA191__SC2| & 3\\\hline
    \verb|MGT_310__CG| & 3 & \verb|MGT_317__CG| & 1.5 & \verb|MGT_370__CG| & 3\\\hline
    \verb|MGT_390__CG| & 1.5 & \verb|MGT_395__CG| & 3 & \verb|MGT_398__CG| & 1.5\\\hline
    \verb|MGT_432__CG| & 3 & \verb|MGT_446__CG| & 1.5 & \verb|MGT_472__CG| & 1.5\\\hline
    \verb|MGT_477__CG| & 1.5 & \verb|MGT_503__CG| & 1.5 & \verb|MGT_555__CG| & 3\\\hline
    \verb|MLLC021__PZ| & 1.5 & \verb|MLLC033__PZ| & 1.5 & \verb|MLLC111__PZ| & 3\\\hline
    \verb|MLLC122__PZ| & 3 & \verb|MLLC144__PZ| & 3 & \verb|MLLC150__PZ| & 3\\\hline
    \verb|MLLC155__PZ| & 3 & \verb|MLLC188__PZ| & 3 & \verb|MOBI188__PO| & 3\\\hline
    \verb|MOBI191A_PO| & 1.5 & \verb|MOBI191A_PO2| & 1.5 & \verb|MOBI191B_PO| & 1.5\\\hline
    \verb|MOBI191B_PO2| & 1.5 & \verb|MOBI194A_PO| & 3 & \verb|MOBI194A_PO2| & 3\\\hline
    \verb|MOBI194B_PO| & 3 & \verb|MOBI194B_PO2| & 3 & \verb|MSL_001A_CM| & 1.5\\\hline
    \verb|MSL_099__CM| & 0.1 & \verb|MSL_101A_CM| & 0.1 & \verb|MSL_101B_CM| & 0.1\\\hline
    \verb|MSL_102A_CM| & 1.5 & \verb|MSL_102B_CM| & 1.5 & \verb|MSL_103A_CM| & 1.5\\\hline
    \verb|MSL_103B_CM| & 1.5 & \verb|MSL_104A_CM| & 1.5 & \verb|MSL_104B_CM| & 1.5\\\hline
    \verb|MS__038__SC| & 3 & \verb|MS__041__SC| & 3 & \verb|MS__044__PO| & 3\\\hline
    \verb|MS__045__PZ| & 3 & \verb|MS__049__PO| & 3 & \verb|MS__049__PZ| & 3\\\hline
    \verb|MS__049__SC| & 3 & \verb|MS__050__PO| & 3 & \verb|MS__050__PZ| & 3\\\hline
    \verb|MS__051__PO| & 3 & \verb|MS__051__PZ| & 3 & \verb|MS__051__SC| & 3\\\hline
    \verb|MS__052__PZ| & 3 & \verb|MS__053__SC| & 3 & \verb|MS__055__PZ| & 3\\\hline
    \verb|MS__057__SC| & 3 & \verb|MS__059__SC| & 3 & \verb|MS__060__PZ| & 3\\\hline
    \verb|MS__070__PZ| & 3 & \verb|MS__071__PO| & 3 & \verb|MS__072__PO| & 3\\\hline
    \verb|MS__073__PO| & 3 & \verb|MS__078__PZ| & 3 & \verb|MS__082__PZ| & 3\\\hline
    \verb|MS__082__SC| & 3 & \verb|MS__085__PO| & 3 & \verb|MS__087__PZ| & 3\\\hline
    \verb|MS__088__PZ| & 3 & \verb|MS__089B_PO| & 3 & \verb|MS__090__PZ| & 3\\\hline
    \verb|MS__092__PO| & 3 & \verb|MS__093__PZ| & 3 & \verb|MS__102__PO| & 3\\\hline
    \verb|MS__102__PZ| & 3 & \verb|MS__120__HM| & 3 & \verb|MS__120__PO| & 3\\\hline
    \verb|MS__120__PZ| & 3 & \verb|MS__120__SC| & 3 & \verb|MS__123__JT| & 3\\\hline
    \verb|MS__124__PZ| & 3 & \verb|MS__125__PZ| & 3 & \verb|MS__125__SC| & 3\\\hline
    \verb|MS__126__PZ| & 3 & \verb|MS__128__PZ| & 3 & \verb|MS__130__SC| & 3\\\hline
    \verb|MS__131B_SC| & 3 & \verb|MS__131__PO| & 3 & \verb|MS__131__SC| & 3\\\hline
    \verb|MS__135__PO| & 3 & \verb|MS__139__SC| & 3 & \verb|MS__140__PO| & 3\\\hline
    \verb|MS__148A_PO| & 3 & \verb|MS__148C_PO| & 3 & \verb|MS__148D_PO| & 3\\\hline
    \verb|MS__149G_PO| & 3 & \verb|MS__149T_PO| & 3 & \verb|MS__150__PO| & 3\\\hline
    \verb|MS__153__PO| & 3 & \verb|MS__160__SC| & 3 & \verb|MS__170__HM| & 3\\\hline
    \verb|MS__172__HM| & 3 & \verb|MS__173__HM| & 3 & \verb|MS__175__PO| & 3\\\hline
    \verb|MS__179I_HM| & 3 & \verb|MS__183__PO| & 3 & \verb|MS__187__SC| & 3\\\hline
    \verb|MS__190F_JT| & 3 & \verb|MS__190F_JT2| & 3 & \verb|MS__190S_JT| & 3\\\hline
    \verb|MS__190S_JT2| & 3 & \verb|MS__190__JT| & 3 & \verb|MS__190__JT2| & 3\\\hline
    \verb|MS__194__PZ| & 3 & \verb|MS__194__PZ2| & 3 & \verb|MS__196__PZ| & 3\\\hline
    \verb|MS__196__PZ2| & 3 & \verb|MS__199__SC| & 3 & \verb|MS__199__SC2| & 3\\\hline
    \verb|MUS_003__HM| & 3 & \verb|MUS_003__SC| & 3 & \verb|MUS_004__PO| & 3\\\hline
    \verb|MUS_007__PO| & 0.75 & \verb|MUS_008__PO| & 0.75 & \verb|MUS_010__PO| & 0.75\\\hline
    \verb|MUS_015VOPO| & 0.75 & \verb|MUS_015__PO| & 0.75 & \verb|MUS_020__PO| & 0.75\\\hline
    \verb|MUS_031__PO| & 1.5 & \verb|MUS_032__PO| & 1.5 & \verb|MUS_033__PO| & 1.5\\\hline
    \verb|MUS_035__PO| & 1.5 & \verb|MUS_037__PO| & 1.5 & \verb|MUS_040__PO| & 0.75\\\hline
    \verb|MUS_041__PO| & 1.5 & \verb|MUS_042B_PO| & 1.5 & \verb|MUS_042C_PO| & 1.5\\\hline
    \verb|MUS_046__HM| & 1 & \verb|MUS_047__PO| & 3 & \verb|MUS_049__HM| & 1\\\hline
    \verb|MUS_051__PO| & 3 & \verb|MUS_057__PO| & 3 & \verb|MUS_060__PO| & 3\\\hline
    \verb|MUS_062__PO| & 3 & \verb|MUS_065__PO| & 3 & \verb|MUS_066__SC| & 3\\\hline
    \verb|MUS_067__HM| & 3 & \verb|MUS_068__PO| & 3 & \verb|MUS_070__PO| & 3\\\hline
    \verb|MUS_071__SC| & 1.5 & \verb|MUS_072__SC| & 1.5 & \verb|MUS_074__PO| & 3\\\hline
    \verb|MUS_077__PO| & 3 & \verb|MUS_080L_PO| & 0.1 & \verb|MUS_080__PO| & 3\\\hline
    \verb|MUS_081L_PO| & 0.1 & \verb|MUS_081__HM| & 3 & \verb|MUS_081__JM| & 3\\\hline
    \verb|MUS_081__PO| & 3 & \verb|MUS_082L_PO| & 0.1 & \verb|MUS_082__PO| & 3\\\hline
    \verb|MUS_085__SC| & 1.5 & \verb|MUS_087__PO| & 3 & \verb|MUS_088__PO| & 3\\\hline
    \verb|MUS_089__SC| & 1.5 & \verb|MUS_091__PO| & 3 & \verb|MUS_092__PO| & 3\\\hline
    \verb|MUS_094__PO| & 3 & \verb|MUS_096A_PO| & 3 & \verb|MUS_096B_PO| & 3\\\hline
    \verb|MUS_099__HM| & 3 & \verb|MUS_100__PO| & 1.5 & \verb|MUS_101__SC| & 3\\\hline
    \verb|MUS_102__SC| & 3 & \verb|MUS_103__SC| & 3 & \verb|MUS_104__SC| & 3\\\hline
    \verb|MUS_110A_SC| & 3 & \verb|MUS_110B_SC| & 3 & \verb|MUS_112__SC| & 3\\\hline
    \verb|MUS_118__PO| & 3 & \verb|MUS_119__SC| & 3 & \verb|MUS_121__PO| & 3\\\hline
    \verb|MUS_121__SC| & 3 & \verb|MUS_122__PO| & 3 & \verb|MUS_126__SC| & 3\\\hline
    \verb|MUS_130__SC| & 3 & \verb|MUS_131__SC| & 3 & \verb|MUS_140__PO| & 1.5\\\hline
    \verb|MUS_170F_SC| & 3 & \verb|MUS_170H_SC| & 1.5 & \verb|MUS_171F_SC| & 3\\\hline
    \verb|MUS_171H_SC| & 1.5 & \verb|MUS_172__SC| & 1.5 & \verb|MUS_173__JM| & 1.5\\\hline
    \verb|MUS_175__JM| & 1.5 & \verb|MUS_176__JM| & 1.5 & \verb|MUS_177F_SC| & 3\\\hline
    \verb|MUS_177H_SC| & 1.5 & \verb|MUS_177__PO| & 3 & \verb|MUS_179D_HM| & 3\\\hline
    \verb|MUS_179E_HM| & 3 & \verb|MUS_179F_HM| & 3 & \verb|MUS_179G_HM| & 3\\\hline
    \verb|MUS_179H_HM| & 3 & \verb|MUS_179I_HM| & 3 & \verb|MUS_180F_SC| & 3\\\hline
    \verb|MUS_180H_SC| & 1.5 & \verb|MUS_184__PO| & 3 & \verb|MUS_190__PO| & 3\\\hline
    \verb|MUS_190__PO2| & 3 & \verb|MUS_191__SC| & 3 & \verb|MUS_191__SC2| & 3\\\hline
    \verb|MUS_192F_PO| & 3 & \verb|MUS_192F_PO2| & 3 & \verb|MUS_192H_PO| & 1.5\\\hline
    \verb|MUS_192H_PO2| & 1.5 & \verb|NEUR080LXKS| & 0.1 & \verb|NEUR080L_KS| & 3\\\hline
    \verb|NEUR086L_KS| & 3 & \verb|NEUR089A_PO| & 3 & \verb|NEUR095L_JT| & 3\\\hline
    \verb|NEUR101ALPO| & 0.1 & \verb|NEUR101A_PO| & 3 & \verb|NEUR101BLPO| & 0.1\\\hline
    \verb|NEUR101B_PO| & 3 & \verb|NEUR102__PO| & 3 & \verb|NEUR103__KS| & 3\\\hline
    \verb|NEUR103__PO| & 3 & \verb|NEUR118__PO| & 3 & \verb|NEUR122L_PO| & 0.1\\\hline
    \verb|NEUR122__PO| & 3 & \verb|NEUR123__PO| & 3 & \verb|NEUR125L_JT| & 3\\\hline
    \verb|NEUR129L_KS| & 3 & \verb|NEUR131__PO| & 3 & \verb|NEUR133L_KS| & 3\\\hline
    \verb|NEUR140__KS| & 1.5 & \verb|NEUR148L_KS| & 3 & \verb|NEUR149__KS| & 3\\\hline
    \verb|NEUR157__PO| & 3 & \verb|NEUR161L_KS| & 3 & \verb|NEUR161__KS| & 3\\\hline
    \verb|NEUR168L_PO| & 0.1 & \verb|NEUR168__PO| & 3 & \verb|NEUR178L_PO| & 0.1\\\hline
    \verb|NEUR178__PO| & 3 & \verb|NEUR181__KS| & 3 & \verb|NEUR188L_KS| & 3\\\hline
    \verb|NEUR189A_PO| & 3 & \verb|NEUR189B_PO| & 3 & \verb|NEUR189L_KS| & 0.1\\\hline
    \verb|NEUR189S_PO| & 3 & \verb|NEUR190L_KS| & 3 & \verb|NEUR190L_KS2| & 3\\\hline
    \verb|NEUR190__PO| & 1.5 & \verb|NEUR190__PO2| & 1.5 & \verb|NEUR191__KS| & 3\\\hline
    \verb|NEUR191__KS2| & 3 & \verb|NEUR191__PO| & 1.5 & \verb|NEUR191__PO2| & 1.5\\\hline
    \verb|NEUR192__PO| & 3 & \verb|NEUR192__PO2| & 3 & \verb|NEUR194A_PO| & 1.5\\\hline
    \verb|NEUR194A_PO2| & 1.5 & \verb|NEUR194B_PO| & 3 & \verb|NEUR194B_PO2| & 3\\\hline
    \verb|NEUR196__KS| & 0.75 & \verb|NEUR196__KS2| & 0.75 & \verb|NEUR197__KS| & 1.5\\\hline
    \verb|NEUR197__KS2| & 1.5 & \verb|ORST050__PZ| & 3 & \verb|ORST060__PZ| & 3\\\hline
    \verb|ORST089__PZ| & 3 & \verb|ORST100A_PZ| & 3 & \verb|ORST100__PZ| & 3\\\hline
    \verb|ORST102__PZ| & 3 & \verb|ORST105__PZ| & 3 & \verb|ORST114__PZ| & 3\\\hline
    \verb|ORST135__PZ| & 3 & \verb|ORST148__PZ| & 3 & \verb|ORST170__PZ| & 3\\\hline
    \verb|ORST173A_PZ| & 3 & \verb|ORST175__PZ| & 3 & \verb|ORST177__PZ| & 3\\\hline
    \verb|ORST180__PZ| & 3 & \verb|ORST191__SC| & 3 & \verb|ORST191__SC2| & 3\\\hline
    \verb|ORST192__PZ| & 3 & \verb|ORST192__PZ2| & 3 & \verb|ORST198C_PZ| & 3\\\hline
    \verb|ORST198C_PZ2| & 3 & \verb|ORST198J_PZ| & 3 & \verb|ORST198J_PZ2| & 3\\\hline
    \verb|ORST198M_PZ| & 3 & \verb|ORST198M_PZ2| & 3 & \verb|PHIL001__PO| & 3\\\hline
    \verb|PHIL007__PO| & 3 & \verb|PHIL007__PZ| & 3 & \verb|PHIL030__CM| & 3\\\hline
    \verb|PHIL030__PO| & 3 & \verb|PHIL030__PZ| & 3 & \verb|PHIL031__CM| & 3\\\hline
    \verb|PHIL031__PO| & 3 & \verb|PHIL031__PZ| & 3 & \verb|PHIL033__CM| & 3\\\hline
    \verb|PHIL033__PO| & 3 & \verb|PHIL033__PZ| & 3 & \verb|PHIL034__CM| & 3\\\hline
    \verb|PHIL034__PO| & 3 & \verb|PHIL035H_CM| & 3 & \verb|PHIL035__CM| & 3\\\hline
    \verb|PHIL035__PO| & 3 & \verb|PHIL036__PO| & 3 & \verb|PHIL037__CM| & 3\\\hline
    \verb|PHIL037__PO| & 3 & \verb|PHIL038__CM| & 3 & \verb|PHIL038__PO| & 3\\\hline
    \verb|PHIL039__CM| & 3 & \verb|PHIL039__PO| & 3 & \verb|PHIL040__PO| & 3\\\hline
    \verb|PHIL042__PO| & 3 & \verb|PHIL043__PO| & 3 & \verb|PHIL044__PO| & 3\\\hline
    \verb|PHIL046__PO| & 3 & \verb|PHIL048__PO| & 3 & \verb|PHIL052__PZ| & 3\\\hline
    \verb|PHIL054__PO| & 3 & \verb|PHIL055__PZ| & 3 & \verb|PHIL057__JT| & 3\\\hline
    \verb|PHIL058__PZ| & 3 & \verb|PHIL060__PO| & 3 & \verb|PHIL070__PO| & 3\\\hline
    \verb|PHIL075__PZ| & 3 & \verb|PHIL080__PO| & 3 & \verb|PHIL084__PZ| & 3\\\hline
    \verb|PHIL090__SC| & 3 & \verb|PHIL095__CM| & 3 & \verb|PHIL100E_CM| & 3\\\hline
    \verb|PHIL100F_CM| & 3 & \verb|PHIL101C_CM| & 3 & \verb|PHIL101D_CM| & 3\\\hline
    \verb|PHIL104__PO| & 3 & \verb|PHIL104__PZ| & 3 & \verb|PHIL106__CM| & 3\\\hline
    \verb|PHIL106__PO| & 3 & \verb|PHIL108__PO| & 3 & \verb|PHIL118__SC| & 3\\\hline
    \verb|PHIL120__PO| & 3 & \verb|PHIL121__CM| & 3 & \verb|PHIL121__HM| & 3\\\hline
    \verb|PHIL122__HM| & 3 & \verb|PHIL123__CM| & 3 & \verb|PHIL124__HM| & 3\\\hline
    \verb|PHIL126__CM| & 3 & \verb|PHIL126__SC| & 3 & \verb|PHIL133__SC| & 3\\\hline
    \verb|PHIL134__CM| & 3 & \verb|PHIL135__CM| & 3 & \verb|PHIL136__SC| & 3\\\hline
    \verb|PHIL138__CM| & 3 & \verb|PHIL139__CM| & 3 & \verb|PHIL140__SC| & 3\\\hline
    \verb|PHIL142__CM| & 3 & \verb|PHIL144__SC| & 3 & \verb|PHIL151__SC| & 3\\\hline
    \verb|PHIL155__SC| & 3 & \verb|PHIL156__PZ| & 3 & \verb|PHIL158__CM| & 3\\\hline
    \verb|PHIL160__CM| & 3 & \verb|PHIL160__SC| & 3 & \verb|PHIL162__SC| & 3\\\hline
    \verb|PHIL163__PO| & 3 & \verb|PHIL163__SC| & 3 & \verb|PHIL164__CM| & 3\\\hline
    \verb|PHIL170__SC| & 3 & \verb|PHIL176__CM| & 3 & \verb|PHIL177__CM| & 3\\\hline
    \verb|PHIL179F_HM| & 3 & \verb|PHIL179G_HM| & 3 & \verb|PHIL181__CM| & 3\\\hline
    \verb|PHIL182__CM| & 3 & \verb|PHIL185C_PO| & 3 & \verb|PHIL185C_PZ| & 3\\\hline
    \verb|PHIL185L_PO| & 3 & \verb|PHIL185N_PZ| & 3 & \verb|PHIL185P_PO| & 3\\\hline
    \verb|PHIL186K_PO| & 3 & \verb|PHIL190__CM| & 3 & \verb|PHIL190__CM2| & 3\\\hline
    \verb|PHIL190__PO| & 3 & \verb|PHIL190__PO2| & 3 & \verb|PHIL190__SC| & 3\\\hline
    \verb|PHIL190__SC2| & 3 & \verb|PHIL191__PO| & 3 & \verb|PHIL191__PO2| & 3\\\hline
    \verb|PHIL191__PZ| & 1.5 & \verb|PHIL191__PZ2| & 1.5 & \verb|PHIL191__SC| & 3\\\hline
    \verb|PHIL191__SC2| & 3 & \verb|PHIL192__CM| & 3 & \verb|PHIL192__CM2| & 3\\\hline
    \verb|PHIL198__CM| & 3 & \verb|PHIL198__CM2| & 3 & \verb|PHIL198__SC| & 1.5\\\hline
    \verb|PHIL198__SC2| & 1.5 & \verb|PHYS003__PO| & 3 & \verb|PHYS016__PO| & 3\\\hline
    \verb|PHYS017__PO| & 3 & \verb|PHYS023__HM| & 1.5 & \verb|PHYS024__HM| & 4\\\hline
    \verb|PHYS030LXKS| & 0.1 & \verb|PHYS030L_KS| & 3 & \verb|PHYS031LXKS| & 0.1\\\hline
    \verb|PHYS031L_KS| & 3 & \verb|PHYS033L_KS| & 3 & \verb|PHYS034L_KS| & 3\\\hline
    \verb|PHYS035__KS| & 3 & \verb|PHYS041L_PO| & 0.1 & \verb|PHYS041__PO| & 3\\\hline
    \verb|PHYS042L_PO| & 0.1 & \verb|PHYS042__PO| & 3 & \verb|PHYS049B_HM| & 1\\\hline
    \verb|PHYS050__HM| & 1 & \verb|PHYS051__HM| & 3 & \verb|PHYS052__HM| & 3\\\hline
    \verb|PHYS054__HM| & 1 & \verb|PHYS064__HM| & 3 & \verb|PHYS070L_PO| & 0.1\\\hline
    \verb|PHYS070__PO| & 3 & \verb|PHYS071__PO| & 1.5 & \verb|PHYS072__PO| & 1.5\\\hline
    \verb|PHYS077L_KS| & 3 & \verb|PHYS084__HM| & 3 & \verb|PHYS100__KS| & 3\\\hline
    \verb|PHYS101L_PO| & 0.1 & \verb|PHYS101__KS| & 3 & \verb|PHYS101__PO| & 3\\\hline
    \verb|PHYS102__KS| & 3 & \verb|PHYS106L_KS| & 3 & \verb|PHYS108__KS| & 3\\\hline
    \verb|PHYS109L_KS| & 3 & \verb|PHYS110__PO| & 1.5 & \verb|PHYS111__HM| & 3\\\hline
    \verb|PHYS114__KS| & 3 & \verb|PHYS115__KS| & 3 & \verb|PHYS116__HM| & 3\\\hline
    \verb|PHYS117__HM| & 3 & \verb|PHYS125__PO| & 3 & \verb|PHYS128__PO| & 3\\\hline
    \verb|PHYS133__HM| & 1 & \verb|PHYS134__HM| & 2 & \verb|PHYS139__PO| & 3\\\hline
    \verb|PHYS142__PO| & 3 & \verb|PHYS150__PO| & 3 & \verb|PHYS151__HM| & 3\\\hline
    \verb|PHYS154__HM| & 3 & \verb|PHYS156__HM| & 3 & \verb|PHYS160__PO| & 3\\\hline
    \verb|PHYS161__HM| & 2 & \verb|PHYS162__HM| & 2 & \verb|PHYS164__HM| & 2\\\hline
    \verb|PHYS166__HM| & 2 & \verb|PHYS170X_HM| & 3 & \verb|PHYS170__HM| & 2\\\hline
    \verb|PHYS170__PO| & 3 & \verb|PHYS172__HM| & 2 & \verb|PHYS174L_PO| & 0.1\\\hline
    \verb|PHYS174__PO| & 3 & \verb|PHYS175__PO| & 3 & \verb|PHYS176__PO| & 3\\\hline
    \verb|PHYS178__KS| & 3 & \verb|PHYS181__HM| & 2 & \verb|PHYS183__HM| & 3\\\hline
    \verb|PHYS187B_KS| & 3 & \verb|PHYS188L_KS| & 3 & \verb|PHYS189L_KS| & 0.1\\\hline
    \verb|PHYS190L_KS| & 3 & \verb|PHYS190L_KS2| & 3 & \verb|PHYS190__PO| & 3\\\hline
    \verb|PHYS190__PO2| & 3 & \verb|PHYS191E_PO| & 3 & \verb|PHYS191E_PO2| & 3\\\hline
    \verb|PHYS191L_PO| & 1.5 & \verb|PHYS191L_PO2| & 1.5 & \verb|PHYS191__HM| & 3\\\hline
    \verb|PHYS191__HM2| & 3 & \verb|PHYS191__KS| & 3 & \verb|PHYS191__KS2| & 3\\\hline
    \verb|PHYS193__HM| & 3 & \verb|PHYS193__HM2| & 3 & \verb|PHYS193__PO| & 0.1\\\hline
    \verb|PHYS193__PO2| & 0.1 & \verb|PHYS194__HM| & 3 & \verb|PHYS194__HM2| & 3\\\hline
    \verb|PHYS195__HM| & 0.1 & \verb|PHYS195__HM2| & 0.1 & \verb|PHYS195__HM3| & 0.1\\\hline
    \verb|PHYS195__HM4| & 0.1 & \verb|PHYS196__KS| & 0.75 & \verb|PHYS196__KS2| & 0.75\\\hline
    \verb|PHYS197__HM| & 3 & \verb|PHYS197__HM2| & 3 & \verb|PHYS197__KS| & 1.5\\\hline
    \verb|PHYS197__KS2| & 1.5 & \verb|PHYS199__HM| & 3 & \verb|PHYS199__HM2| & 3\\\hline
    \verb|PHYS900__HM| & 3 & \verb|PHYS901__HM| & 3 & \verb|PHYS902__HM| & 3\\\hline
    \verb|PHYS903__HM| & 3 & \verb|PHYS904__HM| & 3 & \verb|POLI001A_PO| & 3\\\hline
    \verb|POLI001B_PO| & 3 & \verb|POLI002__PO| & 3 & \verb|POLI003__PO| & 3\\\hline
    \verb|POLI005__PO| & 3 & \verb|POLI007__PO| & 3 & \verb|POLI008__PO| & 3\\\hline
    \verb|POLI010__PO| & 3 & \verb|POLI030__PO| & 3 & \verb|POLI031__PO| & 3\\\hline
    \verb|POLI033A_PO| & 3 & \verb|POLI033B_PO| & 3 & \verb|POLI060__PO| & 3\\\hline
    \verb|POLI061__PO| & 3 & \verb|POLI070__PO| & 3 & \verb|POLI090__PO| & 3\\\hline
    \verb|POLI100__SC| & 3 & \verb|POLI103__SC| & 3 & \verb|POLI106__SC| & 3\\\hline
    \verb|POLI109__SC| & 3 & \verb|POLI110__SC| & 3 & \verb|POLI111__PO| & 3\\\hline
    \verb|POLI112__PO| & 3 & \verb|POLI115__PO| & 3 & \verb|POLI115__SC| & 3\\\hline
    \verb|POLI116__SC| & 3 & \verb|POLI118__SC| & 3 & \verb|POLI120__SC| & 3\\\hline
    \verb|POLI121__SC| & 3 & \verb|POLI122__SC| & 3 & \verb|POLI123__SC| & 3\\\hline
    \verb|POLI124__PO| & 3 & \verb|POLI126__SC| & 3 & \verb|POLI130__SC| & 3\\\hline
    \verb|POLI131__PO| & 3 & \verb|POLI133__JT| & 3 & \verb|POLI134__PO| & 3\\\hline
    \verb|POLI134__SC| & 3 & \verb|POLI135L_SC| & 1.5 & \verb|POLI135__PO| & 3\\\hline
    \verb|POLI135__SC| & 3 & \verb|POLI140__SC| & 3 & \verb|POLI142__SC| & 3\\\hline
    \verb|POLI144__SC| & 3 & \verb|POLI145__SC| & 3 & \verb|POLI146__SC| & 3\\\hline
    \verb|POLI148__SC| & 3 & \verb|POLI150__SC| & 3 & \verb|POLI151__SC| & 3\\\hline
    \verb|POLI152__PO| & 3 & \verb|POLI153__SC| & 3 & \verb|POLI155__SC| & 3\\\hline
    \verb|POLI158__PO| & 3 & \verb|POLI161__PO| & 3 & \verb|POLI162__PO| & 3\\\hline
    \verb|POLI165__PO| & 3 & \verb|POLI167__JT| & 3 & \verb|POLI168__PO| & 3\\\hline
    \verb|POLI171__PO| & 3 & \verb|POLI173__SC| & 3 & \verb|POLI177__PO| & 3\\\hline
    \verb|POLI179A_PO| & 3 & \verb|POLI180__SC| & 3 & \verb|POLI181__PO| & 3\\\hline
    \verb|POLI184__SC| & 3 & \verb|POLI189F_PO| & 3 & \verb|POLI189G_PO| & 3\\\hline
    \verb|POLI189H_PO| & 3 & \verb|POLI190C_PO| & 3 & \verb|POLI190C_PO2| & 3\\\hline
    \verb|POLI190D_PO| & 3 & \verb|POLI190D_PO2| & 3 & \verb|POLI190E_PO| & 3\\\hline
    \verb|POLI190E_PO2| & 3 & \verb|POLI191__PO| & 3 & \verb|POLI191__PO2| & 3\\\hline
    \verb|POLI191__SC| & 3 & \verb|POLI191__SC2| & 3 & \verb|POLI193__PO| & 1.5\\\hline
    \verb|POLI193__PO2| & 1.5 & \verb|POLI195__PO| & 0.1 & \verb|POLI195__PO2| & 0.1\\\hline
    \verb|PONT100__CM| & 3 & \verb|PORT001__PZ| & 3 & \verb|PORT002__PZ| & 3\\\hline
    \verb|PORT022__CM| & 3 & \verb|PORT033__CM| & 3 & \verb|PORT035__PZ| & 1.5\\\hline
    \verb|PORT044__CM| & 3 & \verb|PORT044__PZ| & 3 & \verb|POST016X_PZ| & 3\\\hline
    \verb|POST020A_PZ| & 3 & \verb|POST020__PZ| & 3 & \verb|POST030__PZ| & 3\\\hline
    \verb|POST040__PZ| & 3 & \verb|POST050__PZ| & 3 & \verb|POST060__PZ| & 3\\\hline
    \verb|POST066__PZ| & 3 & \verb|POST070__PZ| & 3 & \verb|POST091__PZ| & 3\\\hline
    \verb|POST101__PZ| & 3 & \verb|POST106__PZ| & 3 & \verb|POST107__CH| & 3\\\hline
    \verb|POST114__HM| & 3 & \verb|POST124__PZ| & 3 & \verb|POST126__PZ| & 3\\\hline
    \verb|POST128__PZ| & 3 & \verb|POST131__PZ| & 3 & \verb|POST138__PZ| & 3\\\hline
    \verb|POST140__HM| & 3 & \verb|POST141__PZ| & 3 & \verb|POST144__PZ| & 3\\\hline
    \verb|POST151__PZ| & 3 & \verb|POST152__PZ| & 3 & \verb|POST153__PZ| & 3\\\hline
    \verb|POST162__PZ| & 3 & \verb|POST172__PZ| & 3 & \verb|POST174__PZ| & 3\\\hline
    \verb|POST179C_HM| & 3 & \verb|POST179D_HM| & 3 & \verb|POST182__PZ| & 3\\\hline
    \verb|POST195__PZ| & 3 & \verb|POST195__PZ2| & 3 & \verb|POST199__PZ| & 3\\\hline
    \verb|POST199__PZ2| & 3 & \verb|PPA_190__PO| & 3 & \verb|PPA_190__PO2| & 3\\\hline
    \verb|PPA_191__PO| & 3 & \verb|PPA_191__PO2| & 3 & \verb|PPA_191__SC| & 3\\\hline
    \verb|PPA_191__SC2| & 3 & \verb|PPA_192__SC| & 3 & \verb|PPA_192__SC2| & 3\\\hline
    \verb|PPA_195__PO| & 3 & \verb|PPA_195__PO2| & 3 & \verb|PPE_001A_CM| & 3\\\hline
    \verb|PPE_001B_CM| & 3 & \verb|PPE_011A_CM| & 3 & \verb|PPE_011B_CM| & 3\\\hline
    \verb|PPE_100__PO| & 3 & \verb|PPE_110A_CM| & 3 & \verb|PPE_110B_CM| & 3\\\hline
    \verb|PPE_160__PO| & 3 & \verb|PPE_195__PO| & 3 & \verb|PPE_195__PO2| & 3\\\hline
    \verb|PSYC010__PZ| & 3 & \verb|PSYC030__CM| & 3 & \verb|PSYC037__CM| & 3\\\hline
    \verb|PSYC040__CM| & 3 & \verb|PSYC051__CM| & 3 & \verb|PSYC051__PO| & 3\\\hline
    \verb|PSYC052__SC| & 3 & \verb|PSYC065__CM| & 3 & \verb|PSYC070__CM| & 3\\\hline
    \verb|PSYC081__CM| & 3 & \verb|PSYC084__CH| & 3 & \verb|PSYC091P_PZ| & 0.1\\\hline
    \verb|PSYC091__PZ| & 3 & \verb|PSYC092P_PZ| & 1.5 & \verb|PSYC092__CM| & 3\\\hline
    \verb|PSYC092__PZ| & 3 & \verb|PSYC097__CM| & 3 & \verb|PSYC101__PZ| & 3\\\hline
    \verb|PSYC103__PZ| & 3 & \verb|PSYC103__SC| & 3 & \verb|PSYC104L_SC| & 1.5\\\hline
    \verb|PSYC104P_PZ| & 1.5 & \verb|PSYC104__PZ| & 3 & \verb|PSYC104__SC| & 3\\\hline
    \verb|PSYC105__PZ| & 3 & \verb|PSYC106__PO| & 3 & \verb|PSYC107__CM| & 3\\\hline
    \verb|PSYC108__PO| & 3 & \verb|PSYC109__CM| & 3 & \verb|PSYC110__CM| & 3\\\hline
    \verb|PSYC111P_PZ| & 1.5 & \verb|PSYC111__CM| & 3 & \verb|PSYC111__PZ| & 3\\\hline
    \verb|PSYC112__PZ| & 3 & \verb|PSYC112__SC| & 3 & \verb|PSYC113__PZ| & 3\\\hline
    \verb|PSYC113__SC| & 3 & \verb|PSYC114__SC| & 3 & \verb|PSYC116__PZ| & 3\\\hline
    \verb|PSYC117__PZ| & 3 & \verb|PSYC119__CM| & 3 & \verb|PSYC120__CM| & 3\\\hline
    \verb|PSYC121__SC| & 3 & \verb|PSYC124__SC| & 3 & \verb|PSYC125__AF| & 3\\\hline
    \verb|PSYC126__PZ| & 3 & \verb|PSYC127__PZ| & 3 & \verb|PSYC128__SC| & 3\\\hline
    \verb|PSYC130P_PZ| & 1.5 & \verb|PSYC130__PZ| & 3 & \verb|PSYC131A_PO| & 3\\\hline
    \verb|PSYC131L_SC| & 1.5 & \verb|PSYC131__CM| & 3 & \verb|PSYC131__PO| & 3\\\hline
    \verb|PSYC131__SC| & 3 & \verb|PSYC132__CM| & 3 & \verb|PSYC133__PO| & 3\\\hline
    \verb|PSYC133__SC| & 3 & \verb|PSYC134__CM| & 3 & \verb|PSYC134__HM| & 3\\\hline
    \verb|PSYC134__SC| & 3 & \verb|PSYC136__CM| & 3 & \verb|PSYC137__PO| & 3\\\hline
    \verb|PSYC139__HM| & 3 & \verb|PSYC140A_PZ| & 3 & \verb|PSYC142__CM| & 3\\\hline
    \verb|PSYC143__CM| & 3 & \verb|PSYC143__SC| & 3 & \verb|PSYC144__CM| & 3\\\hline
    \verb|PSYC147__CH| & 3 & \verb|PSYC147__SC| & 3 & \verb|PSYC148__PZ| & 3\\\hline
    \verb|PSYC150__AF| & 3 & \verb|PSYC151__SC| & 3 & \verb|PSYC153__AA| & 3\\\hline
    \verb|PSYC154__CM| & 3 & \verb|PSYC154__PO| & 3 & \verb|PSYC155__CM| & 3\\\hline
    \verb|PSYC157__PO| & 3 & \verb|PSYC158__PO| & 3 & \verb|PSYC159__CM| & 3\\\hline
    \verb|PSYC160__CM| & 3 & \verb|PSYC160__PO| & 3 & \verb|PSYC162__PO| & 3\\\hline
    \verb|PSYC162__SC| & 3 & \verb|PSYC163__PO| & 3 & \verb|PSYC163__SC| & 3\\\hline
    \verb|PSYC164__CM| & 3 & \verb|PSYC164__SC| & 3 & \verb|PSYC166__CM| & 3\\\hline
    \verb|PSYC167__CM| & 3 & \verb|PSYC168L_SC| & 1.5 & \verb|PSYC168__SC| & 3\\\hline
    \verb|PSYC169__SC| & 3 & \verb|PSYC170__PO| & 3 & \verb|PSYC175__SC| & 3\\\hline
    \verb|PSYC176__PO| & 3 & \verb|PSYC179M_HM| & 3 & \verb|PSYC179N_HM| & 3\\\hline
    \verb|PSYC179O_HM| & 3 & \verb|PSYC179U_HM| & 3 & \verb|PSYC180C_PO| & 3\\\hline
    \verb|PSYC180F_PO| & 3 & \verb|PSYC180H_PO| & 3 & \verb|PSYC180M_PO| & 3\\\hline
    \verb|PSYC180N_CH| & 3 & \verb|PSYC180R_PO| & 3 & \verb|PSYC180S_PO| & 3\\\hline
    \verb|PSYC180W_PO| & 3 & \verb|PSYC180__CM| & 3 & \verb|PSYC181__PZ| & 3\\\hline
    \verb|PSYC183__SC| & 3 & \verb|PSYC184__PZ| & 3 & \verb|PSYC188__CM| & 3\\\hline
    \verb|PSYC189Q_PO| & 3 & \verb|PSYC189__PZ| & 3 & \verb|PSYC190A_PO| & 1.5\\\hline
    \verb|PSYC190A_PO2| & 1.5 & \verb|PSYC190R_PO| & 1.5 & \verb|PSYC190R_PO2| & 1.5\\\hline
    \verb|PSYC190__PO| & 1.5 & \verb|PSYC190__PO2| & 1.5 & \verb|PSYC191__PO| & 3\\\hline
    \verb|PSYC191__PO2| & 3 & \verb|PSYC191__PZ| & 1.5 & \verb|PSYC191__PZ2| & 1.5\\\hline
    \verb|PSYC191__SC| & 3 & \verb|PSYC191__SC2| & 3 & \verb|PSYC193__SC| & 3\\\hline
    \verb|PSYC193__SC2| & 3 & \verb|PSYC194__PZ| & 3 & \verb|PSYC194__PZ2| & 3\\\hline
    \verb|PSYC197A_CM| & 0.75 & \verb|PSYC197A_CM2| & 0.75 & \verb|PSYC197B_CM| & 0.1\\\hline
    \verb|PSYC197B_CM2| & 0.1 & \verb|PSYC197__PZ| & 3 & \verb|PSYC197__PZ2| & 3\\\hline
    \verb|PSYC198__CM| & 3 & \verb|PSYC198__CM2| & 3 & \verb|PSYC308B_CG| & 1.5\\\hline
    \verb|PSYC308C_CG| & 1.5 & \verb|PSYC315F_CG| & 1.5 & \verb|PSYC315H_CG| & 1.5\\\hline
    \verb|PSYC332__CG| & 3 & \verb|PSYC450__CG| & 3 & \verb|RLIT191__PO| & 1.5\\\hline
    \verb|RLIT191__PO2| & 1.5 & \verb|RLST010__CM| & 3 & \verb|RLST020__PO| & 3\\\hline
    \verb|RLST037__CM| & 3 & \verb|RLST038__CM| & 3 & \verb|RLST040__PO| & 3\\\hline
    \verb|RLST045__CM| & 3 & \verb|RLST049__PO| & 3 & \verb|RLST055__CM| & 3\\\hline
    \verb|RLST056__CM| & 3 & \verb|RLST058__CM| & 3 & \verb|RLST060__SC| & 3\\\hline
    \verb|RLST061__SC| & 3 & \verb|RLST064__CM| & 3 & \verb|RLST082__CM| & 3\\\hline
    \verb|RLST083__CM| & 3 & \verb|RLST084__CM| & 3 & \verb|RLST086__SC| & 3\\\hline
    \verb|RLST092__SC| & 3 & \verb|RLST093__CM| & 3 & \verb|RLST098__SC| & 3\\\hline
    \verb|RLST100__PO| & 3 & \verb|RLST102__CM| & 3 & \verb|RLST103__PO| & 3\\\hline
    \verb|RLST107__PO| & 3 & \verb|RLST109__CM| & 3 & \verb|RLST110__PO| & 3\\\hline
    \verb|RLST113__HM| & 3 & \verb|RLST114__HM| & 3 & \verb|RLST125__CM| & 3\\\hline
    \verb|RLST126__CM| & 3 & \verb|RLST128__PO| & 3 & \verb|RLST129__CM| & 3\\\hline
    \verb|RLST130A_CM| & 3 & \verb|RLST130__CM| & 3 & \verb|RLST135__CM| & 3\\\hline
    \verb|RLST139__PO| & 3 & \verb|RLST142__AF| & 3 & \verb|RLST144__PO| & 3\\\hline
    \verb|RLST146__PO| & 3 & \verb|RLST148__PO| & 3 & \verb|RLST150__AF| & 3\\\hline
    \verb|RLST152__PO| & 3 & \verb|RLST157__PO| & 3 & \verb|RLST158__PO| & 3\\\hline
    \verb|RLST160__PO| & 3 & \verb|RLST164__PO| & 3 & \verb|RLST166B_CM| & 3\\\hline
    \verb|RLST168__HM| & 1.5 & \verb|RLST169__CM| & 3 & \verb|RLST171__CM| & 3\\\hline
    \verb|RLST179G_HM| & 3 & \verb|RLST179__SC| & 3 & \verb|RLST180__CM| & 3\\\hline
    \verb|RLST180__PO| & 3 & \verb|RLST181__PO| & 3 & \verb|RLST183__HM| & 3\\\hline
    \verb|RLST184__PO| & 3 & \verb|RLST185__SC| & 3 & \verb|RLST187__PO| & 3\\\hline
    \verb|RLST189C_PO| & 3 & \verb|RLST189L_PO| & 3 & \verb|RLST190__PO| & 3\\\hline
    \verb|RLST190__PO2| & 3 & \verb|RLST191__PO| & 3 & \verb|RLST191__PO2| & 3\\\hline
    \verb|RLST191__SC| & 3 & \verb|RLST191__SC2| & 3 & \verb|RUSS001__PO| & 3\\\hline
    \verb|RUSS002__PO| & 3 & \verb|RUSS011__PO| & 0.75 & \verb|RUSS013__PO| & 0.75\\\hline
    \verb|RUSS033__PO| & 3 & \verb|RUSS044__PO| & 3 & \verb|RUSS180__PO| & 3\\\hline
    \verb|RUSS181__PO| & 3 & \verb|RUSS182__PO| & 3 & \verb|RUSS187__PO| & 3\\\hline
    \verb|RUSS191__PO| & 1.5 & \verb|RUSS191__PO2| & 1.5 & \verb|RUSS191__SC| & 1.5\\\hline
    \verb|RUSS191__SC2| & 1.5 & \verb|RUST079__PO| & 3 & \verb|RUST100__PO| & 3\\\hline
    \verb|RUST110__PO| & 3 & \verb|RUST112__PO| & 3 & \verb|RUST185__PO| & 3\\\hline
    \verb|RUST191__PO| & 1.5 & \verb|RUST191__PO2| & 1.5 & \verb|SCI_010L_CM| & 3\\\hline
    \verb|SCI_030LACM| & 3 & \verb|SCI_030LBCM| & 3 & \verb|SCI_031LBCM| & 3\\\hline
    \verb|SCI_050L_CM| & 3 & \verb|SCI_119__CM| & 3 & \verb|SCI_131__CM| & 3\\\hline
    \verb|SCI_132__CM| & 3 & \verb|SCI_153__CM| & 3 & \verb|SCI_196__CM| & 0.75\\\hline
    \verb|SCI_196__CM2| & 0.75 & \verb|SCI_197__CM| & 1.5 & \verb|SCI_197__CM2| & 1.5\\\hline
    \verb|SOC_001X_PZ| & 3 & \verb|SOC_001__PZ| & 3 & \verb|SOC_012__PZ| & 3\\\hline
    \verb|SOC_030__CH| & 3 & \verb|SOC_051__PO| & 3 & \verb|SOC_061__PZ| & 3\\\hline
    \verb|SOC_069__PZ| & 3 & \verb|SOC_072__PZ| & 3 & \verb|SOC_083__PZ| & 3\\\hline
    \verb|SOC_090__PO| & 3 & \verb|SOC_091__PZ| & 3 & \verb|SOC_092__PZ| & 3\\\hline
    \verb|SOC_095__PZ| & 3 & \verb|SOC_099__PZ| & 3 & \verb|SOC_101__PZ| & 3\\\hline
    \verb|SOC_102__PO| & 3 & \verb|SOC_102__PZ| & 3 & \verb|SOC_103__JT| & 0.75\\\hline
    \verb|SOC_109__PZ| & 3 & \verb|SOC_110__PZ| & 3 & \verb|SOC_118__PZ| & 3\\\hline
    \verb|SOC_119__PO| & 3 & \verb|SOC_119__PZ| & 3 & \verb|SOC_122__PZ| & 3\\\hline
    \verb|SOC_123__PO| & 3 & \verb|SOC_126__AA| & 3 & \verb|SOC_128__JT| & 3\\\hline
    \verb|SOC_130__PO| & 3 & \verb|SOC_138__PO| & 3 & \verb|SOC_146__PO| & 3\\\hline
    \verb|SOC_147__PO| & 3 & \verb|SOC_148__PO| & 3 & \verb|SOC_150__AA| & 3\\\hline
    \verb|SOC_150__CH| & 3 & \verb|SOC_154__PO| & 3 & \verb|SOC_157__PZ| & 3\\\hline
    \verb|SOC_169X_PO| & 3 & \verb|SOC_170__PZ| & 3 & \verb|SOC_179C_HM| & 3\\\hline
    \verb|SOC_179D_HM| & 3 & \verb|SOC_179E_HM| & 3 & \verb|SOC_179F_HM| & 3\\\hline
    \verb|SOC_189E_PO| & 3 & \verb|SOC_189F_PO| & 3 & \verb|SOC_189Z_PO| & 3\\\hline
    \verb|SOC_190__PO| & 3 & \verb|SOC_190__PO2| & 3 & \verb|SOC_191__PO| & 3\\\hline
    \verb|SOC_191__PO2| & 3 & \verb|SOC_199A_PZ| & 3 & \verb|SOC_199A_PZ2| & 3\\\hline
    \verb|SOC_199B_PZ| & 1.5 & \verb|SOC_199B_PZ2| & 1.5 & \verb|SOC_199C_PZ| & 3\\\hline
    \verb|SOC_199C_PZ2| & 3 & \verb|SOSC150__HM| & 3 & \verb|SPAN001__CM| & 3\\\hline
    \verb|SPAN001__PO| & 3 & \verb|SPAN001__PZ| & 3 & \verb|SPAN002__CM| & 3\\\hline
    \verb|SPAN002__PO| & 3 & \verb|SPAN002__PZ| & 3 & \verb|SPAN011__PO| & 0.75\\\hline
    \verb|SPAN013__PO| & 0.75 & \verb|SPAN022__CM| & 3 & \verb|SPAN022__PO| & 3\\\hline
    \verb|SPAN022__PZ| & 3 & \verb|SPAN022__SC| & 3 & \verb|SPAN031__PZ| & 1.5\\\hline
    \verb|SPAN033L_CM| & 0.1 & \verb|SPAN033__CM| & 3 & \verb|SPAN033__PO| & 3\\\hline
    \verb|SPAN033__PZ| & 3 & \verb|SPAN033__SC| & 3 & \verb|SPAN035__PZ| & 1.5\\\hline
    \verb|SPAN044L_CM| & 0.1 & \verb|SPAN044__CM| & 3 & \verb|SPAN044__PO| & 3\\\hline
    \verb|SPAN044__PZ| & 3 & \verb|SPAN044__SC| & 3 & \verb|SPAN050__PZ| & 3\\\hline
    \verb|SPAN055__PZ| & 1.5 & \verb|SPAN100__PZ| & 3 & \verb|SPAN101__PO| & 3\\\hline
    \verb|SPAN101__SC| & 3 & \verb|SPAN102__CM| & 3 & \verb|SPAN103__SC| & 3\\\hline
    \verb|SPAN104__PO| & 3 & \verb|SPAN104__PZ| & 3 & \verb|SPAN106__PO| & 3\\\hline
    \verb|SPAN107__PZ| & 3 & \verb|SPAN108__PO| & 3 & \verb|SPAN109__PO| & 3\\\hline
    \verb|SPAN120A_PO| & 3 & \verb|SPAN125A_CM| & 3 & \verb|SPAN125A_PO| & 3\\\hline
    \verb|SPAN125B_CM| & 3 & \verb|SPAN125B_PO| & 3 & \verb|SPAN126__PO| & 3\\\hline
    \verb|SPAN127__CH| & 3 & \verb|SPAN130__CM| & 3 & \verb|SPAN131__SC| & 3\\\hline
    \verb|SPAN132__PZ| & 3 & \verb|SPAN134__SC| & 3 & \verb|SPAN136__PZ| & 3\\\hline
    \verb|SPAN137__PZ| & 3 & \verb|SPAN139__SC| & 3 & \verb|SPAN140__PO| & 3\\\hline
    \verb|SPAN140__SC| & 3 & \verb|SPAN141__SC| & 3 & \verb|SPAN143__SC| & 3\\\hline
    \verb|SPAN144__PZ| & 3 & \verb|SPAN145__SC| & 3 & \verb|SPAN146__PO| & 3\\\hline
    \verb|SPAN148__PZ| & 3 & \verb|SPAN150__PZ| & 3 & \verb|SPAN153__PO| & 3\\\hline
    \verb|SPAN156__PZ| & 3 & \verb|SPAN157__PZ| & 3 & \verb|SPAN159__PO| & 3\\\hline
    \verb|SPAN159__PZ| & 3 & \verb|SPAN163__PZ| & 3 & \verb|SPAN163__SC| & 3\\\hline
    \verb|SPAN167__CM| & 3 & \verb|SPAN171__PZ| & 3 & \verb|SPAN172__PZ| & 3\\\hline
    \verb|SPAN181__PZ| & 3 & \verb|SPAN181__SC| & 3 & \verb|SPAN183__CM| & 3\\\hline
    \verb|SPAN183__SC| & 3 & \verb|SPAN185__PO| & 3 & \verb|SPAN189A_PO| & 3\\\hline
    \verb|SPAN191__PO| & 1.5 & \verb|SPAN191__PO2| & 1.5 & \verb|SPAN191__SC| & 3\\\hline
    \verb|SPAN191__SC2| & 3 & \verb|SPAN192__PO| & 1.5 & \verb|SPAN192__PO2| & 1.5\\\hline
    \verb|SPAN199__PZ| & 3 & \verb|SPAN199__PZ2| & 3 & \verb|SPCH061A_CM| & 1.5\\\hline
    \verb|SPCH061B_CM| & 0.1 & \verb|STEM000__5C| & 3 & \verb|STEM001__5C| & 3\\\hline
    \verb|STEM002__5C| & 3 & \verb|STEM003__5C| & 3 & \verb|STEM004__5C| & 3\\\hline
    \verb|STEM005__5C| & 3 & \verb|STEM006__5C| & 3 & \verb|STEM007__5C| & 3\\\hline
    \verb|STEM008__5C| & 3 & \verb|STEM009__5C| & 3 & \verb|STEM010__5C| & 3\\\hline
    \verb|STEM011__5C| & 3 & \verb|STS_010__PO| & 3 & \verb|STS_010__PZ| & 3\\\hline
    \verb|STS_080__PO| & 3 & \verb|STS_179N_HM| & 3 & \verb|STS_190__PO| & 3\\\hline
    \verb|STS_190__PO2| & 3 & \verb|STS_191__PO| & 3 & \verb|STS_191__PO2| & 3\\\hline
    \verb|STS_191__SC| & 3 & \verb|STS_191__SC2| & 3 & \verb|TEST001__PZ| & 1.5\\\hline
    \verb|TEST002__PZ| & 3 & \verb|THEA001A_PO| & 3 & \verb|THEA001D_PO| & 3\\\hline
    \verb|THEA001G_PO| & 3 & \verb|THEA002__PO| & 3 & \verb|THEA006__PO| & 3\\\hline
    \verb|THEA010__PO| & 3 & \verb|THEA012__PO| & 3 & \verb|THEA017__PO| & 3\\\hline
    \verb|THEA021__PO| & 3 & \verb|THEA022__PO| & 3 & \verb|THEA023__PO| & 3\\\hline
    \verb|THEA026__PO| & 3 & \verb|THEA030__PO| & 3 & \verb|THEA031__PO| & 3\\\hline
    \verb|THEA041__PO| & 3 & \verb|THEA051C_PO| & 0.75 & \verb|THEA051H_PO| & 1.5\\\hline
    \verb|THEA052C_PO| & 0.75 & \verb|THEA052H_PO| & 1.5 & \verb|THEA053HGPO| & 1.5\\\hline
    \verb|THEA054C_PO| & 3 & \verb|THEA056C_PO| & 0.75 & \verb|THEA061__PO| & 3\\\hline
    \verb|THEA080__PO| & 3 & \verb|THEA081__PO| & 3 & \verb|THEA082__PO| & 3\\\hline
    \verb|THEA083__PO| & 3 & \verb|THEA084__PO| & 3 & \verb|THEA085__PO| & 3\\\hline
    \verb|THEA089C_PO| & 3 & \verb|THEA100C_PO| & 3 & \verb|THEA100G_PO| & 3\\\hline
    \verb|THEA100K_PO| & 3 & \verb|THEA100S_PO| & 3 & \verb|THEA115O_PO| & 3\\\hline
    \verb|THEA130__PO| & 3 & \verb|THEA132__PO| & 3 & \verb|THEA141__PO| & 3\\\hline
    \verb|THEA142__PO| & 3 & \verb|THEA170__PO| & 3 & \verb|THEA171__PO| & 3\\\hline
    \verb|THEA172__PO| & 3 & \verb|THEA187__PO| & 1.5 & \verb|THEA188__PO| & 3\\\hline
    \verb|THEA190__PO| & 1.5 & \verb|THEA190__PO2| & 1.5 & \verb|THEA191H_PO| & 1.5\\\hline
    \verb|THEA191H_PO2| & 1.5 & \verb|THEA192H_PO| & 1.5 & \verb|THEA192H_PO2| & 1.5\\\hline
    \verb|THES191D_SC| & 3 & \verb|THES191D_SC2| & 3 & \verb|THES192D_SC| & 3\\\hline
    \verb|THES192D_SC2| & 3 & \verb|TNDY409D_CG| & 3 & \verb|WGS_321__CG| & 3\\\hline
    \verb|WRIT001E_HM| & 3 & \verb|WRIT001__HM| & 2 & \verb|WRIT015__PZ| & 3\\\hline
    \verb|WRIT022__PZ| & 3 & \verb|WRIT100A_PZ| & 1.5 & \verb|WRIT100B_PZ| & 3\\\hline
    \verb|WRIT101__PZ| & 3 & \verb|WRIT105__SC| & 3 & \verb|WRIT109__SC| & 1.5\\\hline
    \verb|WRIT112__SC| & 1.5 & \verb|WRIT122__SC| & 1.5 & \verb|WRIT125__PZ| & 3\\\hline
    \verb|WRIT132__SC| & 3 & \verb|WRIT137__SC| & 3 & \verb|WRIT140__SC| & 3\\\hline
    \verb|WRIT146__SC| & 3 & \verb|WRIT153IOSC| & 3 & \verb|WRIT155__SC| & 3\\\hline
    \verb|WRIT160__SC| & 1.5 & \verb|WRIT175__SC| & 3 & \verb|WRIT190__PZ| & 3\\\hline
    \verb|WRIT190__PZ2| & 3 & \verb|WRIT191__SC| & 3 & \verb|WRIT191__SC2| & 3\\\hline
    \verb|WRIT197N_SC| & 3 & \verb|WRIT197N_SC2| & 3 & \verb|WRIT197Q_SC| & 3\\\hline
    \verb|WRIT197Q_SC2| & 3 & \verb|XGOV192__CM| & 3 & \verb|XGOV192__CM2| & 3\\\hline
    \caption{Credits by course.}
    \label{tab:credits}
\end{longtable}
\begin{longtable}{|l|l||l|l||l|l|}
    \hline
    Course ID & Offerings & Course ID & Offerings & Course ID & Offerings\\\hline
    \verb|AFRI010ASAF| & 4, 8 & \verb|AFRI010A_AF| & 2, 3, 6, 7, 1, 5 & \verb|AFRI010BSAF| & 4, 8\\\hline
    \verb|AFRI010B_AF| & 2, 3, 6, 7, 1, 5 & \verb|AFRI014__AF| & 1, 5 & \verb|AFRI020__PZ| & 2, 4, 6, 8\\\hline
    \verb|AFRI110__PZ| & 2, 6 & \verb|AFRI114__AF| & 3, 7 & \verb|AFRI116__AF| & 2, 3, 6, 7, 1, 5\\\hline
    \verb|AFRI120__PZ| & 2, 3, 6, 7, 1, 5 & \verb|AFRI121__AF| & 2, 3, 6, 7, 1, 5 & \verb|AFRI124__AF| & 2, 4, 6, 8\\\hline
    \verb|AFRI125__AF| & 3, 7 & \verb|AFRI128__AF| & 1, 5 & \verb|AFRI134__PO| & 2, 6\\\hline
    \verb|AFRI138__PO| & 4, 8 & \verb|AFRI149__AF| & 4, 8 & \verb|AFRI150__PZ| & 1, 3, 5, 7, 4, 8\\\hline
    \verb|AFRI159__AF| & 4, 8 & \verb|AFRI180__AF| & 4, 8 & \verb|AFRI190__AF| & 1, 5\\\hline
    \verb|AFRI190__AF2| & 1, 5 & \verb|AFRI191C_AF| & 2, 6 & \verb|AFRI191C_AF2| & 2, 6\\\hline
    \verb|AFRI191__AF| & 8, 3, 7, 4 & \verb|AFRI191__AF2| & 8, 3, 7, 4 & \verb|AFRI195D_AF| & 2, 6\\\hline
    \verb|AFRI195D_AF2| & 2, 6 & \verb|AMST103__JT| & 2, 3, 6, 7 & \verb|AMST103__PO| & 4, 8\\\hline
    \verb|AMST103__PZ| & 1, 5 & \verb|AMST113__SC| & 2, 6 & \verb|AMST118__PO| & 2, 6\\\hline
    \verb|AMST140__SC| & 2, 6 & \verb|AMST179A_HM| & 3, 7 & \verb|AMST180__PZ| & 4, 8\\\hline
    \verb|AMST180__SC| & 2, 6 & \verb|AMST190__PO| & 1, 3, 5, 7 & \verb|AMST190__PO2| & 1, 3, 5, 7\\\hline
    \verb|AMST191__PO| & 2, 4, 6, 8 & \verb|AMST191__PO2| & 2, 4, 6, 8 & \verb|AMST191__SC| & 4, 8\\\hline
    \verb|AMST191__SC2| & 4, 8 & \verb|ANTH001__PZ| & 2, 3, 6, 7, 4, 8 & \verb|ANTH002__PO| & 2, 3, 6, 7, 1, 5, 4, 8\\\hline
    \verb|ANTH002__PZ| & 2, 3, 6, 7, 1, 5, 4, 8 & \verb|ANTH002__SC| & 2, 3, 6, 7, 1, 5, 4, 8 & \verb|ANTH003__PZ| & 2, 4, 6, 8\\\hline
    \verb|ANTH005__PZ| & 1, 5 & \verb|ANTH009__PZ| & 1, 4, 5, 8 & \verb|ANTH012__PZ| & 2, 6\\\hline
    \verb|ANTH015__PO| & 2, 6 & \verb|ANTH022__SC| & 3, 7 & \verb|ANTH025__SC| & 2, 1, 6, 5\\\hline
    \verb|ANTH027__SC| & 2, 6 & \verb|ANTH053__PO| & 2, 4, 6, 8 & \verb|ANTH055__PZ| & 2, 3, 6, 7, 1, 5\\\hline
    \verb|ANTH070__PZ| & 1, 4, 5, 8 & \verb|ANTH075__PZ| & 1, 5 & \verb|ANTH077__PO| & 2, 4, 6, 8\\\hline
    \verb|ANTH080__PO| & 1, 5 & \verb|ANTH087__SC| & 1, 5 & \verb|ANTH099__PZ| & 4, 8\\\hline
    \verb|ANTH100__PO| & 4, 8 & \verb|ANTH101__PZ| & 4, 8 & \verb|ANTH104__PO| & 1, 3, 5, 7\\\hline
    \verb|ANTH105__PO| & 2, 4, 6, 8 & \verb|ANTH105__PZ| & 1, 3, 5, 7 & \verb|ANTH106__PZ| & 4, 8\\\hline
    \verb|ANTH107__PO| & 2, 4, 6, 8 & \verb|ANTH107__PZ| & 2, 3, 6, 7, 1, 5, 4, 8 & \verb|ANTH107__SC| & 3, 7\\\hline
    \verb|ANTH109__SC| & 3, 7 & \verb|ANTH111__PZ| & 3, 7 & \verb|ANTH112__PZ| & 2, 3, 6, 7\\\hline
    \verb|ANTH113__SC| & 2, 4, 6, 8 & \verb|ANTH114__PZ| & 4, 8 & \verb|ANTH114__SC| & 2, 6\\\hline
    \verb|ANTH115__SC| & 4, 8 & \verb|ANTH116__PO| & 3, 7 & \verb|ANTH123__PO| & 3, 7\\\hline
    \verb|ANTH124__SC| & 4, 8 & \verb|ANTH126__SC| & 1, 5 & \verb|ANTH137__PZ| & 4, 8\\\hline
    \verb|ANTH143__SC| & 1, 5 & \verb|ANTH153__SC| & 1, 3, 5, 7 & \verb|ANTH184__PO| & 3, 7\\\hline
    \verb|ANTH185P_SC| & 2, 6 & \verb|ANTH185__SC| & 4, 8 & \verb|ANTH189P_PO| & 2, 6\\\hline
    \verb|ANTH190__PO| & 1, 3, 5, 7 & \verb|ANTH190__PO2| & 1, 3, 5, 7 & \verb|ANTH190__SC| & 1, 3, 5, 7\\\hline
    \verb|ANTH190__SC2| & 1, 3, 5, 7 & \verb|ANTH191__SC| & 4, 8 & \verb|ANTH191__SC2| & 4, 8\\\hline
    \verb|ARBC001__CM| & 1, 3, 5, 7 & \verb|ARBC002__CM| & 2, 4, 6, 8 & \verb|ARBC033__CM| & 1, 3, 5, 7\\\hline
    \verb|ARBC044__CM| & 2, 4, 6, 8 & \verb|ARBC130__CM| & 2, 4, 6, 8 & \verb|ARBC141__CM| & 4, 8\\\hline
    \verb|ARBC166__CM| & 1, 3, 5, 7 & \verb|ARCN101__SC| & 4, 8 & \verb|ARCN115__SC| & 3, 7\\\hline
    \verb|ARCN120__SC| & 2, 6 & \verb|ARCN144__SC| & 1, 5 & \verb|ARCN191__SC| & 4, 8\\\hline
    \verb|ARCN191__SC2| & 4, 8 & \verb|ARCT179A_HM| & 1, 3, 5, 7 & \verb|ARCT179B_HM| & 4, 8\\\hline
    \verb|ARCT179C_HM| & 1, 5 & \verb|ARHI001A_PO| & 3, 7 & \verb|ARHI001B_PO| & 2, 4, 6, 8\\\hline
    \verb|ARHI114__PO| & 3, 7 & \verb|ARHI120__PO| & 2, 3, 6, 7, 1, 5 & \verb|ARHI123__PO| & 2, 6\\\hline
    \verb|ARHI124__PO| & 4, 8 & \verb|ARHI127__PO| & 4, 8 & \verb|ARHI130__PO| & 2, 1, 6, 5\\\hline
    \verb|ARHI131__PO| & 2, 6 & \verb|ARHI135__PZ| & 4, 8 & \verb|ARHI138__PO| & 3, 7\\\hline
    \verb|ARHI139__PO| & 4, 8 & \verb|ARHI140__PO| & 2, 3, 6, 7 & \verb|ARHI141A_PO| & 4, 8\\\hline
    \verb|ARHI144B_PO| & 1, 5 & \verb|ARHI150__SC| & 1, 3, 5, 7 & \verb|ARHI151__SC| & 4, 8\\\hline
    \verb|ARHI154__SC| & 2, 3, 6, 7, 1, 5 & \verb|ARHI158__SC| & 4, 8 & \verb|ARHI164__SC| & 4, 8\\\hline
    \verb|ARHI174__PO| & 2, 6 & \verb|ARHI178__PO| & 3, 7 & \verb|ARHI179E_HM| & 4, 8\\\hline
    \verb|ARHI179__PO| & 3, 7 & \verb|ARHI183__SC| & 1, 4, 5, 8 & \verb|ARHI185__SC| & 1, 3, 5, 7\\\hline
    \verb|ARHI186A_PZ| & 3, 7 & \verb|ARHI186B_PZ| & 1, 5 & \verb|ARHI186C_SC| & 2, 3, 6, 7\\\hline
    \verb|ARHI186D_PZ| & 2, 6 & \verb|ARHI186G_PO| & 4, 8 & \verb|ARHI186L_PO| & 1, 5\\\hline
    \verb|ARHI186M_SC| & 2, 3, 6, 7, 1, 5 & \verb|ARHI186N_SC| & 2, 6 & \verb|ARHI186W_PO| & 2, 6\\\hline
    \verb|ARHI186Y_PO| & 3, 7 & \verb|ARHI187__SC| & 2, 4, 6, 8 & \verb|ARHI189__SC| & 3, 7\\\hline
    \verb|ARHI190__PO| & 1, 3, 5, 7 & \verb|ARHI190__PO2| & 1, 3, 5, 7 & \verb|ARHI191__PO| & 2, 4, 6, 8\\\hline
    \verb|ARHI191__PO2| & 2, 4, 6, 8 & \verb|ARHI191__SC| & 4, 8 & \verb|ARHI191__SC2| & 4, 8\\\hline
    \verb|ART_005__PO| & 2, 3, 6, 7, 1, 5, 4, 8 & \verb|ART_010__PO| & 2, 3, 6, 7, 1, 5, 4, 8 & \verb|ART_011__PZ| & 1, 3, 5, 7\\\hline
    \verb|ART_012__PZ| & 3, 7 & \verb|ART_015__PZ| & 2, 3, 6, 7 & \verb|ART_016__PZ| & 4, 8\\\hline
    \verb|ART_017X_PZ| & 4, 8 & \verb|ART_020__PO| & 1, 5 & \verb|ART_020__PZ| & 1, 5\\\hline
    \verb|ART_021__PO| & 2, 1, 6, 5, 4, 8 & \verb|ART_022__PO| & 4, 8 & \verb|ART_024__PO| & 2, 6\\\hline
    \verb|ART_025A_PO| & 1, 5 & \verb|ART_026__PO| & 4, 8 & \verb|ART_027A_PO| & 3, 7\\\hline
    \verb|ART_027__PO| & 4, 8 & \verb|ART_028__PO| & 3, 7 & \verb|ART_030__PZ| & 2, 3, 6, 7, 4, 8\\\hline
    \verb|ART_033__PO| & 3, 7 & \verb|ART_060__HM| & 1, 3, 5, 7 & \verb|ART_075__PZ| & 4, 8\\\hline
    \verb|ART_101__SC| & 1, 3, 5, 7 & \verb|ART_102__SC| & 4, 8 & \verb|ART_105A_PO| & 4, 8\\\hline
    \verb|ART_105__SC| & 2, 3, 6, 7, 1, 5, 4, 8 & \verb|ART_106__SC| & 2, 6 & \verb|ART_111__PO| & 3, 7\\\hline
    \verb|ART_112__PZ| & 4, 8 & \verb|ART_113__PZ| & 2, 6 & \verb|ART_114__PO| & 1, 3, 5, 7\\\hline
    \verb|ART_115__PO| & 2, 6 & \verb|ART_116__SC| & 3, 7 & \verb|ART_119__PZ| & 4, 8\\\hline
    \verb|ART_120__PZ| & 2, 3, 6, 7, 1, 5, 4, 8 & \verb|ART_121__SC| & 2, 3, 6, 7, 1, 5, 4, 8 & \verb|ART_122__PO| & 2, 6\\\hline
    \verb|ART_123__PO| & 2, 6 & \verb|ART_123__SC| & 3, 7 & \verb|ART_125__PZ| & 4, 8\\\hline
    \verb|ART_125__SC| & 2, 3, 6, 7, 1, 5, 4, 8 & \verb|ART_126B_PO| & 3, 7 & \verb|ART_126B_SC| & 2, 6\\\hline
    \verb|ART_126__PZ| & 1, 3, 5, 7 & \verb|ART_127__PZ| & 2, 1, 6, 5 & \verb|ART_129__SC| & 2, 4, 6, 8\\\hline
    \verb|ART_132__SC| & 4, 8 & \verb|ART_133__PZ| & 1, 5 & \verb|ART_134__SC| & 2, 3, 6, 7, 1, 5\\\hline
    \verb|ART_135__SC| & 2, 3, 6, 7, 1, 5, 4, 8 & \verb|ART_136__SC| & 2, 3, 6, 7 & \verb|ART_139__PO| & 2, 6\\\hline
    \verb|ART_140__PO| & 2, 4, 6, 8 & \verb|ART_141__SC| & 2, 3, 6, 7, 1, 5, 4, 8 & \verb|ART_142__SC| & 2, 6\\\hline
    \verb|ART_143__SC| & 2, 6 & \verb|ART_144__SC| & 1, 3, 5, 7, 4, 8 & \verb|ART_145__SC| & 2, 3, 6, 7, 1, 5, 4, 8\\\hline
    \verb|ART_146__SC| & 2, 6 & \verb|ART_147__SC| & 3, 7 & \verb|ART_148__SC| & 2, 3, 6, 7\\\hline
    \verb|ART_149__SC| & 4, 8 & \verb|ART_150__SC| & 4, 8 & \verb|ART_151__SC| & 1, 4, 5, 8\\\hline
    \verb|ART_179M_HM| & 2, 6 & \verb|ART_179S_HM| & 3, 7 & \verb|ART_179T_HM| & 1, 3, 5, 7\\\hline
    \verb|ART_179V_HM| & 2, 1, 6, 5 & \verb|ART_179W_HM| & 3, 7 & \verb|ART_179X_HM| & 1, 3, 5, 7\\\hline
    \verb|ART_179Y_HM| & 4, 8 & \verb|ART_181D_SC| & 2, 6 & \verb|ART_181G_SC| & 4, 8\\\hline
    \verb|ART_181M_SC| & 3, 7 & \verb|ART_181__SC| & 1, 5 & \verb|ART_188__PO| & 2, 6\\\hline
    \verb|ART_189__PZ| & 1, 3, 5, 7 & \verb|ART_190__PO| & 1, 3, 5, 7 & \verb|ART_190__PO2| & 1, 3, 5, 7\\\hline
    \verb|ART_192__PO| & 2, 4, 6, 8 & \verb|ART_192__PO2| & 2, 4, 6, 8 & \verb|ART_192__SC| & 1, 3, 5, 7\\\hline
    \verb|ART_192__SC2| & 1, 3, 5, 7 & \verb|ART_193__SC| & 2, 4, 6, 8 & \verb|ART_193__SC2| & 2, 4, 6, 8\\\hline
    \verb|ART_199__PZ| & 2, 4, 6, 8 & \verb|ART_199__PZ2| & 2, 4, 6, 8 & \verb|ASAM055__AA| & 1, 4, 5, 8\\\hline
    \verb|ASAM070__PO| & 3, 7 & \verb|ASAM077B_PZ| & 3, 7 & \verb|ASAM082__PZ| & 3, 7\\\hline
    \verb|ASAM085__PZ| & 2, 6 & \verb|ASAM086__HM| & 4, 8 & \verb|ASAM088__PZ| & 1, 5\\\hline
    \verb|ASAM089__PZ| & 2, 3, 6, 7, 1, 5 & \verb|ASAM090__PZ| & 3, 7 & \verb|ASAM094__PZ| & 4, 8\\\hline
    \verb|ASAM105B_PZ| & 4, 8 & \verb|ASAM115__CM| & 1, 5 & \verb|ASAM115__PO| & 4, 8\\\hline
    \verb|ASAM123__HM| & 2, 3, 6, 7 & \verb|ASAM124__HM| & 2, 6 & \verb|ASAM125__AA| & 1, 3, 5, 7, 4, 8\\\hline
    \verb|ASAM126__HM| & 1, 5 & \verb|ASAM128__PZ| & 1, 5 & \verb|ASAM130__PZ| & 2, 6\\\hline
    \verb|ASAM135B_PZ| & 1, 5 & \verb|ASAM142__PO| & 2, 6 & \verb|ASAM143__PO| & 2, 6\\\hline
    \verb|ASAM144__PO| & 3, 7 & \verb|ASAM151__AA| & 4, 8 & \verb|ASAM160__AA| & 2, 4, 6, 8\\\hline
    \verb|ASAM161__AA| & 3, 7 & \verb|ASAM172__PZ| & 2, 6 & \verb|ASAM175__PZ| & 1, 5\\\hline
    \verb|ASAM179B_AA| & 3, 7 & \verb|ASAM179H_AA| & 3, 7 & \verb|ASAM179I_AA| & 4, 8\\\hline
    \verb|ASAM179P_AA| & 2, 6 & \verb|ASAM179Q_AA| & 2, 6 & \verb|ASAM179R_HM| & 3, 7\\\hline
    \verb|ASAM187__AA| & 4, 8 & \verb|ASAM190A_PO| & 1, 3, 5, 7 & \verb|ASAM190A_PO2| & 1, 3, 5, 7\\\hline
    \verb|ASAM191__PO| & 2, 1, 6, 5, 4, 8 & \verb|ASAM191__PO2| & 2, 1, 6, 5, 4, 8 & \verb|ASIA081__PO| & 2, 6\\\hline
    \verb|ASIA083__PO| & 2, 6 & \verb|ASIA190__PO| & 1, 3, 5, 7 & \verb|ASIA190__PO2| & 1, 3, 5, 7\\\hline
    \verb|ASIA191__PO| & 2, 4, 6, 8 & \verb|ASIA191__PO2| & 2, 4, 6, 8 & \verb|ASIA192__PO| & 2, 6\\\hline
    \verb|ASIA192__PO2| & 2, 6 & \verb|ASTR001L_PO| & 1, 3, 5, 7 & \verb|ASTR001__PO| & 1, 3, 5, 7\\\hline
    \verb|ASTR002__PO| & 2, 4, 6, 8 & \verb|ASTR051L_PO| & 2, 4, 6, 8 & \verb|ASTR051__PO| & 2, 4, 6, 8\\\hline
    \verb|ASTR062__HM| & 2, 4, 6, 8 & \verb|ASTR066LXKS| & 2, 3, 6, 7 & \verb|ASTR066L_KS| & 2, 3, 6, 7, 4, 8\\\hline
    \verb|ASTR101L_PO| & 1, 3, 5, 7 & \verb|ASTR101__PO| & 1, 3, 5, 7 & \verb|ASTR121__PO| & 2, 6\\\hline
    \verb|ASTR122__HM| & 4, 8 & \verb|ASTR123__PO| & 4, 8 & \verb|ASTR125__PO| & 2, 6\\\hline
    \verb|ASTR127__PO| & 4, 8 & \verb|ASTR128__HM| & 2, 6 & \verb|BIOL001D_PO| & 1, 5\\\hline
    \verb|BIOL007L_PO| & 2, 4, 6, 8 & \verb|BIOL007__PO| & 2, 4, 6, 8 & \verb|BIOL023__HM| & 1, 2\\\hline
    \verb|BIOL040__PO| & 1, 3, 5, 7 & \verb|BIOL041C_PO| & 2, 4, 6, 8 & \verb|BIOL041E_PO| & 2, 4, 6, 8\\\hline
    \verb|BIOL042L_KS| & 2, 4, 6, 8 & \verb|BIOL043LXKS| & 1, 3, 5, 7 & \verb|BIOL043L_KS| & 1, 3, 5, 7\\\hline
    \verb|BIOL044LXKS| & 2, 4, 6, 8 & \verb|BIOL044L_KS| & 2, 4, 6, 8 & \verb|BIOL046R_HM| & 2, 3, 6, 7, 1, 5, 4, 8\\\hline
    \verb|BIOL046__HM| & 1, 2 & \verb|BIOL047__PO| & 2, 6 & \verb|BIOL054__HM| & 2, 4, 6, 8\\\hline
    \verb|BIOL057LXKS| & 3, 7 & \verb|BIOL057L_KS| & 2, 3, 6, 7, 1, 5, 4, 8 & \verb|BIOL059L_KS| & 3, 7\\\hline
    \verb|BIOL099__KS| & 1, 3, 5, 7 & \verb|BIOL101__HM| & 2, 4, 6, 8 & \verb|BIOL103__HM| & 4, 8\\\hline
    \verb|BIOL104A_PO| & 4, 8 & \verb|BIOL107__PO| & 3, 7 & \verb|BIOL108__HM| & 2, 4, 6, 8\\\hline
    \verb|BIOL108__PO| & 4, 8 & \verb|BIOL109__HM| & 3, 7 & \verb|BIOL110__HM| & 2, 6\\\hline
    \verb|BIOL111__HM| & 1, 3, 5, 7 & \verb|BIOL111__KS| & 2, 4, 6, 8 & \verb|BIOL112__KS| & 3, 7\\\hline
    \verb|BIOL112__PO| & 2, 1, 6, 5 & \verb|BIOL113L_KS| & 4, 8 & \verb|BIOL113__HM| & 1, 3, 5, 7\\\hline
    \verb|BIOL116__PO| & 4, 8 & \verb|BIOL119__HM| & 1, 5 & \verb|BIOL119__KS| & 2, 3, 6, 7, 4, 8\\\hline
    \verb|BIOL120__KS| & 1, 3, 5, 7 & \verb|BIOL121__HM| & 3, 7 & \verb|BIOL121__PO| & 3, 7\\\hline
    \verb|BIOL122__HM| & 4, 8 & \verb|BIOL131L_KS| & 2, 3, 6, 7, 1, 5, 4, 8 & \verb|BIOL131__KS| & 2, 4, 6, 8\\\hline
    \verb|BIOL133__PO| & 4, 8 & \verb|BIOL135L_KS| & 2, 4, 6, 8 & \verb|BIOL138L_KS| & 2, 4, 6, 8\\\hline
    \verb|BIOL138__KS| & 2, 4, 6, 8 & \verb|BIOL140__PO| & 1, 3, 5, 7 & \verb|BIOL141L_KS| & 1, 3, 5, 7\\\hline
    \verb|BIOL142L_KS| & 2, 6 & \verb|BIOL143__KS| & 2, 3, 6, 7, 1, 5, 4, 8 & \verb|BIOL145__KS| & 1, 3, 5, 7\\\hline
    \verb|BIOL146L_KS| & 1, 3, 5, 7 & \verb|BIOL147__KS| & 3, 7 & \verb|BIOL151L_KS| & 4, 8\\\hline
    \verb|BIOL152__PO| & 2, 6 & \verb|BIOL154__HM| & 2, 4, 6, 8 & \verb|BIOL154__KS| & 3, 7\\\hline
    \verb|BIOL156L_KS| & 2, 6 & \verb|BIOL157L_KS| & 2, 3, 6, 7, 1, 5, 4, 8 & \verb|BIOL160__HM| & 2, 6\\\hline
    \verb|BIOL160__PO| & 3, 7 & \verb|BIOL161__HM| & 2, 3, 6, 7, 1, 5, 4, 8 & \verb|BIOL162__PO| & 4, 8\\\hline
    \verb|BIOL163L_PO| & 1, 5 & \verb|BIOL163__PO| & 1, 3, 5, 7 & \verb|BIOL164__KS| & 2, 4, 6, 8\\\hline
    \verb|BIOL168L_KS| & 2, 4, 6, 8 & \verb|BIOL169L_PO| & 2, 6 & \verb|BIOL169__PO| & 2, 4, 6, 8\\\hline
    \verb|BIOL170L_KS| & 2, 3, 6, 7, 1, 5, 4, 8 & \verb|BIOL170__PO| & 2, 4, 6, 8 & \verb|BIOL173L_KS| & 4, 8\\\hline
    \verb|BIOL173__PO| & 2, 6 & \verb|BIOL174__PO| & 2, 6 & \verb|BIOL175__KS| & 2, 3, 6, 7, 1, 5, 4, 8\\\hline
    \verb|BIOL176__KS| & 2, 4, 6, 8 & \verb|BIOL177__KS| & 2, 3, 6, 7, 1, 5, 4, 8 & \verb|BIOL182L_KS| & 1, 5\\\hline
    \verb|BIOL182__HM| & 2, 4, 6, 8 & \verb|BIOL183__KS| & 1, 3, 5, 7 & \verb|BIOL184L_KS| & 2, 3, 6, 7\\\hline
    \verb|BIOL184__HM| & 2, 4, 6, 8 & \verb|BIOL185G_HM| & 1, 5 & \verb|BIOL187JLKS| & 1, 3, 5, 7\\\hline
    \verb|BIOL187M_KS| & 2, 4, 6, 8 & \verb|BIOL187P_KS| & 1, 5 & \verb|BIOL188L_KS| & 2, 3, 6, 7, 1, 5, 4, 8\\\hline
    \verb|BIOL188__HM| & 3, 7 & \verb|BIOL189L_KS| & 2, 3, 6, 7, 1, 5, 4, 8 & \verb|BIOL189__HM| & 1, 3, 5, 7\\\hline
    \verb|BIOL190L_KS| & 2, 3, 6, 7, 1, 5, 4, 8 & \verb|BIOL190L_KS2| & 2, 3, 6, 7, 1, 5, 4, 8 & \verb|BIOL190__PO| & 2, 3, 6, 7, 1, 5, 4, 8\\\hline
    \verb|BIOL190__PO2| & 2, 3, 6, 7, 1, 5, 4, 8 & \verb|BIOL191F_PO| & 2, 3, 6, 7, 1, 5, 4, 8 & \verb|BIOL191F_PO2| & 2, 3, 6, 7, 1, 5, 4, 8\\\hline
    \verb|BIOL191H_PO| & 2, 3, 6, 7, 1, 5, 4, 8 & \verb|BIOL191H_PO2| & 2, 3, 6, 7, 1, 5, 4, 8 & \verb|BIOL191__HM| & 7, 8, 5, 6\\\hline
    \verb|BIOL191__HM2| & 7, 8, 5, 6 & \verb|BIOL191__HM3| & 7, 8, 5, 6 & \verb|BIOL191__HM4| & 7, 8, 5, 6\\\hline
    \verb|BIOL191__KS| & 2, 3, 6, 7, 1, 5, 4, 8 & \verb|BIOL191__KS2| & 2, 3, 6, 7, 1, 5, 4, 8 & \verb|BIOL193__HM| & 8, 7\\\hline
    \verb|BIOL193__HM2| & 8, 7 & \verb|BIOL194A_PO| & 1, 3, 5, 7 & \verb|BIOL194A_PO2| & 1, 3, 5, 7\\\hline
    \verb|BIOL194B_PO| & 2, 4, 6, 8 & \verb|BIOL194B_PO2| & 2, 4, 6, 8 & \verb|BIOL195__HM| & 8, 7\\\hline
    \verb|BIOL195__HM2| & 8, 7 & \verb|BIOL196__KS| & 2, 3, 6, 7, 1, 5, 4, 8 & \verb|BIOL196__KS2| & 2, 3, 6, 7, 1, 5, 4, 8\\\hline
    \verb|BIOL197__HM| & 2, 3, 6, 7, 1, 5, 4, 8 & \verb|BIOL197__HM2| & 2, 3, 6, 7, 1, 5, 4, 8 & \verb|BIOL197__KS| & 2, 3, 6, 7, 1, 5, 4, 8\\\hline
    \verb|BIOL197__KS2| & 2, 3, 6, 7, 1, 5, 4, 8 & \verb|BIOL900__HM| & 1, 2, 5, 3, 4, 6, 7, 8 & \verb|BIOL901__HM| & 1, 2, 5, 3, 4, 6, 7, 8\\\hline
    \verb|BIOL902__HM| & 1, 2, 5, 3, 4, 6, 7, 8 & \verb|BIOL903__HM| & 1, 2, 5, 3, 4, 6, 7, 8 & \verb|BIOL904__HM| & 1, 2, 5, 3, 4, 6, 7, 8\\\hline
    \verb|CASA101__PZ| & 2, 3, 6, 7, 1, 5, 4, 8 & \verb|CASA105__PZ| & 2, 3, 6, 7, 1, 5, 4, 8 & \verb|CASA200__PZ| & 2, 6\\\hline
    \verb|CGH_100__JT| & 3, 7 & \verb|CGH_301X_CG| & 1, 5 & \verb|CGS_010__PZ| & 1, 3, 5, 7\\\hline
    \verb|CGS_050__PZ| & 2, 4, 6, 8 & \verb|CGS_060__PZ| & 2, 1, 6, 5, 4, 8 & \verb|CGS_075__PZ| & 2, 6\\\hline
    \verb|CGS_080__PZ| & 2, 3, 6, 7 & \verb|CGS_085__PZ| & 4, 8 & \verb|CGS_101__PZ| & 4, 8\\\hline
    \verb|CGS_120__PZ| & 1, 5 & \verb|CGS_122__PZ| & 1, 3, 5, 7 & \verb|CGS_127__PZ| & 3, 7\\\hline
    \verb|CGS_146__PZ| & 1, 3, 5, 7 & \verb|CGS_150__PZ| & 2, 6 & \verb|CGS_151__PZ| & 2, 4, 6, 8\\\hline
    \verb|CGS_190__PZ| & 2, 4, 6, 8 & \verb|CGS_190__PZ2| & 2, 4, 6, 8 & \verb|CGS_199__PZ| & 4, 8\\\hline
    \verb|CGS_199__PZ2| & 4, 8 & \verb|CHEM001ALPO| & 1, 3, 5, 7 & \verb|CHEM001A_PO| & 1, 3, 5, 7\\\hline
    \verb|CHEM001BLPO| & 2, 4, 6, 8 & \verb|CHEM001B_PO| & 2, 4, 6, 8 & \verb|CHEM014L_KS| & 1, 3, 5, 7\\\hline
    \verb|CHEM015LXKS| & 2, 4, 6, 8 & \verb|CHEM015L_KS| & 2, 4, 6, 8 & \verb|CHEM023A_PO| & 1, 3, 5, 7\\\hline
    \verb|CHEM023BLPO| & 2, 4, 6, 8 & \verb|CHEM023B_PO| & 2, 4, 6, 8 & \verb|CHEM024__HM| & 1, 2\\\hline
    \verb|CHEM029L_KS| & 2, 4, 6, 8 & \verb|CHEM042L_KS| & 4, 8 & \verb|CHEM042R_HM| & 2, 3, 6, 7, 4, 8\\\hline
    \verb|CHEM042__HM| & 1, 2 & \verb|CHEM048__HM| & 4, 8 & \verb|CHEM051L_PO| & 1, 3, 5, 7\\\hline
    \verb|CHEM051__HM| & 1, 3, 5, 7 & \verb|CHEM051__PO| & 1, 3, 5, 7 & \verb|CHEM052__HM| & 2, 6\\\hline
    \verb|CHEM053__HM| & 1, 3, 5, 7 & \verb|CHEM056__HM| & 1, 3, 5, 7 & \verb|CHEM058__HM| & 1, 3, 5, 7\\\hline
    \verb|CHEM060__HM| & 2, 6 & \verb|CHEM070L_KS| & 2, 3, 6, 7, 1, 5 & \verb|CHEM080__HM| & 1, 4, 5, 8\\\hline
    \verb|CHEM081LXJT| & 4, 8 & \verb|CHEM081L_JT| & 4, 8 & \verb|CHEM103__HM| & 4, 8\\\hline
    \verb|CHEM104__HM| & 2, 4, 6, 8 & \verb|CHEM105__HM| & 2, 4, 6, 8 & \verb|CHEM106__PO| & 2, 4, 6, 8\\\hline
    \verb|CHEM109__HM| & 4, 8 & \verb|CHEM110ALPO| & 1, 3, 5, 7 & \verb|CHEM110A_PO| & 1, 3, 5, 7\\\hline
    \verb|CHEM110BLPO| & 2, 4, 6, 8 & \verb|CHEM110B_PO| & 2, 4, 6, 8 & \verb|CHEM110__HM| & 2, 4, 6, 8\\\hline
    \verb|CHEM111__HM| & 2, 4, 6, 8 & \verb|CHEM112__HM| & 2, 1, 6, 5 & \verb|CHEM112__PO| & 2, 6\\\hline
    \verb|CHEM114__HM| & 2, 1, 6, 5 & \verb|CHEM115L_PO| & 2, 3, 6, 7, 1, 5, 4, 8 & \verb|CHEM115__PO| & 2, 3, 6, 7, 1, 5, 4, 8\\\hline
    \verb|CHEM116LXKS| & 1, 3, 5, 7 & \verb|CHEM116L_KS| & 1, 3, 5, 7 & \verb|CHEM116__HM| & 1, 5\\\hline
    \verb|CHEM117LXKS| & 2, 4, 6, 8 & \verb|CHEM117L_KS| & 2, 4, 6, 8 & \verb|CHEM121__KS| & 2, 4, 6, 8\\\hline
    \verb|CHEM122__HM| & 3, 7 & \verb|CHEM122__KS| & 1, 3, 5, 7 & \verb|CHEM123__KS| & 2, 6\\\hline
    \verb|CHEM126L_KS| & 1, 3, 5, 7 & \verb|CHEM127L_KS| & 2, 4, 6, 8 & \verb|CHEM128__KS| & 4, 8\\\hline
    \verb|CHEM140__KS| & 3, 7 & \verb|CHEM147L_PO| & 2, 4, 6, 8 & \verb|CHEM147__PO| & 2, 4, 6, 8\\\hline
    \verb|CHEM150__HM| & 2, 3, 6, 7, 1, 5, 4, 8 & \verb|CHEM151__HM| & 7 & \verb|CHEM151__PO| & 4, 8\\\hline
    \verb|CHEM152__HM| & 8 & \verb|CHEM156__PO| & 2, 4, 6, 8 & \verb|CHEM158A_PO| & 1, 3, 5, 7\\\hline
    \verb|CHEM158BLPO| & 2, 4, 6, 8 & \verb|CHEM158B_PO| & 2, 4, 6, 8 & \verb|CHEM161L_PO| & 1, 3, 5, 7\\\hline
    \verb|CHEM161__PO| & 1, 3, 5, 7 & \verb|CHEM164__PO| & 1, 3, 5, 7 & \verb|CHEM170__HM| & 4, 8\\\hline
    \verb|CHEM171__HM| & 3, 7 & \verb|CHEM175__PO| & 4, 8 & \verb|CHEM177__KS| & 2, 3, 6, 7, 1, 5, 4, 8\\\hline
    \verb|CHEM181__KS| & 2, 6 & \verb|CHEM181__PO| & 2, 6 & \verb|CHEM182__HM| & 2, 4, 6, 8\\\hline
    \verb|CHEM184__HM| & 2, 4, 6, 8 & \verb|CHEM187__PO| & 1, 3, 5, 7 & \verb|CHEM188L_KS| & 2, 3, 6, 7, 1, 5, 4, 8\\\hline
    \verb|CHEM189L_KS| & 2, 3, 6, 7, 1, 5 & \verb|CHEM190L_KS| & 2, 3, 6, 7, 1, 5, 4, 8 & \verb|CHEM190L_KS2| & 2, 3, 6, 7, 1, 5, 4, 8\\\hline
    \verb|CHEM191__KS| & 2, 3, 6, 7, 1, 5, 4, 8 & \verb|CHEM191__KS2| & 2, 3, 6, 7, 1, 5, 4, 8 & \verb|CHEM191__PO| & 1, 3, 5, 7, 4, 8\\\hline
    \verb|CHEM191__PO2| & 1, 3, 5, 7, 4, 8 & \verb|CHEM193M_HM| & 3, 7 & \verb|CHEM193M_HM2| & 3, 7\\\hline
    \verb|CHEM193P_HM| & 2, 6 & \verb|CHEM193P_HM2| & 2, 6 & \verb|CHEM194__HM| & 4, 8\\\hline
    \verb|CHEM194__HM2| & 4, 8 & \verb|CHEM194__PO| & 1, 3, 5, 7, 4, 8 & \verb|CHEM194__PO2| & 1, 3, 5, 7, 4, 8\\\hline
    \verb|CHEM196__KS| & 2, 3, 6, 7, 1, 5, 4, 8 & \verb|CHEM196__KS2| & 2, 3, 6, 7, 1, 5, 4, 8 & \verb|CHEM197__HM| & 1, 3, 5, 7\\\hline
    \verb|CHEM197__HM2| & 1, 3, 5, 7 & \verb|CHEM197__KS| & 2, 3, 6, 7, 1, 5, 4, 8 & \verb|CHEM197__KS2| & 2, 3, 6, 7, 1, 5, 4, 8\\\hline
    \verb|CHEM198__HM| & 2, 6 & \verb|CHEM198__HM2| & 2, 6 & \verb|CHEM199__HM| & 7, 8, 5, 6\\\hline
    \verb|CHEM199__HM2| & 7, 8, 5, 6 & \verb|CHEM199__HM3| & 7, 8, 5, 6 & \verb|CHEM199__HM4| & 7, 8, 5, 6\\\hline
    \verb|CHEM900__HM| & 1, 2, 5, 3, 4, 6, 7, 8 & \verb|CHEM901__HM| & 1, 2, 5, 3, 4, 6, 7, 8 & \verb|CHEM902__HM| & 1, 2, 5, 3, 4, 6, 7, 8\\\hline
    \verb|CHEM903__HM| & 1, 2, 5, 3, 4, 6, 7, 8 & \verb|CHEM904__HM| & 1, 2, 5, 3, 4, 6, 7, 8 & \verb|CHIN001A_PO| & 1, 3, 5, 7\\\hline
    \verb|CHIN001B_PO| & 2, 4, 6, 8 & \verb|CHIN002__PO| & 1, 3, 5, 7 & \verb|CHIN011__PO| & 1, 3, 5, 7, 4, 8\\\hline
    \verb|CHIN013__PO| & 1, 3, 5, 7, 4, 8 & \verb|CHIN051A_PO| & 1, 3, 5, 7 & \verb|CHIN051B_PO| & 2, 4, 6, 8\\\hline
    \verb|CHIN051H_PO| & 2, 4, 6, 8 & \verb|CHIN055__PO| & 2, 6 & \verb|CHIN111A_PO| & 1, 3, 5, 7\\\hline
    \verb|CHIN111B_PO| & 2, 4, 6, 8 & \verb|CHIN123__PO| & 2, 6 & \verb|CHIN124__PO| & 4, 8\\\hline
    \verb|CHIN129__PO| & 1, 3, 5, 7 & \verb|CHIN131__PO| & 3, 7 & \verb|CHIN150__PO| & 1, 5\\\hline
    \verb|CHIN153__PO| & 3, 7 & \verb|CHIN192A_PO| & 2, 3, 6, 7 & \verb|CHIN192A_PO2| & 2, 3, 6, 7\\\hline
    \verb|CHLT013__CH| & 4, 8 & \verb|CHLT060__CH| & 2, 6 & \verb|CHLT072__CH| & 1, 3, 5, 7\\\hline
    \verb|CHLT085__PZ| & 2, 4, 6, 8 & \verb|CHLT103__CH| & 1, 5 & \verb|CHLT138__CH| & 1, 5\\\hline
    \verb|CHLT151__CH| & 4, 8 & \verb|CHLT170__CH| & 2, 6 & \verb|CHLT181__CH| & 1, 5\\\hline
    \verb|CHNT168__PO| & 2, 4, 6, 8 & \verb|CHNT178__PO| & 4, 8 & \verb|CHNT180__PO| & 2, 6\\\hline
    \verb|CHST015__CH| & 2, 4, 6, 8 & \verb|CHST028__CH| & 2, 1, 6, 5, 4, 8 & \verb|CHST064__CH| & 1, 3, 5, 7\\\hline
    \verb|CHST066__CH| & 2, 4, 6, 8 & \verb|CHST067__CH| & 4, 8 & \verb|CHST074__CH| & 1, 5\\\hline
    \verb|CHST081__PO| & 1, 5 & \verb|CHST102__PO| & 4, 8 & \verb|CHST120__CH| & 2, 4, 6, 8\\\hline
    \verb|CHST124__PO| & 4, 8 & \verb|CHST125__CH| & 2, 6 & \verb|CHST128__CH| & 8, 3, 7, 4\\\hline
    \verb|CHST130__CH| & 2, 3, 6, 7 & \verb|CHST132__CH| & 2, 1, 6, 5 & \verb|CHST136__CH| & 3, 7\\\hline
    \verb|CHST182__CH| & 2, 4, 6, 8 & \verb|CHST184D_CH| & 3, 7 & \verb|CHST184__CH| & 1, 5\\\hline
    \verb|CHST185C_CH| & 2, 4, 6, 8 & \verb|CHST189B_PO| & 1, 3, 5, 7, 4, 8 & \verb|CHST190__CH| & 1, 3, 5, 7\\\hline
    \verb|CHST190__CH2| & 1, 3, 5, 7 & \verb|CHST191__CH| & 4, 8 & \verb|CHST191__CH2| & 4, 8\\\hline
    \verb|CHST192__CH| & 4, 8 & \verb|CHST192__CH2| & 4, 8 & \verb|CLAS001__PO| & 2, 4, 6, 8\\\hline
    \verb|CLAS010__SC| & 3, 7 & \verb|CLAS012__SC| & 1, 3, 5, 7 & \verb|CLAS019__SC| & 3, 7\\\hline
    \verb|CLAS023__SC| & 2, 6 & \verb|CLAS064__PO| & 1, 5 & \verb|CLAS088__PZ| & 1, 4, 5, 8\\\hline
    \verb|CLAS112__PO| & 1, 5 & \verb|CLAS114__PO| & 4, 8 & \verb|CLAS116A_PO| & 2, 6\\\hline
    \verb|CLAS116__SC| & 2, 6 & \verb|CLAS117__PO| & 3, 7 & \verb|CLAS121__JT| & 2, 6\\\hline
    \verb|CLAS133__SC| & 4, 8 & \verb|CLAS150BEPZ| & 2, 6 & \verb|CLAS161__PZ| & 3, 7\\\hline
    \verb|CLAS162__PZ| & 4, 8 & \verb|CLAS175__PZ| & 4, 8 & \verb|CLAS190__PO| & 1, 3, 5, 7\\\hline
    \verb|CLAS190__PO2| & 1, 3, 5, 7 & \verb|CLAS191__PO| & 2, 6 & \verb|CLAS191__PO2| & 2, 6\\\hline
    \verb|CLAS191__SC| & 4, 8 & \verb|CLAS191__SC2| & 4, 8 & \verb|CLAS192__PO| & 2, 6\\\hline
    \verb|CLAS192__PO2| & 2, 6 & \verb|CLES049A_HM| & 1, 5 & \verb|CLES101__HM| & 1, 3, 5, 7\\\hline
    \verb|CLES120__HM| & 2, 4, 6, 8 & \verb|CLES121__HM| & 4, 8 & \verb|CLES122__HM| & 4, 8\\\hline
    \verb|CLES130__HM| & 1, 5 & \verb|CLES131__HM| & 1, 5 & \verb|CLES179A_HM| & 2, 4, 6, 8\\\hline
    \verb|CLES179B_HM| & 2, 6 & \verb|CLES179C_HM| & 3, 7 & \verb|CLES197__HM| & 1, 4, 5, 8\\\hline
    \verb|CLES197__HM2| & 1, 4, 5, 8 & \verb|CLES199__HM| & 1, 5 & \verb|CLES199__HM2| & 1, 5\\\hline
    \verb|COGS011__PZ| & 2, 4, 6, 8 & \verb|COGS123L_JT| & 2, 6 & \verb|COGS182__PZ| & 3, 7\\\hline
    \verb|CORE001__SC| & 3, 7 & \verb|CORE002__SC| & 2, 4, 6, 8 & \verb|CORE003__SC| & 1, 3, 5, 7\\\hline
    \verb|CORE010A_SC| & 1, 5 & \verb|CORE079R_HM| & 2, 6 & \verb|CORE079__HM| & 2, 6\\\hline
    \verb|CORE099R_HM| & 4, 8 & \verb|CORE099__HM| & 4 & \verb|CSCI004__PZ| & 2, 4, 6, 8\\\hline
    \verb|CSCI005L_HM| & 1, 3, 5, 7 & \verb|CSCI005__HM| & 1 & \verb|CSCI022__SC| & 2, 6\\\hline
    \verb|CSCI035__HM| & 4, 8 & \verb|CSCI036P_PZ| & 3, 7 & \verb|CSCI036__CM| & 2, 3, 6, 7, 1, 5, 4, 8\\\hline
    \verb|CSCI036__PZ| & 2, 4, 6, 8 & \verb|CSCI040__CM| & 1, 3, 5, 7, 4, 8 & \verb|CSCI046L_CM| & 2, 1, 6, 5, 4, 8\\\hline
    \verb|CSCI046__CM| & 2, 3, 6, 7, 1, 5, 4, 8 & \verb|CSCI050__PO| & 1, 5 & \verb|CSCI051ALPO| & 1, 4, 5, 8\\\hline
    \verb|CSCI051A_PO| & 1, 4, 5, 8 & \verb|CSCI051L_PO| & 1, 5 & \verb|CSCI051PLPO| & 2, 3, 6, 7\\\hline
    \verb|CSCI051P_PO| & 2, 3, 6, 7 & \verb|CSCI051__PO| & 1, 5 & \verb|CSCI054__PO| & 2, 3, 6, 7, 1, 5, 4, 8\\\hline
    \verb|CSCI060__HM| & 2, 3, 6, 7, 1, 5, 4, 8 & \verb|CSCI062L_PO| & 2, 3, 6, 7, 1, 5, 4, 8 & \verb|CSCI062__PO| & 2, 3, 6, 7, 1, 5, 4, 8\\\hline
    \verb|CSCI070__HM| & 2, 3, 6, 7, 1, 5, 4, 8 & \verb|CSCI081__HM| & 2, 3, 6, 7, 1, 5, 4, 8 & \verb|CSCI101__PO| & 2, 3, 6, 7, 1, 5, 4, 8\\\hline
    \verb|CSCI105L_PO| & 2, 3, 6, 7, 1, 5, 4, 8 & \verb|CSCI105__HM| & 2, 3, 6, 7, 1, 5, 4, 8 & \verb|CSCI105__PO| & 2, 3, 6, 7, 1, 5, 4, 8\\\hline
    \verb|CSCI120__HM| & 1, 5 & \verb|CSCI123__HM| & 2, 3, 6, 7, 1, 5, 4, 8 & \verb|CSCI131__HM| & 2, 3, 6, 7, 4, 8\\\hline
    \verb|CSCI133__HM| & 3, 4, 8, 7 & \verb|CSCI134__HM| & 3, 7 & \verb|CSCI140__HM| & 2, 3, 6, 7, 1, 5, 4, 8\\\hline
    \verb|CSCI140__PO| & 2, 3, 6, 7, 1, 5, 4, 8 & \verb|CSCI143__CM| & 2, 4, 6, 8 & \verb|CSCI144__HM| & 2, 4, 6, 8\\\hline
    \verb|CSCI145__CM| & 1, 3, 5, 7 & \verb|CSCI148__CM| & 4, 8 & \verb|CSCI151__HM| & 4, 8\\\hline
    \verb|CSCI152__HM| & 2, 4, 6, 8 & \verb|CSCI152__PO| & 4, 8 & \verb|CSCI153__HM| & 1, 3, 5, 7, 4, 8\\\hline
    \verb|CSCI158__HM| & 2, 1, 6, 5, 4, 8 & \verb|CSCI158__PO| & 1, 5 & \verb|CSCI158__PZ| & 2, 4, 6, 8\\\hline
    \verb|CSCI159__HM| & 2, 1, 6, 5 & \verb|CSCI159__PO| & 3, 7 & \verb|CSCI181AAHM| & 2, 6\\\hline
    \verb|CSCI181AAPO| & 1, 5 & \verb|CSCI181ACHM| & 2, 6 & \verb|CSCI181AFHM| & 4, 8\\\hline
    \verb|CSCI181AGHM| & 3, 7 & \verb|CSCI181ALHM| & 3, 7 & \verb|CSCI181AOHM| & 2, 6\\\hline
    \verb|CSCI181APHM| & 1, 3, 5, 7, 4, 8 & \verb|CSCI181AQHM| & 4, 8 & \verb|CSCI181ARHM| & 4, 8\\\hline
    \verb|CSCI181ATHM| & 1, 5 & \verb|CSCI181AUHM| & 1, 5 & \verb|CSCI181CAPO| & 1, 5\\\hline
    \verb|CSCI181DTPO| & 2, 3, 6, 7 & \verb|CSCI181DVPO| & 4, 8 & \verb|CSCI181G_PO| & 2, 6\\\hline
    \verb|CSCI181NWPO| & 1, 5 & \verb|CSCI181RTPO| & 4, 8 & \verb|CSCI181R_PO| & 3, 7\\\hline
    \verb|CSCI181S_PO| & 2, 6 & \verb|CSCI181V_HM| & 1, 5 & \verb|CSCI181V_PZ| & 4, 8\\\hline
    \verb|CSCI181W_PO| & 4, 8 & \verb|CSCI181__CM| & 3, 7 & \verb|CSCI183__HM| & 7\\\hline
    \verb|CSCI184__HM| & 8 & \verb|CSCI186__HM| & 2, 3, 6, 7, 1, 5, 4, 8 & \verb|CSCI188__PO| & 2, 3, 6, 7, 1, 5, 4, 8\\\hline
    \verb|CSCI189__HM| & 2, 3, 6, 7, 4, 8 & \verb|CSCI190__PO| & 2, 3, 6, 7, 1, 5, 4, 8 & \verb|CSCI190__PO2| & 2, 3, 6, 7, 1, 5, 4, 8\\\hline
    \verb|CSCI191__PO| & 2, 3, 6, 7, 1, 5, 4, 8 & \verb|CSCI191__PO2| & 2, 3, 6, 7, 1, 5, 4, 8 & \verb|CSCI192__PO| & 2, 4, 6, 8\\\hline
    \verb|CSCI192__PO2| & 2, 4, 6, 8 & \verb|CSCI195__HM| & 7, 8, 5, 6 & \verb|CSCI195__HM2| & 7, 8, 5, 6\\\hline
    \verb|CSCI195__HM3| & 7, 8, 5, 6 & \verb|CSCI195__HM4| & 7, 8, 5, 6 & \verb|CSCI198__SC| & 4, 8\\\hline
    \verb|CSCI198__SC2| & 4, 8 & \verb|CSCI900__HM| & 1, 2, 5, 3, 4, 6, 7, 8 & \verb|CSCI901__HM| & 1, 2, 5, 3, 4, 6, 7, 8\\\hline
    \verb|CSCI902__HM| & 1, 2, 5, 3, 4, 6, 7, 8 & \verb|CSCI903__HM| & 1, 2, 5, 3, 4, 6, 7, 8 & \verb|CSCI904__HM| & 1, 2, 5, 3, 4, 6, 7, 8\\\hline
    \verb|CSMT183__HM| & 8, 7 & \verb|CSMT184__HM| & 8, 7 & \verb|DANC010__PO| & 2, 3, 6, 7, 1, 5, 4, 8\\\hline
    \verb|DANC012P_PO| & 2, 3, 6, 7, 1, 5, 4, 8 & \verb|DANC012__PO| & 2, 3, 6, 7, 1, 5, 4, 8 & \verb|DANC050P_PO| & 2, 3, 6, 7, 1, 5, 4, 8\\\hline
    \verb|DANC050__PO| & 2, 3, 6, 7, 1, 5, 4, 8 & \verb|DANC051P_PO| & 2, 3, 6, 7, 1, 5, 4, 8 & \verb|DANC051__PO| & 2, 3, 6, 7, 1, 5, 4, 8\\\hline
    \verb|DANC068__SC| & 1, 3, 5, 7 & \verb|DANC071A_SC| & 2, 3, 6, 7, 4, 8 & \verb|DANC071B_SC| & 2, 3, 6, 7, 4, 8\\\hline
    \verb|DANC081A_SC| & 4, 8 & \verb|DANC081B_SC| & 4, 8 & \verb|DANC083A_SC| & 2, 1, 6, 5, 4, 8\\\hline
    \verb|DANC083B_SC| & 2, 6 & \verb|DANC100A_SC| & 2, 6 & \verb|DANC100B_SC| & 2, 3, 6, 7, 1, 5\\\hline
    \verb|DANC102__SC| & 1, 5 & \verb|DANC108A_SC| & 4, 8 & \verb|DANC108B_SC| & 4, 8\\\hline
    \verb|DANC111A_SC| & 2, 3, 6, 7, 1, 5, 4, 8 & \verb|DANC111B_SC| & 2, 3, 6, 7, 1, 5, 4, 8 & \verb|DANC114A_SC| & 2, 1, 6, 5, 4, 8\\\hline
    \verb|DANC114B_SC| & 2, 1, 6, 5, 4, 8 & \verb|DANC117B_SC| & 4, 8 & \verb|DANC119__PO| & 2, 3, 6, 7, 1, 5, 4, 8\\\hline
    \verb|DANC120P_PO| & 2, 3, 6, 7, 1, 5, 4, 8 & \verb|DANC120__PO| & 2, 3, 6, 7, 1, 5, 4, 8 & \verb|DANC122P_PO| & 2, 3, 6, 7, 1, 5, 4, 8\\\hline
    \verb|DANC122__PO| & 2, 3, 6, 7, 1, 5, 4, 8 & \verb|DANC123__PO| & 2, 3, 6, 7, 1, 5, 4, 8 & \verb|DANC124P_PO| & 2, 3, 6, 7, 1, 5, 4, 8\\\hline
    \verb|DANC124__PO| & 2, 3, 6, 7, 1, 5, 4, 8 & \verb|DANC130__PO| & 3, 7 & \verb|DANC131__SC| & 2, 1, 6, 5\\\hline
    \verb|DANC135__PO| & 3, 7 & \verb|DANC136B_SC| & 3, 7 & \verb|DANC136__PO| & 3, 7\\\hline
    \verb|DANC138__PO| & 2, 4, 6, 8 & \verb|DANC139__PO| & 2, 4, 6, 8 & \verb|DANC140__PO| & 2, 3, 6, 7, 1, 5\\\hline
    \verb|DANC142P_PO| & 4, 8 & \verb|DANC142__PO| & 4, 8 & \verb|DANC150A_PO| & 2, 6\\\hline
    \verb|DANC150C_PO| & 2, 3, 6, 7, 1, 5, 4, 8 & \verb|DANC151P_PO| & 2, 3, 6, 7, 1, 5, 4, 8 & \verb|DANC151__PO| & 2, 3, 6, 7, 1, 5, 4, 8\\\hline
    \verb|DANC152IOSC| & 1, 3, 5, 7 & \verb|DANC152P_PO| & 2, 3, 6, 7, 1, 5, 4, 8 & \verb|DANC152__PO| & 2, 3, 6, 7, 1, 5, 4, 8\\\hline
    \verb|DANC153P_PO| & 1, 5 & \verb|DANC153__PO| & 1, 5 & \verb|DANC159__SC| & 3, 7\\\hline
    \verb|DANC160__PO| & 2, 1, 6, 5 & \verb|DANC161__SC| & 2, 6 & \verb|DANC162A_SC| & 2, 3, 6, 7, 1, 5, 4, 8\\\hline
    \verb|DANC162B_SC| & 2, 3, 6, 7, 1, 5, 4, 8 & \verb|DANC163__SC| & 3, 7 & \verb|DANC166P_PO| & 4, 8\\\hline
    \verb|DANC166__PO| & 4, 8 & \verb|DANC170__PO| & 1, 5 & \verb|DANC172__PO| & 2, 3, 6, 7, 1, 5, 4, 8\\\hline
    \verb|DANC173__PO| & 2, 3, 6, 7, 1, 5, 4, 8 & \verb|DANC174__PO| & 2, 3, 6, 7, 1, 5, 4, 8 & \verb|DANC175__PO| & 2, 3, 6, 7, 1, 5, 4, 8\\\hline
    \verb|DANC176__PO| & 2, 3, 6, 7, 1, 5, 4, 8 & \verb|DANC180P_PO| & 2, 3, 6, 7, 1, 5, 4, 8 & \verb|DANC180__PO| & 2, 3, 6, 7, 1, 5, 4, 8\\\hline
    \verb|DANC180__SC| & 4, 8 & \verb|DANC181P_PO| & 2, 3, 6, 7, 1, 5, 4, 8 & \verb|DANC181__PO| & 2, 3, 6, 7, 1, 5, 4, 8\\\hline
    \verb|DANC182__PO| & 1, 4, 5, 8 & \verb|DANC190__SC| & 1, 3, 5, 7 & \verb|DANC190__SC2| & 1, 3, 5, 7\\\hline
    \verb|DANC191__SC| & 2, 3, 6, 7, 4, 8 & \verb|DANC191__SC2| & 2, 3, 6, 7, 4, 8 & \verb|DANC192__PO| & 1, 3, 5, 7, 4, 8\\\hline
    \verb|DANC192__PO2| & 1, 3, 5, 7, 4, 8 & \verb|DANC193__SC| & 2, 3, 6, 7, 1, 5, 4, 8 & \verb|DANC193__SC2| & 2, 3, 6, 7, 1, 5, 4, 8\\\hline
    \verb|DS__001__SC| & 1, 5 & \verb|DS__002R_PO| & 1, 3, 5, 7, 4, 8 & \verb|DS__002__SC| & 2, 4, 6, 8\\\hline
    \verb|DS__108__PZ| & 2, 3, 6, 7 & \verb|DS__180__CM| & 2, 3, 6, 7, 1, 5, 4, 8 & \verb|DS__183__SC| & 4, 8\\\hline
    \verb|DS__190__PO| & 1, 3, 5, 7 & \verb|DS__190__PO2| & 1, 3, 5, 7 & \verb|DS__191__SC| & 1, 4, 5, 8\\\hline
    \verb|DS__191__SC2| & 1, 4, 5, 8 & \verb|EA__010__PO| & 2, 3, 6, 7, 1, 5, 4, 8 & \verb|EA__010__PZ| & 2, 3, 6, 7, 1, 5, 4, 8\\\hline
    \verb|EA__010__SC| & 2, 3, 6, 7, 1, 5 & \verb|EA__020__PO| & 2, 3, 6, 7, 1, 5, 4, 8 & \verb|EA__030L_KS| & 2, 3, 6, 7, 1, 5, 4, 8\\\hline
    \verb|EA__030L_PO| & 2, 1, 6, 5, 4, 8 & \verb|EA__030__PO| & 2, 3, 6, 7, 1, 5, 4, 8 & \verb|EA__031__PZ| & 2, 4, 6, 8\\\hline
    \verb|EA__032__PZ| & 2, 3, 6, 7 & \verb|EA__034__PZ| & 3, 7 & \verb|EA__039__PZ| & 2, 6\\\hline
    \verb|EA__055LCKS| & 1, 5 & \verb|EA__086__PZ| & 2, 3, 6, 7, 4, 8 & \verb|EA__086__SC| & 4, 8\\\hline
    \verb|EA__097__PZ| & 4, 8 & \verb|EA__100L_KS| & 1, 3, 5, 7 & \verb|EA__100__KS| & 1, 3, 5, 7\\\hline
    \verb|EA__101L_PO| & 3, 7 & \verb|EA__101__PO| & 2, 3, 6, 7, 4, 8 & \verb|EA__103L_KS| & 4, 8\\\hline
    \verb|EA__103__KS| & 2, 4, 6, 8 & \verb|EA__103__PZ| & 1, 5 & \verb|EA__105__KS| & 1, 5\\\hline
    \verb|EA__106__KS| & 1, 5 & \verb|EA__106__PZ| & 4, 8 & \verb|EA__112__PZ| & 2, 1, 6, 5, 4, 8\\\hline
    \verb|EA__126__PZ| & 1, 4, 5, 8 & \verb|EA__129__PZ| & 1, 3, 5, 7, 4, 8 & \verb|EA__130__PZ| & 4, 8\\\hline
    \verb|EA__133__PZ| & 3, 7 & \verb|EA__134__PZ| & 2, 4, 6, 8 & \verb|EA__135__PZ| & 4, 8\\\hline
    \verb|EA__138__PZ| & 2, 3, 6, 7, 1, 5, 4, 8 & \verb|EA__141__PZ| & 8, 3, 7, 4 & \verb|EA__142__PZ| & 2, 3, 6, 7, 4, 8\\\hline
    \verb|EA__143__PZ| & 4, 8 & \verb|EA__144__PZ| & 1, 5 & \verb|EA__150__PZ| & 1, 3, 5, 7\\\hline
    \verb|EA__153__PZ| & 2, 6 & \verb|EA__163__PZ| & 3, 7 & \verb|EA__170__PO| & 3, 7\\\hline
    \verb|EA__171__PO| & 1, 5 & \verb|EA__180__PO| & 1, 3, 5, 7 & \verb|EA__182__PO| & 1, 5\\\hline
    \verb|EA__185__PO| & 2, 1, 6, 5, 4, 8 & \verb|EA__188L_KS| & 1, 3, 5, 7, 4, 8 & \verb|EA__189F_PO| & 1, 5\\\hline
    \verb|EA__189G_PO| & 3, 7 & \verb|EA__189L_KS| & 1, 3, 5, 7 & \verb|EA__189M_PO| & 3, 7\\\hline
    \verb|EA__190L_KS| & 1, 3, 5, 7, 4, 8 & \verb|EA__190L_KS2| & 1, 3, 5, 7, 4, 8 & \verb|EA__190__PO| & 2, 4, 6, 8\\\hline
    \verb|EA__190__PO2| & 2, 4, 6, 8 & \verb|EA__191H_PO| & 4, 8 & \verb|EA__191H_PO2| & 4, 8\\\hline
    \verb|EA__191__KS| & 1, 3, 5, 7, 4, 8 & \verb|EA__191__KS2| & 1, 3, 5, 7, 4, 8 & \verb|EA__191__PO| & 1, 3, 5, 7\\\hline
    \verb|EA__191__PO2| & 1, 3, 5, 7 & \verb|EA__196__KS| & 2, 3, 6, 7, 1, 5, 4, 8 & \verb|EA__196__KS2| & 2, 3, 6, 7, 1, 5, 4, 8\\\hline
    \verb|EA__197__KS| & 2, 3, 6, 7, 1, 5, 4, 8 & \verb|EA__197__KS2| & 2, 3, 6, 7, 1, 5, 4, 8 & \verb|EA__197__PZ| & 2, 4, 6, 8\\\hline
    \verb|EA__197__PZ2| & 2, 4, 6, 8 & \verb|ECON020__PZ| & 1, 3, 5, 7, 4, 8 & \verb|ECON030__PZ| & 2, 6\\\hline
    \verb|ECON050__CM| & 2, 3, 6, 7, 1, 5, 4, 8 & \verb|ECON050__SC| & 1, 3, 5, 7, 4, 8 & \verb|ECON051__PO| & 2, 3, 6, 7, 1, 5, 4, 8\\\hline
    \verb|ECON051__PZ| & 2, 3, 6, 7, 1, 5, 4, 8 & \verb|ECON051__SC| & 2, 3, 6, 7 & \verb|ECON052__PO| & 2, 3, 6, 7, 1, 5, 4, 8\\\hline
    \verb|ECON052__PZ| & 2, 3, 6, 7, 1, 5, 4, 8 & \verb|ECON052__SC| & 2, 6 & \verb|ECON053__HM| & 2, 3, 6, 7, 1, 5, 4, 8\\\hline
    \verb|ECON054__HM| & 1, 3, 5, 7 & \verb|ECON057B_PO| & 2, 3, 6, 7, 1, 5 & \verb|ECON057__PO| & 2, 3, 6, 7, 1, 5, 4, 8\\\hline
    \verb|ECON086__CM| & 2, 3, 6, 7, 1, 5, 4, 8 & \verb|ECON089A_PO| & 2, 1, 6, 5 & \verb|ECON091__PZ| & 2, 3, 6, 7, 1, 5, 4, 8\\\hline
    \verb|ECON097__CM| & 1, 3, 5, 7 & \verb|ECON100__CM| & 2, 3, 6, 7, 1, 5, 4, 8 & \verb|ECON101__CM| & 2, 3, 6, 7, 1, 5, 4, 8\\\hline
    \verb|ECON101__PO| & 2, 3, 6, 7, 1, 5, 4, 8 & \verb|ECON101__SC| & 2, 3, 6, 7, 1, 5, 4, 8 & \verb|ECON102__CM| & 2, 3, 6, 7, 1, 5, 4, 8\\\hline
    \verb|ECON102__PO| & 2, 3, 6, 7, 1, 5, 4, 8 & \verb|ECON102__SC| & 2, 3, 6, 7, 1, 5, 4, 8 & \verb|ECON104__PZ| & 2, 1, 6, 5, 4, 8\\\hline
    \verb|ECON105__PZ| & 2, 3, 6, 7, 1, 5, 4, 8 & \verb|ECON107__CM| & 2, 3, 6, 7 & \verb|ECON107__PO| & 2, 3, 6, 7, 1, 5, 4, 8\\\hline
    \verb|ECON111__SC| & 1, 5 & \verb|ECON114__SC| & 2, 4, 6, 8 & \verb|ECON115__PO| & 2, 6\\\hline
    \verb|ECON117__PO| & 2, 4, 6, 8 & \verb|ECON118__CM| & 1, 3, 5, 7 & \verb|ECON120__CM| & 2, 3, 6, 7, 1, 5, 4, 8\\\hline
    \verb|ECON120__PO| & 4, 8 & \verb|ECON120__SC| & 1, 3, 5, 7, 4, 8 & \verb|ECON121__PO| & 2, 4, 6, 8\\\hline
    \verb|ECON122__CM| & 1, 5 & \verb|ECON122__PO| & 2, 6 & \verb|ECON123__PO| & 1, 3, 5, 7\\\hline
    \verb|ECON124__PO| & 3, 7 & \verb|ECON125__CM| & 2, 3, 6, 7, 1, 5, 4, 8 & \verb|ECON125__PO| & 2, 4, 6, 8\\\hline
    \verb|ECON125__PZ| & 2, 3, 6, 7, 1, 5, 4, 8 & \verb|ECON125__SC| & 2, 3, 6, 7, 1, 5, 4, 8 & \verb|ECON126__CM| & 2, 4, 6, 8\\\hline
    \verb|ECON126__PO| & 1, 3, 5, 7 & \verb|ECON127__PO| & 1, 3, 5, 7 & \verb|ECON128__PO| & 2, 4, 6, 8\\\hline
    \verb|ECON129__CM| & 1, 3, 5, 7 & \verb|ECON129__PO| & 4, 8 & \verb|ECON130__PO| & 2, 4, 6, 8\\\hline
    \verb|ECON131__PO| & 1, 5 & \verb|ECON132__PO| & 2, 3, 6, 7 & \verb|ECON133__SC| & 2, 4, 6, 8\\\hline
    \verb|ECON134B_CM| & 2, 4, 6, 8 & \verb|ECON134E_CM| & 4, 8 & \verb|ECON134__CM| & 2, 3, 6, 7, 1, 5, 4, 8\\\hline
    \verb|ECON134__SC| & 1, 3, 5, 7 & \verb|ECON135__CM| & 1, 3, 5, 7 & \verb|ECON135__SC| & 4, 8\\\hline
    \verb|ECON136__CM| & 2, 6 & \verb|ECON138__CM| & 4, 8 & \verb|ECON139__CM| & 1, 3, 5, 7\\\hline
    \verb|ECON140__CM| & 4, 8 & \verb|ECON140__PZ| & 3, 7 & \verb|ECON141__CM| & 2, 3, 6, 7, 1, 5\\\hline
    \verb|ECON142__CM| & 2, 4, 6, 8 & \verb|ECON142__PZ| & 3, 7 & \verb|ECON143__PO| & 3, 7\\\hline
    \verb|ECON145__CM| & 2, 4, 6, 8 & \verb|ECON145__PZ| & 8, 3, 7, 4 & \verb|ECON146__HM| & 3, 7\\\hline
    \verb|ECON150__CM| & 2, 3, 6, 7, 1, 5, 4, 8 & \verb|ECON150__PO| & 4, 8 & \verb|ECON151__CM| & 2, 4, 6, 8\\\hline
    \verb|ECON152__CM| & 1, 3, 5, 7 & \verb|ECON152__PO| & 1, 5 & \verb|ECON154__CM| & 2, 3, 6, 7, 1, 5, 4, 8\\\hline
    \verb|ECON154__HM| & 2, 4, 6, 8 & \verb|ECON154__PO| & 2, 3, 6, 7 & \verb|ECON155__CM| & 2, 4, 6, 8\\\hline
    \verb|ECON155__PO| & 2, 6 & \verb|ECON156__PO| & 1, 3, 5, 7, 4, 8 & \verb|ECON157__CM| & 2, 1, 6, 5, 4, 8\\\hline
    \verb|ECON157__PO| & 1, 3, 5, 7 & \verb|ECON158__CM| & 1, 3, 5, 7 & \verb|ECON159__PO| & 3, 7\\\hline
    \verb|ECON160__CM| & 2, 4, 6, 8 & \verb|ECON161__PO| & 2, 6 & \verb|ECON165__CM| & 2, 6\\\hline
    \verb|ECON165__PO| & 1, 5 & \verb|ECON165__PZ| & 2, 1, 6, 5, 4, 8 & \verb|ECON166__PO| & 4, 8\\\hline
    \verb|ECON167__PO| & 4, 8 & \verb|ECON168__CM| & 1, 5 & \verb|ECON170__SC| & 2, 1, 6, 5, 4, 8\\\hline
    \verb|ECON171__CM| & 2, 3, 6, 7, 1, 5 & \verb|ECON171__PZ| & 4, 8 & \verb|ECON172__CM| & 2, 4, 6, 8\\\hline
    \verb|ECON172__PZ| & 1, 5 & \verb|ECON175__PZ| & 2, 1, 6, 5, 4, 8 & \verb|ECON175__SC| & 2, 3, 6, 7\\\hline
    \verb|ECON176__CM| & 2, 4, 6, 8 & \verb|ECON177__CM| & 2, 4, 6, 8 & \verb|ECON179F_HM| & 2, 6\\\hline
    \verb|ECON179G_HM| & 1, 4, 5, 8 & \verb|ECON179__CM| & 3, 7 & \verb|ECON180__CM| & 2, 3, 6, 7, 1, 5, 4, 8\\\hline
    \verb|ECON180__PZ| & 1, 3, 5, 7, 4, 8 & \verb|ECON181__CM| & 2, 3, 6, 7, 1, 5, 4, 8 & \verb|ECON184__SC| & 2, 6\\\hline
    \verb|ECON185__CM| & 2, 4, 6, 8 & \verb|ECON186__CM| & 1, 3, 5, 7 & \verb|ECON190__PO| & 2, 4, 6, 8\\\hline
    \verb|ECON190__PO2| & 2, 4, 6, 8 & \verb|ECON191H_SC| & 4, 8 & \verb|ECON191H_SC2| & 4, 8\\\hline
    \verb|ECON191__CM| & 2, 4, 6, 8 & \verb|ECON191__CM2| & 2, 4, 6, 8 & \verb|ECON191__SC| & 1, 3, 5, 7, 4, 8\\\hline
    \verb|ECON191__SC2| & 1, 3, 5, 7, 4, 8 & \verb|ECON194A_CM| & 1, 3, 5, 7 & \verb|ECON194A_CM2| & 1, 3, 5, 7\\\hline
    \verb|ECON194B_CM| & 2, 4, 6, 8 & \verb|ECON194B_CM2| & 2, 4, 6, 8 & \verb|ECON195__PO| & 2, 4, 6, 8\\\hline
    \verb|ECON195__PO2| & 2, 4, 6, 8 & \verb|ECON196__CM| & 2, 4, 6, 8 & \verb|ECON196__CM2| & 2, 4, 6, 8\\\hline
    \verb|ECON197SBCM| & 2, 6 & \verb|ECON197SBCM2| & 2, 6 & \verb|ECON197S_CM| & 2, 3, 6, 7, 4, 8\\\hline
    \verb|ECON197S_CM2| & 2, 3, 6, 7, 4, 8 & \verb|ECON198__PZ| & 1, 3, 5, 7 & \verb|ECON198__PZ2| & 1, 3, 5, 7\\\hline
    \verb|ECON199__CM| & 4, 8 & \verb|ECON199__CM2| & 4, 8 & \verb|ECON199__PZ| & 4, 8\\\hline
    \verb|ECON199__PZ2| & 4, 8 & \verb|ECON316__CG| & 1, 5 & \verb|EDUC301G_CG| & 4, 8\\\hline
    \verb|EDUC424__CG| & 4, 8 & \verb|EEP_191__SC| & 8, 3, 7, 4 & \verb|EEP_191__SC2| & 8, 3, 7, 4\\\hline
    \verb|ENGL001__PZ| & 1, 3, 5, 7 & \verb|ENGL002__PZ| & 3, 7 & \verb|ENGL010B_PZ| & 4, 8\\\hline
    \verb|ENGL010__PO| & 3, 7 & \verb|ENGL011A_PZ| & 3, 7 & \verb|ENGL012B_AF| & 1, 3, 5, 7\\\hline
    \verb|ENGL013__PZ| & 2, 1, 6, 5 & \verb|ENGL015__PO| & 1, 5 & \verb|ENGL016__PZ| & 1, 5\\\hline
    \verb|ENGL017__PO| & 4, 8 & \verb|ENGL018__PO| & 1, 5 & \verb|ENGL019__PO| & 2, 3, 6, 7, 1, 5\\\hline
    \verb|ENGL020__PO| & 2, 6 & \verb|ENGL023__PO| & 2, 4, 6, 8 & \verb|ENGL030__PZ| & 2, 3, 6, 7, 1, 5, 4, 8\\\hline
    \verb|ENGL031__PZ| & 1, 5 & \verb|ENGL036__PZ| & 1, 4, 5, 8 & \verb|ENGL037__PZ| & 1, 4, 5, 8\\\hline
    \verb|ENGL040__PO| & 3, 7 & \verb|ENGL041__PO| & 3, 7 & \verb|ENGL043__PO| & 4, 8\\\hline
    \verb|ENGL044__PO| & 2, 6 & \verb|ENGL045__PO| & 4, 8 & \verb|ENGL046__PO| & 1, 5\\\hline
    \verb|ENGL046__PZ| & 3, 7 & \verb|ENGL047__PZ| & 2, 6 & \verb|ENGL048__PO| & 4, 8\\\hline
    \verb|ENGL048__PZ| & 2, 1, 6, 5 & \verb|ENGL049__PZ| & 4, 8 & \verb|ENGL055A_PO| & 1, 5\\\hline
    \verb|ENGL060__PZ| & 2, 6 & \verb|ENGL061__PZ| & 1, 5 & \verb|ENGL064A_PO| & 2, 4, 6, 8\\\hline
    \verb|ENGL064B_PO| & 2, 4, 6, 8 & \verb|ENGL064C_PO| & 2, 4, 6, 8 & \verb|ENGL072__PO| & 2, 6\\\hline
    \verb|ENGL074__PO| & 3, 7 & \verb|ENGL080__PO| & 1, 4, 5, 8 & \verb|ENGL084__PO| & 2, 6\\\hline
    \verb|ENGL086__PO| & 2, 6 & \verb|ENGL087F_PO| & 2, 3, 6, 7, 1, 5 & \verb|ENGL087H_PO| & 2, 1, 6, 5\\\hline
    \verb|ENGL093__PO| & 3, 7 & \verb|ENGL095__PO| & 4, 8 & \verb|ENGL099__PO| & 2, 6\\\hline
    \verb|ENGL101__PZ| & 1, 5 & \verb|ENGL101__SC| & 1, 3, 5, 7 & \verb|ENGL102__PZ| & 3, 7\\\hline
    \verb|ENGL102__SC| & 2, 4, 6, 8 & \verb|ENGL110__PZ| & 2, 6 & \verb|ENGL111__SC| & 4, 8\\\hline
    \verb|ENGL113S_SC| & 1, 5 & \verb|ENGL115__SC| & 2, 6 & \verb|ENGL116__SC| & 2, 6\\\hline
    \verb|ENGL118S_SC| & 3, 7 & \verb|ENGL120__PZ| & 4, 8 & \verb|ENGL121__SC| & 3, 7\\\hline
    \verb|ENGL124__SC| & 1, 3, 5, 7 & \verb|ENGL125C_AF| & 2, 6 & \verb|ENGL125__SC| & 1, 5\\\hline
    \verb|ENGL128__PZ| & 3, 7 & \verb|ENGL130__PZ| & 2, 4, 6, 8 & \verb|ENGL131__SC| & 2, 6\\\hline
    \verb|ENGL133S_SC| & 3, 7 & \verb|ENGL141__PZ| & 2, 4, 6, 8 & \verb|ENGL141__SC| & 3, 7\\\hline
    \verb|ENGL145__SC| & 3, 7 & \verb|ENGL146__PO| & 4, 8 & \verb|ENGL148__SC| & 1, 5\\\hline
    \verb|ENGL149S_SC| & 4, 8 & \verb|ENGL151__PO| & 1, 5 & \verb|ENGL151__PZ| & 1, 5\\\hline
    \verb|ENGL153__PO| & 4, 8 & \verb|ENGL154__SC| & 2, 6 & \verb|ENGL157__PO| & 4, 8\\\hline
    \verb|ENGL160S_SC| & 1, 5 & \verb|ENGL161__PO| & 1, 5 & \verb|ENGL162__PZ| & 1, 5\\\hline
    \verb|ENGL163__PO| & 1, 5 & \verb|ENGL164S_SC| & 4, 8 & \verb|ENGL165__PO| & 1, 3, 5, 7\\\hline
    \verb|ENGL165__SC| & 2, 6 & \verb|ENGL168S_SC| & 2, 3, 6, 7 & \verb|ENGL169__PO| & 2, 1, 6, 5\\\hline
    \verb|ENGL170A_PO| & 4, 8 & \verb|ENGL170H_PO| & 4, 8 & \verb|ENGL170J_PO| & 1, 5\\\hline
    \verb|ENGL170P_PO| & 2, 6 & \verb|ENGL170R_PO| & 1, 5 & \verb|ENGL170U_PO| & 3, 7\\\hline
    \verb|ENGL170X_PO| & 2, 6 & \verb|ENGL174__SC| & 3, 7 & \verb|ENGL175__PZ| & 4, 8\\\hline
    \verb|ENGL177__SC| & 4, 8 & \verb|ENGL180__SC| & 2, 4, 6, 8 & \verb|ENGL181__PZ| & 3, 7\\\hline
    \verb|ENGL181__SC| & 4, 8 & \verb|ENGL183A_PO| & 1, 3, 5, 7 & \verb|ENGL183B_PO| & 1, 3, 5, 7\\\hline
    \verb|ENGL188__PO| & 3, 7 & \verb|ENGL190__SC| & 1, 3, 5, 7 & \verb|ENGL190__SC2| & 1, 3, 5, 7\\\hline
    \verb|ENGL191__PO| & 2, 3, 6, 7, 1, 5, 4, 8 & \verb|ENGL191__PO2| & 2, 3, 6, 7, 1, 5, 4, 8 & \verb|ENGL191__SC| & 1, 3, 5, 7\\\hline
    \verb|ENGL191__SC2| & 1, 3, 5, 7 & \verb|ENGL192__SC| & 4, 8 & \verb|ENGL192__SC2| & 4, 8\\\hline
    \verb|ENGL193__SC| & 2, 6 & \verb|ENGL193__SC2| & 2, 6 & \verb|ENGL194__PZ| & 3, 7\\\hline
    \verb|ENGL194__PZ2| & 3, 7 & \verb|ENGL194__SC| & 1, 5 & \verb|ENGL194__SC2| & 1, 5\\\hline
    \verb|ENGL195B_PO| & 2, 4, 6, 8 & \verb|ENGL195B_PO2| & 2, 4, 6, 8 & \verb|ENGL195__PO| & 2, 3, 6, 7, 1, 5, 4, 8\\\hline
    \verb|ENGL195__PO2| & 2, 3, 6, 7, 1, 5, 4, 8 & \verb|ENGL195__SC| & 3, 7 & \verb|ENGL195__SC2| & 3, 7\\\hline
    \verb|ENGL197N_SC| & 2, 6 & \verb|ENGL197N_SC2| & 2, 6 & \verb|ENGL199T_SC| & 1, 3, 5, 7, 4, 8\\\hline
    \verb|ENGL199T_SC2| & 1, 3, 5, 7, 4, 8 & \verb|ENGR004L_HM| & 2, 3, 6, 7, 1, 5, 4, 8 & \verb|ENGR004__HM| & 1, 2, 3, 4\\\hline
    \verb|ENGR015__HM| & 1, 4, 5, 8 & \verb|ENGR025__HM| & 1, 3, 5, 7 & \verb|ENGR026__HM| & 1, 3, 5, 7\\\hline
    \verb|ENGR072__HM| & 2, 4, 6, 8 & \verb|ENGR079P_HM| & 1, 3, 5, 7 & \verb|ENGR079__HM| & 3\\\hline
    \verb|ENGR080__HM| & 2, 4, 6, 8 & \verb|ENGR082__HM| & 2, 3, 6, 7, 1, 5, 4, 8 & \verb|ENGR083__HM| & 2, 3, 6, 7, 1, 5, 4, 8\\\hline
    \verb|ENGR084__HM| & 2, 3, 6, 7, 1, 5, 4, 8 & \verb|ENGR085A_HM| & 2, 3, 6, 7, 1, 5, 4, 8 & \verb|ENGR085__HM| & 2, 3, 6, 7, 1, 5, 4, 8\\\hline
    \verb|ENGR086__HM| & 2, 3, 6, 7, 1, 5, 4, 8 & \verb|ENGR091__HM| & 2, 3, 6, 7, 1, 5, 4, 8 & \verb|ENGR101__HM| & 1, 3, 5, 7\\\hline
    \verb|ENGR102__HM| & 2, 4, 6, 8 & \verb|ENGR111__HM| & 5, 6 & \verb|ENGR112__HM| & 7\\\hline
    \verb|ENGR113__HM| & 8 & \verb|ENGR114__HM| & 2, 4, 6, 8 & \verb|ENGR122__HM| & 5\\\hline
    \verb|ENGR124__HM| & 6 & \verb|ENGR131__HM| & 1, 3, 5, 7 & \verb|ENGR132__HM| & 2, 6\\\hline
    \verb|ENGR133__HM| & 4, 8 & \verb|ENGR134__HM| & 3, 7 & \verb|ENGR138__HM| & 3, 7\\\hline
    \verb|ENGR151__HM| & 2, 4, 6, 8 & \verb|ENGR154__HM| & 4, 8 & \verb|ENGR155__HM| & 1, 3, 5, 7\\\hline
    \verb|ENGR157__HM| & 3, 7 & \verb|ENGR164__HM| & 4, 8 & \verb|ENGR171__HM| & 3, 7\\\hline
    \verb|ENGR175__HM| & 1, 3, 5, 7 & \verb|ENGR177__HM| & 4, 8 & \verb|ENGR178__HM| & 2, 6\\\hline
    \verb|ENGR180__HM| & 2, 3, 6, 7, 1, 5, 4, 8 & \verb|ENGR181__HM| & 2, 4, 6, 8 & \verb|ENGR183__HM| & 1, 3, 5, 7\\\hline
    \verb|ENGR185A_HM| & 2, 3, 6, 7, 4, 8 & \verb|ENGR185B_HM| & 2, 4, 6, 8 & \verb|ENGR186__HM| & 4, 8\\\hline
    \verb|ENGR187__HM| & 1, 3, 5, 7 & \verb|ENGR190BAHM| & 2, 6 & \verb|ENGR190BAHM2| & 2, 6\\\hline
    \verb|ENGR190BFHM| & 2, 4, 6, 8 & \verb|ENGR190BFHM2| & 2, 4, 6, 8 & \verb|ENGR190BGHM| & 2, 3, 6, 7, 1, 5\\\hline
    \verb|ENGR190BGHM2| & 2, 3, 6, 7, 1, 5 & \verb|ENGR190BHHM| & 2, 6 & \verb|ENGR190BHHM2| & 2, 6\\\hline
    \verb|ENGR190BIHM| & 1, 3, 5, 7 & \verb|ENGR190BIHM2| & 1, 3, 5, 7 & \verb|ENGR190BJHM| & 1, 5\\\hline
    \verb|ENGR190BJHM2| & 1, 5 & \verb|ENGR191__HM| & 2, 3, 6, 7, 1, 5, 4, 8 & \verb|ENGR191__HM2| & 2, 3, 6, 7, 1, 5, 4, 8\\\hline
    \verb|ENGR205__HM| & 4, 8 & \verb|ENGR207__HM| & 2, 4, 6, 8 & \verb|ENGR208__HM| & 3, 7\\\hline
    \verb|ENGR278__HM| & 2, 6 & \verb|ENGR900__HM| & 1, 2, 5, 3, 4, 6, 7, 8 & \verb|ENGR901__HM| & 1, 2, 5, 3, 4, 6, 7, 8\\\hline
    \verb|ENGR902__HM| & 1, 2, 5, 3, 4, 6, 7, 8 & \verb|ENGR903__HM| & 1, 2, 5, 3, 4, 6, 7, 8 & \verb|ENGR904__HM| & 1, 2, 5, 3, 4, 6, 7, 8\\\hline
    \verb|ENG_5300_KG| & 4, 8 & \verb|ENTR179A_HM| & 1, 3, 5, 7 & \verb|ENTR179B_HM| & 4, 8\\\hline
    \verb|FGSS026__SC| & 1, 3, 5, 7, 4, 8 & \verb|FGSS036__SC| & 1, 3, 5, 7 & \verb|FGSS177__SC| & 2, 3, 6, 7\\\hline
    \verb|FGSS183__SC| & 2, 4, 6, 8 & \verb|FGSS184__SC| & 1, 3, 5, 7 & \verb|FGSS186__SC| & 2, 6\\\hline
    \verb|FGSS188E_SC| & 2, 4, 6, 8 & \verb|FGSS188__SC| & 3, 7 & \verb|FGSS191__SC| & 4, 8\\\hline
    \verb|FGSS191__SC2| & 4, 8 & \verb|FGSS192__SC| & 1, 4, 5, 8 & \verb|FGSS192__SC2| & 1, 4, 5, 8\\\hline
    \verb|FGSS198__SC| & 4, 8 & \verb|FGSS198__SC2| & 4, 8 & \verb|FHS_010__CM| & 2, 3, 6, 7, 1, 5, 4, 8\\\hline
    \verb|FLAN191__SC| & 1, 3, 5, 7, 4, 8 & \verb|FLAN191__SC2| & 1, 3, 5, 7, 4, 8 & \verb|FREN001L_CM| & 1, 4, 5, 8\\\hline
    \verb|FREN001L_SC| & 8, 3, 7, 4 & \verb|FREN001__PO| & 2, 3, 6, 7, 1, 5, 4, 8 & \verb|FREN001__PZ| & 4, 8\\\hline
    \verb|FREN001__SC| & 2, 3, 6, 7, 1, 5 & \verb|FREN002L_CM| & 2, 1, 6, 5, 4, 8 & \verb|FREN002L_SC| & 1, 3, 5, 7, 4, 8\\\hline
    \verb|FREN002__PO| & 2, 4, 6, 8 & \verb|FREN002__PZ| & 2, 4, 6, 8 & \verb|FREN002__SC| & 2, 3, 6, 7, 1, 5\\\hline
    \verb|FREN011__PO| & 1, 3, 5, 7, 4, 8 & \verb|FREN013__PO| & 1, 3, 5, 7, 4, 8 & \verb|FREN022__PO| & 1, 3, 5, 7\\\hline
    \verb|FREN033L_CM| & 1, 4, 5, 8 & \verb|FREN033L_SC| & 1, 3, 5, 7, 4, 8 & \verb|FREN033__CM| & 1, 5\\\hline
    \verb|FREN033__PO| & 2, 3, 6, 7, 1, 5, 4, 8 & \verb|FREN033__SC| & 2, 3, 6, 7, 1, 5, 4, 8 & \verb|FREN044L_CM| & 1, 4, 5, 8\\\hline
    \verb|FREN044L_SC| & 1, 4, 5, 8 & \verb|FREN044__CM| & 1, 4, 5, 8 & \verb|FREN044__PO| & 2, 3, 6, 7, 1, 5, 4, 8\\\hline
    \verb|FREN044__SC| & 2, 3, 6, 7, 1, 5, 4, 8 & \verb|FREN045L_SC| & 3, 7 & \verb|FREN045__SC| & 3, 7\\\hline
    \verb|FREN100__CM| & 2, 3, 6, 7 & \verb|FREN100__PZ| & 1, 5 & \verb|FREN101A_PO| & 2, 3, 6, 7\\\hline
    \verb|FREN101B_PO| & 1, 4, 5, 8 & \verb|FREN103__PO| & 4, 8 & \verb|FREN104__PO| & 3, 7\\\hline
    \verb|FREN105__PO| & 3, 7 & \verb|FREN107__PO| & 2, 6 & \verb|FREN110__PO| & 1, 5\\\hline
    \verb|FREN112__PO| & 3, 7 & \verb|FREN113__SC| & 4, 8 & \verb|FREN115__CM| & 3, 7\\\hline
    \verb|FREN125__SC| & 3, 7 & \verb|FREN126__SC| & 4, 8 & \verb|FREN127__SC| & 1, 5\\\hline
    \verb|FREN129__PO| & 1, 5 & \verb|FREN133__CM| & 1, 3, 5, 7 & \verb|FREN135__CM| & 2, 6\\\hline
    \verb|FREN150C_PO| & 2, 6 & \verb|FREN153__SC| & 2, 6 & \verb|FREN161__PZ| & 3, 7\\\hline
    \verb|FREN179__SC| & 1, 5 & \verb|FREN180__PO| & 4, 8 & \verb|FREN181__SC| & 2, 6\\\hline
    \verb|FREN185__PO| & 2, 6 & \verb|FREN187A_SC| & 4, 8 & \verb|FREN188__PO| & 1, 5\\\hline
    \verb|FREN191__PO| & 1, 3, 5, 7, 4, 8 & \verb|FREN191__PO2| & 1, 3, 5, 7, 4, 8 & \verb|FREN191__SC| & 4, 8\\\hline
    \verb|FREN191__SC2| & 4, 8 & \verb|FREN192__PO| & 2, 3, 6, 7, 1, 5, 4, 8 & \verb|FREN192__PO2| & 2, 3, 6, 7, 1, 5, 4, 8\\\hline
    \verb|FREN193__PO| & 2, 3, 6, 7, 1, 5, 4, 8 & \verb|FREN193__PO2| & 2, 3, 6, 7, 1, 5, 4, 8 & \verb|FWS_010__CM| & 2, 3, 6, 7, 1, 5, 4, 8\\\hline
    \verb|GEOG145__HM| & 1, 5 & \verb|GEOG179H_HM| & 4, 8 & \verb|GEOG179I_HM| & 2, 6\\\hline
    \verb|GEOG179J_HM| & 2, 6 & \verb|GEOG179K_HM| & 3, 7 & \verb|GEOL015__PO| & 4, 8\\\hline
    \verb|GEOL020A_PO| & 2, 4, 6, 8 & \verb|GEOL020C_PO| & 2, 4, 6, 8 & \verb|GEOL020E_PO| & 1, 3, 5, 7\\\hline
    \verb|GEOL025__PO| & 4, 8 & \verb|GEOL115__PO| & 2, 6 & \verb|GEOL125__PO| & 2, 4, 6, 8\\\hline
    \verb|GEOL127__PO| & 1, 3, 5, 7 & \verb|GEOL129__PO| & 2, 4, 6, 8 & \verb|GEOL131__PO| & 3, 7\\\hline
    \verb|GEOL133__PO| & 2, 6 & \verb|GEOL143__PO| & 3, 7 & \verb|GEOL181__PO| & 2, 6\\\hline
    \verb|GEOL183__PO| & 1, 3, 5, 7 & \verb|GEOL185__PO| & 2, 4, 6, 8 & \verb|GEOL189C_PO| & 4, 8\\\hline
    \verb|GEOL189G_PO| & 4, 8 & \verb|GEOL189I_PO| & 1, 5 & \verb|GEOL192__PO| & 2, 3, 6, 7, 1, 5, 4, 8\\\hline
    \verb|GEOL192__PO2| & 2, 3, 6, 7, 1, 5, 4, 8 & \verb|GERM001__PO| & 1, 3, 5, 7 & \verb|GERM001__SC| & 1, 3, 5, 7\\\hline
    \verb|GERM002__PO| & 2, 4, 6, 8 & \verb|GERM002__SC| & 2, 4, 6, 8 & \verb|GERM011__PO| & 1, 3, 5, 7, 4, 8\\\hline
    \verb|GERM012__SC| & 1, 5 & \verb|GERM013__PO| & 1, 3, 5, 7, 4, 8 & \verb|GERM033__PO| & 1, 3, 5, 7\\\hline
    \verb|GERM033__SC| & 1, 3, 5, 7 & \verb|GERM044__SC| & 2, 4, 6, 8 & \verb|GERM102__PO| & 4, 8\\\hline
    \verb|GERM107__SC| & 3, 7 & \verb|GERM112__SC| & 1, 4, 5, 8 & \verb|GERM152__PO| & 3, 7\\\hline
    \verb|GERM154F_PO| & 2, 1, 6, 5 & \verb|GERM191__PO| & 2, 3, 6, 7, 1, 5, 4, 8 & \verb|GERM191__PO2| & 2, 3, 6, 7, 1, 5, 4, 8\\\hline
    \verb|GERM191__SC| & 4, 8 & \verb|GERM191__SC2| & 4, 8 & \verb|GERM193__PO| & 3, 7\\\hline
    \verb|GERM193__PO2| & 3, 7 & \verb|GFS_036__PZ| & 1, 5 & \verb|GFS_170__PZ| & 1, 5\\\hline
    \verb|GLAS001__PZ| & 8, 3, 7, 4 & \verb|GLAS180IOPZ| & 4, 8 & \verb|GLAS194A_PZ| & 1, 3, 5, 7\\\hline
    \verb|GLAS194A_PZ2| & 1, 3, 5, 7 & \verb|GLAS194B_PZ| & 1, 3, 5, 7 & \verb|GLAS194B_PZ2| & 1, 3, 5, 7\\\hline
    \verb|GOVT020H_CM| & 1, 3, 5, 7 & \verb|GOVT020__CM| & 2, 3, 6, 7, 1, 5, 4, 8 & \verb|GOVT041B_CM| & 2, 3, 6, 7, 1, 5, 4, 8\\\hline
    \verb|GOVT041__CM| & 2, 3, 6, 7, 1, 5, 4, 8 & \verb|GOVT054__CM| & 8, 3, 7, 4 & \verb|GOVT055__CM| & 1, 3, 5, 7, 4, 8\\\hline
    \verb|GOVT060__CM| & 2, 3, 6, 7, 1, 5, 4, 8 & \verb|GOVT070H_CM| & 2, 4, 6, 8 & \verb|GOVT070__CM| & 2, 3, 6, 7, 1, 5, 4, 8\\\hline
    \verb|GOVT071__CM| & 2, 3, 6, 7, 4, 8 & \verb|GOVT080__CM| & 2, 3, 6, 7, 1, 5, 4, 8 & \verb|GOVT094__CM| & 4, 8\\\hline
    \verb|GOVT097__CM| & 1, 3, 5, 7 & \verb|GOVT100__CM| & 2, 3, 6, 7, 1, 5, 4, 8 & \verb|GOVT101__CM| & 2, 4, 6, 8\\\hline
    \verb|GOVT102__CM| & 2, 4, 6, 8 & \verb|GOVT103__CM| & 4, 8 & \verb|GOVT107__CM| & 1, 5\\\hline
    \verb|GOVT110__CM| & 2, 3, 6, 7, 1, 5 & \verb|GOVT112A_CM| & 1, 3, 5, 7 & \verb|GOVT112B_CM| & 2, 3, 6, 7, 1, 5\\\hline
    \verb|GOVT115__CM| & 3, 7 & \verb|GOVT116__CM| & 1, 5 & \verb|GOVT117__CM| & 2, 6\\\hline
    \verb|GOVT118__CM| & 1, 3, 5, 7 & \verb|GOVT121__CM| & 1, 5 & \verb|GOVT123__CM| & 2, 6\\\hline
    \verb|GOVT128__JT| & 2, 6 & \verb|GOVT129A_CM| & 2, 6 & \verb|GOVT129E_CM| & 4, 8\\\hline
    \verb|GOVT129__CM| & 1, 4, 5, 8 & \verb|GOVT131__CM| & 2, 4, 6, 8 & \verb|GOVT132E_CM| & 2, 4, 6, 8\\\hline
    \verb|GOVT134__CM| & 3, 7 & \verb|GOVT136C_CM| & 1, 5 & \verb|GOVT137A_CM| & 2, 4, 6, 8\\\hline
    \verb|GOVT137B_CM| & 2, 3, 6, 7 & \verb|GOVT137__CM| & 2, 3, 6, 7, 1, 5, 4, 8 & \verb|GOVT138C_CM| & 2, 4, 6, 8\\\hline
    \verb|GOVT139__CM| & 1, 3, 5, 7 & \verb|GOVT140__CM| & 2, 3, 6, 7, 1, 5 & \verb|GOVT141__CM| & 1, 4, 5, 8\\\hline
    \verb|GOVT142C_CM| & 4, 8 & \verb|GOVT142E_CM| & 2, 3, 6, 7, 1, 5 & \verb|GOVT142__CM| & 2, 6\\\hline
    \verb|GOVT143C_CM| & 2, 4, 6, 8 & \verb|GOVT144D_CM| & 1, 3, 5, 7 & \verb|GOVT145E_CM| & 3, 7\\\hline
    \verb|GOVT146A_CM| & 2, 3, 6, 7, 1, 5 & \verb|GOVT146__CM| & 4, 8 & \verb|GOVT147__CM| & 2, 1, 6, 5, 4, 8\\\hline
    \verb|GOVT148__CM| & 2, 3, 6, 7 & \verb|GOVT149__CM| & 2, 3, 6, 7, 1, 5, 4, 8 & \verb|GOVT150C_CM| & 3, 7\\\hline
    \verb|GOVT151__CM| & 1, 5 & \verb|GOVT152C_CM| & 4, 8 & \verb|GOVT152__CM| & 1, 3, 5, 7\\\hline
    \verb|GOVT154E_CM| & 4, 8 & \verb|GOVT156E_CM| & 2, 1, 6, 5 & \verb|GOVT157S_CM| & 1, 5\\\hline
    \verb|GOVT164__CM| & 2, 6 & \verb|GOVT166__CM| & 2, 3, 6, 7 & \verb|GOVT173C_CM| & 1, 3, 5, 7\\\hline
    \verb|GOVT174C_CM| & 3, 7 & \verb|GOVT177C_CM| & 3, 7 & \verb|GOVT177__CM| & 2, 1, 6, 5\\\hline
    \verb|GOVT178__CM| & 2, 4, 6, 8 & \verb|GOVT179C_CM| & 4, 8 & \verb|GOVT179__CM| & 3, 7\\\hline
    \verb|GOVT182__CM| & 4, 8 & \verb|GOVT191__CM| & 1, 5 & \verb|GOVT191__CM2| & 1, 5\\\hline
    \verb|GOVT192__CM| & 1, 5 & \verb|GOVT192__CM2| & 1, 5 & \verb|GOVT999__CM| & 2, 6\\\hline
    \verb|GREK022__PO| & 2, 3, 6, 7, 4, 8 & \verb|GREK022__SC| & 1, 5 & \verb|GREK033__PO| & 3, 7\\\hline
    \verb|GREK033__SC| & 2, 1, 6, 5, 4, 8 & \verb|GREK044__PO| & 3, 7 & \verb|GREK044__SC| & 2, 1, 6, 5, 4, 8\\\hline
    \verb|GRMT014__PO| & 4, 8 & \verb|GRMT114__SC| & 4, 8 & \verb|GRMT115__SC| & 1, 5\\\hline
    \verb|GRMT131__PO| & 1, 5 & \verb|GRMT164__PO| & 4, 8 & \verb|GWS_026__PO| & 2, 3, 6, 7, 1, 5, 4, 8\\\hline
    \verb|GWS_036__PO| & 4, 8 & \verb|GWS_070__PO| & 2, 6 & \verb|GWS_081__PO| & 2, 4, 6, 8\\\hline
    \verb|GWS_140__PO| & 4, 8 & \verb|GWS_162__PO| & 1, 4, 5, 8 & \verb|GWS_166__PO| & 2, 3, 6, 7, 4, 8\\\hline
    \verb|GWS_170__PO| & 3, 7 & \verb|GWS_172__PO| & 3, 7 & \verb|GWS_173__PO| & 2, 1, 6, 5\\\hline
    \verb|GWS_179A_HM| & 4, 8 & \verb|GWS_179B_HM| & 1, 5 & \verb|GWS_180__PO| & 1, 3, 5, 7\\\hline
    \verb|GWS_183__PO| & 2, 4, 6, 8 & \verb|GWS_184__PO| & 1, 3, 5, 7 & \verb|GWS_190__PO| & 1, 3, 5, 7\\\hline
    \verb|GWS_190__PO2| & 1, 3, 5, 7 & \verb|GWS_191__PO| & 2, 4, 6, 8 & \verb|GWS_191__PO2| & 2, 4, 6, 8\\\hline
    \verb|HIST008__PO| & 1, 3, 5, 7 & \verb|HIST009__SC| & 1, 5 & \verb|HIST011__PO| & 2, 4, 6, 8\\\hline
    \verb|HIST011__PZ| & 2, 3, 6, 7, 1, 5 & \verb|HIST012__PZ| & 1, 3, 5, 7 & \verb|HIST014__PO| & 3, 7\\\hline
    \verb|HIST017__CH| & 1, 3, 5, 7 & \verb|HIST020__PO| & 1, 3, 5, 7 & \verb|HIST021__PO| & 2, 4, 6, 8\\\hline
    \verb|HIST025__CH| & 4, 8 & \verb|HIST025__PZ| & 4, 8 & \verb|HIST031__CH| & 3, 7\\\hline
    \verb|HIST032__CH| & 2, 4, 6, 8 & \verb|HIST034__CH| & 1, 3, 5, 7 & \verb|HIST036__PO| & 2, 6\\\hline
    \verb|HIST037__PO| & 4, 8 & \verb|HIST040__AF| & 3, 7 & \verb|HIST041__AF| & 2, 4, 6, 8\\\hline
    \verb|HIST043__PO| & 4, 8 & \verb|HIST044__PZ| & 4, 8 & \verb|HIST046__PZ| & 2, 4, 6, 8\\\hline
    \verb|HIST047__PO| & 2, 6 & \verb|HIST048__SC| & 2, 1, 6, 5 & \verb|HIST054__CM| & 3, 4, 8, 7\\\hline
    \verb|HIST055__CM| & 1, 3, 5, 7 & \verb|HIST056__CM| & 2, 4, 6, 8 & \verb|HIST059__CM| & 1, 5\\\hline
    \verb|HIST060__PO| & 1, 3, 5, 7 & \verb|HIST060__PZ| & 3, 7 & \verb|HIST061__PO| & 4, 8\\\hline
    \verb|HIST062__PO| & 2, 6 & \verb|HIST064__PO| & 4, 8 & \verb|HIST065__PZ| & 2, 1, 6, 5, 4, 8\\\hline
    \verb|HIST066__PZ| & 2, 6 & \verb|HIST069__PZ| & 2, 6 & \verb|HIST070B_SC| & 1, 5\\\hline
    \verb|HIST070__PO| & 1, 3, 5, 7 & \verb|HIST071__PO| & 2, 4, 6, 8 & \verb|HIST072__CM| & 3, 7\\\hline
    \verb|HIST072__PO| & 4, 8 & \verb|HIST081__HM| & 2, 6 & \verb|HIST082__HM| & 3, 7\\\hline
    \verb|HIST083__PZ| & 3, 7 & \verb|HIST090__CM| & 4, 8 & \verb|HIST096__CM| & 2, 1, 6, 5\\\hline
    \verb|HIST097__CM| & 3, 7 & \verb|HIST099__CM| & 4, 8 & \verb|HIST100__CM| & 3, 7\\\hline
    \verb|HIST101ADPO| & 2, 6 & \verb|HIST101A_PO| & 2, 6 & \verb|HIST101B_PO| & 4, 8\\\hline
    \verb|HIST101D_PO| & 2, 6 & \verb|HIST101F_PO| & 1, 3, 5, 7 & \verb|HIST101HAPO| & 2, 6\\\hline
    \verb|HIST101H_PO| & 2, 6 & \verb|HIST101N_PO| & 1, 5 & \verb|HIST101Q_PO| & 2, 6\\\hline
    \verb|HIST101S_CH| & 1, 5 & \verb|HIST101T_CH| & 4, 8 & \verb|HIST101U_PO| & 1, 5\\\hline
    \verb|HIST101Y_PO| & 1, 4, 5, 8 & \verb|HIST101__CM| & 2, 6 & \verb|HIST102__PO| & 2, 3, 6, 7, 1, 5\\\hline
    \verb|HIST105__CM| & 2, 6 & \verb|HIST105__PO| & 4, 8 & \verb|HIST106__CM| & 1, 5\\\hline
    \verb|HIST108__CM| & 1, 5 & \verb|HIST108__SC| & 4, 8 & \verb|HIST109__CM| & 1, 5\\\hline
    \verb|HIST111__SC| & 3, 7 & \verb|HIST112A_CM| & 3, 7 & \verb|HIST112__CM| & 1, 5\\\hline
    \verb|HIST113__CM| & 4, 8 & \verb|HIST113__PO| & 1, 5 & \verb|HIST113__SC| & 1, 5\\\hline
    \verb|HIST114C_CM| & 4, 8 & \verb|HIST114__CM| & 4, 8 & \verb|HIST115__CM| & 3, 7\\\hline
    \verb|HIST116IOCM| & 2, 6 & \verb|HIST116__CM| & 4, 8 & \verb|HIST117__SC| & 2, 6\\\hline
    \verb|HIST118__PO| & 2, 6 & \verb|HIST119__CM| & 3, 7 & \verb|HIST120E_CM| & 2, 6\\\hline
    \verb|HIST120__CM| & 1, 5 & \verb|HIST121__CM| & 1, 5 & \verb|HIST122__CM| & 1, 5\\\hline
    \verb|HIST123__CM| & 3, 7 & \verb|HIST123__PO| & 4, 8 & \verb|HIST125__CM| & 4, 8\\\hline
    \verb|HIST125__PO| & 3, 7 & \verb|HIST126__CM| & 2, 1, 6, 5 & \verb|HIST128__CM| & 4, 8\\\hline
    \verb|HIST128__PO| & 3, 7 & \verb|HIST129__CM| & 1, 5 & \verb|HIST130__CH| & 2, 4, 6, 8\\\hline
    \verb|HIST132__CM| & 3, 7 & \verb|HIST133A_CM| & 1, 5 & \verb|HIST133B_CM| & 4, 8\\\hline
    \verb|HIST133__AF| & 3, 7 & \verb|HIST134__CM| & 2, 4, 6, 8 & \verb|HIST134__SC| & 1, 3, 5, 7\\\hline
    \verb|HIST135__CM| & 2, 6 & \verb|HIST136C_CM| & 1, 5 & \verb|HIST139E_CM| & 3, 7\\\hline
    \verb|HIST139__PO| & 3, 7 & \verb|HIST140C_CM| & 4, 8 & \verb|HIST140__AF| & 2, 4, 6, 8\\\hline
    \verb|HIST141__SC| & 3, 7 & \verb|HIST142__CM| & 2, 6 & \verb|HIST142__PZ| & 2, 6\\\hline
    \verb|HIST143__SC| & 2, 4, 6, 8 & \verb|HIST144__SC| & 3, 7 & \verb|HIST146__AF| & 3, 7\\\hline
    \verb|HIST146__PO| & 4, 8 & \verb|HIST146__SC| & 4, 8 & \verb|HIST147__CM| & 1, 5\\\hline
    \verb|HIST147__PO| & 1, 5 & \verb|HIST148__PO| & 2, 6 & \verb|HIST149A_SC| & 4, 8\\\hline
    \verb|HIST149__SC| & 2, 6 & \verb|HIST150C_CM| & 3, 7 & \verb|HIST150D_CM| & 4, 8\\\hline
    \verb|HIST150__CM| & 3, 4, 8, 7 & \verb|HIST150__PZ| & 2, 1, 6, 5 & \verb|HIST151__CM| & 1, 3, 5, 7\\\hline
    \verb|HIST151__HM| & 4, 8 & \verb|HIST151__SC| & 2, 6 & \verb|HIST154__CM| & 2, 4, 6, 8\\\hline
    \verb|HIST154__SC| & 4, 8 & \verb|HIST155__SC| & 1, 4, 5, 8 & \verb|HIST157__CM| & 4, 8\\\hline
    \verb|HIST157__PZ| & 2, 6 & \verb|HIST158__CM| & 3, 4, 8, 7 & \verb|HIST158__PZ| & 2, 6\\\hline
    \verb|HIST159__SC| & 4, 8 & \verb|HIST160__PO| & 4, 8 & \verb|HIST161__CM| & 2, 3, 6, 7\\\hline
    \verb|HIST164__SC| & 1, 5 & \verb|HIST165__PO| & 1, 5 & \verb|HIST166__CM| & 2, 1, 6, 5\\\hline
    \verb|HIST167__PO| & 2, 6 & \verb|HIST168__CM| & 3, 7 & \verb|HIST168__PO| & 2, 1, 6, 5\\\hline
    \verb|HIST169__CM| & 3, 7 & \verb|HIST169__PO| & 1, 5 & \verb|HIST170__PO| & 4, 8\\\hline
    \verb|HIST171__PO| & 2, 3, 6, 7, 1, 5 & \verb|HIST173__CM| & 1, 5 & \verb|HIST174__CM| & 2, 1, 6, 5\\\hline
    \verb|HIST174__PO| & 1, 5 & \verb|HIST175__CM| & 4, 8 & \verb|HIST175__PO| & 1, 3, 5, 7\\\hline
    \verb|HIST176__CM| & 3, 7 & \verb|HIST178__CM| & 4, 8 & \verb|HIST178__PZ| & 4, 8\\\hline
    \verb|HIST179L_HM| & 4, 8 & \verb|HIST179__CM| & 2, 3, 6, 7, 1, 5 & \verb|HIST180__PO| & 2, 6\\\hline
    \verb|HIST180__SC| & 2, 4, 6, 8 & \verb|HIST181__PZ| & 3, 7 & \verb|HIST182__CM| & 4, 8\\\hline
    \verb|HIST184__PO| & 2, 4, 6, 8 & \verb|HIST185__PO| & 3, 7 & \verb|HIST186__PO| & 8, 3, 7, 4\\\hline
    \verb|HIST188__CM| & 3, 7 & \verb|HIST189F_PO| & 1, 5 & \verb|HIST189G_PO| & 1, 5\\\hline
    \verb|HIST190__CM| & 2, 4, 6, 8 & \verb|HIST190__CM2| & 2, 4, 6, 8 & \verb|HIST190__PO| & 1, 3, 5, 7\\\hline
    \verb|HIST190__PO2| & 1, 3, 5, 7 & \verb|HIST191__PO| & 2, 1, 6, 5, 4, 8 & \verb|HIST191__PO2| & 2, 1, 6, 5, 4, 8\\\hline
    \verb|HIST191__SC| & 2, 3, 6, 7, 1, 5, 4, 8 & \verb|HIST191__SC2| & 2, 3, 6, 7, 1, 5, 4, 8 & \verb|HIST192__PO| & 2, 6\\\hline
    \verb|HIST192__PO2| & 2, 6 & \verb|HIST192__SC| & 4, 8 & \verb|HIST192__SC2| & 4, 8\\\hline
    \verb|HIST196__CM| & 3, 7 & \verb|HIST196__CM2| & 3, 7 & \verb|HIST197__CM| & 2, 6\\\hline
    \verb|HIST197__CM2| & 2, 6 & \verb|HIST197__PZ| & 1, 3, 5, 7 & \verb|HIST197__PZ2| & 1, 3, 5, 7\\\hline
    \verb|HIST199__PZ| & 2, 4, 6, 8 & \verb|HIST199__PZ2| & 2, 4, 6, 8 & \verb|HMSC133__SC| & 1, 5\\\hline
    \verb|HMSC138__SC| & 3, 7 & \verb|HMSC191__SC| & 1, 3, 5, 7, 4, 8 & \verb|HMSC191__SC2| & 1, 3, 5, 7, 4, 8\\\hline
    \verb|HMSC192__SC| & 4, 8 & \verb|HMSC192__SC2| & 4, 8 & \verb|HSA_000__5C| & 1, 2, 5, 3, 4, 6, 7, 8\\\hline
    \verb|HSA_000__HM| & 1, 2, 5, 3, 4, 6, 7, 8 & \verb|HSA_001__5C| & 1, 2, 5, 3, 4, 6, 7, 8 & \verb|HSA_001__HM| & 1, 2, 5, 3, 4, 6, 7, 8\\\hline
    \verb|HSA_002__5C| & 1, 2, 5, 3, 4, 6, 7, 8 & \verb|HSA_002__HM| & 1, 2, 5, 3, 4, 6, 7, 8 & \verb|HSA_003__5C| & 1, 2, 5, 3, 4, 6, 7, 8\\\hline
    \verb|HSA_003__HM| & 1, 2, 5, 3, 4, 6, 7, 8 & \verb|HSA_004__5C| & 1, 2, 5, 3, 4, 6, 7, 8 & \verb|HSA_005__5C| & 1, 2, 5, 3, 4, 6, 7, 8\\\hline
    \verb|HSA_006__5C| & 1, 2, 5, 3, 4, 6, 7, 8 & \verb|HSA_007__5C| & 1, 2, 5, 3, 4, 6, 7, 8 & \verb|HSA_008__5C| & 1, 2, 5, 3, 4, 6, 7, 8\\\hline
    \verb|HSA_009__5C| & 1, 2, 5, 3, 4, 6, 7, 8 & \verb|HSA_010__HM| & 2 & \verb|HSA_100__HM| & 1, 2, 5, 3, 4, 6, 7, 8\\\hline
    \verb|HUM_195J_SC| & 2, 4, 6, 8 & \verb|HUM_195J_SC2| & 2, 4, 6, 8 & \verb|HUM_196__PO| & 1, 3, 5, 7, 4, 8\\\hline
    \verb|HUM_196__PO2| & 1, 3, 5, 7, 4, 8 & \verb|ID__009__PO| & 2, 4, 6, 8 & \verb|ID__076__JT| & 2, 1, 6, 5, 4, 8\\\hline
    \verb|ID__077__JT| & 1, 5 & \verb|ID__080__CM| & 2, 3, 6, 7, 4, 8 & \verb|ID__081__CM| & 4, 8\\\hline
    \verb|ID__191D_SC| & 4, 8 & \verb|ID__191D_SC2| & 4, 8 & \verb|ID__191H_SC| & 4, 8\\\hline
    \verb|ID__191H_SC2| & 4, 8 & \verb|ID__191S_SC| & 4, 8 & \verb|ID__191S_SC2| & 4, 8\\\hline
    \verb|ID__199P1PO| & 1, 3, 5, 7 & \verb|ID__199P1PO2| & 1, 3, 5, 7 & \verb|ID__199P2PO| & 4, 8\\\hline
    \verb|ID__199P2PO2| & 4, 8 & \verb|ID__199P3PO| & 1, 3, 5, 7 & \verb|ID__199P3PO2| & 1, 3, 5, 7\\\hline
    \verb|ID__199P4PO| & 4, 8 & \verb|ID__199P4PO2| & 4, 8 & \verb|ID__199S1PO| & 3, 7\\\hline
    \verb|ID__199S1PO2| & 3, 7 & \verb|ID__199S2PO| & 8, 3, 7, 4 & \verb|ID__199S2PO2| & 8, 3, 7, 4\\\hline
    \verb|ID__199S3PO| & 1, 4, 5, 8 & \verb|ID__199S3PO2| & 1, 4, 5, 8 & \verb|ID__199S4PO| & 1, 5\\\hline
    \verb|ID__199S4PO2| & 1, 5 & \verb|IR__100__PO| & 2, 3, 6, 7, 1, 5, 4, 8 & \verb|IR__101__PO| & 1, 5\\\hline
    \verb|IR__118__PO| & 2, 6 & \verb|IR__190__PO| & 1, 3, 5, 7 & \verb|IR__190__PO2| & 1, 3, 5, 7\\\hline
    \verb|IR__191__PO| & 2, 4, 6, 8 & \verb|IR__191__PO2| & 2, 4, 6, 8 & \verb|IR__192__SC| & 4, 8\\\hline
    \verb|IR__192__SC2| & 4, 8 & \verb|IST_302__CG| & 1, 5 & \verb|ITAL001__SC| & 1, 3, 5, 7\\\hline
    \verb|ITAL002__SC| & 2, 4, 6, 8 & \verb|ITAL022__SC| & 2, 4, 6, 8 & \verb|ITAL033__SC| & 1, 3, 5, 7\\\hline
    \verb|ITAL044__SC| & 2, 4, 6, 8 & \verb|ITAL131__SC| & 2, 6 & \verb|ITAL134__SC| & 3, 7\\\hline
    \verb|ITAL140__SC| & 3, 7 & \verb|ITAL141__SC| & 1, 5 & \verb|ITAL150__SC| & 4, 8\\\hline
    \verb|ITAL187C_SC| & 2, 6 & \verb|ITAL191__SC| & 4, 8 & \verb|ITAL191__SC2| & 4, 8\\\hline
    \verb|JAPN001A_PO| & 1, 3, 5, 7 & \verb|JAPN001B_PO| & 2, 4, 6, 8 & \verb|JAPN011__PO| & 1, 3, 5, 7, 4, 8\\\hline
    \verb|JAPN012A_PO| & 1, 3, 5, 7 & \verb|JAPN012B_PO| & 2, 4, 6, 8 & \verb|JAPN013__PO| & 1, 3, 5, 7, 4, 8\\\hline
    \verb|JAPN014A_PO| & 1, 3, 5, 7 & \verb|JAPN014B_PO| & 2, 4, 6, 8 & \verb|JAPN051A_PO| & 1, 3, 5, 7\\\hline
    \verb|JAPN051B_PO| & 2, 4, 6, 8 & \verb|JAPN111A_PO| & 1, 3, 5, 7 & \verb|JAPN111B_PO| & 2, 4, 6, 8\\\hline
    \verb|JAPN125__PO| & 2, 6 & \verb|JAPN127__PO| & 4, 8 & \verb|JPNT171__PO| & 2, 6\\\hline
    \verb|JPNT175__PO| & 3, 7 & \verb|KORE001__CM| & 1, 3, 5, 7 & \verb|KORE002__CM| & 2, 4, 6, 8\\\hline
    \verb|KORE033__CM| & 1, 3, 5, 7 & \verb|KORE044__CM| & 2, 4, 6, 8 & \verb|KORE120__CM| & 3, 7\\\hline
    \verb|KRNT130__CM| & 2, 4, 6, 8 & \verb|LAMS191__PO| & 4, 8 & \verb|LAMS191__PO2| & 4, 8\\\hline
    \verb|LAST184__PO| & 3, 7 & \verb|LAST190__PO| & 2, 3, 6, 7, 4, 8 & \verb|LAST190__PO2| & 2, 3, 6, 7, 4, 8\\\hline
    \verb|LAST191__PO| & 2, 4, 6, 8 & \verb|LAST191__PO2| & 2, 4, 6, 8 & \verb|LAST191__SC| & 4, 8\\\hline
    \verb|LAST191__SC2| & 4, 8 & \verb|LAS_180__PO| & 4, 8 & \verb|LATN022__PO| & 2, 3, 6, 7, 1, 5, 4, 8\\\hline
    \verb|LATN033__PO| & 1, 4, 5, 8 & \verb|LATN033__SC| & 2, 3, 6, 7 & \verb|LATN044__PO| & 1, 3, 5, 7\\\hline
    \verb|LATN044__SC| & 2, 4, 6, 8 & \verb|LATN103__PO| & 2, 3, 6, 7, 1, 5, 4, 8 & \verb|LEAD010__CM| & 2, 3, 6, 7, 1, 5, 4, 8\\\hline
    \verb|LEAD041__CM| & 2, 1, 6, 5, 4, 8 & \verb|LEAD101__HM| & 2, 3, 6, 7, 1, 5, 4, 8 & \verb|LEAD134__CM| & 1, 4, 5, 8\\\hline
    \verb|LEAD143__CM| & 2, 6 & \verb|LEAD144__CM| & 2, 3, 6, 7, 1, 5, 4, 8 & \verb|LEAD149__CM| & 3, 7\\\hline
    \verb|LEAD151__HM| & 2, 3, 6, 7, 1, 5, 4, 8 & \verb|LGCS010__PO| & 2, 3, 6, 7, 1, 5, 4, 8 & \verb|LGCS010__PZ| & 1, 5\\\hline
    \verb|LGCS011__PO| & 2, 3, 6, 7, 1, 5, 4, 8 & \verb|LGCS101__PO| & 2, 4, 6, 8 & \verb|LGCS104__PO| & 4, 8\\\hline
    \verb|LGCS105__PO| & 1, 3, 5, 7 & \verb|LGCS106__PO| & 1, 5 & \verb|LGCS108__PO| & 1, 3, 5, 7\\\hline
    \verb|LGCS110__PZ| & 4, 8 & \verb|LGCS112__PZ| & 2, 3, 6, 7, 1, 5 & \verb|LGCS113__PO| & 3, 7\\\hline
    \verb|LGCS114__PO| & 2, 6 & \verb|LGCS115__PZ| & 2, 6 & \verb|LGCS117__PO| & 2, 4, 6, 8\\\hline
    \verb|LGCS118__PO| & 4, 8 & \verb|LGCS119__PO| & 2, 4, 6, 8 & \verb|LGCS120__PO| & 2, 4, 6, 8\\\hline
    \verb|LGCS121__PO| & 1, 3, 5, 7 & \verb|LGCS123__PO| & 2, 4, 6, 8 & \verb|LGCS124__PO| & 3, 7\\\hline
    \verb|LGCS125__PO| & 2, 3, 6, 7 & \verb|LGCS128__PO| & 4, 8 & \verb|LGCS132__PO| & 1, 3, 5, 7\\\hline
    \verb|LGCS135__PO| & 4, 8 & \verb|LGCS145__PO| & 2, 6 & \verb|LGCS166__PZ| & 1, 4, 5, 8\\\hline
    \verb|LGCS183__PO| & 3, 7 & \verb|LGCS185__PO| & 2, 3, 6, 7, 1, 5 & \verb|LGCS188__PO| & 2, 1, 6, 5\\\hline
    \verb|LGCS190__PO| & 1, 3, 5, 7 & \verb|LGCS190__PO2| & 1, 3, 5, 7 & \verb|LGCS191__PO| & 2, 3, 6, 7, 1, 5, 4, 8\\\hline
    \verb|LGCS191__PO2| & 2, 3, 6, 7, 1, 5, 4, 8 & \verb|LIT_031__CM| & 2, 3, 6, 7, 1, 5, 4, 8 & \verb|LIT_034__CM| & 3, 7\\\hline
    \verb|LIT_035__HM| & 3, 7 & \verb|LIT_037__CM| & 2, 3, 6, 7, 1, 5 & \verb|LIT_038__CM| & 2, 3, 6, 7, 1, 5, 4, 8\\\hline
    \verb|LIT_057__CM| & 1, 3, 5, 7 & \verb|LIT_058__CM| & 2, 4, 6, 8 & \verb|LIT_060__CM| & 3, 7\\\hline
    \verb|LIT_061__CM| & 2, 4, 6, 8 & \verb|LIT_062__CM| & 2, 4, 6, 8 & \verb|LIT_066__CM| & 3, 7\\\hline
    \verb|LIT_067__CM| & 4, 8 & \verb|LIT_068__CM| & 2, 6 & \verb|LIT_081__CM| & 1, 5\\\hline
    \verb|LIT_090__CM| & 3, 7 & \verb|LIT_093__CM| & 3, 7 & \verb|LIT_097__CM| & 3, 7\\\hline
    \verb|LIT_099A_CM| & 2, 1, 6, 5, 4, 8 & \verb|LIT_099B_CM| & 2, 1, 6, 5, 4, 8 & \verb|LIT_099C_CM| & 2, 1, 6, 5, 4, 8\\\hline
    \verb|LIT_099__CM| & 2, 3, 6, 7, 4, 8 & \verb|LIT_100__CM| & 1, 3, 5, 7 & \verb|LIT_102__CM| & 1, 5\\\hline
    \verb|LIT_107__HM| & 3, 7 & \verb|LIT_110__HM| & 4, 8 & \verb|LIT_111__CM| & 2, 1, 6, 5\\\hline
    \verb|LIT_118__CM| & 1, 3, 5, 7 & \verb|LIT_119__CM| & 4, 8 & \verb|LIT_120__CM| & 1, 5\\\hline
    \verb|LIT_125__CM| & 4, 8 & \verb|LIT_129IOCM| & 2, 1, 6, 5, 4, 8 & \verb|LIT_129__CM| & 3, 7\\\hline
    \verb|LIT_130__CM| & 2, 4, 6, 8 & \verb|LIT_131__CM| & 2, 1, 6, 5 & \verb|LIT_132__CM| & 4, 8\\\hline
    \verb|LIT_134C_CM| & 2, 4, 6, 8 & \verb|LIT_134__CM| & 1, 5 & \verb|LIT_135__CM| & 1, 5\\\hline
    \verb|LIT_138__CM| & 3, 7 & \verb|LIT_139__CM| & 3, 7 & \verb|LIT_141__HM| & 2, 4, 6, 8\\\hline
    \verb|LIT_143__CM| & 4, 8 & \verb|LIT_144__HM| & 1, 5 & \verb|LIT_148__CM| & 3, 7\\\hline
    \verb|LIT_152__CM| & 2, 6 & \verb|LIT_154__CM| & 3, 7 & \verb|LIT_156__HM| & 1, 3, 5, 7\\\hline
    \verb|LIT_161__AF| & 2, 6 & \verb|LIT_162__AF| & 1, 3, 5, 7 & \verb|LIT_163__AF| & 4, 8\\\hline
    \verb|LIT_165__AF| & 2, 4, 6, 8 & \verb|LIT_179ABHM| & 2, 6 & \verb|LIT_179AHHM| & 2, 6\\\hline
    \verb|LIT_179AIHM| & 2, 6 & \verb|LIT_179AJHM| & 3, 7 & \verb|LIT_179AKHM| & 3, 7\\\hline
    \verb|LIT_179Y_HM| & 1, 5 & \verb|LIT_182__CM| & 3, 7 & \verb|LIT_183__CM| & 1, 5\\\hline
    \verb|LIT_185__CM| & 4, 8 & \verb|LIT_199__CM| & 4, 8 & \verb|LIT_199__CM2| & 4, 8\\\hline
    \verb|MATH001__PO| & 2, 6 & \verb|MATH015B_PZ| & 2, 6 & \verb|MATH019L_HM| & 1, 3, 5, 7\\\hline
    \verb|MATH019__HM| & 1 & \verb|MATH022__SC| & 2, 4, 6, 8 & \verb|MATH023__SC| & 1, 3, 5, 7\\\hline
    \verb|MATH025__PZ| & 1, 3, 5, 7 & \verb|MATH030__CM| & 2, 3, 6, 7, 1, 5, 4, 8 & \verb|MATH030__PO| & 2, 3, 6, 7, 1, 5, 4, 8\\\hline
    \verb|MATH030__PZ| & 2, 3, 6, 7, 1, 5, 4, 8 & \verb|MATH030__SC| & 2, 3, 6, 7, 1, 5, 4, 8 & \verb|MATH031A_CM| & 1, 3, 5, 7\\\hline
    \verb|MATH031H_PO| & 3, 7 & \verb|MATH031__CM| & 2, 3, 6, 7, 1, 5, 4, 8 & \verb|MATH031__PO| & 2, 3, 6, 7, 1, 5, 4, 8\\\hline
    \verb|MATH031__PZ| & 2, 3, 6, 7, 1, 5, 4, 8 & \verb|MATH031__SC| & 2, 3, 6, 7, 1, 5, 4, 8 & \verb|MATH032H_CM| & 1, 3, 5, 7\\\hline
    \verb|MATH032S_PO| & 4, 8 & \verb|MATH032__CM| & 2, 3, 6, 7, 1, 5, 4, 8 & \verb|MATH032__PO| & 2, 3, 6, 7, 1, 5\\\hline
    \verb|MATH032__PZ| & 2, 4, 6, 8 & \verb|MATH032__SC| & 2, 3, 6, 7, 1, 5, 4, 8 & \verb|MATH052__CM| & 2, 3, 6, 7, 4, 8\\\hline
    \verb|MATH052__PZ| & 2, 4, 6, 8 & \verb|MATH055__CM| & 2, 3, 6, 7 & \verb|MATH055__HM| & 2, 3, 6, 7, 1, 5, 4, 8\\\hline
    \verb|MATH055__PZ| & 2, 3, 6, 7, 1, 5, 4, 8 & \verb|MATH055__SC| & 4, 8 & \verb|MATH056__HM| & 2, 3, 6, 7, 1, 5, 4, 8\\\hline
    \verb|MATH058B_PO| & 2, 4, 6, 8 & \verb|MATH058__PO| & 2, 3, 6, 7, 1, 5, 4, 8 & \verb|MATH060C_CM| & 1, 3, 5, 7\\\hline
    \verb|MATH060__CM| & 2, 1, 6, 5, 4, 8 & \verb|MATH060__PO| & 2, 3, 6, 7, 1, 5, 4, 8 & \verb|MATH060__PZ| & 2, 3, 6, 7, 1, 5, 4, 8\\\hline
    \verb|MATH060__SC| & 2, 3, 6, 7, 1, 5, 4, 8 & \verb|MATH061__PO| & 3, 7 & \verb|MATH062__HM| & 2, 4, 6, 8\\\hline
    \verb|MATH067__PO| & 2, 4, 6, 8 & \verb|MATH073__HM| & 2 & \verb|MATH082__HM| & 1, 3, 5, 7\\\hline
    \verb|MATH093__HM| & 1, 3, 5, 7 & \verb|MATH101__PO| & 2, 1, 6, 5, 4, 8 & \verb|MATH102__PO| & 2, 3, 6, 7, 1, 5, 4, 8\\\hline
    \verb|MATH102__PZ| & 1, 3, 5, 7, 4, 8 & \verb|MATH102__SC| & 2, 4, 6, 8 & \verb|MATH103__PO| & 2, 3, 6, 7, 1, 5, 4, 8\\\hline
    \verb|MATH104__HM| & 4, 8 & \verb|MATH106__HM| & 2, 6 & \verb|MATH108B_PZ| & 1, 5\\\hline
    \verb|MATH111__CM| & 2, 3, 6, 7, 1, 5 & \verb|MATH112__PO| & 2, 6 & \verb|MATH113__PO| & 3, 7\\\hline
    \verb|MATH119__HM| & 1, 5 & \verb|MATH131__CM| & 2, 3, 6, 7, 4, 8 & \verb|MATH131__HM| & 2, 3, 6, 7, 1, 5, 4, 8\\\hline
    \verb|MATH131__PO| & 1, 3, 5, 7 & \verb|MATH131__SC| & 2, 6 & \verb|MATH132__HM| & 1, 5\\\hline
    \verb|MATH132__PO| & 4, 8 & \verb|MATH135__CM| & 1, 4, 5, 8 & \verb|MATH135__PO| & 2, 6\\\hline
    \verb|MATH135__SC| & 3, 7 & \verb|MATH136__HM| & 4, 8 & \verb|MATH137__CM| & 1, 5\\\hline
    \verb|MATH137__PO| & 3, 7 & \verb|MATH138__PO| & 4, 8 & \verb|MATH139__SC| & 4, 8\\\hline
    \verb|MATH140__CM| & 1, 5 & \verb|MATH142__PO| & 2, 6 & \verb|MATH144__CM| & 2, 1, 6, 5\\\hline
    \verb|MATH144__PZ| & 4, 8 & \verb|MATH145A_PZ| & 2, 6 & \verb|MATH145__PO| & 4, 8\\\hline
    \verb|MATH147__HM| & 4, 8 & \verb|MATH147__PO| & 2, 6 & \verb|MATH148__CM| & 3, 7\\\hline
    \verb|MATH150__PO| & 4, 8 & \verb|MATH151__CM| & 2, 3, 6, 7, 1, 5, 4, 8 & \verb|MATH151__PO| & 2, 3, 6, 7, 1, 5, 4, 8\\\hline
    \verb|MATH152__CM| & 2, 1, 6, 5, 4, 8 & \verb|MATH152__PO| & 1, 3, 5, 7 & \verb|MATH153__PO| & 3, 7\\\hline
    \verb|MATH154__PO| & 1, 3, 5, 7 & \verb|MATH155__PO| & 2, 6 & \verb|MATH156__CM| & 1, 5\\\hline
    \verb|MATH156__HM| & 4, 8 & \verb|MATH157__HM| & 2, 3, 6, 7, 1, 5, 4, 8 & \verb|MATH158__HM| & 1, 3, 5, 7\\\hline
    \verb|MATH158__PO| & 2, 4, 6, 8 & \verb|MATH160__CM| & 4, 8 & \verb|MATH162__PZ| & 2, 6\\\hline
    \verb|MATH163__CM| & 4, 8 & \verb|MATH164__HM| & 2, 4, 6, 8 & \verb|MATH165__CM| & 2, 6\\\hline
    \verb|MATH165__HM| & 1, 3, 5, 7 & \verb|MATH171__CM| & 1, 3, 5, 7 & \verb|MATH171__HM| & 2, 3, 6, 7, 1, 5, 4, 8\\\hline
    \verb|MATH171__PO| & 2, 4, 6, 8 & \verb|MATH172__CM| & 4, 8 & \verb|MATH173__PO| & 2, 4, 6, 8\\\hline
    \verb|MATH174__HM| & 4, 8 & \verb|MATH175__CM| & 2, 6 & \verb|MATH175__HM| & 4, 8\\\hline
    \verb|MATH175__SC| & 1, 5 & \verb|MATH176__HM| & 2, 6 & \verb|MATH176__PO| & 3, 7\\\hline
    \verb|MATH177__HM| & 1, 3, 5, 7 & \verb|MATH177__PO| & 1, 5 & \verb|MATH179__HM| & 2, 6\\\hline
    \verb|MATH180__CM| & 4, 8 & \verb|MATH180__HM| & 1, 3, 5, 7 & \verb|MATH180__PO| & 2, 6\\\hline
    \verb|MATH181__HM| & 2, 6 & \verb|MATH181__PO| & 4, 8 & \verb|MATH182__PZ| & 1, 5\\\hline
    \verb|MATH183__PO| & 1, 5 & \verb|MATH183__SC| & 3, 7 & \verb|MATH185__PO| & 4, 8\\\hline
    \verb|MATH185__SC| & 2, 6 & \verb|MATH186__CM| & 2, 6 & \verb|MATH187__HM| & 1, 3, 5, 7\\\hline
    \verb|MATH187__PO| & 2, 4, 6, 8 & \verb|MATH189AAHM| & 1, 5 & \verb|MATH189AHHM| & 2, 6\\\hline
    \verb|MATH189AIHM| & 3, 7 & \verb|MATH189AJHM| & 1, 5 & \verb|MATH189B_PO| & 3, 7\\\hline
    \verb|MATH190__CM| & 2, 4, 6, 8 & \verb|MATH190__CM2| & 2, 4, 6, 8 & \verb|MATH190__PO| & 1, 3, 5, 7\\\hline
    \verb|MATH190__PO2| & 1, 3, 5, 7 & \verb|MATH191__CM| & 2, 4, 6, 8 & \verb|MATH191__CM2| & 2, 4, 6, 8\\\hline
    \verb|MATH191__PO| & 2, 3, 6, 7, 4, 8 & \verb|MATH191__PO2| & 2, 3, 6, 7, 4, 8 & \verb|MATH191__SC| & 1, 3, 5, 7\\\hline
    \verb|MATH191__SC2| & 1, 3, 5, 7 & \verb|MATH193__HM| & 8, 7 & \verb|MATH193__HM2| & 8, 7\\\hline
    \verb|MATH194__PZ| & 1, 5 & \verb|MATH194__PZ2| & 1, 5 & \verb|MATH195__CM| & 3, 7\\\hline
    \verb|MATH195__CM2| & 3, 7 & \verb|MATH196__HM| & 2, 3, 6, 7, 1, 5 & \verb|MATH196__HM2| & 2, 3, 6, 7, 1, 5\\\hline
    \verb|MATH197__HM| & 8, 7 & \verb|MATH197__HM2| & 8, 7 & \verb|MATH198__HM| & 5, 6\\\hline
    \verb|MATH198__HM2| & 5, 6 & \verb|MATH198__PZ| & 2, 3, 6, 7, 1, 5, 4, 8 & \verb|MATH198__PZ2| & 2, 3, 6, 7, 1, 5, 4, 8\\\hline
    \verb|MATH294__CG| & 1, 5 & \verb|MATH389L_CG| & 1, 5 & \verb|MATH398__CG| & 1, 4, 5, 8\\\hline
    \verb|MATH454__CG| & 1, 5 & \verb|MATH458__CG| & 4, 8 & \verb|MATH900__HM| & 1, 2, 5, 3, 4, 6, 7, 8\\\hline
    \verb|MATH901__HM| & 1, 2, 5, 3, 4, 6, 7, 8 & \verb|MATH902__HM| & 1, 2, 5, 3, 4, 6, 7, 8 & \verb|MATH903__HM| & 1, 2, 5, 3, 4, 6, 7, 8\\\hline
    \verb|MATH904__HM| & 1, 2, 5, 3, 4, 6, 7, 8 & \verb|MCBI117__HM| & 4, 8 & \verb|MCBI118A_HM| & 2, 4, 6, 8\\\hline
    \verb|MCBI118B_HM| & 2, 4, 6, 8 & \verb|MCBI199__HM| & 7, 8, 5, 6 & \verb|MCBI199__HM2| & 7, 8, 5, 6\\\hline
    \verb|MCBI199__HM3| & 7, 8, 5, 6 & \verb|MCBI199__HM4| & 7, 8, 5, 6 & \verb|MCSI195__PZ| & 1, 3, 5, 7\\\hline
    \verb|MCSI195__PZ2| & 1, 3, 5, 7 & \verb|MENA191__SC| & 4, 8 & \verb|MENA191__SC2| & 4, 8\\\hline
    \verb|MGT_310__CG| & 3, 7 & \verb|MGT_317__CG| & 4, 8 & \verb|MGT_370__CG| & 4, 8\\\hline
    \verb|MGT_390__CG| & 4, 8 & \verb|MGT_395__CG| & 4, 8 & \verb|MGT_398__CG| & 4, 8\\\hline
    \verb|MGT_432__CG| & 4, 8 & \verb|MGT_446__CG| & 4, 8 & \verb|MGT_472__CG| & 4, 8\\\hline
    \verb|MGT_477__CG| & 4, 8 & \verb|MGT_503__CG| & 4, 8 & \verb|MGT_555__CG| & 4, 8\\\hline
    \verb|MLLC021__PZ| & 2, 6 & \verb|MLLC033__PZ| & 2, 6 & \verb|MLLC111__PZ| & 1, 3, 5, 7\\\hline
    \verb|MLLC122__PZ| & 1, 3, 5, 7 & \verb|MLLC144__PZ| & 2, 4, 6, 8 & \verb|MLLC150__PZ| & 1, 3, 5, 7\\\hline
    \verb|MLLC155__PZ| & 2, 4, 6, 8 & \verb|MLLC188__PZ| & 2, 4, 6, 8 & \verb|MOBI188__PO| & 2, 3, 6, 7, 1, 5, 4, 8\\\hline
    \verb|MOBI191A_PO| & 1, 3, 5, 7 & \verb|MOBI191A_PO2| & 1, 3, 5, 7 & \verb|MOBI191B_PO| & 2, 4, 6, 8\\\hline
    \verb|MOBI191B_PO2| & 2, 4, 6, 8 & \verb|MOBI194A_PO| & 1, 3, 5, 7 & \verb|MOBI194A_PO2| & 1, 3, 5, 7\\\hline
    \verb|MOBI194B_PO| & 2, 4, 6, 8 & \verb|MOBI194B_PO2| & 2, 4, 6, 8 & \verb|MSL_001A_CM| & 2, 3, 6, 7, 4, 8\\\hline
    \verb|MSL_099__CM| & 2, 3, 6, 7, 1, 5, 4, 8 & \verb|MSL_101A_CM| & 1, 3, 5, 7 & \verb|MSL_101B_CM| & 2, 4, 6, 8\\\hline
    \verb|MSL_102A_CM| & 1, 3, 5, 7 & \verb|MSL_102B_CM| & 2, 4, 6, 8 & \verb|MSL_103A_CM| & 1, 3, 5, 7\\\hline
    \verb|MSL_103B_CM| & 2, 4, 6, 8 & \verb|MSL_104A_CM| & 1, 3, 5, 7 & \verb|MSL_104B_CM| & 2, 4, 6, 8\\\hline
    \verb|MS__038__SC| & 8, 3, 7, 4 & \verb|MS__041__SC| & 1, 3, 5, 7 & \verb|MS__044__PO| & 1, 5\\\hline
    \verb|MS__045__PZ| & 3, 7 & \verb|MS__049__PO| & 2, 3, 6, 7, 1, 5, 4, 8 & \verb|MS__049__PZ| & 1, 5\\\hline
    \verb|MS__049__SC| & 1, 4, 5, 8 & \verb|MS__050__PO| & 1, 3, 5, 7, 4, 8 & \verb|MS__050__PZ| & 2, 1, 6, 5, 4, 8\\\hline
    \verb|MS__051__PO| & 2, 6 & \verb|MS__051__PZ| & 2, 3, 6, 7, 4, 8 & \verb|MS__051__SC| & 3, 7\\\hline
    \verb|MS__052__PZ| & 3, 7 & \verb|MS__053__SC| & 2, 6 & \verb|MS__055__PZ| & 1, 3, 5, 7\\\hline
    \verb|MS__057__SC| & 8, 3, 7, 4 & \verb|MS__059__SC| & 3, 7 & \verb|MS__060__PZ| & 4, 8\\\hline
    \verb|MS__070__PZ| & 4, 8 & \verb|MS__071__PO| & 3, 7 & \verb|MS__072__PO| & 2, 4, 6, 8\\\hline
    \verb|MS__073__PO| & 1, 5 & \verb|MS__078__PZ| & 3, 7 & \verb|MS__082__PZ| & 2, 3, 6, 7, 1, 5, 4, 8\\\hline
    \verb|MS__082__SC| & 1, 4, 5, 8 & \verb|MS__085__PO| & 2, 6 & \verb|MS__087__PZ| & 2, 1, 6, 5\\\hline
    \verb|MS__088__PZ| & 2, 6 & \verb|MS__089B_PO| & 1, 5 & \verb|MS__090__PZ| & 2, 4, 6, 8\\\hline
    \verb|MS__092__PO| & 4, 8 & \verb|MS__093__PZ| & 1, 4, 5, 8 & \verb|MS__102__PO| & 4, 8\\\hline
    \verb|MS__102__PZ| & 2, 6 & \verb|MS__120__HM| & 3, 7 & \verb|MS__120__PO| & 1, 5\\\hline
    \verb|MS__120__PZ| & 3, 7 & \verb|MS__120__SC| & 3, 7 & \verb|MS__123__JT| & 1, 5\\\hline
    \verb|MS__124__PZ| & 2, 1, 6, 5, 4, 8 & \verb|MS__125__PZ| & 2, 6 & \verb|MS__125__SC| & 2, 4, 6, 8\\\hline
    \verb|MS__126__PZ| & 1, 3, 5, 7 & \verb|MS__128__PZ| & 4, 8 & \verb|MS__130__SC| & 4, 8\\\hline
    \verb|MS__131B_SC| & 4, 8 & \verb|MS__131__PO| & 2, 6 & \verb|MS__131__SC| & 2, 3, 6, 7\\\hline
    \verb|MS__135__PO| & 4, 8 & \verb|MS__139__SC| & 2, 6 & \verb|MS__140__PO| & 4, 8\\\hline
    \verb|MS__148A_PO| & 1, 5 & \verb|MS__148C_PO| & 1, 5 & \verb|MS__148D_PO| & 2, 4, 6, 8\\\hline
    \verb|MS__149G_PO| & 2, 6 & \verb|MS__149T_PO| & 2, 4, 6, 8 & \verb|MS__150__PO| & 3, 7\\\hline
    \verb|MS__153__PO| & 2, 6 & \verb|MS__160__SC| & 2, 6 & \verb|MS__170__HM| & 1, 5\\\hline
    \verb|MS__172__HM| & 1, 3, 5, 7 & \verb|MS__173__HM| & 4, 8 & \verb|MS__175__PO| & 2, 1, 6, 5\\\hline
    \verb|MS__179I_HM| & 1, 5 & \verb|MS__183__PO| & 2, 6 & \verb|MS__187__SC| & 2, 6\\\hline
    \verb|MS__190F_JT| & 1, 3, 5, 7 & \verb|MS__190F_JT2| & 1, 3, 5, 7 & \verb|MS__190S_JT| & 4, 8\\\hline
    \verb|MS__190S_JT2| & 4, 8 & \verb|MS__190__JT| & 2, 6 & \verb|MS__190__JT2| & 2, 6\\\hline
    \verb|MS__194__PZ| & 2, 4, 6, 8 & \verb|MS__194__PZ2| & 2, 4, 6, 8 & \verb|MS__196__PZ| & 2, 3, 6, 7, 1, 5, 4, 8\\\hline
    \verb|MS__196__PZ2| & 2, 3, 6, 7, 1, 5, 4, 8 & \verb|MS__199__SC| & 4, 8 & \verb|MS__199__SC2| & 4, 8\\\hline
    \verb|MUS_003__HM| & 2, 4, 6, 8 & \verb|MUS_003__SC| & 1, 3, 5, 7 & \verb|MUS_004__PO| & 2, 3, 6, 7, 1, 5, 4, 8\\\hline
    \verb|MUS_007__PO| & 2, 3, 6, 7, 1, 5, 4, 8 & \verb|MUS_008__PO| & 2, 3, 6, 7, 1, 5, 4, 8 & \verb|MUS_010__PO| & 1, 3, 5, 7, 4, 8\\\hline
    \verb|MUS_015VOPO| & 3, 7 & \verb|MUS_015__PO| & 1, 3, 5, 7, 4, 8 & \verb|MUS_020__PO| & 1, 3, 5, 7, 4, 8\\\hline
    \verb|MUS_031__PO| & 2, 3, 6, 7, 1, 5, 4, 8 & \verb|MUS_032__PO| & 2, 4, 6, 8 & \verb|MUS_033__PO| & 2, 3, 6, 7, 1, 5, 4, 8\\\hline
    \verb|MUS_035__PO| & 2, 3, 6, 7, 1, 5, 4, 8 & \verb|MUS_037__PO| & 2, 3, 6, 7, 1, 5, 4, 8 & \verb|MUS_040__PO| & 2, 3, 6, 7, 1, 5, 4, 8\\\hline
    \verb|MUS_041__PO| & 2, 3, 6, 7, 1, 5, 4, 8 & \verb|MUS_042B_PO| & 2, 4, 6, 8 & \verb|MUS_042C_PO| & 1, 3, 5, 7\\\hline
    \verb|MUS_046__HM| & 2, 4, 6, 8 & \verb|MUS_047__PO| & 1, 3, 5, 7 & \verb|MUS_049__HM| & 1, 3, 5, 7\\\hline
    \verb|MUS_051__PO| & 2, 3, 6, 7, 1, 5, 4, 8 & \verb|MUS_057__PO| & 3, 7 & \verb|MUS_060__PO| & 2, 4, 6, 8\\\hline
    \verb|MUS_062__PO| & 2, 6 & \verb|MUS_065__PO| & 2, 4, 6, 8 & \verb|MUS_066__SC| & 1, 3, 5, 7\\\hline
    \verb|MUS_067__HM| & 1, 4, 5, 8 & \verb|MUS_068__PO| & 2, 6 & \verb|MUS_070__PO| & 1, 3, 5, 7\\\hline
    \verb|MUS_071__SC| & 1, 3, 5, 7 & \verb|MUS_072__SC| & 2, 4, 6, 8 & \verb|MUS_074__PO| & 4, 8\\\hline
    \verb|MUS_077__PO| & 2, 6 & \verb|MUS_080L_PO| & 2, 3, 6, 7, 1, 5, 4, 8 & \verb|MUS_080__PO| & 2, 3, 6, 7, 1, 5, 4, 8\\\hline
    \verb|MUS_081L_PO| & 1, 3, 5, 7 & \verb|MUS_081__HM| & 3, 7 & \verb|MUS_081__JM| & 2, 6\\\hline
    \verb|MUS_081__PO| & 1, 3, 5, 7 & \verb|MUS_082L_PO| & 2, 4, 6, 8 & \verb|MUS_082__PO| & 2, 4, 6, 8\\\hline
    \verb|MUS_085__SC| & 2, 3, 6, 7, 1, 5, 4, 8 & \verb|MUS_087__PO| & 1, 3, 5, 7 & \verb|MUS_088__PO| & 3, 7\\\hline
    \verb|MUS_089__SC| & 2, 3, 6, 7, 1, 5, 4, 8 & \verb|MUS_091__PO| & 3, 7 & \verb|MUS_092__PO| & 4, 8\\\hline
    \verb|MUS_094__PO| & 4, 8 & \verb|MUS_096A_PO| & 1, 3, 5, 7 & \verb|MUS_096B_PO| & 2, 4, 6, 8\\\hline
    \verb|MUS_099__HM| & 1, 5 & \verb|MUS_100__PO| & 1, 3, 5, 7, 4, 8 & \verb|MUS_101__SC| & 1, 3, 5, 7\\\hline
    \verb|MUS_102__SC| & 2, 4, 6, 8 & \verb|MUS_103__SC| & 3, 7 & \verb|MUS_104__SC| & 1, 5\\\hline
    \verb|MUS_110A_SC| & 1, 3, 5, 7 & \verb|MUS_110B_SC| & 2, 4, 6, 8 & \verb|MUS_112__SC| & 2, 4, 6, 8\\\hline
    \verb|MUS_118__PO| & 4, 8 & \verb|MUS_119__SC| & 4, 8 & \verb|MUS_121__PO| & 2, 4, 6, 8\\\hline
    \verb|MUS_121__SC| & 3, 7 & \verb|MUS_122__PO| & 1, 3, 5, 7 & \verb|MUS_126__SC| & 2, 6\\\hline
    \verb|MUS_130__SC| & 1, 3, 5, 7 & \verb|MUS_131__SC| & 2, 4, 6, 8 & \verb|MUS_140__PO| & 2, 3, 6, 7, 1, 5, 4, 8\\\hline
    \verb|MUS_170F_SC| & 2, 3, 6, 7, 1, 5, 4, 8 & \verb|MUS_170H_SC| & 2, 3, 6, 7, 1, 5, 4, 8 & \verb|MUS_171F_SC| & 2, 3, 6, 7, 1, 5, 4, 8\\\hline
    \verb|MUS_171H_SC| & 2, 3, 6, 7, 1, 5, 4, 8 & \verb|MUS_172__SC| & 2, 3, 6, 7, 1, 5, 4, 8 & \verb|MUS_173__JM| & 2, 3, 6, 7, 1, 5, 4, 8\\\hline
    \verb|MUS_175__JM| & 2, 3, 6, 7, 1, 5, 4, 8 & \verb|MUS_176__JM| & 2, 3, 6, 7, 1, 5 & \verb|MUS_177F_SC| & 2, 3, 6, 7, 1, 5, 4, 8\\\hline
    \verb|MUS_177H_SC| & 2, 3, 6, 7, 1, 5, 4, 8 & \verb|MUS_177__PO| & 2, 6 & \verb|MUS_179D_HM| & 3, 7\\\hline
    \verb|MUS_179E_HM| & 2, 6 & \verb|MUS_179F_HM| & 3, 7 & \verb|MUS_179G_HM| & 1, 5\\\hline
    \verb|MUS_179H_HM| & 4, 8 & \verb|MUS_179I_HM| & 1, 5 & \verb|MUS_180F_SC| & 2, 3, 6, 7, 1, 5, 4, 8\\\hline
    \verb|MUS_180H_SC| & 2, 3, 6, 7, 1, 5, 4, 8 & \verb|MUS_184__PO| & 1, 3, 5, 7 & \verb|MUS_190__PO| & 1, 3, 5, 7\\\hline
    \verb|MUS_190__PO2| & 1, 3, 5, 7 & \verb|MUS_191__SC| & 4, 8 & \verb|MUS_191__SC2| & 4, 8\\\hline
    \verb|MUS_192F_PO| & 4, 8 & \verb|MUS_192F_PO2| & 4, 8 & \verb|MUS_192H_PO| & 4, 8\\\hline
    \verb|MUS_192H_PO2| & 4, 8 & \verb|NEUR080LXKS| & 3, 7 & \verb|NEUR080L_KS| & 3, 7\\\hline
    \verb|NEUR086L_KS| & 3, 7 & \verb|NEUR089A_PO| & 2, 6 & \verb|NEUR095L_JT| & 2, 4, 6, 8\\\hline
    \verb|NEUR101ALPO| & 1, 3, 5, 7 & \verb|NEUR101A_PO| & 1, 3, 5, 7 & \verb|NEUR101BLPO| & 2, 4, 6, 8\\\hline
    \verb|NEUR101B_PO| & 2, 4, 6, 8 & \verb|NEUR102__PO| & 2, 6 & \verb|NEUR103__KS| & 2, 1, 6, 5, 4, 8\\\hline
    \verb|NEUR103__PO| & 2, 4, 6, 8 & \verb|NEUR118__PO| & 1, 4, 5, 8 & \verb|NEUR122L_PO| & 1, 5\\\hline
    \verb|NEUR122__PO| & 1, 5 & \verb|NEUR123__PO| & 2, 4, 6, 8 & \verb|NEUR125L_JT| & 2, 6\\\hline
    \verb|NEUR129L_KS| & 4, 8 & \verb|NEUR131__PO| & 3, 7 & \verb|NEUR133L_KS| & 1, 3, 5, 7\\\hline
    \verb|NEUR140__KS| & 2, 3, 6, 7, 1, 5 & \verb|NEUR148L_KS| & 2, 3, 6, 7, 1, 5, 4, 8 & \verb|NEUR149__KS| & 2, 3, 6, 7, 1, 5, 4, 8\\\hline
    \verb|NEUR157__PO| & 4, 8 & \verb|NEUR161L_KS| & 4, 8 & \verb|NEUR161__KS| & 2, 6\\\hline
    \verb|NEUR168L_PO| & 2, 4, 6, 8 & \verb|NEUR168__PO| & 2, 4, 6, 8 & \verb|NEUR178L_PO| & 1, 3, 5, 7\\\hline
    \verb|NEUR178__PO| & 1, 3, 5, 7 & \verb|NEUR181__KS| & 2, 3, 6, 7 & \verb|NEUR188L_KS| & 1, 3, 5, 7, 4, 8\\\hline
    \verb|NEUR189A_PO| & 3, 7 & \verb|NEUR189B_PO| & 2, 4, 6, 8 & \verb|NEUR189L_KS| & 1, 3, 5, 7\\\hline
    \verb|NEUR189S_PO| & 4, 8 & \verb|NEUR190L_KS| & 1, 3, 5, 7, 4, 8 & \verb|NEUR190L_KS2| & 1, 3, 5, 7, 4, 8\\\hline
    \verb|NEUR190__PO| & 1, 3, 5, 7 & \verb|NEUR190__PO2| & 1, 3, 5, 7 & \verb|NEUR191__KS| & 1, 3, 5, 7, 4, 8\\\hline
    \verb|NEUR191__KS2| & 1, 3, 5, 7, 4, 8 & \verb|NEUR191__PO| & 2, 3, 6, 7, 1, 5, 4, 8 & \verb|NEUR191__PO2| & 2, 3, 6, 7, 1, 5, 4, 8\\\hline
    \verb|NEUR192__PO| & 2, 4, 6, 8 & \verb|NEUR192__PO2| & 2, 4, 6, 8 & \verb|NEUR194A_PO| & 1, 3, 5, 7\\\hline
    \verb|NEUR194A_PO2| & 1, 3, 5, 7 & \verb|NEUR194B_PO| & 2, 4, 6, 8 & \verb|NEUR194B_PO2| & 2, 4, 6, 8\\\hline
    \verb|NEUR196__KS| & 2, 3, 6, 7, 1, 5, 4, 8 & \verb|NEUR196__KS2| & 2, 3, 6, 7, 1, 5, 4, 8 & \verb|NEUR197__KS| & 2, 3, 6, 7, 1, 5, 4, 8\\\hline
    \verb|NEUR197__KS2| & 2, 3, 6, 7, 1, 5, 4, 8 & \verb|ORST050__PZ| & 1, 4, 5, 8 & \verb|ORST060__PZ| & 8, 3, 7, 4\\\hline
    \verb|ORST089__PZ| & 2, 3, 6, 7 & \verb|ORST100A_PZ| & 3, 7 & \verb|ORST100__PZ| & 1, 3, 5, 7\\\hline
    \verb|ORST102__PZ| & 4, 8 & \verb|ORST105__PZ| & 1, 3, 5, 7 & \verb|ORST114__PZ| & 2, 4, 6, 8\\\hline
    \verb|ORST135__PZ| & 2, 1, 6, 5, 4, 8 & \verb|ORST148__PZ| & 1, 4, 5, 8 & \verb|ORST170__PZ| & 2, 6\\\hline
    \verb|ORST173A_PZ| & 1, 5 & \verb|ORST175__PZ| & 3, 7 & \verb|ORST177__PZ| & 2, 6\\\hline
    \verb|ORST180__PZ| & 2, 4, 6, 8 & \verb|ORST191__SC| & 4, 8 & \verb|ORST191__SC2| & 4, 8\\\hline
    \verb|ORST192__PZ| & 1, 3, 5, 7 & \verb|ORST192__PZ2| & 1, 3, 5, 7 & \verb|ORST198C_PZ| & 2, 6\\\hline
    \verb|ORST198C_PZ2| & 2, 6 & \verb|ORST198J_PZ| & 1, 5 & \verb|ORST198J_PZ2| & 1, 5\\\hline
    \verb|ORST198M_PZ| & 1, 3, 5, 7 & \verb|ORST198M_PZ2| & 1, 3, 5, 7 & \verb|PHIL001__PO| & 2, 3, 6, 7, 1, 5, 4, 8\\\hline
    \verb|PHIL007__PO| & 1, 3, 5, 7 & \verb|PHIL007__PZ| & 1, 3, 5, 7 & \verb|PHIL030__CM| & 2, 3, 6, 7, 1, 5, 4, 8\\\hline
    \verb|PHIL030__PO| & 2, 4, 6, 8 & \verb|PHIL030__PZ| & 2, 3, 6, 7, 1, 5 & \verb|PHIL031__CM| & 3, 7\\\hline
    \verb|PHIL031__PO| & 1, 3, 5, 7 & \verb|PHIL031__PZ| & 3, 7 & \verb|PHIL033__CM| & 2, 1, 6, 5\\\hline
    \verb|PHIL033__PO| & 2, 3, 6, 7 & \verb|PHIL033__PZ| & 1, 5 & \verb|PHIL034__CM| & 2, 3, 6, 7, 1, 5, 4, 8\\\hline
    \verb|PHIL034__PO| & 2, 3, 6, 7 & \verb|PHIL035H_CM| & 3, 7 & \verb|PHIL035__CM| & 2, 1, 6, 5, 4, 8\\\hline
    \verb|PHIL035__PO| & 2, 4, 6, 8 & \verb|PHIL036__PO| & 2, 4, 6, 8 & \verb|PHIL037__CM| & 2, 6\\\hline
    \verb|PHIL037__PO| & 2, 1, 6, 5 & \verb|PHIL038__CM| & 2, 3, 6, 7, 1, 5, 4, 8 & \verb|PHIL038__PO| & 3, 7\\\hline
    \verb|PHIL039__CM| & 1, 5 & \verb|PHIL039__PO| & 2, 6 & \verb|PHIL040__PO| & 1, 3, 5, 7\\\hline
    \verb|PHIL042__PO| & 2, 4, 6, 8 & \verb|PHIL043__PO| & 4, 8 & \verb|PHIL044__PO| & 4, 8\\\hline
    \verb|PHIL046__PO| & 4, 8 & \verb|PHIL048__PO| & 1, 5 & \verb|PHIL052__PZ| & 4, 8\\\hline
    \verb|PHIL054__PO| & 3, 7 & \verb|PHIL055__PZ| & 4, 8 & \verb|PHIL057__JT| & 2, 6\\\hline
    \verb|PHIL058__PZ| & 1, 4, 5, 8 & \verb|PHIL060__PO| & 1, 3, 5, 7 & \verb|PHIL070__PO| & 4, 8\\\hline
    \verb|PHIL075__PZ| & 4, 8 & \verb|PHIL080__PO| & 2, 6 & \verb|PHIL084__PZ| & 2, 6\\\hline
    \verb|PHIL090__SC| & 2, 3, 6, 7, 1, 5 & \verb|PHIL095__CM| & 2, 4, 6, 8 & \verb|PHIL100E_CM| & 2, 3, 6, 7, 1, 5, 4, 8\\\hline
    \verb|PHIL100F_CM| & 1, 5 & \verb|PHIL101C_CM| & 2, 6 & \verb|PHIL101D_CM| & 4, 8\\\hline
    \verb|PHIL104__PO| & 2, 6 & \verb|PHIL104__PZ| & 4, 8 & \verb|PHIL106__CM| & 4, 8\\\hline
    \verb|PHIL106__PO| & 4, 8 & \verb|PHIL108__PO| & 2, 6 & \verb|PHIL118__SC| & 2, 1, 6, 5, 4, 8\\\hline
    \verb|PHIL120__PO| & 3, 7 & \verb|PHIL121__CM| & 2, 1, 6, 5 & \verb|PHIL121__HM| & 4, 8\\\hline
    \verb|PHIL122__HM| & 1, 5 & \verb|PHIL123__CM| & 1, 5 & \verb|PHIL124__HM| & 3, 7\\\hline
    \verb|PHIL126__CM| & 4, 8 & \verb|PHIL126__SC| & 2, 1, 6, 5 & \verb|PHIL133__SC| & 1, 5\\\hline
    \verb|PHIL134__CM| & 2, 3, 6, 7 & \verb|PHIL135__CM| & 2, 1, 6, 5, 4, 8 & \verb|PHIL136__SC| & 2, 6\\\hline
    \verb|PHIL138__CM| & 2, 1, 6, 5 & \verb|PHIL139__CM| & 3, 7 & \verb|PHIL140__SC| & 3, 7\\\hline
    \verb|PHIL142__CM| & 1, 5 & \verb|PHIL144__SC| & 2, 3, 6, 7, 1, 5, 4, 8 & \verb|PHIL151__SC| & 1, 5\\\hline
    \verb|PHIL155__SC| & 1, 3, 5, 7 & \verb|PHIL156__PZ| & 4, 8 & \verb|PHIL158__CM| & 1, 3, 5, 7\\\hline
    \verb|PHIL160__CM| & 1, 3, 5, 7, 4, 8 & \verb|PHIL160__SC| & 2, 4, 6, 8 & \verb|PHIL162__SC| & 3, 7\\\hline
    \verb|PHIL163__PO| & 2, 6 & \verb|PHIL163__SC| & 3, 7 & \verb|PHIL164__CM| & 2, 3, 6, 7, 1, 5\\\hline
    \verb|PHIL170__SC| & 3, 7 & \verb|PHIL176__CM| & 1, 5 & \verb|PHIL177__CM| & 4, 8\\\hline
    \verb|PHIL179F_HM| & 2, 4, 6, 8 & \verb|PHIL179G_HM| & 4, 8 & \verb|PHIL181__CM| & 4, 8\\\hline
    \verb|PHIL182__CM| & 1, 5 & \verb|PHIL185C_PO| & 3, 7 & \verb|PHIL185C_PZ| & 4, 8\\\hline
    \verb|PHIL185L_PO| & 4, 8 & \verb|PHIL185N_PZ| & 4, 8 & \verb|PHIL185P_PO| & 4, 8\\\hline
    \verb|PHIL186K_PO| & 1, 5 & \verb|PHIL190__CM| & 3, 7 & \verb|PHIL190__CM2| & 3, 7\\\hline
    \verb|PHIL190__PO| & 1, 3, 5, 7 & \verb|PHIL190__PO2| & 1, 3, 5, 7 & \verb|PHIL190__SC| & 1, 3, 5, 7\\\hline
    \verb|PHIL190__SC2| & 1, 3, 5, 7 & \verb|PHIL191__PO| & 2, 4, 6, 8 & \verb|PHIL191__PO2| & 2, 4, 6, 8\\\hline
    \verb|PHIL191__PZ| & 2, 3, 6, 7, 1, 5, 4, 8 & \verb|PHIL191__PZ2| & 2, 3, 6, 7, 1, 5, 4, 8 & \verb|PHIL191__SC| & 4, 8\\\hline
    \verb|PHIL191__SC2| & 4, 8 & \verb|PHIL192__CM| & 2, 6 & \verb|PHIL192__CM2| & 2, 6\\\hline
    \verb|PHIL198__CM| & 2, 3, 6, 7, 1, 5, 4, 8 & \verb|PHIL198__CM2| & 2, 3, 6, 7, 1, 5, 4, 8 & \verb|PHIL198__SC| & 4, 8\\\hline
    \verb|PHIL198__SC2| & 4, 8 & \verb|PHYS003__PO| & 3, 7 & \verb|PHYS016__PO| & 3, 7\\\hline
    \verb|PHYS017__PO| & 2, 4, 6, 8 & \verb|PHYS023__HM| & 1 & \verb|PHYS024__HM| & 2\\\hline
    \verb|PHYS030LXKS| & 1, 5 & \verb|PHYS030L_KS| & 1, 3, 5, 7 & \verb|PHYS031LXKS| & 2, 4, 6, 8\\\hline
    \verb|PHYS031L_KS| & 2, 4, 6, 8 & \verb|PHYS033L_KS| & 1, 3, 5, 7 & \verb|PHYS034L_KS| & 2, 4, 6, 8\\\hline
    \verb|PHYS035__KS| & 1, 3, 5, 7 & \verb|PHYS041L_PO| & 1, 3, 5, 7 & \verb|PHYS041__PO| & 1, 3, 5, 7\\\hline
    \verb|PHYS042L_PO| & 2, 4, 6, 8 & \verb|PHYS042__PO| & 2, 4, 6, 8 & \verb|PHYS049B_HM| & 2, 4, 6, 8\\\hline
    \verb|PHYS050__HM| & 3, 4 & \verb|PHYS051__HM| & 1, 3, 5, 7 & \verb|PHYS052__HM| & 2, 4, 6, 8\\\hline
    \verb|PHYS054__HM| & 2, 4, 6, 8 & \verb|PHYS064__HM| & 2, 4, 6, 8 & \verb|PHYS070L_PO| & 1, 3, 5, 7\\\hline
    \verb|PHYS070__PO| & 1, 3, 5, 7 & \verb|PHYS071__PO| & 2, 4, 6, 8 & \verb|PHYS072__PO| & 2, 4, 6, 8\\\hline
    \verb|PHYS077L_KS| & 4, 8 & \verb|PHYS084__HM| & 2, 6 & \verb|PHYS100__KS| & 2, 4, 6, 8\\\hline
    \verb|PHYS101L_PO| & 1, 3, 5, 7 & \verb|PHYS101__KS| & 1, 3, 5, 7 & \verb|PHYS101__PO| & 1, 3, 5, 7\\\hline
    \verb|PHYS102__KS| & 2, 4, 6, 8 & \verb|PHYS106L_KS| & 2, 6 & \verb|PHYS108__KS| & 2, 6\\\hline
    \verb|PHYS109L_KS| & 4, 8 & \verb|PHYS110__PO| & 4, 8 & \verb|PHYS111__HM| & 1, 3, 5, 7\\\hline
    \verb|PHYS114__KS| & 2, 4, 6, 8 & \verb|PHYS115__KS| & 1, 3, 5, 7 & \verb|PHYS116__HM| & 2, 4, 6, 8\\\hline
    \verb|PHYS117__HM| & 1, 3, 5, 7 & \verb|PHYS125__PO| & 2, 4, 6, 8 & \verb|PHYS128__PO| & 2, 4, 6, 8\\\hline
    \verb|PHYS133__HM| & 1, 3, 5, 7 & \verb|PHYS134__HM| & 2, 4, 6, 8 & \verb|PHYS139__PO| & 1, 3, 5, 7\\\hline
    \verb|PHYS142__PO| & 2, 6 & \verb|PHYS150__PO| & 2, 4, 6, 8 & \verb|PHYS151__HM| & 1, 3, 5, 7\\\hline
    \verb|PHYS154__HM| & 4, 8 & \verb|PHYS156__HM| & 2, 6 & \verb|PHYS160__PO| & 1, 5\\\hline
    \verb|PHYS161__HM| & 1, 3, 5, 7 & \verb|PHYS162__HM| & 2, 6 & \verb|PHYS164__HM| & 4, 8\\\hline
    \verb|PHYS166__HM| & 4, 8 & \verb|PHYS170X_HM| & 4, 8 & \verb|PHYS170__HM| & 4, 8\\\hline
    \verb|PHYS170__PO| & 1, 3, 5, 7 & \verb|PHYS172__HM| & 2, 6 & \verb|PHYS174L_PO| & 2, 4, 6, 8\\\hline
    \verb|PHYS174__PO| & 2, 4, 6, 8 & \verb|PHYS175__PO| & 2, 4, 6, 8 & \verb|PHYS176__PO| & 3, 7\\\hline
    \verb|PHYS178__KS| & 4, 8 & \verb|PHYS181__HM| & 1, 3, 5, 7 & \verb|PHYS183__HM| & 1, 3, 5, 7, 4, 8\\\hline
    \verb|PHYS187B_KS| & 2, 6 & \verb|PHYS188L_KS| & 2, 3, 6, 7, 1, 5, 4, 8 & \verb|PHYS189L_KS| & 2, 3, 6, 7, 1, 5\\\hline
    \verb|PHYS190L_KS| & 2, 3, 6, 7, 1, 5, 4, 8 & \verb|PHYS190L_KS2| & 2, 3, 6, 7, 1, 5, 4, 8 & \verb|PHYS190__PO| & 1, 3, 5, 7\\\hline
    \verb|PHYS190__PO2| & 1, 3, 5, 7 & \verb|PHYS191E_PO| & 1, 4, 5, 8 & \verb|PHYS191E_PO2| & 1, 4, 5, 8\\\hline
    \verb|PHYS191L_PO| & 4, 8 & \verb|PHYS191L_PO2| & 4, 8 & \verb|PHYS191__HM| & 2, 3, 6, 7, 1, 5, 4, 8\\\hline
    \verb|PHYS191__HM2| & 2, 3, 6, 7, 1, 5, 4, 8 & \verb|PHYS191__KS| & 2, 3, 6, 7, 1, 5, 4, 8 & \verb|PHYS191__KS2| & 2, 3, 6, 7, 1, 5, 4, 8\\\hline
    \verb|PHYS193__HM| & 7 & \verb|PHYS193__HM2| & 7 & \verb|PHYS193__PO| & 1, 3, 5, 7, 4, 8\\\hline
    \verb|PHYS193__PO2| & 1, 3, 5, 7, 4, 8 & \verb|PHYS194__HM| & 8 & \verb|PHYS194__HM2| & 8\\\hline
    \verb|PHYS195__HM| & 7, 8, 5, 6 & \verb|PHYS195__HM2| & 7, 8, 5, 6 & \verb|PHYS195__HM3| & 7, 8, 5, 6\\\hline
    \verb|PHYS195__HM4| & 7, 8, 5, 6 & \verb|PHYS196__KS| & 2, 3, 6, 7, 1, 5, 4, 8 & \verb|PHYS196__KS2| & 2, 3, 6, 7, 1, 5, 4, 8\\\hline
    \verb|PHYS197__HM| & 2, 3, 6, 7, 1, 5, 4, 8 & \verb|PHYS197__HM2| & 2, 3, 6, 7, 1, 5, 4, 8 & \verb|PHYS197__KS| & 2, 3, 6, 7, 1, 5, 4, 8\\\hline
    \verb|PHYS197__KS2| & 2, 3, 6, 7, 1, 5, 4, 8 & \verb|PHYS199__HM| & 8, 7 & \verb|PHYS199__HM2| & 8, 7\\\hline
    \verb|PHYS900__HM| & 1, 2, 5, 3, 4, 6, 7, 8 & \verb|PHYS901__HM| & 1, 2, 5, 3, 4, 6, 7, 8 & \verb|PHYS902__HM| & 1, 2, 5, 3, 4, 6, 7, 8\\\hline
    \verb|PHYS903__HM| & 1, 2, 5, 3, 4, 6, 7, 8 & \verb|PHYS904__HM| & 1, 2, 5, 3, 4, 6, 7, 8 & \verb|POLI001A_PO| & 1, 3, 5, 7\\\hline
    \verb|POLI001B_PO| & 2, 4, 6, 8 & \verb|POLI002__PO| & 2, 6 & \verb|POLI003__PO| & 2, 3, 6, 7, 1, 5, 4, 8\\\hline
    \verb|POLI005__PO| & 2, 3, 6, 7, 1, 5, 4, 8 & \verb|POLI007__PO| & 2, 3, 6, 7, 1, 5, 4, 8 & \verb|POLI008__PO| & 2, 3, 6, 7, 1, 5, 4, 8\\\hline
    \verb|POLI010__PO| & 1, 5 & \verb|POLI030__PO| & 2, 6 & \verb|POLI031__PO| & 4, 8\\\hline
    \verb|POLI033A_PO| & 3, 7 & \verb|POLI033B_PO| & 2, 4, 6, 8 & \verb|POLI060__PO| & 4, 8\\\hline
    \verb|POLI061__PO| & 1, 5 & \verb|POLI070__PO| & 1, 5 & \verb|POLI090__PO| & 2, 3, 6, 7, 1, 5, 4, 8\\\hline
    \verb|POLI100__SC| & 1, 3, 5, 7 & \verb|POLI103__SC| & 3, 7 & \verb|POLI106__SC| & 2, 4, 6, 8\\\hline
    \verb|POLI109__SC| & 2, 4, 6, 8 & \verb|POLI110__SC| & 2, 3, 6, 7, 1, 5 & \verb|POLI111__PO| & 3, 7\\\hline
    \verb|POLI112__PO| & 2, 6 & \verb|POLI115__PO| & 4, 8 & \verb|POLI115__SC| & 4, 8\\\hline
    \verb|POLI116__SC| & 2, 6 & \verb|POLI118__SC| & 2, 6 & \verb|POLI120__SC| & 1, 4, 5, 8\\\hline
    \verb|POLI121__SC| & 4, 8 & \verb|POLI122__SC| & 3, 7 & \verb|POLI123__SC| & 2, 1, 6, 5, 4, 8\\\hline
    \verb|POLI124__PO| & 4, 8 & \verb|POLI126__SC| & 4, 8 & \verb|POLI130__SC| & 4, 8\\\hline
    \verb|POLI131__PO| & 3, 7 & \verb|POLI133__JT| & 2, 6 & \verb|POLI134__PO| & 1, 5\\\hline
    \verb|POLI134__SC| & 3, 7 & \verb|POLI135L_SC| & 4, 8 & \verb|POLI135__PO| & 1, 3, 5, 7\\\hline
    \verb|POLI135__SC| & 4, 8 & \verb|POLI140__SC| & 4, 8 & \verb|POLI142__SC| & 2, 6\\\hline
    \verb|POLI144__SC| & 3, 7 & \verb|POLI145__SC| & 1, 3, 5, 7 & \verb|POLI146__SC| & 2, 1, 6, 5\\\hline
    \verb|POLI148__SC| & 2, 6 & \verb|POLI150__SC| & 3, 7 & \verb|POLI151__SC| & 3, 7\\\hline
    \verb|POLI152__PO| & 3, 7 & \verb|POLI153__SC| & 2, 6 & \verb|POLI155__SC| & 4, 8\\\hline
    \verb|POLI158__PO| & 2, 4, 6, 8 & \verb|POLI161__PO| & 1, 3, 5, 7 & \verb|POLI162__PO| & 2, 1, 6, 5\\\hline
    \verb|POLI165__PO| & 4, 8 & \verb|POLI167__JT| & 4, 8 & \verb|POLI168__PO| & 2, 4, 6, 8\\\hline
    \verb|POLI171__PO| & 2, 6 & \verb|POLI173__SC| & 3, 7 & \verb|POLI177__PO| & 2, 6\\\hline
    \verb|POLI179A_PO| & 3, 7 & \verb|POLI180__SC| & 2, 1, 6, 5 & \verb|POLI181__PO| & 1, 3, 5, 7\\\hline
    \verb|POLI184__SC| & 1, 5 & \verb|POLI189F_PO| & 4, 8 & \verb|POLI189G_PO| & 4, 8\\\hline
    \verb|POLI189H_PO| & 2, 6 & \verb|POLI190C_PO| & 2, 6 & \verb|POLI190C_PO2| & 2, 6\\\hline
    \verb|POLI190D_PO| & 8, 3, 7, 4 & \verb|POLI190D_PO2| & 8, 3, 7, 4 & \verb|POLI190E_PO| & 1, 5\\\hline
    \verb|POLI190E_PO2| & 1, 5 & \verb|POLI191__PO| & 2, 3, 6, 7, 1, 5, 4, 8 & \verb|POLI191__PO2| & 2, 3, 6, 7, 1, 5, 4, 8\\\hline
    \verb|POLI191__SC| & 1, 3, 5, 7, 4, 8 & \verb|POLI191__SC2| & 1, 3, 5, 7, 4, 8 & \verb|POLI193__PO| & 2, 3, 6, 7, 1, 5, 4, 8\\\hline
    \verb|POLI193__PO2| & 2, 3, 6, 7, 1, 5, 4, 8 & \verb|POLI195__PO| & 2, 3, 6, 7, 1, 5, 4, 8 & \verb|POLI195__PO2| & 2, 3, 6, 7, 1, 5, 4, 8\\\hline
    \verb|PONT100__CM| & 2, 3, 6, 7, 1, 5 & \verb|PORT001__PZ| & 1, 3, 5, 7 & \verb|PORT002__PZ| & 2, 4, 6, 8\\\hline
    \verb|PORT022__CM| & 2, 4, 6, 8 & \verb|PORT033__CM| & 1, 3, 5, 7 & \verb|PORT035__PZ| & 2, 3, 6, 7, 1, 5, 4, 8\\\hline
    \verb|PORT044__CM| & 4, 8 & \verb|PORT044__PZ| & 2, 6 & \verb|POST016X_PZ| & 3, 7\\\hline
    \verb|POST020A_PZ| & 3, 7 & \verb|POST020__PZ| & 1, 5 & \verb|POST030__PZ| & 2, 3, 6, 7, 1, 5, 4, 8\\\hline
    \verb|POST040__PZ| & 3, 7 & \verb|POST050__PZ| & 2, 3, 6, 7, 1, 5, 4, 8 & \verb|POST060__PZ| & 2, 6\\\hline
    \verb|POST066__PZ| & 4, 8 & \verb|POST070__PZ| & 4, 8 & \verb|POST091__PZ| & 1, 5\\\hline
    \verb|POST101__PZ| & 3, 7 & \verb|POST106__PZ| & 2, 3, 6, 7 & \verb|POST107__CH| & 2, 6\\\hline
    \verb|POST114__HM| & 1, 5 & \verb|POST124__PZ| & 3, 7 & \verb|POST126__PZ| & 2, 4, 6, 8\\\hline
    \verb|POST128__PZ| & 4, 8 & \verb|POST131__PZ| & 2, 6 & \verb|POST138__PZ| & 1, 5\\\hline
    \verb|POST140__HM| & 3, 7 & \verb|POST141__PZ| & 2, 1, 6, 5 & \verb|POST144__PZ| & 2, 6\\\hline
    \verb|POST151__PZ| & 4, 8 & \verb|POST152__PZ| & 4, 8 & \verb|POST153__PZ| & 4, 8\\\hline
    \verb|POST162__PZ| & 2, 4, 6, 8 & \verb|POST172__PZ| & 1, 5 & \verb|POST174__PZ| & 4, 8\\\hline
    \verb|POST179C_HM| & 4, 8 & \verb|POST179D_HM| & 1, 5 & \verb|POST182__PZ| & 1, 5\\\hline
    \verb|POST195__PZ| & 1, 3, 5, 7 & \verb|POST195__PZ2| & 1, 3, 5, 7 & \verb|POST199__PZ| & 2, 4, 6, 8\\\hline
    \verb|POST199__PZ2| & 2, 4, 6, 8 & \verb|PPA_190__PO| & 1, 3, 5, 7 & \verb|PPA_190__PO2| & 1, 3, 5, 7\\\hline
    \verb|PPA_191__PO| & 2, 4, 6, 8 & \verb|PPA_191__PO2| & 2, 4, 6, 8 & \verb|PPA_191__SC| & 4, 8\\\hline
    \verb|PPA_191__SC2| & 4, 8 & \verb|PPA_192__SC| & 4, 8 & \verb|PPA_192__SC2| & 4, 8\\\hline
    \verb|PPA_195__PO| & 1, 3, 5, 7 & \verb|PPA_195__PO2| & 1, 3, 5, 7 & \verb|PPE_001A_CM| & 2, 4, 6, 8\\\hline
    \verb|PPE_001B_CM| & 2, 4, 6, 8 & \verb|PPE_011A_CM| & 2, 4, 6, 8 & \verb|PPE_011B_CM| & 2, 4, 6, 8\\\hline
    \verb|PPE_100__PO| & 2, 4, 6, 8 & \verb|PPE_110A_CM| & 2, 3, 6, 7, 1, 5, 4, 8 & \verb|PPE_110B_CM| & 2, 3, 6, 7, 1, 5, 4, 8\\\hline
    \verb|PPE_160__PO| & 1, 3, 5, 7 & \verb|PPE_195__PO| & 2, 1, 6, 5, 4, 8 & \verb|PPE_195__PO2| & 2, 1, 6, 5, 4, 8\\\hline
    \verb|PSYC010__PZ| & 2, 3, 6, 7, 1, 5, 4, 8 & \verb|PSYC030__CM| & 2, 3, 6, 7, 1, 5, 4, 8 & \verb|PSYC037__CM| & 2, 3, 6, 7, 1, 5, 4, 8\\\hline
    \verb|PSYC040__CM| & 2, 3, 6, 7, 1, 5, 4, 8 & \verb|PSYC051__CM| & 1, 5 & \verb|PSYC051__PO| & 2, 3, 6, 7, 1, 5, 4, 8\\\hline
    \verb|PSYC052__SC| & 2, 1, 6, 5, 4, 8 & \verb|PSYC065__CM| & 2, 3, 6, 7, 1, 5, 4, 8 & \verb|PSYC070__CM| & 2, 3, 6, 7, 1, 5, 4, 8\\\hline
    \verb|PSYC081__CM| & 2, 3, 6, 7, 1, 5, 4, 8 & \verb|PSYC084__CH| & 3, 7 & \verb|PSYC091P_PZ| & 2, 3, 6, 7, 1, 5, 4, 8\\\hline
    \verb|PSYC091__PZ| & 2, 3, 6, 7, 1, 5, 4, 8 & \verb|PSYC092P_PZ| & 2, 3, 6, 7, 1, 5, 4, 8 & \verb|PSYC092__CM| & 2, 3, 6, 7, 1, 5, 4, 8\\\hline
    \verb|PSYC092__PZ| & 2, 3, 6, 7, 1, 5, 4, 8 & \verb|PSYC097__CM| & 2, 3, 6, 7 & \verb|PSYC101__PZ| & 2, 3, 6, 7, 4, 8\\\hline
    \verb|PSYC103__PZ| & 2, 3, 6, 7, 1, 5 & \verb|PSYC103__SC| & 2, 3, 6, 7, 1, 5, 4, 8 & \verb|PSYC104L_SC| & 2, 3, 6, 7, 1, 5, 4, 8\\\hline
    \verb|PSYC104P_PZ| & 4, 8 & \verb|PSYC104__PZ| & 4, 8 & \verb|PSYC104__SC| & 2, 3, 6, 7, 1, 5, 4, 8\\\hline
    \verb|PSYC105__PZ| & 2, 3, 6, 7, 1, 5, 4, 8 & \verb|PSYC106__PO| & 2, 6 & \verb|PSYC107__CM| & 2, 3, 6, 7\\\hline
    \verb|PSYC108__PO| & 1, 3, 5, 7 & \verb|PSYC109__CM| & 2, 3, 6, 7, 1, 5, 4, 8 & \verb|PSYC110__CM| & 2, 3, 6, 7, 1, 5, 4, 8\\\hline
    \verb|PSYC111P_PZ| & 2, 3, 6, 7, 1, 5 & \verb|PSYC111__CM| & 2, 3, 6, 7, 1, 5, 4, 8 & \verb|PSYC111__PZ| & 2, 3, 6, 7, 1, 5\\\hline
    \verb|PSYC112__PZ| & 1, 3, 5, 7 & \verb|PSYC112__SC| & 3, 7 & \verb|PSYC113__PZ| & 1, 4, 5, 8\\\hline
    \verb|PSYC113__SC| & 2, 1, 6, 5, 4, 8 & \verb|PSYC114__SC| & 2, 4, 6, 8 & \verb|PSYC116__PZ| & 2, 6\\\hline
    \verb|PSYC117__PZ| & 4, 8 & \verb|PSYC119__CM| & 1, 3, 5, 7 & \verb|PSYC120__CM| & 2, 4, 6, 8\\\hline
    \verb|PSYC121__SC| & 2, 6 & \verb|PSYC124__SC| & 2, 6 & \verb|PSYC125__AF| & 2, 4, 6, 8\\\hline
    \verb|PSYC126__PZ| & 1, 5 & \verb|PSYC127__PZ| & 2, 6 & \verb|PSYC128__SC| & 3, 7\\\hline
    \verb|PSYC130P_PZ| & 3, 7 & \verb|PSYC130__PZ| & 3, 7 & \verb|PSYC131A_PO| & 4, 8\\\hline
    \verb|PSYC131L_SC| & 2, 6 & \verb|PSYC131__CM| & 1, 5 & \verb|PSYC131__PO| & 2, 6\\\hline
    \verb|PSYC131__SC| & 2, 3, 6, 7, 1, 5 & \verb|PSYC132__CM| & 2, 6 & \verb|PSYC133__PO| & 1, 3, 5, 7\\\hline
    \verb|PSYC133__SC| & 2, 6 & \verb|PSYC134__CM| & 1, 4, 5, 8 & \verb|PSYC134__HM| & 2, 6\\\hline
    \verb|PSYC134__SC| & 4, 8 & \verb|PSYC136__CM| & 1, 5 & \verb|PSYC137__PO| & 3, 7\\\hline
    \verb|PSYC139__HM| & 3, 7 & \verb|PSYC140A_PZ| & 1, 5 & \verb|PSYC142__CM| & 1, 3, 5, 7, 4, 8\\\hline
    \verb|PSYC143__CM| & 2, 6 & \verb|PSYC143__SC| & 3, 7 & \verb|PSYC144__CM| & 2, 3, 6, 7, 1, 5, 4, 8\\\hline
    \verb|PSYC147__CH| & 4, 8 & \verb|PSYC147__SC| & 1, 4, 5, 8 & \verb|PSYC148__PZ| & 4, 8\\\hline
    \verb|PSYC150__AF| & 1, 3, 5, 7 & \verb|PSYC151__SC| & 4, 8 & \verb|PSYC153__AA| & 2, 4, 6, 8\\\hline
    \verb|PSYC154__CM| & 2, 4, 6, 8 & \verb|PSYC154__PO| & 2, 3, 6, 7 & \verb|PSYC155__CM| & 3, 7\\\hline
    \verb|PSYC157__PO| & 2, 4, 6, 8 & \verb|PSYC158__PO| & 1, 3, 5, 7 & \verb|PSYC159__CM| & 1, 3, 5, 7\\\hline
    \verb|PSYC160__CM| & 1, 5 & \verb|PSYC160__PO| & 1, 3, 5, 7 & \verb|PSYC162__PO| & 2, 4, 6, 8\\\hline
    \verb|PSYC162__SC| & 3, 7 & \verb|PSYC163__PO| & 2, 6 & \verb|PSYC163__SC| & 4, 8\\\hline
    \verb|PSYC164__CM| & 3, 7 & \verb|PSYC164__SC| & 2, 6 & \verb|PSYC166__CM| & 2, 6\\\hline
    \verb|PSYC167__CM| & 1, 3, 5, 7 & \verb|PSYC168L_SC| & 2, 6 & \verb|PSYC168__SC| & 2, 4, 6, 8\\\hline
    \verb|PSYC169__SC| & 3, 7 & \verb|PSYC170__PO| & 4, 8 & \verb|PSYC175__SC| & 2, 6\\\hline
    \verb|PSYC176__PO| & 1, 3, 5, 7 & \verb|PSYC179M_HM| & 4, 8 & \verb|PSYC179N_HM| & 3, 7\\\hline
    \verb|PSYC179O_HM| & 2, 6 & \verb|PSYC179U_HM| & 1, 5 & \verb|PSYC180C_PO| & 4, 8\\\hline
    \verb|PSYC180F_PO| & 2, 6 & \verb|PSYC180H_PO| & 2, 4, 6, 8 & \verb|PSYC180M_PO| & 2, 6\\\hline
    \verb|PSYC180N_CH| & 2, 6 & \verb|PSYC180R_PO| & 4, 8 & \verb|PSYC180S_PO| & 3, 7\\\hline
    \verb|PSYC180W_PO| & 3, 7 & \verb|PSYC180__CM| & 2, 3, 6, 7, 1, 5 & \verb|PSYC181__PZ| & 4, 8\\\hline
    \verb|PSYC183__SC| & 2, 4, 6, 8 & \verb|PSYC184__PZ| & 3, 7 & \verb|PSYC188__CM| & 2, 6\\\hline
    \verb|PSYC189Q_PO| & 2, 6 & \verb|PSYC189__PZ| & 2, 3, 6, 7, 1, 5 & \verb|PSYC190A_PO| & 3, 7\\\hline
    \verb|PSYC190A_PO2| & 3, 7 & \verb|PSYC190R_PO| & 1, 3, 5, 7 & \verb|PSYC190R_PO2| & 1, 3, 5, 7\\\hline
    \verb|PSYC190__PO| & 1, 5 & \verb|PSYC190__PO2| & 1, 5 & \verb|PSYC191__PO| & 2, 4, 6, 8\\\hline
    \verb|PSYC191__PO2| & 2, 4, 6, 8 & \verb|PSYC191__PZ| & 2, 4, 6, 8 & \verb|PSYC191__PZ2| & 2, 4, 6, 8\\\hline
    \verb|PSYC191__SC| & 1, 3, 5, 7 & \verb|PSYC191__SC2| & 1, 3, 5, 7 & \verb|PSYC193__SC| & 4, 8\\\hline
    \verb|PSYC193__SC2| & 4, 8 & \verb|PSYC194__PZ| & 3, 7 & \verb|PSYC194__PZ2| & 3, 7\\\hline
    \verb|PSYC197A_CM| & 2, 3, 6, 7, 1, 5, 4, 8 & \verb|PSYC197A_CM2| & 2, 3, 6, 7, 1, 5, 4, 8 & \verb|PSYC197B_CM| & 2, 3, 6, 7, 1, 5, 4, 8\\\hline
    \verb|PSYC197B_CM2| & 2, 3, 6, 7, 1, 5, 4, 8 & \verb|PSYC197__PZ| & 3, 7 & \verb|PSYC197__PZ2| & 3, 7\\\hline
    \verb|PSYC198__CM| & 1, 3, 5, 7 & \verb|PSYC198__CM2| & 1, 3, 5, 7 & \verb|PSYC308B_CG| & 3, 7\\\hline
    \verb|PSYC308C_CG| & 4, 8 & \verb|PSYC315F_CG| & 4, 8 & \verb|PSYC315H_CG| & 4, 8\\\hline
    \verb|PSYC332__CG| & 4, 8 & \verb|PSYC450__CG| & 4, 8 & \verb|RLIT191__PO| & 8, 3, 7, 4\\\hline
    \verb|RLIT191__PO2| & 8, 3, 7, 4 & \verb|RLST010__CM| & 2, 6 & \verb|RLST020__PO| & 4, 8\\\hline
    \verb|RLST037__CM| & 2, 3, 6, 7, 1, 5, 4, 8 & \verb|RLST038__CM| & 2, 3, 6, 7 & \verb|RLST040__PO| & 4, 8\\\hline
    \verb|RLST045__CM| & 3, 7 & \verb|RLST049__PO| & 4, 8 & \verb|RLST055__CM| & 4, 8\\\hline
    \verb|RLST056__CM| & 1, 3, 5, 7 & \verb|RLST058__CM| & 2, 3, 6, 7, 1, 5, 4, 8 & \verb|RLST060__SC| & 4, 8\\\hline
    \verb|RLST061__SC| & 2, 6 & \verb|RLST064__CM| & 1, 5 & \verb|RLST082__CM| & 1, 5\\\hline
    \verb|RLST083__CM| & 2, 4, 6, 8 & \verb|RLST084__CM| & 4, 8 & \verb|RLST086__SC| & 1, 5\\\hline
    \verb|RLST092__SC| & 3, 7 & \verb|RLST093__CM| & 1, 3, 5, 7 & \verb|RLST098__SC| & 4, 8\\\hline
    \verb|RLST100__PO| & 8, 3, 7, 4 & \verb|RLST102__CM| & 1, 3, 5, 7 & \verb|RLST103__PO| & 8, 3, 7, 4\\\hline
    \verb|RLST107__PO| & 2, 6 & \verb|RLST109__CM| & 2, 6 & \verb|RLST110__PO| & 2, 1, 6, 5\\\hline
    \verb|RLST113__HM| & 1, 5 & \verb|RLST114__HM| & 3, 7 & \verb|RLST125__CM| & 2, 1, 6, 5\\\hline
    \verb|RLST126__CM| & 2, 6 & \verb|RLST128__PO| & 1, 5 & \verb|RLST129__CM| & 2, 6\\\hline
    \verb|RLST130A_CM| & 4, 8 & \verb|RLST130__CM| & 1, 4, 5, 8 & \verb|RLST135__CM| & 2, 4, 6, 8\\\hline
    \verb|RLST139__PO| & 4, 8 & \verb|RLST142__AF| & 2, 3, 6, 7, 4, 8 & \verb|RLST144__PO| & 4, 8\\\hline
    \verb|RLST146__PO| & 2, 6 & \verb|RLST148__PO| & 3, 7 & \verb|RLST150__AF| & 2, 3, 6, 7, 4, 8\\\hline
    \verb|RLST152__PO| & 1, 5 & \verb|RLST157__PO| & 3, 7 & \verb|RLST158__PO| & 2, 6\\\hline
    \verb|RLST160__PO| & 4, 8 & \verb|RLST164__PO| & 2, 6 & \verb|RLST166B_CM| & 3, 7\\\hline
    \verb|RLST168__HM| & 3, 7 & \verb|RLST169__CM| & 4, 8 & \verb|RLST171__CM| & 3, 7\\\hline
    \verb|RLST179G_HM| & 1, 3, 5, 7 & \verb|RLST179__SC| & 4, 8 & \verb|RLST180__CM| & 4, 8\\\hline
    \verb|RLST180__PO| & 2, 6 & \verb|RLST181__PO| & 1, 3, 5, 7 & \verb|RLST183__HM| & 1, 5\\\hline
    \verb|RLST184__PO| & 2, 6 & \verb|RLST185__SC| & 2, 6 & \verb|RLST187__PO| & 4, 8\\\hline
    \verb|RLST189C_PO| & 2, 6 & \verb|RLST189L_PO| & 2, 6 & \verb|RLST190__PO| & 1, 3, 5, 7\\\hline
    \verb|RLST190__PO2| & 1, 3, 5, 7 & \verb|RLST191__PO| & 2, 4, 6, 8 & \verb|RLST191__PO2| & 2, 4, 6, 8\\\hline
    \verb|RLST191__SC| & 4, 8 & \verb|RLST191__SC2| & 4, 8 & \verb|RUSS001__PO| & 1, 3, 5, 7\\\hline
    \verb|RUSS002__PO| & 2, 4, 6, 8 & \verb|RUSS011__PO| & 1, 3, 5, 7, 4, 8 & \verb|RUSS013__PO| & 1, 3, 5, 7, 4, 8\\\hline
    \verb|RUSS033__PO| & 1, 3, 5, 7 & \verb|RUSS044__PO| & 2, 4, 6, 8 & \verb|RUSS180__PO| & 1, 5\\\hline
    \verb|RUSS181__PO| & 2, 6 & \verb|RUSS182__PO| & 4, 8 & \verb|RUSS187__PO| & 3, 7\\\hline
    \verb|RUSS191__PO| & 8, 3, 7, 4 & \verb|RUSS191__PO2| & 8, 3, 7, 4 & \verb|RUSS191__SC| & 4, 8\\\hline
    \verb|RUSS191__SC2| & 4, 8 & \verb|RUST079__PO| & 4, 8 & \verb|RUST100__PO| & 3, 7\\\hline
    \verb|RUST110__PO| & 4, 8 & \verb|RUST112__PO| & 2, 6 & \verb|RUST185__PO| & 2, 1, 6, 5\\\hline
    \verb|RUST191__PO| & 2, 6 & \verb|RUST191__PO2| & 2, 6 & \verb|SCI_010L_CM| & 2, 3, 6, 7, 1, 5, 4, 8\\\hline
    \verb|SCI_030LACM| & 1, 5 & \verb|SCI_030LBCM| & 1, 5 & \verb|SCI_031LBCM| & 1, 5\\\hline
    \verb|SCI_050L_CM| & 2, 1, 6, 5 & \verb|SCI_119__CM| & 2, 6 & \verb|SCI_131__CM| & 3, 7\\\hline
    \verb|SCI_132__CM| & 4, 8 & \verb|SCI_153__CM| & 3, 7 & \verb|SCI_196__CM| & 1, 4, 5, 8\\\hline
    \verb|SCI_196__CM2| & 1, 4, 5, 8 & \verb|SCI_197__CM| & 1, 4, 5, 8 & \verb|SCI_197__CM2| & 1, 4, 5, 8\\\hline
    \verb|SOC_001X_PZ| & 2, 1, 6, 5, 4, 8 & \verb|SOC_001__PZ| & 3, 7 & \verb|SOC_012__PZ| & 3, 7\\\hline
    \verb|SOC_030__CH| & 4, 8 & \verb|SOC_051__PO| & 2, 3, 6, 7, 1, 5, 4, 8 & \verb|SOC_061__PZ| & 2, 3, 6, 7\\\hline
    \verb|SOC_069__PZ| & 2, 6 & \verb|SOC_072__PZ| & 2, 3, 6, 7 & \verb|SOC_083__PZ| & 2, 4, 6, 8\\\hline
    \verb|SOC_090__PO| & 2, 6 & \verb|SOC_091__PZ| & 8, 3, 7, 4 & \verb|SOC_092__PZ| & 1, 5\\\hline
    \verb|SOC_095__PZ| & 2, 1, 6, 5, 4, 8 & \verb|SOC_099__PZ| & 1, 5 & \verb|SOC_101__PZ| & 2, 3, 6, 7, 1, 5, 4, 8\\\hline
    \verb|SOC_102__PO| & 1, 3, 5, 7 & \verb|SOC_102__PZ| & 2, 3, 6, 7, 1, 5, 4, 8 & \verb|SOC_103__JT| & 2, 1, 6, 5, 4, 8\\\hline
    \verb|SOC_109__PZ| & 3, 7 & \verb|SOC_110__PZ| & 1, 3, 5, 7 & \verb|SOC_118__PZ| & 4, 8\\\hline
    \verb|SOC_119__PO| & 3, 7 & \verb|SOC_119__PZ| & 2, 4, 6, 8 & \verb|SOC_122__PZ| & 3, 7\\\hline
    \verb|SOC_123__PO| & 4, 8 & \verb|SOC_126__AA| & 3, 7 & \verb|SOC_128__JT| & 2, 6\\\hline
    \verb|SOC_130__PO| & 4, 8 & \verb|SOC_138__PO| & 1, 4, 5, 8 & \verb|SOC_146__PO| & 2, 6\\\hline
    \verb|SOC_147__PO| & 2, 4, 6, 8 & \verb|SOC_148__PO| & 3, 7 & \verb|SOC_150__AA| & 2, 4, 6, 8\\\hline
    \verb|SOC_150__CH| & 2, 4, 6, 8 & \verb|SOC_154__PO| & 2, 4, 6, 8 & \verb|SOC_157__PZ| & 2, 4, 6, 8\\\hline
    \verb|SOC_169X_PO| & 3, 7 & \verb|SOC_170__PZ| & 2, 4, 6, 8 & \verb|SOC_179C_HM| & 3, 7\\\hline
    \verb|SOC_179D_HM| & 4, 8 & \verb|SOC_179E_HM| & 1, 5 & \verb|SOC_179F_HM| & 1, 5\\\hline
    \verb|SOC_189E_PO| & 4, 8 & \verb|SOC_189F_PO| & 3, 7 & \verb|SOC_189Z_PO| & 2, 6\\\hline
    \verb|SOC_190__PO| & 1, 3, 5, 7 & \verb|SOC_190__PO2| & 1, 3, 5, 7 & \verb|SOC_191__PO| & 2, 4, 6, 8\\\hline
    \verb|SOC_191__PO2| & 2, 4, 6, 8 & \verb|SOC_199A_PZ| & 2, 3, 6, 7, 1, 5 & \verb|SOC_199A_PZ2| & 2, 3, 6, 7, 1, 5\\\hline
    \verb|SOC_199B_PZ| & 2, 4, 6, 8 & \verb|SOC_199B_PZ2| & 2, 4, 6, 8 & \verb|SOC_199C_PZ| & 4, 8\\\hline
    \verb|SOC_199C_PZ2| & 4, 8 & \verb|SOSC150__HM| & 1, 3, 5, 7 & \verb|SPAN001__CM| & 1, 3, 5, 7\\\hline
    \verb|SPAN001__PO| & 1, 3, 5, 7 & \verb|SPAN001__PZ| & 1, 3, 5, 7 & \verb|SPAN002__CM| & 2, 4, 6, 8\\\hline
    \verb|SPAN002__PO| & 2, 4, 6, 8 & \verb|SPAN002__PZ| & 2, 4, 6, 8 & \verb|SPAN011__PO| & 1, 3, 5, 7, 4, 8\\\hline
    \verb|SPAN013__PO| & 1, 3, 5, 7, 4, 8 & \verb|SPAN022__CM| & 2, 3, 6, 7, 1, 5, 4, 8 & \verb|SPAN022__PO| & 1, 3, 5, 7\\\hline
    \verb|SPAN022__PZ| & 2, 3, 6, 7, 1, 5, 4, 8 & \verb|SPAN022__SC| & 1, 3, 5, 7 & \verb|SPAN031__PZ| & 2, 3, 6, 7, 1, 5, 4, 8\\\hline
    \verb|SPAN033L_CM| & 1, 3, 5, 7, 4, 8 & \verb|SPAN033__CM| & 2, 3, 6, 7, 1, 5, 4, 8 & \verb|SPAN033__PO| & 2, 3, 6, 7, 1, 5, 4, 8\\\hline
    \verb|SPAN033__PZ| & 2, 3, 6, 7, 1, 5, 4, 8 & \verb|SPAN033__SC| & 2, 3, 6, 7, 1, 5, 4, 8 & \verb|SPAN035__PZ| & 2, 3, 6, 7, 1, 5, 4, 8\\\hline
    \verb|SPAN044L_CM| & 1, 3, 5, 7, 4, 8 & \verb|SPAN044__CM| & 2, 3, 6, 7, 1, 5 & \verb|SPAN044__PO| & 2, 3, 6, 7, 1, 5, 4, 8\\\hline
    \verb|SPAN044__PZ| & 2, 3, 6, 7, 1, 5, 4, 8 & \verb|SPAN044__SC| & 2, 3, 6, 7, 1, 5, 4, 8 & \verb|SPAN050__PZ| & 1, 3, 5, 7\\\hline
    \verb|SPAN055__PZ| & 2, 3, 6, 7, 1, 5, 4, 8 & \verb|SPAN100__PZ| & 2, 3, 6, 7, 1, 5, 4, 8 & \verb|SPAN101__PO| & 2, 3, 6, 7, 1, 5, 4, 8\\\hline
    \verb|SPAN101__SC| & 2, 3, 6, 7, 1, 5, 4, 8 & \verb|SPAN102__CM| & 1, 3, 5, 7, 4, 8 & \verb|SPAN103__SC| & 1, 5\\\hline
    \verb|SPAN104__PO| & 2, 6 & \verb|SPAN104__PZ| & 1, 3, 5, 7 & \verb|SPAN106__PO| & 2, 1, 6, 5\\\hline
    \verb|SPAN107__PZ| & 4, 8 & \verb|SPAN108__PO| & 2, 6 & \verb|SPAN109__PO| & 3, 7\\\hline
    \verb|SPAN120A_PO| & 3, 7 & \verb|SPAN125A_CM| & 3, 7 & \verb|SPAN125A_PO| & 1, 5\\\hline
    \verb|SPAN125B_CM| & 4, 8 & \verb|SPAN125B_PO| & 2, 6 & \verb|SPAN126__PO| & 8, 3, 7, 4\\\hline
    \verb|SPAN127__CH| & 2, 4, 6, 8 & \verb|SPAN130__CM| & 2, 1, 6, 5 & \verb|SPAN131__SC| & 2, 6\\\hline
    \verb|SPAN132__PZ| & 4, 8 & \verb|SPAN134__SC| & 4, 8 & \verb|SPAN136__PZ| & 2, 3, 6, 7\\\hline
    \verb|SPAN137__PZ| & 2, 6 & \verb|SPAN139__SC| & 3, 7 & \verb|SPAN140__PO| & 3, 7\\\hline
    \verb|SPAN140__SC| & 2, 6 & \verb|SPAN141__SC| & 2, 4, 6, 8 & \verb|SPAN143__SC| & 1, 5\\\hline
    \verb|SPAN144__PZ| & 2, 3, 6, 7, 1, 5 & \verb|SPAN145__SC| & 4, 8 & \verb|SPAN146__PO| & 4, 8\\\hline
    \verb|SPAN148__PZ| & 1, 5 & \verb|SPAN150__PZ| & 4, 8 & \verb|SPAN153__PO| & 2, 6\\\hline
    \verb|SPAN156__PZ| & 1, 5 & \verb|SPAN157__PZ| & 4, 8 & \verb|SPAN159__PO| & 1, 5\\\hline
    \verb|SPAN159__PZ| & 2, 6 & \verb|SPAN163__PZ| & 3, 7 & \verb|SPAN163__SC| & 1, 3, 5, 7\\\hline
    \verb|SPAN167__CM| & 4, 8 & \verb|SPAN171__PZ| & 4, 8 & \verb|SPAN172__PZ| & 4, 8\\\hline
    \verb|SPAN181__PZ| & 2, 6 & \verb|SPAN181__SC| & 3, 7 & \verb|SPAN183__CM| & 1, 3, 5, 7\\\hline
    \verb|SPAN183__SC| & 2, 6 & \verb|SPAN185__PO| & 2, 6 & \verb|SPAN189A_PO| & 2, 4, 6, 8\\\hline
    \verb|SPAN191__PO| & 2, 3, 6, 7, 1, 5, 4, 8 & \verb|SPAN191__PO2| & 2, 3, 6, 7, 1, 5, 4, 8 & \verb|SPAN191__SC| & 1, 3, 5, 7, 4, 8\\\hline
    \verb|SPAN191__SC2| & 1, 3, 5, 7, 4, 8 & \verb|SPAN192__PO| & 1, 4, 5, 8 & \verb|SPAN192__PO2| & 1, 4, 5, 8\\\hline
    \verb|SPAN199__PZ| & 2, 4, 6, 8 & \verb|SPAN199__PZ2| & 2, 4, 6, 8 & \verb|SPCH061A_CM| & 2, 3, 6, 7, 1, 5, 4, 8\\\hline
    \verb|SPCH061B_CM| & 2, 3, 6, 7, 1, 5, 4, 8 & \verb|STEM000__5C| & 1, 2, 5, 3, 4, 6, 7, 8 & \verb|STEM001__5C| & 1, 2, 5, 3, 4, 6, 7, 8\\\hline
    \verb|STEM002__5C| & 1, 2, 5, 3, 4, 6, 7, 8 & \verb|STEM003__5C| & 1, 2, 5, 3, 4, 6, 7, 8 & \verb|STEM004__5C| & 1, 2, 5, 3, 4, 6, 7, 8\\\hline
    \verb|STEM005__5C| & 1, 2, 5, 3, 4, 6, 7, 8 & \verb|STEM006__5C| & 1, 2, 5, 3, 4, 6, 7, 8 & \verb|STEM007__5C| & 1, 2, 5, 3, 4, 6, 7, 8\\\hline
    \verb|STEM008__5C| & 1, 2, 5, 3, 4, 6, 7, 8 & \verb|STEM009__5C| & 1, 2, 5, 3, 4, 6, 7, 8 & \verb|STEM010__5C| & 1, 2, 5, 3, 4, 6, 7, 8\\\hline
    \verb|STEM011__5C| & 1, 2, 5, 3, 4, 6, 7, 8 & \verb|STS_010__PO| & 4, 8 & \verb|STS_010__PZ| & 1, 5\\\hline
    \verb|STS_080__PO| & 3, 7 & \verb|STS_179N_HM| & 2, 6 & \verb|STS_190__PO| & 1, 3, 5, 7\\\hline
    \verb|STS_190__PO2| & 1, 3, 5, 7 & \verb|STS_191__PO| & 2, 6 & \verb|STS_191__PO2| & 2, 6\\\hline
    \verb|STS_191__SC| & 1, 4, 5, 8 & \verb|STS_191__SC2| & 1, 4, 5, 8 & \verb|TEST001__PZ| & 3, 7\\\hline
    \verb|TEST002__PZ| & 3, 7 & \verb|THEA001A_PO| & 2, 3, 6, 7, 1, 5, 4, 8 & \verb|THEA001D_PO| & 2, 4, 6, 8\\\hline
    \verb|THEA001G_PO| & 2, 3, 6, 7 & \verb|THEA002__PO| & 2, 3, 6, 7, 1, 5, 4, 8 & \verb|THEA006__PO| & 3, 7\\\hline
    \verb|THEA010__PO| & 2, 6 & \verb|THEA012__PO| & 2, 3, 6, 7, 1, 5, 4, 8 & \verb|THEA017__PO| & 2, 1, 6, 5, 4, 8\\\hline
    \verb|THEA021__PO| & 3, 7 & \verb|THEA022__PO| & 2, 4, 6, 8 & \verb|THEA023__PO| & 1, 3, 5, 7\\\hline
    \verb|THEA026__PO| & 2, 4, 6, 8 & \verb|THEA030__PO| & 1, 3, 5, 7 & \verb|THEA031__PO| & 2, 4, 6, 8\\\hline
    \verb|THEA041__PO| & 1, 3, 5, 7 & \verb|THEA051C_PO| & 2, 3, 6, 7, 1, 5, 4, 8 & \verb|THEA051H_PO| & 2, 3, 6, 7, 1, 5, 4, 8\\\hline
    \verb|THEA052C_PO| & 2, 3, 6, 7, 1, 5, 4, 8 & \verb|THEA052H_PO| & 2, 3, 6, 7, 1, 5, 4, 8 & \verb|THEA053HGPO| & 2, 6\\\hline
    \verb|THEA054C_PO| & 1, 3, 5, 7 & \verb|THEA056C_PO| & 1, 5 & \verb|THEA061__PO| & 2, 6\\\hline
    \verb|THEA080__PO| & 1, 3, 5, 7 & \verb|THEA081__PO| & 2, 4, 6, 8 & \verb|THEA082__PO| & 1, 3, 5, 7\\\hline
    \verb|THEA083__PO| & 2, 6 & \verb|THEA084__PO| & 4, 8 & \verb|THEA085__PO| & 2, 4, 6, 8\\\hline
    \verb|THEA089C_PO| & 1, 3, 5, 7 & \verb|THEA100C_PO| & 2, 6 & \verb|THEA100G_PO| & 3, 7\\\hline
    \verb|THEA100K_PO| & 2, 4, 6, 8 & \verb|THEA100S_PO| & 1, 5 & \verb|THEA115O_PO| & 4, 8\\\hline
    \verb|THEA130__PO| & 1, 3, 5, 7 & \verb|THEA132__PO| & 2, 4, 6, 8 & \verb|THEA141__PO| & 1, 3, 5, 7\\\hline
    \verb|THEA142__PO| & 4, 8 & \verb|THEA170__PO| & 1, 3, 5, 7 & \verb|THEA171__PO| & 2, 6\\\hline
    \verb|THEA172__PO| & 4, 8 & \verb|THEA187__PO| & 1, 5 & \verb|THEA188__PO| & 2, 4, 6, 8\\\hline
    \verb|THEA190__PO| & 1, 3, 5, 7 & \verb|THEA190__PO2| & 1, 3, 5, 7 & \verb|THEA191H_PO| & 1, 4, 5, 8\\\hline
    \verb|THEA191H_PO2| & 1, 4, 5, 8 & \verb|THEA192H_PO| & 1, 4, 5, 8 & \verb|THEA192H_PO2| & 1, 4, 5, 8\\\hline
    \verb|THES191D_SC| & 1, 3, 5, 7, 4, 8 & \verb|THES191D_SC2| & 1, 3, 5, 7, 4, 8 & \verb|THES192D_SC| & 1, 3, 5, 7, 4, 8\\\hline
    \verb|THES192D_SC2| & 1, 3, 5, 7, 4, 8 & \verb|TNDY409D_CG| & 4, 8 & \verb|WGS_321__CG| & 1, 5\\\hline
    \verb|WRIT001E_HM| & 2, 4, 6, 8 & \verb|WRIT001__HM| & 1 & \verb|WRIT015__PZ| & 3, 7\\\hline
    \verb|WRIT022__PZ| & 4, 8 & \verb|WRIT100A_PZ| & 1, 3, 5, 7 & \verb|WRIT100B_PZ| & 1, 3, 5, 7\\\hline
    \verb|WRIT101__PZ| & 2, 6 & \verb|WRIT105__SC| & 3, 7 & \verb|WRIT109__SC| & 2, 4, 6, 8\\\hline
    \verb|WRIT112__SC| & 2, 4, 6, 8 & \verb|WRIT122__SC| & 3, 7 & \verb|WRIT125__PZ| & 4, 8\\\hline
    \verb|WRIT132__SC| & 3, 7 & \verb|WRIT137__SC| & 1, 5 & \verb|WRIT140__SC| & 2, 1, 6, 5\\\hline
    \verb|WRIT146__SC| & 4, 8 & \verb|WRIT153IOSC| & 3, 7 & \verb|WRIT155__SC| & 3, 7\\\hline
    \verb|WRIT160__SC| & 1, 3, 5, 7 & \verb|WRIT175__SC| & 4, 8 & \verb|WRIT190__PZ| & 4, 8\\\hline
    \verb|WRIT190__PZ2| & 4, 8 & \verb|WRIT191__SC| & 2, 1, 6, 5, 4, 8 & \verb|WRIT191__SC2| & 2, 1, 6, 5, 4, 8\\\hline
    \verb|WRIT197N_SC| & 2, 6 & \verb|WRIT197N_SC2| & 2, 6 & \verb|WRIT197Q_SC| & 4, 8\\\hline
    \verb|WRIT197Q_SC2| & 4, 8 & \verb|XGOV192__CM| & 2, 6 & \verb|XGOV192__CM2| & 2, 6\\\hline
    \caption{Course offerings. Semesters are labeled sequentially, with semester 1 indicating the fall semester of a student's first year.}
    \label{tab:offerings}
\end{longtable}
\begin{longtable}{|p{0.5\textwidth}|c|}
    \hline
    Requirement Group Name & Required Groups\\\hline
    \verb|AFRI190__AF2| & 3\\\hline
    \verb|AFRI191C_AF2| & 3\\\hline
    \verb|AFRI191__AF2| & 3\\\hline
    \verb|AFRI195D_AF2| & 3\\\hline
    \verb|AMST190__PO2| & 3\\\hline
    \verb|AMST191__PO2| & 3\\\hline
    \verb|AMST191__SC2| & 3\\\hline
    \verb|ANTH190__PO2| & 3\\\hline
    \verb|ANTH190__SC2| & 3\\\hline
    \verb|ANTH191__SC2| & 3\\\hline
    \verb|ARCN191__SC2| & 3\\\hline
    \verb|ARHI190__PO2| & 3\\\hline
    \verb|ARHI191__PO2| & 3\\\hline
    \verb|ARHI191__SC2| & 3\\\hline
    \verb|ART_190__PO2| & 3\\\hline
    \verb|ART_192__PO2| & 3\\\hline
    \verb|ART_192__SC2| & 3\\\hline
    \verb|ART_193__SC2| & 3\\\hline
    \verb|ART_199__PZ2| & 3\\\hline
    \verb|ASAM190A_PO2| & 3\\\hline
    \verb|ASAM191__PO2| & 3\\\hline
    \verb|ASIA190__PO2| & 3\\\hline
    \verb|ASIA191__PO2| & 3\\\hline
    \verb|ASIA192__PO2| & 3\\\hline
    \verb|ASTR062__HM| & 3\\\hline
    \verb|ASTR122__HM| & 6\\\hline
    \verb|ASTR128__HM| & 10\\\hline
    \verb|BIOL023__HM| & 3\\\hline
    \verb|BIOL046__HM| & 2\\\hline
    \verb|BIOL046__HM_0| & 0.1\\\hline
    \verb|BIOL054__HM| & 4\\\hline
    \verb|BIOL101__HM| & 3\\\hline
    \verb|BIOL103__HM| & 4\\\hline
    \verb|BIOL108__HM| & 7\\\hline
    \verb|BIOL109__HM| & 7\\\hline
    \verb|BIOL110__HM| & 2\\\hline
    \verb|BIOL110__HM_co| & 3\\\hline
    \verb|BIOL110__HM_pre| & 4\\\hline
    \verb|BIOL111__HM| & 2\\\hline
    \verb|BIOL111__HM_co| & 3\\\hline
    \verb|BIOL111__HM_pre| & 4\\\hline
    \verb|BIOL113__HM| & 7\\\hline
    \verb|BIOL119__HM| & 1.5\\\hline
    \verb|BIOL121__HM| & 3\\\hline
    \verb|BIOL122__HM| & 3\\\hline
    \verb|BIOL154__HM| & 2\\\hline
    \verb|BIOL154__HM_co| & 3\\\hline
    \verb|BIOL154__HM_pre| & 0.1\\\hline
    \verb|BIOL160__HM| & 3\\\hline
    \verb|BIOL182__HM| & 3\\\hline
    \verb|BIOL184__HM| & 0.1\\\hline
    \verb|BIOL188__HM| & 1.5\\\hline
    \verb|BIOL189__HM| & 3\\\hline
    \verb|BIOL190L_KS2| & 3\\\hline
    \verb|BIOL190__PO2| & 1.5\\\hline
    \verb|BIOL191F_PO2| & 3\\\hline
    \verb|BIOL191H_PO2| & 1.5\\\hline
    \verb|BIOL191__HM2| & 0.1\\\hline
    \verb|BIOL191__HM3| & 0.1\\\hline
    \verb|BIOL191__HM4| & 0.1\\\hline
    \verb|BIOL191__KS2| & 3\\\hline
    \verb|BIOL193__HM2| & 3\\\hline
    \verb|BIOL194A_PO2| & 3\\\hline
    \verb|BIOL194B_PO2| & 3\\\hline
    \verb|BIOL195__HM| & 1\\\hline
    \verb|BIOL195__HM2| & 6\\\hline
    \verb|BIOL196__KS2| & 0.75\\\hline
    \verb|BIOL197__HM2| & 1\\\hline
    \verb|BIOL197__KS2| & 1.5\\\hline
    \verb|CGS_190__PZ2| & 3\\\hline
    \verb|CGS_199__PZ2| & 3\\\hline
    \verb|CHEM024__HM| & 4\\\hline
    \verb|CHEM042__HM| & 1\\\hline
    \verb|CHEM048__HM| & 2\\\hline
    \verb|CHEM048__HM_co| & 4\\\hline
    \verb|CHEM048__HM_pre| & 3\\\hline
    \verb|CHEM051__HM| & 5\\\hline
    \verb|CHEM053__HM| & 3\\\hline
    \verb|CHEM056__HM| & 5\\\hline
    \verb|CHEM058__HM| & 2\\\hline
    \verb|CHEM058__HM_co| & 3\\\hline
    \verb|CHEM058__HM_pre| & 1\\\hline
    \verb|CHEM080__HM| & 8\\\hline
    \verb|CHEM080__HM_0| & 0.1\\\hline
    \verb|CHEM103__HM| & 5\\\hline
    \verb|CHEM104__HM| & 3\\\hline
    \verb|CHEM105__HM| & 3\\\hline
    \verb|CHEM109__HM| & 3\\\hline
    \verb|CHEM110__HM| & 2\\\hline
    \verb|CHEM110__HM_co| & 3\\\hline
    \verb|CHEM110__HM_pre| & 1\\\hline
    \verb|CHEM111__HM| & 2\\\hline
    \verb|CHEM111__HM_co| & 3\\\hline
    \verb|CHEM111__HM_pre| & 1\\\hline
    \verb|CHEM112__HM| & 2\\\hline
    \verb|CHEM112__HM_co| & 3\\\hline
    \verb|CHEM112__HM_pre| & 4\\\hline
    \verb|CHEM114__HM| & 3\\\hline
    \verb|CHEM116__HM| & 3\\\hline
    \verb|CHEM122__HM| & 0.1\\\hline
    \verb|CHEM170__HM| & 6\\\hline
    \verb|CHEM171__HM| & 6\\\hline
    \verb|CHEM182__HM| & 3\\\hline
    \verb|CHEM184__HM| & 0.1\\\hline
    \verb|CHEM190L_KS2| & 3\\\hline
    \verb|CHEM191__KS2| & 3\\\hline
    \verb|CHEM191__PO2| & 1.5\\\hline
    \verb|CHEM193M_HM2| & 2\\\hline
    \verb|CHEM193P_HM2| & 2\\\hline
    \verb|CHEM194__HM| & 0.1\\\hline
    \verb|CHEM194__HM2| & 2\\\hline
    \verb|CHEM194__PO2| & 1.5\\\hline
    \verb|CHEM196__KS2| & 0.75\\\hline
    \verb|CHEM197__HM2| & 3\\\hline
    \verb|CHEM197__KS2| & 1.5\\\hline
    \verb|CHEM198__HM2| & 3\\\hline
    \verb|CHEM199__HM2| & 0.1\\\hline
    \verb|CHEM199__HM3| & 0.1\\\hline
    \verb|CHEM199__HM4| & 0.1\\\hline
    \verb|CHIN192A_PO2| & 1.5\\\hline
    \verb|CHST190__CH2| & 3\\\hline
    \verb|CHST191__CH2| & 3\\\hline
    \verb|CHST192__CH2| & 3\\\hline
    \verb|CLAS190__PO2| & 3\\\hline
    \verb|CLAS191__PO2| & 3\\\hline
    \verb|CLAS191__SC2| & 3\\\hline
    \verb|CLAS192__PO2| & 3\\\hline
    \verb|CLES197__HM2| & 1\\\hline
    \verb|CLES199__HM2| & 0.1\\\hline
    \verb|CSCI035__HM| & 0.1\\\hline
    \verb|CSCI060__HM| & 0.1\\\hline
    \verb|CSCI070__HM| & 0.1\\\hline
    \verb|CSCI081__HM| & 4\\\hline
    \verb|CSCI081__HM_0| & 0.1\\\hline
    \verb|CSCI081__HM_1| & 0.1\\\hline
    \verb|CSCI081__HM_2| & 0.1\\\hline
    \verb|CSCI081__HM_3| & 0.1\\\hline
    \verb|CSCI105__HM| & 3\\\hline
    \verb|CSCI123__HM| & 3\\\hline
    \verb|CSCI131__HM| & 6\\\hline
    \verb|CSCI133__HM| & 6\\\hline
    \verb|CSCI134__HM| & 3\\\hline
    \verb|CSCI140__HM| & 0.1\\\hline
    \verb|CSCI140__HM_0| & 0.1\\\hline
    \verb|CSCI140__HM_0_0| & 6\\\hline
    \verb|CSCI140__HM_0_0_0| & 0.1\\\hline
    \verb|CSCI140__HM_0_0_1| & 0.1\\\hline
    \verb|CSCI140__HM_0_0_2| & 0.1\\\hline
    \verb|CSCI140__HM_0_1| & 4\\\hline
    \verb|CSCI140__HM_0_1_0| & 0.1\\\hline
    \verb|CSCI140__HM_1| & 3\\\hline
    \verb|CSCI144__HM| & 7\\\hline
    \verb|CSCI144__HM_0| & 0.1\\\hline
    \verb|CSCI151__HM| & 4\\\hline
    \verb|CSCI151__HM_0| & 0.1\\\hline
    \verb|CSCI152__HM| & 7\\\hline
    \verb|CSCI152__HM_0| & 0.1\\\hline
    \verb|CSCI153__HM| & 3\\\hline
    \verb|CSCI158__HM| & 10\\\hline
    \verb|CSCI158__HM_0| & 0.1\\\hline
    \verb|CSCI159__HM| & 4\\\hline
    \verb|CSCI159__HM_0| & 0.1\\\hline
    \verb|CSCI183__HM| & 3\\\hline
    \verb|CSCI184__HM| & 3\\\hline
    \verb|CSCI190__PO2| & 0.1\\\hline
    \verb|CSCI191__PO2| & 3\\\hline
    \verb|CSCI192__PO2| & 1.5\\\hline
    \verb|CSCI195__HM2| & 0.1\\\hline
    \verb|CSCI195__HM3| & 0.1\\\hline
    \verb|CSCI195__HM4| & 0.1\\\hline
    \verb|CSCI198__SC2| & 3\\\hline
    \verb|CSMT184__HM| & 3\\\hline
    \verb|Core| & 32.5\\\hline
    \verb|Core_IntroCS| & 0.1\\\hline
    \verb|Core_Mechanics| & 0.1\\\hline
    \verb|DANC190__SC2| & 1.5\\\hline
    \verb|DANC191__SC2| & 3\\\hline
    \verb|DANC192__PO2| & 1.5\\\hline
    \verb|DANC193__SC2| & 0.1\\\hline
    \verb|DS__190__PO2| & 1.5\\\hline
    \verb|DS__191__SC2| & 3\\\hline
    \verb|EA__190L_KS2| & 3\\\hline
    \verb|EA__190__PO2| & 3\\\hline
    \verb|EA__191H_PO2| & 1.5\\\hline
    \verb|EA__191__KS2| & 3\\\hline
    \verb|EA__191__PO2| & 3\\\hline
    \verb|EA__196__KS2| & 0.75\\\hline
    \verb|EA__197__KS2| & 1.5\\\hline
    \verb|EA__197__PZ2| & 3\\\hline
    \verb|ECON190__PO2| & 3\\\hline
    \verb|ECON191H_SC2| & 3\\\hline
    \verb|ECON191__CM2| & 3\\\hline
    \verb|ECON191__SC2| & 3\\\hline
    \verb|ECON194A_CM2| & 1.5\\\hline
    \verb|ECON194B_CM2| & 1.5\\\hline
    \verb|ECON195__PO2| & 0.1\\\hline
    \verb|ECON196__CM2| & 3\\\hline
    \verb|ECON197SBCM2| & 1.5\\\hline
    \verb|ECON197S_CM2| & 3\\\hline
    \verb|ECON198__PZ2| & 3\\\hline
    \verb|ECON199__CM2| & 1.5\\\hline
    \verb|ECON199__PZ2| & 0.1\\\hline
    \verb|EEP_191__SC2| & 3\\\hline
    \verb|ENGL190__SC2| & 3\\\hline
    \verb|ENGL191__PO2| & 1.5\\\hline
    \verb|ENGL191__SC2| & 3\\\hline
    \verb|ENGL192__SC2| & 3\\\hline
    \verb|ENGL193__SC2| & 3\\\hline
    \verb|ENGL194__PZ2| & 3\\\hline
    \verb|ENGL194__SC2| & 3\\\hline
    \verb|ENGL195B_PO2| & 1.5\\\hline
    \verb|ENGL195__PO2| & 3\\\hline
    \verb|ENGL195__SC2| & 3\\\hline
    \verb|ENGL197N_SC2| & 3\\\hline
    \verb|ENGL199T_SC2| & 3\\\hline
    \verb|ENGR004__HM| & 2\\\hline
    \verb|ENGR004__HM_co| & 4\\\hline
    \verb|ENGR004__HM_pre| & 2\\\hline
    \verb|ENGR072__HM| & 2\\\hline
    \verb|ENGR072__HM_co| & 3\\\hline
    \verb|ENGR072__HM_pre| & 11\\\hline
    \verb|ENGR079__HM| & 4\\\hline
    \verb|ENGR080__HM| & 2\\\hline
    \verb|ENGR080__HM_co| & 6\\\hline
    \verb|ENGR080__HM_pre| & 4\\\hline
    \verb|ENGR082__HM| & 4\\\hline
    \verb|ENGR083__HM| & 8\\\hline
    \verb|ENGR084__HM| & 4\\\hline
    \verb|ENGR085A_HM| & 0.1\\\hline
    \verb|ENGR085__HM| & 0.1\\\hline
    \verb|ENGR086__HM| & 15\\\hline
    \verb|ENGR101__HM| & 10\\\hline
    \verb|ENGR102__HM| & 3\\\hline
    \verb|ENGR111__HM| & 3\\\hline
    \verb|ENGR112__HM| & 0.1\\\hline
    \verb|ENGR113__HM| & 0.1\\\hline
    \verb|ENGR131__HM| & 3\\\hline
    \verb|ENGR132__HM| & 3\\\hline
    \verb|ENGR133__HM| & 3\\\hline
    \verb|ENGR134__HM| & 3\\\hline
    \verb|ENGR138__HM| & 3\\\hline
    \verb|ENGR151__HM| & 7\\\hline
    \verb|ENGR154__HM| & 0.1\\\hline
    \verb|ENGR154__HM_0| & 4.5\\\hline
    \verb|ENGR155__HM| & 0.1\\\hline
    \verb|ENGR155__HM_0| & 4.5\\\hline
    \verb|ENGR157__HM| & 2\\\hline
    \verb|ENGR157__HM_co| & 3\\\hline
    \verb|ENGR157__HM_pre| & 3\\\hline
    \verb|ENGR164__HM| & 7\\\hline
    \verb|ENGR171__HM| & 3\\\hline
    \verb|ENGR175__HM| & 3\\\hline
    \verb|ENGR177__HM| & 7\\\hline
    \verb|ENGR178__HM| & 3\\\hline
    \verb|ENGR183__HM| & 4\\\hline
    \verb|ENGR185A_HM| & 4\\\hline
    \verb|ENGR185B_HM| & 1.5\\\hline
    \verb|ENGR186__HM| & 3\\\hline
    \verb|ENGR187__HM| & 3\\\hline
    \verb|ENGR190BAHM2| & 3\\\hline
    \verb|ENGR190BFHM2| & 1\\\hline
    \verb|ENGR190BGHM2| & 3\\\hline
    \verb|ENGR190BHHM2| & 3\\\hline
    \verb|ENGR190BIHM2| & 3\\\hline
    \verb|ENGR190BJHM2| & 3\\\hline
    \verb|ENGR191__HM2| & 3\\\hline
    \verb|ENGR205__HM| & 3\\\hline
    \verb|ENGR207__HM| & 0.1\\\hline
    \verb|ENGR208__HM| & 0.1\\\hline
    \verb|ENGR278__HM| & 3\\\hline
    \verb|FGSS191__SC2| & 3\\\hline
    \verb|FGSS192__SC2| & 3\\\hline
    \verb|FGSS198__SC2| & 3\\\hline
    \verb|FLAN191__SC2| & 3\\\hline
    \verb|FREN191__PO2| & 1.5\\\hline
    \verb|FREN191__SC2| & 3\\\hline
    \verb|FREN192__PO2| & 1.5\\\hline
    \verb|FREN193__PO2| & 0.1\\\hline
    \verb|GEOL192__PO2| & 1.5\\\hline
    \verb|GERM191__PO2| & 1.5\\\hline
    \verb|GERM191__SC2| & 3\\\hline
    \verb|GERM193__PO2| & 1.5\\\hline
    \verb|GLAS194A_PZ2| & 1.5\\\hline
    \verb|GLAS194B_PZ2| & 1.5\\\hline
    \verb|GOVT191__CM2| & 3\\\hline
    \verb|GOVT192__CM2| & 3\\\hline
    \verb|GWS_190__PO2| & 3\\\hline
    \verb|GWS_191__PO2| & 3\\\hline
    \verb|HIST190__CM2| & 3\\\hline
    \verb|HIST190__PO2| & 3\\\hline
    \verb|HIST191__PO2| & 3\\\hline
    \verb|HIST191__SC2| & 3\\\hline
    \verb|HIST192__PO2| & 3\\\hline
    \verb|HIST192__SC2| & 3\\\hline
    \verb|HIST196__CM2| & 3\\\hline
    \verb|HIST197__CM2| & 3\\\hline
    \verb|HIST197__PZ2| & 3\\\hline
    \verb|HIST199__PZ2| & 3\\\hline
    \verb|HMSC191__SC2| & 3\\\hline
    \verb|HMSC192__SC2| & 3\\\hline
    \verb|HSAs| & 3\\\hline
    \verb|HSAs_MuddHums| & 12\\\hline
    \verb|HSAs_TenHSAs| & 33\\\hline
    \verb|HSAs_WritingIntensive| & 0.1\\\hline
    \verb|HUM_195J_SC2| & 3\\\hline
    \verb|HUM_196__PO2| & 1.5\\\hline
    \verb|ID__191D_SC2| & 3\\\hline
    \verb|ID__191H_SC2| & 3\\\hline
    \verb|ID__191S_SC2| & 3\\\hline
    \verb|ID__199P1PO2| & 1.5\\\hline
    \verb|ID__199P2PO2| & 1.5\\\hline
    \verb|ID__199P3PO2| & 1.5\\\hline
    \verb|ID__199P4PO2| & 1.5\\\hline
    \verb|ID__199S1PO2| & 1.5\\\hline
    \verb|ID__199S2PO2| & 1.5\\\hline
    \verb|ID__199S3PO2| & 1.5\\\hline
    \verb|ID__199S4PO2| & 1.5\\\hline
    \verb|IR__190__PO2| & 3\\\hline
    \verb|IR__191__PO2| & 3\\\hline
    \verb|IR__192__SC2| & 3\\\hline
    \verb|ITAL191__SC2| & 3\\\hline
    \verb|LAMS191__PO2| & 3\\\hline
    \verb|LAST190__PO2| & 3\\\hline
    \verb|LAST191__PO2| & 3\\\hline
    \verb|LAST191__SC2| & 3\\\hline
    \verb|LGCS190__PO2| & 3\\\hline
    \verb|LGCS191__PO2| & 1.5\\\hline
    \verb|LIT_199__CM2| & 1.5\\\hline
    \verb|MATH055__HM| & 3\\\hline
    \verb|MATH056__HM| & 7\\\hline
    \verb|MATH062__HM| & 2\\\hline
    \verb|MATH062__HM_co| & 3\\\hline
    \verb|MATH062__HM_pre| & 4\\\hline
    \verb|MATH073__HM| & 0.1\\\hline
    \verb|MATH082__HM| & 7\\\hline
    \verb|MATH104__HM| & 6\\\hline
    \verb|MATH106__HM| & 3\\\hline
    \verb|MATH119__HM| & 1.5\\\hline
    \verb|MATH131__HM| & 3\\\hline
    \verb|MATH132__HM| & 3\\\hline
    \verb|MATH136__HM| & 6\\\hline
    \verb|MATH147__HM| & 3\\\hline
    \verb|MATH156__HM| & 4.5\\\hline
    \verb|MATH157__HM| & 0.1\\\hline
    \verb|MATH158__HM| & 6\\\hline
    \verb|MATH164__HM| & 9\\\hline
    \verb|MATH165__HM| & 6\\\hline
    \verb|MATH171__HM| & 6\\\hline
    \verb|MATH174__HM| & 3\\\hline
    \verb|MATH175__HM| & 3\\\hline
    \verb|MATH176__HM| & 3\\\hline
    \verb|MATH179__HM| & 5\\\hline
    \verb|MATH179__HM_0| & 0.1\\\hline
    \verb|MATH179__HM_1| & 0.1\\\hline
    \verb|MATH180__HM| & 6\\\hline
    \verb|MATH181__HM| & 2\\\hline
    \verb|MATH181__HM_co| & 3\\\hline
    \verb|MATH181__HM_pre| & 3\\\hline
    \verb|MATH187__HM| & 3\\\hline
    \verb|MATH190__CM2| & 1.5\\\hline
    \verb|MATH190__PO2| & 1.5\\\hline
    \verb|MATH191__CM2| & 1.5\\\hline
    \verb|MATH191__PO2| & 1.5\\\hline
    \verb|MATH191__SC2| & 3\\\hline
    \verb|MATH193__HM2| & 3\\\hline
    \verb|MATH194__PZ2| & 1.5\\\hline
    \verb|MATH195__CM2| & 3\\\hline
    \verb|MATH196__HM2| & 1\\\hline
    \verb|MATH197__HM2| & 3\\\hline
    \verb|MATH198__HM2| & 1\\\hline
    \verb|MATH198__PZ2| & 0.1\\\hline
    \verb|MCBI118A_HM| & 9\\\hline
    \verb|MCBI118B_HM| & 6\\\hline
    \verb|MCBI199__HM2| & 0.1\\\hline
    \verb|MCBI199__HM3| & 0.1\\\hline
    \verb|MCBI199__HM4| & 0.1\\\hline
    \verb|MCSI195__PZ2| & 3\\\hline
    \verb|MENA191__SC2| & 3\\\hline
    \verb|MOBI191A_PO2| & 1.5\\\hline
    \verb|MOBI191B_PO2| & 1.5\\\hline
    \verb|MOBI194A_PO2| & 3\\\hline
    \verb|MOBI194B_PO2| & 3\\\hline
    \verb|MS__190F_JT2| & 3\\\hline
    \verb|MS__190S_JT2| & 3\\\hline
    \verb|MS__190__JT2| & 3\\\hline
    \verb|MS__194__PZ2| & 3\\\hline
    \verb|MS__196__PZ2| & 3\\\hline
    \verb|MS__199__SC2| & 3\\\hline
    \verb|MUS_190__PO2| & 3\\\hline
    \verb|MUS_191__SC2| & 3\\\hline
    \verb|MUS_192F_PO2| & 3\\\hline
    \verb|MUS_192H_PO2| & 1.5\\\hline
    \verb|Maj_Bio| & 28\\\hline
    \verb|Maj_BioChem| & 29\\\hline
    \verb|Maj_BioChem_BioChem| & 0.1\\\hline
    \verb|Maj_BioChem_BioChemLab| & 0.1\\\hline
    \verb|Maj_BioChem_BioColloqium| & 0.2\\\hline
    \verb|Maj_BioChem_BioElective| & 3\\\hline
    \verb|Maj_BioChem_Capstone| & 0.1\\\hline
    \verb|Maj_BioChem_Capstone_BioResearch| & 12\\\hline
    \verb|Maj_BioChem_Capstone_BioThesis| & 6\\\hline
    \verb|Maj_BioChem_Capstone_ChemThesis| & 6\\\hline
    \verb|Maj_BioChem_Capstone_EngrClinic| & 6\\\hline
    \verb|Maj_BioChem_ChemColloqium| & 0.2\\\hline
    \verb|Maj_BioChem_ChemLab| & 2\\\hline
    \verb|Maj_BioChem_ChemWithoutCells| & 0.1\\\hline
    \verb|Maj_BioChem_OrganicBio| & 0.1\\\hline
    \verb|Maj_BioChem_Topics| & 0.1\\\hline
    \verb|Maj_BioClim| & 28\\\hline
    \verb|Maj_BioClim_BioClimateCapstone| & 0.1\\\hline
    \verb|Maj_BioClim_BioClimateCapstone_CSClinic| & 6\\\hline
    \verb|Maj_BioClim_BioClimateCapstone_EngClinic| & 6\\\hline
    \verb|Maj_BioClim_BioClimateCapstone_MathClinic| & 6\\\hline
    \verb|Maj_BioClim_BioClimateCapstone_PhysClinic| & 6\\\hline
    \verb|Maj_BioClim_BioClimateCapstone_Thesis| & 6\\\hline
    \verb|Maj_BioClim_BioClimateElective| & 0.1\\\hline
    \verb|Maj_BioClim_BioColloqium| & 0.2\\\hline
    \verb|Maj_BioClim_BioComputational| & 0.1\\\hline
    \verb|Maj_BioClim_BioElec| & 3\\\hline
    \verb|Maj_BioClim_BioGroup| & 0.1\\\hline
    \verb|Maj_BioClim_BioSeminar| & 0.1\\\hline
    \verb|Maj_BioClim_ClimateColloqium| & 0.2\\\hline
    \verb|Maj_BioClim_ClimateContexts| & 0.1\\\hline
    \verb|Maj_BioClim_ClimateContexts_Justice| & 3.0\\\hline
    \verb|Maj_BioClim_ClimateInterventions| & 0.1\\\hline
    \verb|Maj_BioClim_ClimateInterventions_ClimateInterventionCredits| & 3\\\hline
    \verb|Maj_BioClim_ClimateInterventions_Justice| & 3.0\\\hline
    \verb|Maj_BioClim_ClimateStats| & 0.1\\\hline
    \verb|Maj_BioClim_Thermodynamics| & 0.1\\\hline
    \verb|Maj_Bio_BioColloquium| & 0.2\\\hline
    \verb|Maj_Bio_BioElecs| & 26\\\hline
    \verb|Maj_Bio_BioLab| & 0.1\\\hline
    \verb|Maj_Bio_BioLab_NoBioChem| & 4\\\hline
    \verb|Maj_Bio_BioLab_WithBioChem| & 2\\\hline
    \verb|Maj_Bio_BioLab_WithBioChem_Other| & 0.1\\\hline
    \verb|Maj_Bio_BioSeminar| & 0.1\\\hline
    \verb|Maj_Bio_ThesisOrClinic| & 0.1\\\hline
    \verb|Maj_Bio_ThesisOrClinic_BioThesis| & 6\\\hline
    \verb|Maj_Bio_ThesisOrClinic_CSClinic| & 6\\\hline
    \verb|Maj_Bio_ThesisOrClinic_EngrClinic| & 6\\\hline
    \verb|Maj_Bio_ThesisOrClinic_IntensiveBioThesis| & 12\\\hline
    \verb|Maj_Bio_ThesisOrClinic_MathClinic| & 6\\\hline
    \verb|Maj_Bio_ThesisOrClinic_PhysClinic| & 6\\\hline
    \verb|Maj_CS| & 27\\\hline
    \verb|Maj_CSClimate| & 29\\\hline
    \verb|Maj_CSClimate_Areas| & 6\\\hline
    \verb|Maj_CSClimate_Areas_Computational| & 0.1\\\hline
    \verb|Maj_CSClimate_Areas_ProbStats| & 0.1\\\hline
    \verb|Maj_CSClimate_CSColloquium| & 0.2\\\hline
    \verb|Maj_CSClimate_CSElective| & 6\\\hline
    \verb|Maj_CSClimate_ClimateColloquium| & 0.2\\\hline
    \verb|Maj_CSClimate_ClimateContexts| & 0.1\\\hline
    \verb|Maj_CSClimate_ClimateImpacts| & 0.1\\\hline
    \verb|Maj_CSClimate_ClimateImpacts_ClimateJustice| & 3.0\\\hline
    \verb|Maj_CSClimate_ClimateInterventions| & 3\\\hline
    \verb|Maj_CSClimate_ClimateInterventions_ClimateJustice| & 3.0\\\hline
    \verb|Maj_CSClimate_HumanCentered| & 0.1\\\hline
    \verb|Maj_CSClimate_Principles| & 0.1\\\hline
    \verb|Maj_CSClimate_SysAlgs| & 0.1\\\hline
    \verb|Maj_CSClimate_Thermodynamics| & 0.1\\\hline
    \verb|Maj_CS_CSColloquium| & 0.2\\\hline
    \verb|Maj_CS_CSCore| & 0.1\\\hline
    \verb|Maj_CS_CSElec| & 9\\\hline
    \verb|Maj_Chem| & 23\\\hline
    \verb|Maj_ChemClimate| & 37\\\hline
    \verb|Maj_ChemClimate_ChemClass| & 6\\\hline
    \verb|Maj_ChemClimate_ChemClimateComputational| & 3\\\hline
    \verb|Maj_ChemClimate_ChemClimateProbStats| & 3\\\hline
    \verb|Maj_ChemClimate_ChemSeminar| & 0.2\\\hline
    \verb|Maj_ChemClimate_ChemUpperDiv| & 3\\\hline
    \verb|Maj_ChemClimate_ClimateColloqium| & 0.2\\\hline
    \verb|Maj_ChemClimate_ClimateContexts| & 0.1\\\hline
    \verb|Maj_ChemClimate_ClimateContexts_Justice| & 3.0\\\hline
    \verb|Maj_ChemClimate_ClimateInterventions| & 0.1\\\hline
    \verb|Maj_ChemClimate_ClimateInterventions_ClimateInterventionCredits| & 3\\\hline
    \verb|Maj_ChemClimate_ClimateInterventions_Justice| & 3.0\\\hline
    \verb|Maj_ChemClimate_Or1| & 0.1\\\hline
    \verb|Maj_Chem_CHEMElec| & 8\\\hline
    \verb|Maj_Chem_ChemLab| & 4\\\hline
    \verb|Maj_Chem_ChemSeminar| & 0.2\\\hline
    \verb|Maj_Chem_Math| & 0.1\\\hline
    \verb|Maj_Chem_ThesisOrClinic| & 0.1\\\hline
    \verb|Maj_Chem_ThesisOrClinic_Clinic| & 6\\\hline
    \verb|Maj_Chem_ThesisOrClinic_Thesis| & 6\\\hline
    \verb|Maj_CsPhys| & 37\\\hline
    \verb|Maj_CsPhys_CSColloqium| & 0.2\\\hline
    \verb|Maj_CsPhys_CSCoreElective| & 0.1\\\hline
    \verb|Maj_CsPhys_CSElective| & 6\\\hline
    \verb|Maj_CsPhys_Capstone| & 0.1\\\hline
    \verb|Maj_CsPhys_Capstone_CSMTClinic| & 6\\\hline
    \verb|Maj_CsPhys_Capstone_PhysClinic| & 6\\\hline
    \verb|Maj_CsPhys_Capstone_PhysThesis| & 6\\\hline
    \verb|Maj_CsPhys_IntroCS| & 0.1\\\hline
    \verb|Maj_CsPhys_Lab| & 0.1\\\hline
    \verb|Maj_CsPhys_PhysCSElective| & 9\\\hline
    \verb|Maj_CsPhys_PhysColloqium| & 0.2\\\hline
    \verb|Maj_CsPhys_Quantum| & 0.1\\\hline
    \verb|Maj_Engr| & 48.2\\\hline
    \verb|Maj_Engr_EngrElec| & 9\\\hline
    \verb|Maj_Math| & 22\\\hline
    \verb|Maj_MathCS| & 32.2\\\hline
    \verb|Maj_MathCS_Algs| & 0.1\\\hline
    \verb|Maj_MathCS_CSElec| & 6\\\hline
    \verb|Maj_MathCS_IntroCS| & 0.1\\\hline
    \verb|Maj_MathCS_MathElec| & 6\\\hline
    \verb|Maj_MathCompBio| & 23.1\\\hline
    \verb|Maj_MathCompBio_AdvancedBio| & 0.1\\\hline
    \verb|Maj_MathCompBio_BioColloquium| & 0.2\\\hline
    \verb|Maj_MathCompBio_BioFoundations| & 6\\\hline
    \verb|Maj_MathCompBio_BioLab| & 0.1\\\hline
    \verb|Maj_MathCompBio_BioSeminar| & 0.1\\\hline
    \verb|Maj_MathCompBio_CSElec| & 3\\\hline
    \verb|Maj_MathCompBio_MathElec| & 3\\\hline
    \verb|Maj_MathCompBio_MathOrCSElec| & 11\\\hline
    \verb|Maj_MathCompBio_ThesisOrClinic| & 0.1\\\hline
    \verb|Maj_MathCompBio_ThesisOrClinic_BioThesis| & 6\\\hline
    \verb|Maj_MathCompBio_ThesisOrClinic_CSClinic| & 6\\\hline
    \verb|Maj_MathCompBio_ThesisOrClinic_IntensiveBioThesis| & 12\\\hline
    \verb|Maj_MathCompBio_ThesisOrClinic_MathClinic| & 6\\\hline
    \verb|Maj_MathCompBio_ThesisOrClinic_MathThesis| & 6\\\hline
    \verb|Maj_MathPhys| & 39.5\\\hline
    \verb|Maj_MathPhys_Fields| & 0.1\\\hline
    \verb|Maj_MathPhys_PhysColloquium| & 0.2\\\hline
    \verb|Maj_MathPhys_ScientificComp| & 0.1\\\hline
    \verb|Maj_MathPhys_ThesisOrClinic| & 0.1\\\hline
    \verb|Maj_MathPhys_ThesisOrClinic_MathClinic| & 6\\\hline
    \verb|Maj_MathPhys_ThesisOrClinic_MathThesis| & 6\\\hline
    \verb|Maj_MathPhys_ThesisOrClinic_PhysClinic| & 6\\\hline
    \verb|Maj_MathPhys_ThesisOrClinic_PhysThesis| & 6\\\hline
    \verb|Maj_Math_Computational| & 0.1\\\hline
    \verb|Maj_Math_MathElec| & 7\\\hline
    \verb|Maj_Math_ThesisOrClinic| & 0.1\\\hline
    \verb|Maj_Math_ThesisOrClinic_BioThesis| & 6\\\hline
    \verb|Maj_Math_ThesisOrClinic_MathThesis| & 6\\\hline
    \verb|Maj_MolBio| & 35\\\hline
    \verb|Maj_MolBio_BioElecs| & 26\\\hline
    \verb|Maj_MolBio_BioLab| & 0.1\\\hline
    \verb|Maj_MolBio_Capstone| & 0.1\\\hline
    \verb|Maj_MolBio_Capstone_BioThesis| & 6\\\hline
    \verb|Maj_MolBio_Capstone_CSClinic| & 6\\\hline
    \verb|Maj_MolBio_Capstone_ChemThesis| & 6\\\hline
    \verb|Maj_MolBio_Capstone_EngrClinic| & 6\\\hline
    \verb|Maj_MolBio_Capstone_IntensiveResearch| & 12\\\hline
    \verb|Maj_MolBio_Capstone_MathClinic| & 6\\\hline
    \verb|Maj_MolBio_Capstone_PhysClinic| & 6\\\hline
    \verb|Maj_MolBio_Colloquium| & 0.4\\\hline
    \verb|Maj_MolBio_MolBio| & 0.1\\\hline
    \verb|Maj_MolBio_SeminarCourse| & 0.1\\\hline
    \verb|Maj_Phys| & 27\\\hline
    \verb|Maj_Phys_PhysColloquium| & 0.2\\\hline
    \verb|Maj_Phys_Thesis| & 0.1\\\hline
    \verb|Maj_Phys_Thesis_Clinic| & 6\\\hline
    \verb|Maj_Phys_Thesis_PhysThesis| & 3\\\hline
    \verb|NEUR190L_KS2| & 3\\\hline
    \verb|NEUR190__PO2| & 1.5\\\hline
    \verb|NEUR191__KS2| & 3\\\hline
    \verb|NEUR191__PO2| & 1.5\\\hline
    \verb|NEUR192__PO2| & 3\\\hline
    \verb|NEUR194A_PO2| & 1.5\\\hline
    \verb|NEUR194B_PO2| & 3\\\hline
    \verb|NEUR196__KS2| & 0.75\\\hline
    \verb|NEUR197__KS2| & 1.5\\\hline
    \verb|ORST191__SC2| & 3\\\hline
    \verb|ORST192__PZ2| & 3\\\hline
    \verb|ORST198C_PZ2| & 3\\\hline
    \verb|ORST198J_PZ2| & 3\\\hline
    \verb|ORST198M_PZ2| & 3\\\hline
    \verb|PHIL190__CM2| & 3\\\hline
    \verb|PHIL190__PO2| & 3\\\hline
    \verb|PHIL190__SC2| & 3\\\hline
    \verb|PHIL191__PO2| & 3\\\hline
    \verb|PHIL191__PZ2| & 1.5\\\hline
    \verb|PHIL191__SC2| & 3\\\hline
    \verb|PHIL192__CM2| & 3\\\hline
    \verb|PHIL198__CM2| & 3\\\hline
    \verb|PHIL198__SC2| & 1.5\\\hline
    \verb|PHYS050__HM| & 4\\\hline
    \verb|PHYS051__HM| & 2\\\hline
    \verb|PHYS051__HM_co| & 0.1\\\hline
    \verb|PHYS051__HM_pre| & 5.5\\\hline
    \verb|PHYS052__HM| & 4\\\hline
    \verb|PHYS052__HM_0| & 0.1\\\hline
    \verb|PHYS054__HM| & 4\\\hline
    \verb|PHYS064__HM| & 2\\\hline
    \verb|PHYS064__HM_0| & 0.1\\\hline
    \verb|PHYS064__HM_1| & 0.1\\\hline
    \verb|PHYS084__HM| & 8\\\hline
    \verb|PHYS084__HM_0| & 0.1\\\hline
    \verb|PHYS111__HM| & 6.5\\\hline
    \verb|PHYS111__HM_0| & 0.1\\\hline
    \verb|PHYS116__HM| & 3\\\hline
    \verb|PHYS117__HM| & 3\\\hline
    \verb|PHYS133__HM| & 1\\\hline
    \verb|PHYS134__HM| & 4\\\hline
    \verb|PHYS151__HM| & 5\\\hline
    \verb|PHYS151__HM_0| & 0.1\\\hline
    \verb|PHYS151__HM_1| & 0.1\\\hline
    \verb|PHYS154__HM| & 0.1\\\hline
    \verb|PHYS156__HM| & 4\\\hline
    \verb|PHYS156__HM_0| & 0.1\\\hline
    \verb|PHYS161__HM| & 3\\\hline
    \verb|PHYS162__HM| & 2\\\hline
    \verb|PHYS162__HM_co| & 3\\\hline
    \verb|PHYS162__HM_pre| & 3\\\hline
    \verb|PHYS164__HM| & 3\\\hline
    \verb|PHYS166__HM| & 5.5\\\hline
    \verb|PHYS170X_HM| & 9\\\hline
    \verb|PHYS170__HM| & 9\\\hline
    \verb|PHYS172__HM| & 3\\\hline
    \verb|PHYS181__HM| & 2\\\hline
    \verb|PHYS183__HM| & 0.1\\\hline
    \verb|PHYS190L_KS2| & 3\\\hline
    \verb|PHYS190__PO2| & 3\\\hline
    \verb|PHYS191E_PO2| & 3\\\hline
    \verb|PHYS191L_PO2| & 1.5\\\hline
    \verb|PHYS191__HM2| & 3\\\hline
    \verb|PHYS191__KS2| & 3\\\hline
    \verb|PHYS193__HM2| & 3\\\hline
    \verb|PHYS193__PO2| & 0.1\\\hline
    \verb|PHYS194__HM2| & 3\\\hline
    \verb|PHYS195__HM2| & 0.1\\\hline
    \verb|PHYS195__HM3| & 0.1\\\hline
    \verb|PHYS195__HM4| & 0.1\\\hline
    \verb|PHYS196__KS2| & 0.75\\\hline
    \verb|PHYS197__HM2| & 3\\\hline
    \verb|PHYS197__KS2| & 1.5\\\hline
    \verb|PHYS199__HM2| & 3\\\hline
    \verb|POLI190C_PO2| & 3\\\hline
    \verb|POLI190D_PO2| & 3\\\hline
    \verb|POLI190E_PO2| & 3\\\hline
    \verb|POLI191__PO2| & 3\\\hline
    \verb|POLI191__SC2| & 3\\\hline
    \verb|POLI193__PO2| & 1.5\\\hline
    \verb|POLI195__PO2| & 0.1\\\hline
    \verb|POST195__PZ2| & 3\\\hline
    \verb|POST199__PZ2| & 3\\\hline
    \verb|PPA_190__PO2| & 3\\\hline
    \verb|PPA_191__PO2| & 3\\\hline
    \verb|PPA_191__SC2| & 3\\\hline
    \verb|PPA_192__SC2| & 3\\\hline
    \verb|PPA_195__PO2| & 3\\\hline
    \verb|PPE_195__PO2| & 3\\\hline
    \verb|PSYC190A_PO2| & 1.5\\\hline
    \verb|PSYC190R_PO2| & 1.5\\\hline
    \verb|PSYC190__PO2| & 1.5\\\hline
    \verb|PSYC191__PO2| & 3\\\hline
    \verb|PSYC191__PZ2| & 1.5\\\hline
    \verb|PSYC191__SC2| & 3\\\hline
    \verb|PSYC193__SC2| & 3\\\hline
    \verb|PSYC194__PZ2| & 3\\\hline
    \verb|PSYC197A_CM2| & 0.75\\\hline
    \verb|PSYC197B_CM2| & 0.1\\\hline
    \verb|PSYC197__PZ2| & 3\\\hline
    \verb|PSYC198__CM2| & 3\\\hline
    \verb|RLIT191__PO2| & 1.5\\\hline
    \verb|RLST190__PO2| & 3\\\hline
    \verb|RLST191__PO2| & 3\\\hline
    \verb|RLST191__SC2| & 3\\\hline
    \verb|RUSS191__PO2| & 1.5\\\hline
    \verb|RUSS191__SC2| & 1.5\\\hline
    \verb|RUST191__PO2| & 1.5\\\hline
    \verb|SCI_196__CM2| & 0.75\\\hline
    \verb|SCI_197__CM2| & 1.5\\\hline
    \verb|SOC_190__PO2| & 3\\\hline
    \verb|SOC_191__PO2| & 3\\\hline
    \verb|SOC_199A_PZ2| & 3\\\hline
    \verb|SOC_199B_PZ2| & 1.5\\\hline
    \verb|SOC_199C_PZ2| & 3\\\hline
    \verb|SPAN191__PO2| & 1.5\\\hline
    \verb|SPAN191__SC2| & 3\\\hline
    \verb|SPAN192__PO2| & 1.5\\\hline
    \verb|SPAN199__PZ2| & 3\\\hline
    \verb|STS_190__PO2| & 3\\\hline
    \verb|STS_191__PO2| & 3\\\hline
    \verb|STS_191__SC2| & 3\\\hline
    \verb|THEA190__PO2| & 1.5\\\hline
    \verb|THEA191H_PO2| & 1.5\\\hline
    \verb|THEA192H_PO2| & 1.5\\\hline
    \verb|THES191D_SC2| & 3\\\hline
    \verb|THES192D_SC2| & 3\\\hline
    \verb|WRIT190__PZ2| & 3\\\hline
    \verb|WRIT191__SC2| & 3\\\hline
    \verb|WRIT197N_SC2| & 3\\\hline
    \verb|WRIT197Q_SC2| & 3\\\hline
    \verb|XGOV192__CM2| & 3\\\hline
    \caption{Number of requirements and credits needed to satisfy each requirement group.}
    \label{tab:reqs-nums}
\end{longtable}


\begin{landscape}
\begin{longtable}{|p{4.4in}|p{4.4in}|}
    \hline
    Requirement Group Name & Required Groups\\\hline
    \verb|BIOL046__HM| & \verb|BIOL046__HM_0|\\\hline
    \verb|BIOL110__HM| & \verb|BIOL110__HM_pre|, \verb|BIOL110__HM_co|\\\hline
    \verb|BIOL111__HM| & \verb|BIOL111__HM_pre|, \verb|BIOL111__HM_co|\\\hline
    \verb|BIOL154__HM| & \verb|BIOL154__HM_pre|, \verb|BIOL154__HM_co|\\\hline
    \verb|CHEM048__HM| & \verb|CHEM048__HM_pre|, \verb|CHEM048__HM_co|\\\hline
    \verb|CHEM058__HM| & \verb|CHEM058__HM_pre|, \verb|CHEM058__HM_co|\\\hline
    \verb|CHEM080__HM| & \verb|CHEM080__HM_0|\\\hline
    \verb|CHEM110__HM| & \verb|CHEM110__HM_pre|, \verb|CHEM110__HM_co|\\\hline
    \verb|CHEM111__HM| & \verb|CHEM111__HM_pre|, \verb|CHEM111__HM_co|\\\hline
    \verb|CHEM112__HM| & \verb|CHEM112__HM_pre|, \verb|CHEM112__HM_co|\\\hline
    \verb|CSCI081__HM| & \verb|CSCI081__HM_0|, \verb|CSCI081__HM_1|, \verb|CSCI081__HM_2|, \verb|CSCI081__HM_3|\\\hline
    \verb|CSCI140__HM| & \verb|CSCI140__HM_0|, \verb|CSCI140__HM_1|\\\hline
    \verb|CSCI140__HM_0| & \verb|CSCI140__HM_0_0|, \verb|CSCI140__HM_0_1|\\\hline
    \verb|CSCI140__HM_0_0| & \verb|CSCI140__HM_0_0_0|, \verb|CSCI140__HM_0_0_1|, \verb|CSCI140__HM_0_0_2|\\\hline
    \verb|CSCI140__HM_0_1| & \verb|CSCI140__HM_0_1_0|\\\hline
    \verb|CSCI144__HM| & \verb|CSCI144__HM_0|\\\hline
    \verb|CSCI151__HM| & \verb|CSCI151__HM_0|\\\hline
    \verb|CSCI152__HM| & \verb|CSCI152__HM_0|\\\hline
    \verb|CSCI158__HM| & \verb|CSCI158__HM_0|\\\hline
    \verb|CSCI159__HM| & \verb|CSCI159__HM_0|\\\hline
    \verb|Core| & \verb|Core_IntroCS|, \verb|Core_Mechanics|\\\hline
    \verb|ENGR004__HM| & \verb|ENGR004__HM_pre|, \verb|ENGR004__HM_co|\\\hline
    \verb|ENGR072__HM| & \verb|ENGR072__HM_pre|, \verb|ENGR072__HM_co|\\\hline
    \verb|ENGR080__HM| & \verb|ENGR080__HM_pre|, \verb|ENGR080__HM_co|\\\hline
    \verb|ENGR154__HM| & \verb|ENGR154__HM_0|\\\hline
    \verb|ENGR155__HM| & \verb|ENGR155__HM_0|\\\hline
    \verb|ENGR157__HM| & \verb|ENGR157__HM_pre|, \verb|ENGR157__HM_co|\\\hline
    \verb|HSAs| & \verb|HSAs_TenHSAs|, \verb|HSAs_MuddHums|, \verb|HSAs_WritingIntensive|\\\hline
    \verb|MATH062__HM| & \verb|MATH062__HM_pre|, \verb|MATH062__HM_co|\\\hline
    \verb|MATH179__HM| & \verb|MATH179__HM_0|, \verb|MATH179__HM_1|\\\hline
    \verb|MATH181__HM| & \verb|MATH181__HM_pre|, \verb|MATH181__HM_co|\\\hline
    \verb|Maj_Bio| & \verb|Maj_Bio_BioLab|, \verb|Maj_Bio_BioSeminar|, \verb|Maj_Bio_BioElecs|, \verb|Maj_Bio_BioColloquium|, \verb|Maj_Bio_ThesisOrClinic|\\\hline
    \verb|Maj_BioChem| & \verb|Maj_BioChem_ChemWithoutCells|, \verb|Maj_BioChem_BioChem|, \verb|Maj_BioChem_Topics|, \verb|Maj_BioChem_BioChemLab|, \verb|Maj_BioChem_ChemLab|, \verb|Maj_BioChem_OrganicBio|, \verb|Maj_BioChem_BioElective|, \verb|Maj_BioChem_BioColloqium|, \verb|Maj_BioChem_ChemColloqium|, \verb|Maj_BioChem_Capstone|\\\hline
    \verb|Maj_BioChem_Capstone| & \verb|Maj_BioChem_Capstone_BioThesis|, \verb|Maj_BioChem_Capstone_BioResearch|, \verb|Maj_BioChem_Capstone_ChemThesis|, \verb|Maj_BioChem_Capstone_EngrClinic|\\\hline
    \verb|Maj_BioClim| & \verb|Maj_BioClim_BioGroup|, \verb|Maj_BioClim_BioSeminar|, \verb|Maj_BioClim_BioElec|, \verb|Maj_BioClim_Thermodynamics|, \verb|Maj_BioClim_BioComputational|, \verb|Maj_BioClim_BioClimateElective|, \verb|Maj_BioClim_ClimateInterventions|, \verb|Maj_BioClim_ClimateContexts|, \verb|Maj_BioClim_ClimateStats|, \verb|Maj_BioClim_BioColloqium|, \verb|Maj_BioClim_ClimateColloqium|, \verb|Maj_BioClim_BioClimateCapstone|\\\hline
    \verb|Maj_BioClim_BioClimateCapstone| & \verb|Maj_BioClim_BioClimateCapstone_Thesis|, \verb|Maj_BioClim_BioClimateCapstone_CSClinic|, \verb|Maj_BioClim_BioClimateCapstone_EngClinic|, \verb|Maj_BioClim_BioClimateCapstone_MathClinic|, \verb|Maj_BioClim_BioClimateCapstone_PhysClinic|\\\hline
    \verb|Maj_BioClim_ClimateContexts| & \verb|Maj_BioClim_ClimateContexts_Justice|\\\hline
    \verb|Maj_BioClim_ClimateInterventions| & \verb|Maj_BioClim_ClimateInterventions_ClimateInterventionCredits|, \verb|Maj_BioClim_ClimateInterventions_Justice|\\\hline
    \verb|Maj_Bio_BioLab| & \verb|Maj_Bio_BioLab_WithBioChem|, \verb|Maj_Bio_BioLab_NoBioChem|\\\hline
    \verb|Maj_Bio_BioLab_WithBioChem| & \verb|Maj_Bio_BioLab_WithBioChem_Other|\\\hline
    \verb|Maj_Bio_ThesisOrClinic| & \verb|Maj_Bio_ThesisOrClinic_BioThesis|, \verb|Maj_Bio_ThesisOrClinic_IntensiveBioThesis|, \verb|Maj_Bio_ThesisOrClinic_CSClinic|, \verb|Maj_Bio_ThesisOrClinic_EngrClinic|, \verb|Maj_Bio_ThesisOrClinic_PhysClinic|, \verb|Maj_Bio_ThesisOrClinic_MathClinic|\\\hline
    \verb|Maj_CS| & \verb|Maj_CS_CSCore|, \verb|Maj_CS_CSElec|, \verb|Maj_CS_CSColloquium|\\\hline
    \verb|Maj_CSClimate| & \verb|Maj_CSClimate_Principles|, \verb|Maj_CSClimate_SysAlgs|, \verb|Maj_CSClimate_CSElective|, \verb|Maj_CSClimate_Thermodynamics|, \verb|Maj_CSClimate_Areas|, \verb|Maj_CSClimate_ClimateImpacts|, \verb|Maj_CSClimate_ClimateInterventions|, \verb|Maj_CSClimate_ClimateContexts|, \verb|Maj_CSClimate_HumanCentered|, \verb|Maj_CSClimate_CSColloquium|, \verb|Maj_CSClimate_ClimateColloquium|\\\hline
    \verb|Maj_CSClimate_Areas| & \verb|Maj_CSClimate_Areas_ProbStats|, \verb|Maj_CSClimate_Areas_Computational|\\\hline
    \verb|Maj_CSClimate_ClimateImpacts| & \verb|Maj_CSClimate_ClimateImpacts_ClimateJustice|\\\hline
    \verb|Maj_CSClimate_ClimateInterventions| & \verb|Maj_CSClimate_ClimateInterventions_ClimateJustice|\\\hline
    \verb|Maj_Chem| & \verb|Maj_Chem_Math|, \verb|Maj_Chem_ChemLab|, \verb|Maj_Chem_CHEMElec|, \verb|Maj_Chem_ChemSeminar|, \verb|Maj_Chem_ThesisOrClinic|\\\hline
    \verb|Maj_ChemClimate| & \verb|Maj_ChemClimate_Or1|, \verb|Maj_ChemClimate_ChemClass|, \verb|Maj_ChemClimate_ChemUpperDiv|, \verb|Maj_ChemClimate_ChemClimateProbStats|, \verb|Maj_ChemClimate_ChemClimateComputational|, \verb|Maj_ChemClimate_ClimateInterventions|, \verb|Maj_ChemClimate_ClimateContexts|, \verb|Maj_ChemClimate_ChemSeminar|, \verb|Maj_ChemClimate_ClimateColloqium|\\\hline
    \verb|Maj_ChemClimate_ClimateContexts| & \verb|Maj_ChemClimate_ClimateContexts_Justice|\\\hline
    \verb|Maj_ChemClimate_ClimateInterventions| & \verb|Maj_ChemClimate_ClimateInterventions_ClimateInterventionCredits|, \verb|Maj_ChemClimate_ClimateInterventions_Justice|\\\hline
    \verb|Maj_Chem_ThesisOrClinic| & \verb|Maj_Chem_ThesisOrClinic_Thesis|, \verb|Maj_Chem_ThesisOrClinic_Clinic|\\\hline
    \verb|Maj_CsPhys| & \verb|Maj_CsPhys_IntroCS|, \verb|Maj_CsPhys_CSCoreElective|, \verb|Maj_CsPhys_CSElective|, \verb|Maj_CsPhys_CSColloqium|, \verb|Maj_CsPhys_Quantum|, \verb|Maj_CsPhys_Lab|, \verb|Maj_CsPhys_PhysColloqium|, \verb|Maj_CsPhys_PhysCSElective|, \verb|Maj_CsPhys_Capstone|\\\hline
    \verb|Maj_CsPhys_Capstone| & \verb|Maj_CsPhys_Capstone_CSMTClinic|, \verb|Maj_CsPhys_Capstone_PhysClinic|, \verb|Maj_CsPhys_Capstone_PhysThesis|\\\hline
    \verb|Maj_Engr| & \verb|Maj_Engr_EngrElec|\\\hline
    \verb|Maj_Math| & \verb|Maj_Math_Computational|, \verb|Maj_Math_MathElec|, \verb|Maj_Math_ThesisOrClinic|\\\hline
    \verb|Maj_MathCS| & \verb|Maj_MathCS_IntroCS|, \verb|Maj_MathCS_Algs|, \verb|Maj_MathCS_CSElec|, \verb|Maj_MathCS_MathElec|\\\hline
    \verb|Maj_MathCompBio| & \verb|Maj_MathCompBio_BioFoundations|, \verb|Maj_MathCompBio_BioLab|, \verb|Maj_MathCompBio_BioSeminar|, \verb|Maj_MathCompBio_AdvancedBio|, \verb|Maj_MathCompBio_MathElec|, \verb|Maj_MathCompBio_CSElec|, \verb|Maj_MathCompBio_MathOrCSElec|, \verb|Maj_MathCompBio_BioColloquium|, \verb|Maj_MathCompBio_ThesisOrClinic|\\\hline
    \verb|Maj_MathCompBio_ThesisOrClinic| & \verb|Maj_MathCompBio_ThesisOrClinic_CSClinic|, \verb|Maj_MathCompBio_ThesisOrClinic_MathThesis|, \verb|Maj_MathCompBio_ThesisOrClinic_MathClinic|, \verb|Maj_MathCompBio_ThesisOrClinic_BioThesis|, \verb|Maj_MathCompBio_ThesisOrClinic_IntensiveBioThesis|\\\hline
    \verb|Maj_MathPhys| & \verb|Maj_MathPhys_Fields|, \verb|Maj_MathPhys_PhysColloquium|, \verb|Maj_MathPhys_ScientificComp|, \verb|Maj_MathPhys_ThesisOrClinic|\\\hline
    \verb|Maj_MathPhys_ThesisOrClinic| & \verb|Maj_MathPhys_ThesisOrClinic_PhysThesis|, \verb|Maj_MathPhys_ThesisOrClinic_MathThesis|, \verb|Maj_MathPhys_ThesisOrClinic_PhysClinic|, \verb|Maj_MathPhys_ThesisOrClinic_MathClinic|\\\hline
    \verb|Maj_Math_ThesisOrClinic| & \verb|Maj_Math_ThesisOrClinic_MathThesis|, \verb|Maj_Math_ThesisOrClinic_BioThesis|\\\hline
    \verb|Maj_MolBio| & \verb|Maj_MolBio_MolBio|, \verb|Maj_MolBio_BioLab|, \verb|Maj_MolBio_SeminarCourse|, \verb|Maj_MolBio_BioElecs|, \verb|Maj_MolBio_Colloquium|, \verb|Maj_MolBio_Capstone|\\\hline
    \verb|Maj_MolBio_Capstone| & \verb|Maj_MolBio_Capstone_BioThesis|, \verb|Maj_MolBio_Capstone_IntensiveResearch|, \verb|Maj_MolBio_Capstone_ChemThesis|, \verb|Maj_MolBio_Capstone_CSClinic|, \verb|Maj_MolBio_Capstone_EngrClinic|, \verb|Maj_MolBio_Capstone_MathClinic|, \verb|Maj_MolBio_Capstone_PhysClinic|\\\hline
    \verb|Maj_Phys| & \verb|Maj_Phys_PhysColloquium|, \verb|Maj_Phys_Thesis|\\\hline
    \verb|Maj_Phys_Thesis| & \verb|Maj_Phys_Thesis_Clinic|, \verb|Maj_Phys_Thesis_PhysThesis|\\\hline
    \verb|PHYS051__HM| & \verb|PHYS051__HM_pre|, \verb|PHYS051__HM_co|\\\hline
    \verb|PHYS052__HM| & \verb|PHYS052__HM_0|\\\hline
    \verb|PHYS064__HM| & \verb|PHYS064__HM_0|, \verb|PHYS064__HM_1|\\\hline
    \verb|PHYS084__HM| & \verb|PHYS084__HM_0|\\\hline
    \verb|PHYS111__HM| & \verb|PHYS111__HM_0|\\\hline
    \verb|PHYS151__HM| & \verb|PHYS151__HM_0|, \verb|PHYS151__HM_1|\\\hline
    \verb|PHYS156__HM| & \verb|PHYS156__HM_0|\\\hline
    \verb|PHYS162__HM| & \verb|PHYS162__HM_pre|, \verb|PHYS162__HM_co|\\\hline
    \caption{Requirement group contents: other requirement groups.}
    \label{tab:reqs-groups}
\end{longtable}
\begin{longtable}{|p{4.4in}|p{4.4in}|}
    \hline
    Requirement Group Name & Course Prequirequisites\\\hline
    \verb|AFRI190__AF2| & \verb|AFRI190__AF|\\\hline
    \verb|AFRI191C_AF2| & \verb|AFRI191C_AF|\\\hline
    \verb|AFRI191__AF2| & \verb|AFRI191__AF|\\\hline
    \verb|AFRI195D_AF2| & \verb|AFRI195D_AF|\\\hline
    \verb|AMST190__PO2| & \verb|AMST190__PO|\\\hline
    \verb|AMST191__PO2| & \verb|AMST191__PO|\\\hline
    \verb|AMST191__SC2| & \verb|AMST191__SC|\\\hline
    \verb|ANTH190__PO2| & \verb|ANTH190__PO|\\\hline
    \verb|ANTH190__SC2| & \verb|ANTH190__SC|\\\hline
    \verb|ANTH191__SC2| & \verb|ANTH191__SC|\\\hline
    \verb|ARCN191__SC2| & \verb|ARCN191__SC|\\\hline
    \verb|ARHI190__PO2| & \verb|ARHI190__PO|\\\hline
    \verb|ARHI191__PO2| & \verb|ARHI191__PO|\\\hline
    \verb|ARHI191__SC2| & \verb|ARHI191__SC|\\\hline
    \verb|ART_190__PO2| & \verb|ART_190__PO|\\\hline
    \verb|ART_192__PO2| & \verb|ART_192__PO|\\\hline
    \verb|ART_192__SC2| & \verb|ART_192__SC|\\\hline
    \verb|ART_193__SC2| & \verb|ART_193__SC|\\\hline
    \verb|ART_199__PZ2| & \verb|ART_199__PZ|\\\hline
    \verb|ASAM190A_PO2| & \verb|ASAM190A_PO|\\\hline
    \verb|ASAM191__PO2| & \verb|ASAM191__PO|\\\hline
    \verb|ASIA190__PO2| & \verb|ASIA190__PO|\\\hline
    \verb|ASIA191__PO2| & \verb|ASIA191__PO|\\\hline
    \verb|ASIA192__PO2| & \verb|ASIA192__PO|\\\hline
    \verb|ASTR062__HM| & \verb|PHYS051__HM|\\\hline
    \verb|ASTR122__HM| & \verb|ASTR062__HM|, \verb|PHYS052__HM|\\\hline
    \verb|ASTR128__HM| & \verb|CSCI005__HM|, \verb|MATH073__HM|, \verb|PHYS024__HM|\\\hline
    \verb|BIOL101__HM| & \verb|BIOL046__HM|\\\hline
    \verb|BIOL103__HM| & \verb|BIOL054__HM|, \verb|BIOL101__HM|\\\hline
    \verb|BIOL108__HM| & \verb|BIOL046__HM|, \verb|MATH019__HM|\\\hline
    \verb|BIOL109__HM| & \verb|BIOL046__HM|, \verb|MATH019__HM|\\\hline
    \verb|BIOL110__HM_pre| & \verb|BIOL054__HM|, \verb|BIOL154__HM|\\\hline
    \verb|BIOL111__HM_pre| & \verb|BIOL054__HM|, \verb|BIOL154__HM|\\\hline
    \verb|BIOL113__HM| & \verb|BIOL046__HM|, \verb|CHEM042__HM|\\\hline
    \verb|BIOL119__HM| & \verb|MCBI118A_HM|\\\hline
    \verb|BIOL121__HM| & \verb|BIOL154__HM|\\\hline
    \verb|BIOL122__HM| & \verb|BIOL113__HM|\\\hline
    \verb|BIOL154__HM_pre| & \verb|CSCI005__HM|\\\hline
    \verb|BIOL160__HM| & \verb|BIOL113__HM|\\\hline
    \verb|BIOL182__HM| & \verb|CHEM056__HM|\\\hline
    \verb|BIOL188__HM| & \verb|MCBI118B_HM|\\\hline
    \verb|BIOL189__HM| & \verb|BIOL113__HM|\\\hline
    \verb|BIOL190L_KS2| & \verb|BIOL190L_KS|\\\hline
    \verb|BIOL190__PO2| & \verb|BIOL190__PO|\\\hline
    \verb|BIOL191F_PO2| & \verb|BIOL191F_PO|\\\hline
    \verb|BIOL191H_PO2| & \verb|BIOL191H_PO|\\\hline
    \verb|BIOL191__HM2| & \verb|BIOL191__HM|\\\hline
    \verb|BIOL191__HM3| & \verb|BIOL191__HM2|\\\hline
    \verb|BIOL191__HM4| & \verb|BIOL191__HM3|\\\hline
    \verb|BIOL191__KS2| & \verb|BIOL191__KS|\\\hline
    \verb|BIOL193__HM2| & \verb|BIOL193__HM|\\\hline
    \verb|BIOL194A_PO2| & \verb|BIOL194A_PO|\\\hline
    \verb|BIOL194B_PO2| & \verb|BIOL194B_PO|\\\hline
    \verb|BIOL195__HM| & \verb|BIOL161__HM|\\\hline
    \verb|BIOL195__HM2| & \verb|BIOL195__HM|\\\hline
    \verb|BIOL196__KS2| & \verb|BIOL196__KS|\\\hline
    \verb|BIOL197__HM2| & \verb|BIOL197__HM|\\\hline
    \verb|BIOL197__KS2| & \verb|BIOL197__KS|\\\hline
    \verb|CGS_190__PZ2| & \verb|CGS_190__PZ|\\\hline
    \verb|CGS_199__PZ2| & \verb|CGS_199__PZ|\\\hline
    \verb|CHEM048__HM_pre| & \verb|CSCI005__HM|\\\hline
    \verb|CHEM051__HM| & \verb|CHEM024__HM|, \verb|CHEM042__HM|\\\hline
    \verb|CHEM056__HM| & \verb|CHEM042__HM|, \verb|CHEM024__HM|\\\hline
    \verb|CHEM058__HM_pre| & \verb|CHEM024__HM|\\\hline
    \verb|CHEM080__HM| & \verb|CHEM042__HM|, \verb|MATH073__HM|\\\hline
    \verb|CHEM080__HM_0| & \verb|CSCI005__HM|\\\hline
    \verb|CHEM103__HM| & \verb|CHEM042__HM|, \verb|CHEM024__HM|\\\hline
    \verb|CHEM104__HM| & \verb|CHEM056__HM|\\\hline
    \verb|CHEM105__HM| & \verb|CHEM056__HM|\\\hline
    \verb|CHEM110__HM_pre| & \verb|CHEM058__HM|\\\hline
    \verb|CHEM111__HM_pre| & \verb|CHEM058__HM|\\\hline
    \verb|CHEM112__HM_pre| & \verb|CHEM103__HM|, \verb|CHEM109__HM|\\\hline
    \verb|CHEM114__HM| & \verb|CHEM103__HM|\\\hline
    \verb|CHEM116__HM| & \verb|CHEM051__HM|\\\hline
    \verb|CHEM122__HM| & \verb|ENGR086__HM|, \verb|PHYS051__HM|\\\hline
    \verb|CHEM170__HM| & \verb|CHEM051__HM|, \verb|CHEM056__HM|\\\hline
    \verb|CHEM171__HM| & \verb|CHEM056__HM|, \verb|CHEM105__HM|\\\hline
    \verb|CHEM182__HM| & \verb|CHEM056__HM|\\\hline
    \verb|CHEM190L_KS2| & \verb|CHEM190L_KS|\\\hline
    \verb|CHEM191__KS2| & \verb|CHEM191__KS|\\\hline
    \verb|CHEM191__PO2| & \verb|CHEM191__PO|\\\hline
    \verb|CHEM193M_HM2| & \verb|CHEM193M_HM|\\\hline
    \verb|CHEM193P_HM2| & \verb|CHEM193P_HM|\\\hline
    \verb|CHEM194__HM| & \verb|CHEM056__HM|\\\hline
    \verb|CHEM194__HM2| & \verb|CHEM194__HM|\\\hline
    \verb|CHEM194__PO2| & \verb|CHEM194__PO|\\\hline
    \verb|CHEM196__KS2| & \verb|CHEM196__KS|\\\hline
    \verb|CHEM197__HM2| & \verb|CHEM197__HM|\\\hline
    \verb|CHEM197__KS2| & \verb|CHEM197__KS|\\\hline
    \verb|CHEM198__HM2| & \verb|CHEM198__HM|\\\hline
    \verb|CHEM199__HM2| & \verb|CHEM199__HM|\\\hline
    \verb|CHEM199__HM3| & \verb|CHEM199__HM2|\\\hline
    \verb|CHEM199__HM4| & \verb|CHEM199__HM3|\\\hline
    \verb|CHIN192A_PO2| & \verb|CHIN192A_PO|\\\hline
    \verb|CHST190__CH2| & \verb|CHST190__CH|\\\hline
    \verb|CHST191__CH2| & \verb|CHST191__CH|\\\hline
    \verb|CHST192__CH2| & \verb|CHST192__CH|\\\hline
    \verb|CLAS190__PO2| & \verb|CLAS190__PO|\\\hline
    \verb|CLAS191__PO2| & \verb|CLAS191__PO|\\\hline
    \verb|CLAS191__SC2| & \verb|CLAS191__SC|\\\hline
    \verb|CLAS192__PO2| & \verb|CLAS192__PO|\\\hline
    \verb|CLES197__HM2| & \verb|CLES197__HM|\\\hline
    \verb|CLES199__HM2| & \verb|CLES199__HM|\\\hline
    \verb|CSCI035__HM| & \verb|CSCI005__HM|\\\hline
    \verb|CSCI060__HM| & \verb|CSCI005__HM|\\\hline
    \verb|CSCI070__HM| & \verb|CSCI060__HM|\\\hline
    \verb|CSCI081__HM_0| & \verb|MATH055__HM|\\\hline
    \verb|CSCI081__HM_1| & \verb|CSCI060__HM|\\\hline
    \verb|CSCI081__HM_2| & \verb|MATH019__HM|, \verb|MATH032__CM|, \verb|MATH032__PO|, \verb|MATH032__PZ|, \verb|MATH032__SC|, \verb|MATH032S_PO|\\\hline
    \verb|CSCI081__HM_3| & \verb|MATH073__HM|, \verb|MATH060__CM|, \verb|MATH060__PO|, \verb|MATH060__PZ|, \verb|MATH060__SC|\\\hline
    \verb|CSCI105__HM| & \verb|CSCI070__HM|\\\hline
    \verb|CSCI123__HM| & \verb|CSCI070__HM|\\\hline
    \verb|CSCI131__HM| & \verb|CSCI070__HM|, \verb|CSCI081__HM|\\\hline
    \verb|CSCI133__HM| & \verb|CSCI070__HM|, \verb|CSCI081__HM|\\\hline
    \verb|CSCI134__HM| & \verb|CSCI105__HM|\\\hline
    \verb|CSCI140__HM| & \verb|CSCI081__HM|\\\hline
    \verb|CSCI140__HM_0_0| & \verb|CSCI070__HM|\\\hline
    \verb|CSCI140__HM_0_0_0| & \verb|MATH055__HM|, \verb|MATH055__CM|, \verb|MATH055__PZ|, \verb|MATH055__SC|\\\hline
    \verb|CSCI140__HM_0_0_1| & \verb|MATH019__HM|, \verb|MATH032S_PO|, \verb|MATH067__PO|\\\hline
    \verb|CSCI140__HM_0_0_2| & \verb|MATH073__HM|, \verb|MATH060C_CM|\\\hline
    \verb|CSCI140__HM_0_1| & \verb|MATH131__HM|\\\hline
    \verb|CSCI140__HM_0_1_0| & \verb|CSCI060__HM|\\\hline
    \verb|CSCI140__HM_1| & \verb|CSCI062__PO|\\\hline
    \verb|CSCI144__HM| & \verb|MATH073__HM|, \verb|MATH082__HM|\\\hline
    \verb|CSCI144__HM_0| & \verb|CSCI060__HM|\\\hline
    \verb|CSCI151__HM| & \verb|CSCI070__HM|\\\hline
    \verb|CSCI151__HM_0| & \verb|MATH056__HM|, \verb|MATH062__HM|, \verb|BIOL154__HM|, \verb|MATH151__CM|, \verb|MATH151__PO|\\\hline
    \verb|CSCI152__HM| & \verb|CSCI070__HM|, \verb|MATH073__HM|\\\hline
    \verb|CSCI152__HM_0| & \verb|MATH056__HM|, \verb|MATH062__HM|, \verb|BIOL154__HM|, \verb|MATH151__CM|, \verb|MATH151__PO|\\\hline
    \verb|CSCI153__HM| & \verb|CSCI070__HM|\\\hline
    \verb|CSCI158__HM| & \verb|CSCI070__HM|, \verb|MATH073__HM|, \verb|CSCI151__HM|\\\hline
    \verb|CSCI158__HM_0| & \verb|MATH056__HM|, \verb|MATH062__HM|, \verb|BIOL154__HM|, \verb|MATH151__CM|, \verb|MATH151__PO|\\\hline
    \verb|CSCI159__HM| & \verb|CSCI081__HM|\\\hline
    \verb|CSCI159__HM_0| & \verb|MATH056__HM|, \verb|MATH062__HM|, \verb|BIOL154__HM|, \verb|MATH151__CM|, \verb|MATH151__PO|\\\hline
    \verb|CSCI183__HM| & \verb|CSCI123__HM|\\\hline
    \verb|CSCI184__HM| & \verb|CSCI183__HM|\\\hline
    \verb|CSCI190__PO2| & \verb|CSCI190__PO|\\\hline
    \verb|CSCI191__PO2| & \verb|CSCI191__PO|\\\hline
    \verb|CSCI192__PO2| & \verb|CSCI192__PO|\\\hline
    \verb|CSCI195__HM2| & \verb|CSCI195__HM|\\\hline
    \verb|CSCI195__HM3| & \verb|CSCI195__HM2|\\\hline
    \verb|CSCI195__HM4| & \verb|CSCI195__HM3|\\\hline
    \verb|CSCI198__SC2| & \verb|CSCI198__SC|\\\hline
    \verb|CSMT184__HM| & \verb|CSMT183__HM|\\\hline
    \verb|DANC190__SC2| & \verb|DANC190__SC|\\\hline
    \verb|DANC191__SC2| & \verb|DANC191__SC|\\\hline
    \verb|DANC192__PO2| & \verb|DANC192__PO|\\\hline
    \verb|DANC193__SC2| & \verb|DANC193__SC|\\\hline
    \verb|DS__190__PO2| & \verb|DS__190__PO|\\\hline
    \verb|DS__191__SC2| & \verb|DS__191__SC|\\\hline
    \verb|EA__190L_KS2| & \verb|EA__190L_KS|\\\hline
    \verb|EA__190__PO2| & \verb|EA__190__PO|\\\hline
    \verb|EA__191H_PO2| & \verb|EA__191H_PO|\\\hline
    \verb|EA__191__KS2| & \verb|EA__191__KS|\\\hline
    \verb|EA__191__PO2| & \verb|EA__191__PO|\\\hline
    \verb|EA__196__KS2| & \verb|EA__196__KS|\\\hline
    \verb|EA__197__KS2| & \verb|EA__197__KS|\\\hline
    \verb|EA__197__PZ2| & \verb|EA__197__PZ|\\\hline
    \verb|ECON190__PO2| & \verb|ECON190__PO|\\\hline
    \verb|ECON191H_SC2| & \verb|ECON191H_SC|\\\hline
    \verb|ECON191__CM2| & \verb|ECON191__CM|\\\hline
    \verb|ECON191__SC2| & \verb|ECON191__SC|\\\hline
    \verb|ECON194A_CM2| & \verb|ECON194A_CM|\\\hline
    \verb|ECON194B_CM2| & \verb|ECON194B_CM|\\\hline
    \verb|ECON195__PO2| & \verb|ECON195__PO|\\\hline
    \verb|ECON196__CM2| & \verb|ECON196__CM|\\\hline
    \verb|ECON197SBCM2| & \verb|ECON197SBCM|\\\hline
    \verb|ECON197S_CM2| & \verb|ECON197S_CM|\\\hline
    \verb|ECON198__PZ2| & \verb|ECON198__PZ|\\\hline
    \verb|ECON199__CM2| & \verb|ECON199__CM|\\\hline
    \verb|ECON199__PZ2| & \verb|ECON199__PZ|\\\hline
    \verb|EEP_191__SC2| & \verb|EEP_191__SC|\\\hline
    \verb|ENGL190__SC2| & \verb|ENGL190__SC|\\\hline
    \verb|ENGL191__PO2| & \verb|ENGL191__PO|\\\hline
    \verb|ENGL191__SC2| & \verb|ENGL191__SC|\\\hline
    \verb|ENGL192__SC2| & \verb|ENGL192__SC|\\\hline
    \verb|ENGL193__SC2| & \verb|ENGL193__SC|\\\hline
    \verb|ENGL194__PZ2| & \verb|ENGL194__PZ|\\\hline
    \verb|ENGL194__SC2| & \verb|ENGL194__SC|\\\hline
    \verb|ENGL195B_PO2| & \verb|ENGL195B_PO|\\\hline
    \verb|ENGL195__PO2| & \verb|ENGL195__PO|\\\hline
    \verb|ENGL195__SC2| & \verb|ENGL195__SC|\\\hline
    \verb|ENGL197N_SC2| & \verb|ENGL197N_SC|\\\hline
    \verb|ENGL199T_SC2| & \verb|ENGL199T_SC|\\\hline
    \verb|ENGR004__HM_pre| & \verb|WRIT001__HM|\\\hline
    \verb|ENGR072__HM_pre| & \verb|MATH019__HM|, \verb|MATH073__HM|, \verb|ENGR079__HM|\\\hline
    \verb|ENGR079__HM| & \verb|PHYS024__HM|\\\hline
    \verb|ENGR080__HM_pre| & \verb|ENGR079__HM|\\\hline
    \verb|ENGR082__HM| & \verb|CHEM042__HM|\\\hline
    \verb|ENGR083__HM| & \verb|ENGR079__HM|, \verb|PHYS024__HM|\\\hline
    \verb|ENGR084__HM| & \verb|ENGR079__HM|\\\hline
    \verb|ENGR085A_HM| & \verb|CSCI005__HM|\\\hline
    \verb|ENGR085__HM| & \verb|CSCI005__HM|\\\hline
    \verb|ENGR086__HM| & \verb|CHEM042__HM|, \verb|MATH019__HM|, \verb|MATH073__HM|, \verb|PHYS024__HM|\\\hline
    \verb|ENGR101__HM| & \verb|ENGR072__HM|, \verb|ENGR079__HM|, \verb|ENGR080__HM|\\\hline
    \verb|ENGR102__HM| & \verb|ENGR101__HM|\\\hline
    \verb|ENGR111__HM| & \verb|ENGR080__HM|\\\hline
    \verb|ENGR112__HM| & \verb|ENGR004__HM|, \verb|ENGR080__HM|, \verb|ENGR111__HM|\\\hline
    \verb|ENGR113__HM| & \verb|ENGR004__HM|, \verb|ENGR080__HM|, \verb|ENGR112__HM|\\\hline
    \verb|ENGR131__HM| & \verb|ENGR083__HM|\\\hline
    \verb|ENGR132__HM| & \verb|ENGR082__HM|\\\hline
    \verb|ENGR133__HM| & \verb|ENGR082__HM|\\\hline
    \verb|ENGR134__HM| & \verb|ENGR082__HM|\\\hline
    \verb|ENGR138__HM| & \verb|ENGR082__HM|\\\hline
    \verb|ENGR151__HM| & \verb|ENGR079__HM|, \verb|ENGR084__HM|\\\hline
    \verb|ENGR154__HM| & \verb|ENGR085__HM|\\\hline
    \verb|ENGR154__HM_0| & \verb|ENGR085A_HM|, \verb|CSCI105__HM|\\\hline
    \verb|ENGR155__HM| & \verb|ENGR085__HM|\\\hline
    \verb|ENGR155__HM_0| & \verb|ENGR085A_HM|, \verb|CSCI060__HM|\\\hline
    \verb|ENGR157__HM_pre| & \verb|ENGR084__HM|\\\hline
    \verb|ENGR164__HM| & \verb|BIOL046__HM|, \verb|ENGR079__HM|\\\hline
    \verb|ENGR171__HM| & \verb|ENGR083__HM|\\\hline
    \verb|ENGR177__HM| & \verb|ENGR004__HM|, \verb|ENGR083__HM|\\\hline
    \verb|ENGR178__HM| & \verb|ENGR080__HM|\\\hline
    \verb|ENGR183__HM| & \verb|ENGR004__HM|\\\hline
    \verb|ENGR185A_HM| & \verb|ENGR004__HM|\\\hline
    \verb|ENGR185B_HM| & \verb|ENGR185A_HM|\\\hline
    \verb|ENGR186__HM| & \verb|ENGR180__HM|\\\hline
    \verb|ENGR187__HM| & \verb|MATH073__HM|\\\hline
    \verb|ENGR190BAHM2| & \verb|ENGR190BAHM|\\\hline
    \verb|ENGR190BFHM2| & \verb|ENGR190BFHM|\\\hline
    \verb|ENGR190BGHM2| & \verb|ENGR190BGHM|\\\hline
    \verb|ENGR190BHHM2| & \verb|ENGR190BHHM|\\\hline
    \verb|ENGR190BIHM2| & \verb|ENGR190BIHM|\\\hline
    \verb|ENGR190BJHM2| & \verb|ENGR190BJHM|\\\hline
    \verb|ENGR191__HM2| & \verb|ENGR191__HM|\\\hline
    \verb|ENGR205__HM| & \verb|ENGR102__HM|\\\hline
    \verb|ENGR207__HM| & \verb|ENGR101__HM|, \verb|CSCI060__HM|\\\hline
    \verb|ENGR208__HM| & \verb|ENGR101__HM|, \verb|CSCI060__HM|\\\hline
    \verb|ENGR278__HM| & \verb|ENGR171__HM|\\\hline
    \verb|FGSS191__SC2| & \verb|FGSS191__SC|\\\hline
    \verb|FGSS192__SC2| & \verb|FGSS192__SC|\\\hline
    \verb|FGSS198__SC2| & \verb|FGSS198__SC|\\\hline
    \verb|FLAN191__SC2| & \verb|FLAN191__SC|\\\hline
    \verb|FREN191__PO2| & \verb|FREN191__PO|\\\hline
    \verb|FREN191__SC2| & \verb|FREN191__SC|\\\hline
    \verb|FREN192__PO2| & \verb|FREN192__PO|\\\hline
    \verb|FREN193__PO2| & \verb|FREN193__PO|\\\hline
    \verb|GEOL192__PO2| & \verb|GEOL192__PO|\\\hline
    \verb|GERM191__PO2| & \verb|GERM191__PO|\\\hline
    \verb|GERM191__SC2| & \verb|GERM191__SC|\\\hline
    \verb|GERM193__PO2| & \verb|GERM193__PO|\\\hline
    \verb|GLAS194A_PZ2| & \verb|GLAS194A_PZ|\\\hline
    \verb|GLAS194B_PZ2| & \verb|GLAS194B_PZ|\\\hline
    \verb|GOVT191__CM2| & \verb|GOVT191__CM|\\\hline
    \verb|GOVT192__CM2| & \verb|GOVT192__CM|\\\hline
    \verb|GWS_190__PO2| & \verb|GWS_190__PO|\\\hline
    \verb|GWS_191__PO2| & \verb|GWS_191__PO|\\\hline
    \verb|HIST190__CM2| & \verb|HIST190__CM|\\\hline
    \verb|HIST190__PO2| & \verb|HIST190__PO|\\\hline
    \verb|HIST191__PO2| & \verb|HIST191__PO|\\\hline
    \verb|HIST191__SC2| & \verb|HIST191__SC|\\\hline
    \verb|HIST192__PO2| & \verb|HIST192__PO|\\\hline
    \verb|HIST192__SC2| & \verb|HIST192__SC|\\\hline
    \verb|HIST196__CM2| & \verb|HIST196__CM|\\\hline
    \verb|HIST197__CM2| & \verb|HIST197__CM|\\\hline
    \verb|HIST197__PZ2| & \verb|HIST197__PZ|\\\hline
    \verb|HIST199__PZ2| & \verb|HIST199__PZ|\\\hline
    \verb|HMSC191__SC2| & \verb|HMSC191__SC|\\\hline
    \verb|HMSC192__SC2| & \verb|HMSC192__SC|\\\hline
    \verb|HUM_195J_SC2| & \verb|HUM_195J_SC|\\\hline
    \verb|HUM_196__PO2| & \verb|HUM_196__PO|\\\hline
    \verb|ID__191D_SC2| & \verb|ID__191D_SC|\\\hline
    \verb|ID__191H_SC2| & \verb|ID__191H_SC|\\\hline
    \verb|ID__191S_SC2| & \verb|ID__191S_SC|\\\hline
    \verb|ID__199P1PO2| & \verb|ID__199P1PO|\\\hline
    \verb|ID__199P2PO2| & \verb|ID__199P2PO|\\\hline
    \verb|ID__199P3PO2| & \verb|ID__199P3PO|\\\hline
    \verb|ID__199P4PO2| & \verb|ID__199P4PO|\\\hline
    \verb|ID__199S1PO2| & \verb|ID__199S1PO|\\\hline
    \verb|ID__199S2PO2| & \verb|ID__199S2PO|\\\hline
    \verb|ID__199S3PO2| & \verb|ID__199S3PO|\\\hline
    \verb|ID__199S4PO2| & \verb|ID__199S4PO|\\\hline
    \verb|IR__190__PO2| & \verb|IR__190__PO|\\\hline
    \verb|IR__191__PO2| & \verb|IR__191__PO|\\\hline
    \verb|IR__192__SC2| & \verb|IR__192__SC|\\\hline
    \verb|ITAL191__SC2| & \verb|ITAL191__SC|\\\hline
    \verb|LAMS191__PO2| & \verb|LAMS191__PO|\\\hline
    \verb|LAST190__PO2| & \verb|LAST190__PO|\\\hline
    \verb|LAST191__PO2| & \verb|LAST191__PO|\\\hline
    \verb|LAST191__SC2| & \verb|LAST191__SC|\\\hline
    \verb|LGCS190__PO2| & \verb|LGCS190__PO|\\\hline
    \verb|LGCS191__PO2| & \verb|LGCS191__PO|\\\hline
    \verb|LIT_199__CM2| & \verb|LIT_199__CM|\\\hline
    \verb|MATH056__HM| & \verb|MATH019__HM|, \verb|MATH073__HM|\\\hline
    \verb|MATH062__HM_pre| & \verb|MATH019__HM|\\\hline
    \verb|MATH073__HM| & \verb|MATH019__HM|\\\hline
    \verb|MATH082__HM| & \verb|MATH019__HM|, \verb|MATH073__HM|\\\hline
    \verb|MATH104__HM| & \verb|MATH073__HM|, \verb|MATH055__HM|\\\hline
    \verb|MATH106__HM| & \verb|MATH055__HM|\\\hline
    \verb|MATH119__HM| & \verb|MCBI118A_HM|\\\hline
    \verb|MATH131__HM| & \verb|MATH055__HM|\\\hline
    \verb|MATH132__HM| & \verb|MATH131__HM|\\\hline
    \verb|MATH136__HM| & \verb|MATH073__HM|, \verb|MATH082__HM|\\\hline
    \verb|MATH147__HM| & \verb|MATH131__HM|\\\hline
    \verb|MATH156__HM| & \verb|MATH073__HM|, \verb|MATH157__HM|\\\hline
    \verb|MATH157__HM| & \verb|BIOL154__HM|, \verb|MATH062__HM|, \verb|PHYS052__HM|\\\hline
    \verb|MATH158__HM| & \verb|MATH062__HM|, \verb|MATH073__HM|\\\hline
    \verb|MATH164__HM| & \verb|MATH073__HM|, \verb|MATH082__HM|, \verb|CSCI060__HM|\\\hline
    \verb|MATH165__HM| & \verb|MATH073__HM|, \verb|MATH082__HM|\\\hline
    \verb|MATH171__HM| & \verb|MATH073__HM|, \verb|MATH055__HM|\\\hline
    \verb|MATH174__HM| & \verb|MATH171__HM|\\\hline
    \verb|MATH175__HM| & \verb|MATH055__HM|\\\hline
    \verb|MATH176__HM| & \verb|MATH171__HM|\\\hline
    \verb|MATH179__HM| & \verb|CSCI005__HM|\\\hline
    \verb|MATH179__HM_0| & \verb|MATH019__HM|\\\hline
    \verb|MATH179__HM_1| & \verb|MATH073__HM|\\\hline
    \verb|MATH180__HM| & \verb|MATH082__HM|, \verb|MATH131__HM|\\\hline
    \verb|MATH181__HM_pre| & \verb|MATH082__HM|\\\hline
    \verb|MATH187__HM| & \verb|MATH073__HM|\\\hline
    \verb|MATH190__CM2| & \verb|MATH190__CM|\\\hline
    \verb|MATH190__PO2| & \verb|MATH190__PO|\\\hline
    \verb|MATH191__CM2| & \verb|MATH191__CM|\\\hline
    \verb|MATH191__PO2| & \verb|MATH191__PO|\\\hline
    \verb|MATH191__SC2| & \verb|MATH191__SC|\\\hline
    \verb|MATH193__HM2| & \verb|MATH193__HM|\\\hline
    \verb|MATH194__PZ2| & \verb|MATH194__PZ|\\\hline
    \verb|MATH195__CM2| & \verb|MATH195__CM|\\\hline
    \verb|MATH196__HM2| & \verb|MATH196__HM|\\\hline
    \verb|MATH197__HM2| & \verb|MATH197__HM|\\\hline
    \verb|MATH198__HM2| & \verb|MATH198__HM|\\\hline
    \verb|MATH198__PZ2| & \verb|MATH198__PZ|\\\hline
    \verb|MCBI118A_HM| & \verb|MATH073__HM|, \verb|MATH082__HM|, \verb|BIOL046__HM|\\\hline
    \verb|MCBI118B_HM| & \verb|CSCI005__HM|, \verb|BIOL046__HM|\\\hline
    \verb|MCBI199__HM2| & \verb|MCBI199__HM|\\\hline
    \verb|MCBI199__HM3| & \verb|MCBI199__HM2|\\\hline
    \verb|MCBI199__HM4| & \verb|MCBI199__HM3|\\\hline
    \verb|MCSI195__PZ2| & \verb|MCSI195__PZ|\\\hline
    \verb|MENA191__SC2| & \verb|MENA191__SC|\\\hline
    \verb|MOBI191A_PO2| & \verb|MOBI191A_PO|\\\hline
    \verb|MOBI191B_PO2| & \verb|MOBI191B_PO|\\\hline
    \verb|MOBI194A_PO2| & \verb|MOBI194A_PO|\\\hline
    \verb|MOBI194B_PO2| & \verb|MOBI194B_PO|\\\hline
    \verb|MS__190F_JT2| & \verb|MS__190F_JT|\\\hline
    \verb|MS__190S_JT2| & \verb|MS__190S_JT|\\\hline
    \verb|MS__190__JT2| & \verb|MS__190__JT|\\\hline
    \verb|MS__194__PZ2| & \verb|MS__194__PZ|\\\hline
    \verb|MS__196__PZ2| & \verb|MS__196__PZ|\\\hline
    \verb|MS__199__SC2| & \verb|MS__199__SC|\\\hline
    \verb|MUS_190__PO2| & \verb|MUS_190__PO|\\\hline
    \verb|MUS_191__SC2| & \verb|MUS_191__SC|\\\hline
    \verb|MUS_192F_PO2| & \verb|MUS_192F_PO|\\\hline
    \verb|MUS_192H_PO2| & \verb|MUS_192H_PO|\\\hline
    \verb|NEUR190L_KS2| & \verb|NEUR190L_KS|\\\hline
    \verb|NEUR190__PO2| & \verb|NEUR190__PO|\\\hline
    \verb|NEUR191__KS2| & \verb|NEUR191__KS|\\\hline
    \verb|NEUR191__PO2| & \verb|NEUR191__PO|\\\hline
    \verb|NEUR192__PO2| & \verb|NEUR192__PO|\\\hline
    \verb|NEUR194A_PO2| & \verb|NEUR194A_PO|\\\hline
    \verb|NEUR194B_PO2| & \verb|NEUR194B_PO|\\\hline
    \verb|NEUR196__KS2| & \verb|NEUR196__KS|\\\hline
    \verb|NEUR197__KS2| & \verb|NEUR197__KS|\\\hline
    \verb|ORST191__SC2| & \verb|ORST191__SC|\\\hline
    \verb|ORST192__PZ2| & \verb|ORST192__PZ|\\\hline
    \verb|ORST198C_PZ2| & \verb|ORST198C_PZ|\\\hline
    \verb|ORST198J_PZ2| & \verb|ORST198J_PZ|\\\hline
    \verb|ORST198M_PZ2| & \verb|ORST198M_PZ|\\\hline
    \verb|PHIL190__CM2| & \verb|PHIL190__CM|\\\hline
    \verb|PHIL190__PO2| & \verb|PHIL190__PO|\\\hline
    \verb|PHIL190__SC2| & \verb|PHIL190__SC|\\\hline
    \verb|PHIL191__PO2| & \verb|PHIL191__PO|\\\hline
    \verb|PHIL191__PZ2| & \verb|PHIL191__PZ|\\\hline
    \verb|PHIL191__SC2| & \verb|PHIL191__SC|\\\hline
    \verb|PHIL192__CM2| & \verb|PHIL192__CM|\\\hline
    \verb|PHIL198__CM2| & \verb|PHIL198__CM|\\\hline
    \verb|PHIL198__SC2| & \verb|PHIL198__SC|\\\hline
    \verb|PHYS050__HM| & \verb|PHYS024__HM|\\\hline
    \verb|PHYS051__HM_pre| & \verb|PHYS023__HM|, \verb|PHYS024__HM|\\\hline
    \verb|PHYS052__HM| & \verb|PHYS051__HM|\\\hline
    \verb|PHYS052__HM_0| & \verb|MATH056__HM|, \verb|MATH082__HM|\\\hline
    \verb|PHYS064__HM_0| & \verb|CSCI005__HM|\\\hline
    \verb|PHYS064__HM_1| & \verb|MATH056__HM|, \verb|MATH082__HM|\\\hline
    \verb|PHYS084__HM| & \verb|PHYS024__HM|, \verb|MATH073__HM|\\\hline
    \verb|PHYS084__HM_0| & \verb|CSCI005__HM|\\\hline
    \verb|PHYS111__HM| & \verb|PHYS023__HM|, \verb|PHYS024__HM|\\\hline
    \verb|PHYS111__HM_0| & \verb|MATH082__HM|, \verb|PHYS064__HM|\\\hline
    \verb|PHYS116__HM| & \verb|PHYS052__HM|\\\hline
    \verb|PHYS117__HM| & \verb|PHYS052__HM|\\\hline
    \verb|PHYS133__HM| & \verb|PHYS054__HM|\\\hline
    \verb|PHYS134__HM| & \verb|PHYS051__HM|, \verb|PHYS054__HM|\\\hline
    \verb|PHYS151__HM| & \verb|PHYS051__HM|\\\hline
    \verb|PHYS151__HM_0| & \verb|PHYS111__HM|, \verb|PHYS116__HM|\\\hline
    \verb|PHYS151__HM_1| & \verb|MATH180__HM|, \verb|PHYS064__HM|\\\hline
    \verb|PHYS154__HM| & \verb|MATH180__HM|, \verb|PHYS064__HM|\\\hline
    \verb|PHYS156__HM| & \verb|PHYS111__HM|\\\hline
    \verb|PHYS156__HM_0| & \verb|MATH180__HM|, \verb|PHYS064__HM|\\\hline
    \verb|PHYS161__HM| & \verb|PHYS116__HM|\\\hline
    \verb|PHYS162__HM_pre| & \verb|PHYS117__HM|\\\hline
    \verb|PHYS166__HM| & \verb|PHYS023__HM|, \verb|PHYS024__HM|\\\hline
    \verb|PHYS170X_HM| & \verb|PHYS052__HM|, \verb|PHYS064__HM|, \verb|PHYS111__HM|\\\hline
    \verb|PHYS170__HM| & \verb|PHYS052__HM|, \verb|PHYS064__HM|, \verb|PHYS111__HM|\\\hline
    \verb|PHYS172__HM| & \verb|PHYS111__HM|\\\hline
    \verb|PHYS181__HM| & \verb|PHYS134__HM|\\\hline
    \verb|PHYS190L_KS2| & \verb|PHYS190L_KS|\\\hline
    \verb|PHYS190__PO2| & \verb|PHYS190__PO|\\\hline
    \verb|PHYS191E_PO2| & \verb|PHYS191E_PO|\\\hline
    \verb|PHYS191L_PO2| & \verb|PHYS191L_PO|\\\hline
    \verb|PHYS191__HM2| & \verb|PHYS191__HM|\\\hline
    \verb|PHYS191__KS2| & \verb|PHYS191__KS|\\\hline
    \verb|PHYS193__HM2| & \verb|PHYS193__HM|\\\hline
    \verb|PHYS193__PO2| & \verb|PHYS193__PO|\\\hline
    \verb|PHYS194__HM2| & \verb|PHYS194__HM|\\\hline
    \verb|PHYS195__HM2| & \verb|PHYS195__HM|\\\hline
    \verb|PHYS195__HM3| & \verb|PHYS195__HM2|\\\hline
    \verb|PHYS195__HM4| & \verb|PHYS195__HM3|\\\hline
    \verb|PHYS196__KS2| & \verb|PHYS196__KS|\\\hline
    \verb|PHYS197__HM2| & \verb|PHYS197__HM|\\\hline
    \verb|PHYS197__KS2| & \verb|PHYS197__KS|\\\hline
    \verb|PHYS199__HM2| & \verb|PHYS199__HM|\\\hline
    \verb|POLI190C_PO2| & \verb|POLI190C_PO|\\\hline
    \verb|POLI190D_PO2| & \verb|POLI190D_PO|\\\hline
    \verb|POLI190E_PO2| & \verb|POLI190E_PO|\\\hline
    \verb|POLI191__PO2| & \verb|POLI191__PO|\\\hline
    \verb|POLI191__SC2| & \verb|POLI191__SC|\\\hline
    \verb|POLI193__PO2| & \verb|POLI193__PO|\\\hline
    \verb|POLI195__PO2| & \verb|POLI195__PO|\\\hline
    \verb|POST195__PZ2| & \verb|POST195__PZ|\\\hline
    \verb|POST199__PZ2| & \verb|POST199__PZ|\\\hline
    \verb|PPA_190__PO2| & \verb|PPA_190__PO|\\\hline
    \verb|PPA_191__PO2| & \verb|PPA_191__PO|\\\hline
    \verb|PPA_191__SC2| & \verb|PPA_191__SC|\\\hline
    \verb|PPA_192__SC2| & \verb|PPA_192__SC|\\\hline
    \verb|PPA_195__PO2| & \verb|PPA_195__PO|\\\hline
    \verb|PPE_195__PO2| & \verb|PPE_195__PO|\\\hline
    \verb|PSYC190A_PO2| & \verb|PSYC190A_PO|\\\hline
    \verb|PSYC190R_PO2| & \verb|PSYC190R_PO|\\\hline
    \verb|PSYC190__PO2| & \verb|PSYC190__PO|\\\hline
    \verb|PSYC191__PO2| & \verb|PSYC191__PO|\\\hline
    \verb|PSYC191__PZ2| & \verb|PSYC191__PZ|\\\hline
    \verb|PSYC191__SC2| & \verb|PSYC191__SC|\\\hline
    \verb|PSYC193__SC2| & \verb|PSYC193__SC|\\\hline
    \verb|PSYC194__PZ2| & \verb|PSYC194__PZ|\\\hline
    \verb|PSYC197A_CM2| & \verb|PSYC197A_CM|\\\hline
    \verb|PSYC197B_CM2| & \verb|PSYC197B_CM|\\\hline
    \verb|PSYC197__PZ2| & \verb|PSYC197__PZ|\\\hline
    \verb|PSYC198__CM2| & \verb|PSYC198__CM|\\\hline
    \verb|RLIT191__PO2| & \verb|RLIT191__PO|\\\hline
    \verb|RLST190__PO2| & \verb|RLST190__PO|\\\hline
    \verb|RLST191__PO2| & \verb|RLST191__PO|\\\hline
    \verb|RLST191__SC2| & \verb|RLST191__SC|\\\hline
    \verb|RUSS191__PO2| & \verb|RUSS191__PO|\\\hline
    \verb|RUSS191__SC2| & \verb|RUSS191__SC|\\\hline
    \verb|RUST191__PO2| & \verb|RUST191__PO|\\\hline
    \verb|SCI_196__CM2| & \verb|SCI_196__CM|\\\hline
    \verb|SCI_197__CM2| & \verb|SCI_197__CM|\\\hline
    \verb|SOC_190__PO2| & \verb|SOC_190__PO|\\\hline
    \verb|SOC_191__PO2| & \verb|SOC_191__PO|\\\hline
    \verb|SOC_199A_PZ2| & \verb|SOC_199A_PZ|\\\hline
    \verb|SOC_199B_PZ2| & \verb|SOC_199B_PZ|\\\hline
    \verb|SOC_199C_PZ2| & \verb|SOC_199C_PZ|\\\hline
    \verb|SPAN191__PO2| & \verb|SPAN191__PO|\\\hline
    \verb|SPAN191__SC2| & \verb|SPAN191__SC|\\\hline
    \verb|SPAN192__PO2| & \verb|SPAN192__PO|\\\hline
    \verb|SPAN199__PZ2| & \verb|SPAN199__PZ|\\\hline
    \verb|STS_190__PO2| & \verb|STS_190__PO|\\\hline
    \verb|STS_191__PO2| & \verb|STS_191__PO|\\\hline
    \verb|STS_191__SC2| & \verb|STS_191__SC|\\\hline
    \verb|THEA190__PO2| & \verb|THEA190__PO|\\\hline
    \verb|THEA191H_PO2| & \verb|THEA191H_PO|\\\hline
    \verb|THEA192H_PO2| & \verb|THEA192H_PO|\\\hline
    \verb|THES191D_SC2| & \verb|THES191D_SC|\\\hline
    \verb|THES192D_SC2| & \verb|THES192D_SC|\\\hline
    \verb|WRIT190__PZ2| & \verb|WRIT190__PZ|\\\hline
    \verb|WRIT191__SC2| & \verb|WRIT191__SC|\\\hline
    \verb|WRIT197N_SC2| & \verb|WRIT197N_SC|\\\hline
    \verb|WRIT197Q_SC2| & \verb|WRIT197Q_SC|\\\hline
    \verb|XGOV192__CM2| & \verb|XGOV192__CM|\\\hline
    \caption{Requirement group contents: course prerequisites.}
    \label{tab:reqs-prereqs}
\end{longtable}
\begin{longtable}{|p{4.4in}|p{4.4in}|}
    \hline
    Requirement Group Name & Course Coquirequisites\\\hline
    \verb|BIOL023__HM| & \verb|BIOL046__HM|\\\hline
    \verb|BIOL046__HM| & \verb|BIOL023__HM|\\\hline
    \verb|BIOL046__HM_0| & \verb|CSCI005__HM|\\\hline
    \verb|BIOL054__HM| & \verb|BIOL023__HM|, \verb|BIOL046__HM|\\\hline
    \verb|BIOL110__HM_co| & \verb|BIOL108__HM|\\\hline
    \verb|BIOL111__HM_co| & \verb|BIOL113__HM|\\\hline
    \verb|BIOL154__HM_co| & \verb|BIOL046__HM|\\\hline
    \verb|BIOL184__HM| & \verb|BIOL182__HM|, \verb|CHEM182__HM|\\\hline
    \verb|CHEM024__HM| & \verb|CHEM042__HM|\\\hline
    \verb|CHEM042__HM| & \verb|CHEM024__HM|\\\hline
    \verb|CHEM048__HM_co| & \verb|CHEM042__HM|\\\hline
    \verb|CHEM053__HM| & \verb|CHEM051__HM|\\\hline
    \verb|CHEM058__HM_co| & \verb|CHEM056__HM|\\\hline
    \verb|CHEM109__HM| & \verb|CHEM103__HM|\\\hline
    \verb|CHEM110__HM_co| & \verb|CHEM104__HM|\\\hline
    \verb|CHEM111__HM_co| & \verb|CHEM105__HM|\\\hline
    \verb|CHEM112__HM_co| & \verb|CHEM114__HM|\\\hline
    \verb|CHEM184__HM| & \verb|CHEM182__HM|, \verb|BIOL182__HM|\\\hline
    \verb|Core| & \verb|BIOL023__HM|, \verb|BIOL046__HM|, \verb|CHEM024__HM|, \verb|CHEM042__HM|, \verb|CORE099__HM|, \verb|ENGR079__HM|, \verb|HSA_010__HM|, \verb|MATH019__HM|, \verb|MATH073__HM|, \verb|PHYS023__HM|, \verb|PHYS050__HM|, \verb|WRIT001__HM|\\\hline
    \verb|Core_IntroCS| & \verb|CSCI005__HM|\\\hline
    \verb|Core_Mechanics| & \verb|PHYS024__HM|\\\hline
    \verb|ENGR004__HM_co| & \verb|PHYS024__HM|\\\hline
    \verb|ENGR072__HM_co| & \verb|ENGR080__HM|\\\hline
    \verb|ENGR080__HM_co| & \verb|MATH056__HM|, \verb|ENGR072__HM|\\\hline
    \verb|ENGR157__HM_co| & \verb|ENGR101__HM|\\\hline
    \verb|ENGR175__HM| & \verb|ENGR083__HM|\\\hline
    \verb|HSAs_MuddHums| & \verb|ASAM123__HM|, \verb|ASAM179R_HM|, \verb|PSYC139__HM|, \verb|PSYC179N_HM|, \verb|MUS_049__HM|, \verb|MUS_081__HM|, \verb|MUS_173__JM|, \verb|MUS_175__JM|, \verb|MUS_176__JM|, \verb|MUS_179D_HM|, \verb|MUS_179F_HM|, \verb|RLST114__HM|, \verb|RLST168__HM|, \verb|RLST179G_HM|, \verb|SOC_179C_HM|, \verb|MS__078__PZ|, \verb|MS__120__HM|, \verb|MS__172__HM|, \verb|HIST082__HM|, \verb|ECON053__HM|, \verb|ECON054__HM|, \verb|ECON146__HM|, \verb|POST140__HM|, \verb|LIT_035__HM|, \verb|LIT_107__HM|, \verb|LIT_156__HM|, \verb|LIT_179AJHM|, \verb|LIT_179AKHM|, \verb|AMST179A_HM|, \verb|GEOG179K_HM|, \verb|SOSC150__HM|, \verb|ARCT179A_HM|, \verb|ART_060__HM|, \verb|ART_179S_HM|, \verb|ART_179T_HM|, \verb|ART_179W_HM|, \verb|ART_179X_HM|, \verb|CSCI181ALHM|, \verb|PHIL124__HM|, \verb|ECON052__SC|, \verb|MUS_081__JM|, \verb|LIT_179AIHM|, \verb|PSYC134__HM|, \verb|ASAM124__HM|, \verb|GEOG179J_HM|, \verb|ART_179M_HM|, \verb|ECON154__HM|, \verb|LIT_141__HM|, \verb|STS_179N_HM|, \verb|ART_179V_HM|, \verb|ECON179F_HM|, \verb|GEOG179I_HM|, \verb|HIST081__HM|, \verb|LIT_179ABHM|, \verb|MUS_003__HM|, \verb|MUS_046__HM|, \verb|MUS_179E_HM|, \verb|PHIL179F_HM|, \verb|PSYC179O_HM|, \verb|LIT_179AHHM|, \verb|ASAM126__HM|, \verb|PSYC179U_HM|, \verb|MUS_067__HM|, \verb|MUS_099__HM|, \verb|MUS_179G_HM|, \verb|MUS_179I_HM|, \verb|RLST113__HM|, \verb|RLST183__HM|, \verb|SOC_179E_HM|, \verb|SOC_179F_HM|, \verb|MS__170__HM|, \verb|MS__179I_HM|, \verb|ECON179G_HM|, \verb|POST114__HM|, \verb|POST179D_HM|, \verb|LIT_144__HM|, \verb|LIT_179Y_HM|, \verb|GEOG145__HM|, \verb|ARCT179C_HM|, \verb|GWS_179B_HM|, \verb|PHIL122__HM|, \verb|PSYC179M_HM|, \verb|MUS_179H_HM|, \verb|SOC_179D_HM|, \verb|MS__173__HM|, \verb|HIST151__HM|, \verb|HIST179L_HM|, \verb|ARHI179E_HM|, \verb|POST179C_HM|, \verb|LIT_110__HM|, \verb|GEOG179H_HM|, \verb|ARCT179B_HM|, \verb|GWS_179A_HM|, \verb|PHIL121__HM|, \verb|PHIL179G_HM|, \verb|HSA_000__HM|, \verb|HSA_001__HM|, \verb|HSA_002__HM|, \verb|HSA_003__HM|, \verb|HSA_100__HM|\\\hline
    \verb|HSAs_TenHSAs| & \verb|RUST100__PO|, \verb|ITAL001__SC|, \verb|ITAL033__SC|, \verb|ITAL134__SC|, \verb|ITAL140__SC|, \verb|CGS_010__PZ|, \verb|CGS_080__PZ|, \verb|CGS_122__PZ|, \verb|CGS_127__PZ|, \verb|CGS_146__PZ|, \verb|ASAM070__PO|, \verb|ASAM077B_PZ|, \verb|ASAM082__PZ|, \verb|ASAM089__PZ|, \verb|ASAM090__PZ|, \verb|ASAM123__HM|, \verb|ASAM125__AA|, \verb|ASAM144__PO|, \verb|ASAM161__AA|, \verb|ASAM179B_AA|, \verb|ASAM179H_AA|, \verb|ASAM179R_HM|, \verb|ASAM190A_PO|, \verb|KORE001__CM|, \verb|KORE033__CM|, \verb|KORE120__CM|, \verb|PSYC010__PZ|, \verb|PSYC030__CM|, \verb|PSYC037__CM|, \verb|PSYC040__CM|, \verb|PSYC051__PO|, \verb|PSYC065__CM|, \verb|PSYC070__CM|, \verb|PSYC081__CM|, \verb|PSYC084__CH|, \verb|PSYC092__CM|, \verb|PSYC103__PZ|, \verb|PSYC105__PZ|, \verb|PSYC108__PO|, \verb|PSYC111P_PZ|, \verb|PSYC128__SC|, \verb|PSYC130P_PZ|, \verb|PSYC137__PO|, \verb|PSYC139__HM|, \verb|PSYC144__CM|, \verb|PSYC150__AF|, \verb|PSYC154__PO|, \verb|PSYC155__CM|, \verb|PSYC159__CM|, \verb|PSYC160__PO|, \verb|PSYC162__SC|, \verb|PSYC164__CM|, \verb|PSYC169__SC|, \verb|PSYC176__PO|, \verb|PSYC179N_HM|, \verb|PSYC180S_PO|, \verb|PSYC189__PZ|, \verb|PSYC190A_PO|, \verb|PSYC190R_PO|, \verb|PSYC191__SC|, \verb|PSYC197A_CM|, \verb|PSYC197B_CM|, \verb|PSYC198__CM|, \verb|FREN001__PO|, \verb|FREN001__SC|, \verb|FREN001L_SC|, \verb|FREN002__SC|, \verb|FREN002L_SC|, \verb|FREN011__PO|, \verb|FREN013__PO|, \verb|FREN022__PO|, \verb|FREN033__PO|, \verb|FREN033__SC|, \verb|FREN033L_SC|, \verb|FREN044__PO|, \verb|FREN044__SC|, \verb|FREN045__SC|, \verb|FREN045L_SC|, \verb|FREN100__CM|, \verb|FREN101A_PO|, \verb|FREN104__PO|, \verb|FREN105__PO|, \verb|FREN112__PO|, \verb|FREN115__CM|, \verb|FREN125__SC|, \verb|FREN133__CM|, \verb|FREN161__PZ|, \verb|FREN191__PO|, \verb|FREN192__PO|, \verb|FREN193__PO|, \verb|MUS_003__SC|, \verb|MUS_004__PO|, \verb|MUS_007__PO|, \verb|MUS_008__PO|, \verb|MUS_010__PO|, \verb|MUS_015__PO|, \verb|MUS_015VOPO|, \verb|MUS_020__PO|, \verb|MUS_031__PO|, \verb|MUS_033__PO|, \verb|MUS_035__PO|, \verb|MUS_037__PO|, \verb|MUS_040__PO|, \verb|MUS_041__PO|, \verb|MUS_042C_PO|, \verb|MUS_047__PO|, \verb|MUS_049__HM|, \verb|MUS_051__PO|, \verb|MUS_057__PO|, \verb|MUS_066__SC|, \verb|MUS_070__PO|, \verb|MUS_071__SC|, \verb|MUS_080__PO|, \verb|MUS_080L_PO|, \verb|MUS_081__HM|, \verb|MUS_081__PO|, \verb|MUS_081L_PO|, \verb|MUS_085__SC|, \verb|MUS_087__PO|, \verb|MUS_088__PO|, \verb|MUS_089__SC|, \verb|MUS_091__PO|, \verb|MUS_096A_PO|, \verb|MUS_100__PO|, \verb|MUS_101__SC|, \verb|MUS_103__SC|, \verb|MUS_110A_SC|, \verb|MUS_121__SC|, \verb|MUS_122__PO|, \verb|MUS_130__SC|, \verb|MUS_140__PO|, \verb|MUS_170F_SC|, \verb|MUS_170H_SC|, \verb|MUS_171F_SC|, \verb|MUS_171H_SC|, \verb|MUS_172__SC|, \verb|MUS_173__JM|, \verb|MUS_175__JM|, \verb|MUS_176__JM|, \verb|MUS_177F_SC|, \verb|MUS_177H_SC|, \verb|MUS_179D_HM|, \verb|MUS_179F_HM|, \verb|MUS_180F_SC|, \verb|MUS_180H_SC|, \verb|MUS_184__PO|, \verb|MUS_190__PO|, \verb|GLAS194A_PZ|, \verb|GLAS194B_PZ|, \verb|ORST060__PZ|, \verb|ORST100__PZ|, \verb|ORST100A_PZ|, \verb|ORST175__PZ|, \verb|ORST198M_PZ|, \verb|GERM001__PO|, \verb|GERM001__SC|, \verb|GERM011__PO|, \verb|GERM013__PO|, \verb|GERM033__PO|, \verb|GERM033__SC|, \verb|GERM107__SC|, \verb|GERM152__PO|, \verb|GERM191__PO|, \verb|GERM193__PO|, \verb|WRIT100A_PZ|, \verb|WRIT100B_PZ|, \verb|WRIT105__SC|, \verb|WRIT132__SC|, \verb|WRIT153IOSC|, \verb|WRIT160__SC|, \verb|AFRI010A_AF|, \verb|AFRI010B_AF|, \verb|AFRI114__AF|, \verb|AFRI116__AF|, \verb|AFRI120__PZ|, \verb|AFRI121__AF|, \verb|AFRI125__AF|, \verb|AFRI150__PZ|, \verb|AFRI191__AF|, \verb|CLAS010__SC|, \verb|CLAS012__SC|, \verb|CLAS019__SC|, \verb|CLAS117__PO|, \verb|CLAS161__PZ|, \verb|CLAS190__PO|, \verb|THEA001A_PO|, \verb|THEA001G_PO|, \verb|THEA002__PO|, \verb|THEA006__PO|, \verb|THEA012__PO|, \verb|THEA021__PO|, \verb|THEA023__PO|, \verb|THEA030__PO|, \verb|THEA041__PO|, \verb|THEA051C_PO|, \verb|THEA051H_PO|, \verb|THEA052C_PO|, \verb|THEA052H_PO|, \verb|THEA054C_PO|, \verb|THEA080__PO|, \verb|THEA082__PO|, \verb|THEA089C_PO|, \verb|THEA100G_PO|, \verb|THEA130__PO|, \verb|THEA141__PO|, \verb|THEA170__PO|, \verb|THEA190__PO|, \verb|CHLT072__CH|, \verb|RUSS001__PO|, \verb|RUSS011__PO|, \verb|RUSS013__PO|, \verb|RUSS033__PO|, \verb|RUSS187__PO|, \verb|RUSS191__PO|, \verb|PORT001__PZ|, \verb|PORT033__CM|, \verb|PORT035__PZ|, \verb|GREK022__PO|, \verb|GREK033__PO|, \verb|GREK044__PO|, \verb|IR__100__PO|, \verb|IR__190__PO|, \verb|FLAN191__SC|, \verb|LGCS010__PO|, \verb|LGCS105__PO|, \verb|LGCS108__PO|, \verb|LGCS112__PZ|, \verb|LGCS121__PO|, \verb|LGCS124__PO|, \verb|LGCS125__PO|, \verb|LGCS183__PO|, \verb|LGCS190__PO|, \verb|RLST037__CM|, \verb|RLST038__CM|, \verb|RLST045__CM|, \verb|RLST056__CM|, \verb|RLST058__CM|, \verb|RLST092__SC|, \verb|RLST093__CM|, \verb|RLST100__PO|, \verb|RLST102__CM|, \verb|RLST103__PO|, \verb|RLST114__HM|, \verb|RLST142__AF|, \verb|RLST148__PO|, \verb|RLST150__AF|, \verb|RLST157__PO|, \verb|RLST166B_CM|, \verb|RLST168__HM|, \verb|RLST171__CM|, \verb|RLST179G_HM|, \verb|RLST181__PO|, \verb|RLST190__PO|, \verb|SOC_001__PZ|, \verb|SOC_051__PO|, \verb|SOC_061__PZ|, \verb|SOC_072__PZ|, \verb|SOC_091__PZ|, \verb|SOC_101__PZ|, \verb|SOC_102__PO|, \verb|SOC_102__PZ|, \verb|SOC_109__PZ|, \verb|SOC_110__PZ|, \verb|SOC_119__PO|, \verb|SOC_122__PZ|, \verb|SOC_126__AA|, \verb|SOC_148__PO|, \verb|SOC_169X_PO|, \verb|SOC_179C_HM|, \verb|SOC_189F_PO|, \verb|SOC_190__PO|, \verb|SOC_199A_PZ|, \verb|ARBC001__CM|, \verb|ARBC033__CM|, \verb|ARBC166__CM|, \verb|MS__038__SC|, \verb|MS__041__SC|, \verb|MS__045__PZ|, \verb|MS__049__PO|, \verb|MS__050__PO|, \verb|MS__051__PZ|, \verb|MS__051__SC|, \verb|MS__052__PZ|, \verb|MS__055__PZ|, \verb|MS__057__SC|, \verb|MS__059__SC|, \verb|MS__078__PZ|, \verb|MS__082__PZ|, \verb|MS__120__HM|, \verb|MS__120__PZ|, \verb|MS__120__SC|, \verb|MS__126__PZ|, \verb|MS__131__SC|, \verb|MS__150__PO|, \verb|MS__172__HM|, \verb|MS__190F_JT|, \verb|ANTH002__PO|, \verb|ANTH002__PZ|, \verb|ANTH002__SC|, \verb|ANTH022__SC|, \verb|ANTH055__PZ|, \verb|ANTH104__PO|, \verb|ANTH105__PZ|, \verb|ANTH107__PZ|, \verb|ANTH107__SC|, \verb|ANTH109__SC|, \verb|ANTH111__PZ|, \verb|ANTH116__PO|, \verb|ANTH123__PO|, \verb|ANTH153__SC|, \verb|ANTH184__PO|, \verb|ANTH190__PO|, \verb|ANTH190__SC|, \verb|SPCH061A_CM|, \verb|SPCH061B_CM|, \verb|MLLC111__PZ|, \verb|MLLC122__PZ|, \verb|MLLC150__PZ|, \verb|HIST008__PO|, \verb|HIST011__PZ|, \verb|HIST012__PZ|, \verb|HIST014__PO|, \verb|HIST017__CH|, \verb|HIST020__PO|, \verb|HIST031__CH|, \verb|HIST034__CH|, \verb|HIST040__AF|, \verb|HIST054__CM|, \verb|HIST055__CM|, \verb|HIST060__PO|, \verb|HIST060__PZ|, \verb|HIST070__PO|, \verb|HIST072__CM|, \verb|HIST082__HM|, \verb|HIST083__PZ|, \verb|HIST097__CM|, \verb|HIST100__CM|, \verb|HIST101F_PO|, \verb|HIST102__PO|, \verb|HIST111__SC|, \verb|HIST112A_CM|, \verb|HIST115__CM|, \verb|HIST119__CM|, \verb|HIST123__CM|, \verb|HIST125__PO|, \verb|HIST128__PO|, \verb|HIST132__CM|, \verb|HIST133__AF|, \verb|HIST134__SC|, \verb|HIST139__PO|, \verb|HIST139E_CM|, \verb|HIST141__SC|, \verb|HIST144__SC|, \verb|HIST146__AF|, \verb|HIST150__CM|, \verb|HIST151__CM|, \verb|HIST158__CM|, \verb|HIST161__CM|, \verb|HIST168__CM|, \verb|HIST169__CM|, \verb|HIST171__PO|, \verb|HIST175__PO|, \verb|HIST176__CM|, \verb|HIST179__CM|, \verb|HIST181__PZ|, \verb|HIST185__PO|, \verb|HIST186__PO|, \verb|HIST188__CM|, \verb|HIST190__PO|, \verb|HIST191__SC|, \verb|HIST196__CM|, \verb|HIST197__PZ|, \verb|ECON050__CM|, \verb|ECON050__SC|, \verb|ECON051__PO|, \verb|ECON051__PZ|, \verb|ECON051__SC|, \verb|ECON052__PO|, \verb|ECON052__PZ|, \verb|ECON053__HM|, \verb|ECON054__HM|, \verb|ECON097__CM|, \verb|ECON100__CM|, \verb|ECON101__CM|, \verb|ECON101__PO|, \verb|ECON101__SC|, \verb|ECON102__CM|, \verb|ECON102__PO|, \verb|ECON102__SC|, \verb|ECON105__PZ|, \verb|ECON118__CM|, \verb|ECON123__PO|, \verb|ECON124__PO|, \verb|ECON126__PO|, \verb|ECON127__PO|, \verb|ECON129__CM|, \verb|ECON135__CM|, \verb|ECON139__CM|, \verb|ECON140__PZ|, \verb|ECON141__CM|, \verb|ECON142__PZ|, \verb|ECON143__PO|, \verb|ECON145__PZ|, \verb|ECON146__HM|, \verb|ECON154__PO|, \verb|ECON157__PO|, \verb|ECON159__PO|, \verb|ECON171__CM|, \verb|ECON175__SC|, \verb|ECON179__CM|, \verb|ECON180__PZ|, \verb|ECON186__CM|, \verb|ECON191__SC|, \verb|ECON198__PZ|, \verb|ARHI001A_PO|, \verb|ARHI120__PO|, \verb|ARHI138__PO|, \verb|ARHI140__PO|, \verb|ARHI150__SC|, \verb|ARHI154__SC|, \verb|ARHI178__PO|, \verb|ARHI179__PO|, \verb|ARHI185__SC|, \verb|ARHI186A_PZ|, \verb|ARHI186C_SC|, \verb|ARHI186M_SC|, \verb|ARHI186Y_PO|, \verb|ARHI189__SC|, \verb|ARHI190__PO|, \verb|FGSS026__SC|, \verb|FGSS036__SC|, \verb|FGSS177__SC|, \verb|FGSS184__SC|, \verb|FGSS188__SC|, \verb|GOVT020__CM|, \verb|GOVT020H_CM|, \verb|GOVT041__CM|, \verb|GOVT041B_CM|, \verb|GOVT054__CM|, \verb|GOVT055__CM|, \verb|GOVT060__CM|, \verb|GOVT070__CM|, \verb|GOVT071__CM|, \verb|GOVT080__CM|, \verb|GOVT097__CM|, \verb|GOVT100__CM|, \verb|GOVT110__CM|, \verb|GOVT112A_CM|, \verb|GOVT112B_CM|, \verb|GOVT115__CM|, \verb|GOVT118__CM|, \verb|GOVT134__CM|, \verb|GOVT137__CM|, \verb|GOVT137B_CM|, \verb|GOVT139__CM|, \verb|GOVT140__CM|, \verb|GOVT142E_CM|, \verb|GOVT144D_CM|, \verb|GOVT145E_CM|, \verb|GOVT146A_CM|, \verb|GOVT148__CM|, \verb|GOVT149__CM|, \verb|GOVT150C_CM|, \verb|GOVT152__CM|, \verb|GOVT166__CM|, \verb|GOVT173C_CM|, \verb|GOVT174C_CM|, \verb|GOVT177C_CM|, \verb|GOVT179__CM|, \verb|HUM_196__PO|, \verb|EA__010__PO|, \verb|EA__010__PZ|, \verb|EA__010__SC|, \verb|EA__020__PO|, \verb|EA__032__PZ|, \verb|EA__034__PZ|, \verb|EA__086__PZ|, \verb|EA__101__PO|, \verb|EA__101L_PO|, \verb|EA__129__PZ|, \verb|EA__133__PZ|, \verb|EA__138__PZ|, \verb|EA__141__PZ|, \verb|EA__142__PZ|, \verb|EA__150__PZ|, \verb|EA__163__PZ|, \verb|EA__170__PO|, \verb|EA__180__PO|, \verb|EA__191__PO|, \verb|POST016X_PZ|, \verb|POST020A_PZ|, \verb|POST030__PZ|, \verb|POST040__PZ|, \verb|POST050__PZ|, \verb|POST101__PZ|, \verb|POST106__PZ|, \verb|POST124__PZ|, \verb|POST140__HM|, \verb|POST195__PZ|, \verb|ENGL001__PZ|, \verb|ENGL002__PZ|, \verb|ENGL010__PO|, \verb|ENGL011A_PZ|, \verb|ENGL012B_AF|, \verb|ENGL019__PO|, \verb|ENGL030__PZ|, \verb|ENGL040__PO|, \verb|ENGL041__PO|, \verb|ENGL074__PO|, \verb|ENGL087F_PO|, \verb|ENGL093__PO|, \verb|ENGL101__SC|, \verb|ENGL102__PZ|, \verb|ENGL118S_SC|, \verb|ENGL121__SC|, \verb|ENGL124__SC|, \verb|ENGL128__PZ|, \verb|ENGL133S_SC|, \verb|ENGL141__SC|, \verb|ENGL145__SC|, \verb|ENGL165__PO|, \verb|ENGL170U_PO|, \verb|ENGL174__SC|, \verb|ENGL181__PZ|, \verb|ENGL183A_PO|, \verb|ENGL183B_PO|, \verb|ENGL188__PO|, \verb|ENGL190__SC|, \verb|ENGL191__PO|, \verb|ENGL191__SC|, \verb|ENGL194__PZ|, \verb|ENGL195__PO|, \verb|ENGL199T_SC|, \verb|HMSC138__SC|, \verb|HMSC191__SC|, \verb|LIT_031__CM|, \verb|LIT_034__CM|, \verb|LIT_035__HM|, \verb|LIT_037__CM|, \verb|LIT_038__CM|, \verb|LIT_057__CM|, \verb|LIT_060__CM|, \verb|LIT_066__CM|, \verb|LIT_090__CM|, \verb|LIT_097__CM|, \verb|LIT_099__CM|, \verb|LIT_100__CM|, \verb|LIT_107__HM|, \verb|LIT_118__CM|, \verb|LIT_129__CM|, \verb|LIT_138__CM|, \verb|LIT_139__CM|, \verb|LIT_148__CM|, \verb|LIT_156__HM|, \verb|LIT_162__AF|, \verb|LIT_179AJHM|, \verb|LIT_179AKHM|, \verb|LIT_182__CM|, \verb|CHST064__CH|, \verb|CHST128__CH|, \verb|CHST130__CH|, \verb|CHST136__CH|, \verb|CHST184D_CH|, \verb|CHST189B_PO|, \verb|CHST190__CH|, \verb|AMST103__JT|, \verb|AMST179A_HM|, \verb|AMST190__PO|, \verb|GEOG179K_HM|, \verb|JAPN001A_PO|, \verb|JAPN011__PO|, \verb|JAPN012A_PO|, \verb|JAPN013__PO|, \verb|JAPN014A_PO|, \verb|JAPN051A_PO|, \verb|JAPN111A_PO|, \verb|POLI001A_PO|, \verb|POLI003__PO|, \verb|POLI005__PO|, \verb|POLI007__PO|, \verb|POLI008__PO|, \verb|POLI033A_PO|, \verb|POLI090__PO|, \verb|POLI100__SC|, \verb|POLI103__SC|, \verb|POLI110__SC|, \verb|POLI111__PO|, \verb|POLI122__SC|, \verb|POLI131__PO|, \verb|POLI134__SC|, \verb|POLI135__PO|, \verb|POLI144__SC|, \verb|POLI145__SC|, \verb|POLI150__SC|, \verb|POLI151__SC|, \verb|POLI152__PO|, \verb|POLI161__PO|, \verb|POLI173__SC|, \verb|POLI179A_PO|, \verb|POLI181__PO|, \verb|POLI190D_PO|, \verb|POLI191__PO|, \verb|POLI191__SC|, \verb|POLI193__PO|, \verb|POLI195__PO|, \verb|CHIN001A_PO|, \verb|CHIN002__PO|, \verb|CHIN011__PO|, \verb|CHIN013__PO|, \verb|CHIN051A_PO|, \verb|CHIN111A_PO|, \verb|CHIN129__PO|, \verb|CHIN131__PO|, \verb|CHIN153__PO|, \verb|CHIN192A_PO|, \verb|ID__199P1PO|, \verb|ID__199P3PO|, \verb|ID__199S1PO|, \verb|ID__199S2PO|, \verb|PONT100__CM|, \verb|SOSC150__HM|, \verb|PPA_190__PO|, \verb|PPA_195__PO|, \verb|RLIT191__PO|, \verb|ARCN115__SC|, \verb|STS_080__PO|, \verb|STS_190__PO|, \verb|ARCT179A_HM|, \verb|ART_005__PO|, \verb|ART_010__PO|, \verb|ART_011__PZ|, \verb|ART_012__PZ|, \verb|ART_015__PZ|, \verb|ART_027A_PO|, \verb|ART_028__PO|, \verb|ART_030__PZ|, \verb|ART_033__PO|, \verb|ART_060__HM|, \verb|ART_101__SC|, \verb|ART_105__SC|, \verb|ART_111__PO|, \verb|ART_114__PO|, \verb|ART_116__SC|, \verb|ART_120__PZ|, \verb|ART_121__SC|, \verb|ART_123__SC|, \verb|ART_125__SC|, \verb|ART_126__PZ|, \verb|ART_126B_PO|, \verb|ART_134__SC|, \verb|ART_135__SC|, \verb|ART_136__SC|, \verb|ART_141__SC|, \verb|ART_144__SC|, \verb|ART_145__SC|, \verb|ART_147__SC|, \verb|ART_148__SC|, \verb|ART_179S_HM|, \verb|ART_179T_HM|, \verb|ART_179W_HM|, \verb|ART_179X_HM|, \verb|ART_181M_SC|, \verb|ART_189__PZ|, \verb|ART_190__PO|, \verb|ART_192__SC|, \verb|DANC010__PO|, \verb|DANC012__PO|, \verb|DANC012P_PO|, \verb|DANC050__PO|, \verb|DANC050P_PO|, \verb|DANC051__PO|, \verb|DANC051P_PO|, \verb|DANC068__SC|, \verb|DANC071A_SC|, \verb|DANC071B_SC|, \verb|DANC100B_SC|, \verb|DANC111A_SC|, \verb|DANC111B_SC|, \verb|DANC119__PO|, \verb|DANC120__PO|, \verb|DANC120P_PO|, \verb|DANC122__PO|, \verb|DANC122P_PO|, \verb|DANC123__PO|, \verb|DANC124__PO|, \verb|DANC124P_PO|, \verb|DANC130__PO|, \verb|DANC135__PO|, \verb|DANC136__PO|, \verb|DANC136B_SC|, \verb|DANC140__PO|, \verb|DANC150C_PO|, \verb|DANC151__PO|, \verb|DANC151P_PO|, \verb|DANC152__PO|, \verb|DANC152IOSC|, \verb|DANC152P_PO|, \verb|DANC159__SC|, \verb|DANC162A_SC|, \verb|DANC162B_SC|, \verb|DANC163__SC|, \verb|DANC172__PO|, \verb|DANC173__PO|, \verb|DANC174__PO|, \verb|DANC175__PO|, \verb|DANC176__PO|, \verb|DANC180__PO|, \verb|DANC180P_PO|, \verb|DANC181__PO|, \verb|DANC181P_PO|, \verb|DANC190__SC|, \verb|DANC191__SC|, \verb|DANC192__PO|, \verb|DANC193__SC|, \verb|ASIA190__PO|, \verb|PPE_110A_CM|, \verb|PPE_110B_CM|, \verb|PPE_160__PO|, \verb|SPAN001__CM|, \verb|SPAN001__PO|, \verb|SPAN001__PZ|, \verb|SPAN011__PO|, \verb|SPAN013__PO|, \verb|SPAN022__CM|, \verb|SPAN022__PO|, \verb|SPAN022__PZ|, \verb|SPAN022__SC|, \verb|SPAN031__PZ|, \verb|SPAN033__CM|, \verb|SPAN033__PO|, \verb|SPAN033__PZ|, \verb|SPAN033__SC|, \verb|SPAN033L_CM|, \verb|SPAN035__PZ|, \verb|SPAN044__CM|, \verb|SPAN044__PO|, \verb|SPAN044__PZ|, \verb|SPAN044__SC|, \verb|SPAN044L_CM|, \verb|SPAN050__PZ|, \verb|SPAN055__PZ|, \verb|SPAN100__PZ|, \verb|SPAN101__PO|, \verb|SPAN101__SC|, \verb|SPAN102__CM|, \verb|SPAN104__PZ|, \verb|SPAN109__PO|, \verb|SPAN120A_PO|, \verb|SPAN125A_CM|, \verb|SPAN126__PO|, \verb|SPAN136__PZ|, \verb|SPAN139__SC|, \verb|SPAN140__PO|, \verb|SPAN144__PZ|, \verb|SPAN163__PZ|, \verb|SPAN163__SC|, \verb|SPAN181__SC|, \verb|SPAN183__CM|, \verb|SPAN191__PO|, \verb|SPAN191__SC|, \verb|CASA101__PZ|, \verb|CASA105__PZ|, \verb|GWS_026__PO|, \verb|GWS_166__PO|, \verb|GWS_170__PO|, \verb|GWS_172__PO|, \verb|GWS_180__PO|, \verb|GWS_184__PO|, \verb|GWS_190__PO|, \verb|JPNT175__PO|, \verb|CSCI181ALHM|, \verb|LAST184__PO|, \verb|LAST190__PO|, \verb|PHIL001__PO|, \verb|PHIL007__PO|, \verb|PHIL007__PZ|, \verb|PHIL030__CM|, \verb|PHIL030__PZ|, \verb|PHIL031__CM|, \verb|PHIL031__PO|, \verb|PHIL031__PZ|, \verb|PHIL033__PO|, \verb|PHIL034__CM|, \verb|PHIL034__PO|, \verb|PHIL038__CM|, \verb|PHIL038__PO|, \verb|PHIL040__PO|, \verb|PHIL054__PO|, \verb|PHIL060__PO|, \verb|PHIL090__SC|, \verb|PHIL100E_CM|, \verb|PHIL120__PO|, \verb|PHIL124__HM|, \verb|PHIL134__CM|, \verb|PHIL139__CM|, \verb|PHIL140__SC|, \verb|PHIL144__SC|, \verb|PHIL155__SC|, \verb|PHIL158__CM|, \verb|PHIL160__CM|, \verb|PHIL162__SC|, \verb|PHIL163__SC|, \verb|PHIL164__CM|, \verb|PHIL170__SC|, \verb|PHIL185C_PO|, \verb|PHIL190__CM|, \verb|PHIL190__PO|, \verb|PHIL190__SC|, \verb|PHIL191__PZ|, \verb|PHIL198__CM|, \verb|LATN022__PO|, \verb|LATN033__SC|, \verb|LATN044__PO|, \verb|LATN103__PO|, \verb|ECON165__PZ|, \verb|ECON175__PZ|, \verb|GOVT132E_CM|, \verb|ECON142__CM|, \verb|EA__112__PZ|, \verb|MS__190__JT|, \verb|SOC_128__JT|, \verb|CGS_050__PZ|, \verb|CGS_150__PZ|, \verb|CGS_190__PZ|, \verb|SOC_069__PZ|, \verb|HIST069__PZ|, \verb|ANTH003__PZ|, \verb|ANTH012__PZ|, \verb|MS__072__PO|, \verb|AMST191__PO|, \verb|ASAM085__PZ|, \verb|ASAM130__PZ|, \verb|ASAM172__PZ|, \verb|STS_191__PO|, \verb|ECON104__PZ|, \verb|LGCS115__PZ|, \verb|MLLC144__PZ|, \verb|MLLC155__PZ|, \verb|MLLC188__PZ|, \verb|ORST135__PZ|, \verb|ORST180__PZ|, \verb|PORT002__PZ|, \verb|SOC_083__PZ|, \verb|SOC_095__PZ|, \verb|SOC_170__PZ|, \verb|SOC_199B_PZ|, \verb|SPAN002__PZ|, \verb|AMST113__SC|, \verb|AMST180__SC|, \verb|ANTH025__SC|, \verb|ANTH027__SC|, \verb|ANTH114__SC|, \verb|ANTH185P_SC|, \verb|ARHI186N_SC|, \verb|ART_106__SC|, \verb|ART_126B_SC|, \verb|ART_142__SC|, \verb|ART_143__SC|, \verb|ART_146__SC|, \verb|ART_193__SC|, \verb|ASAM160__AA|, \verb|CHST015__CH|, \verb|CHST066__CH|, \verb|CHST120__CH|, \verb|CLAS116__SC|, \verb|DANC083A_SC|, \verb|DANC083B_SC|, \verb|DANC100A_SC|, \verb|DANC114A_SC|, \verb|DANC114B_SC|, \verb|DANC131__SC|, \verb|DANC161__SC|, \verb|ECON052__SC|, \verb|ECON114__SC|, \verb|ECON133__SC|, \verb|ECON170__SC|, \verb|ECON184__SC|, \verb|ENGL102__SC|, \verb|ENGL115__SC|, \verb|ENGL116__SC|, \verb|ENGL131__SC|, \verb|ENGL180__SC|, \verb|ENGL193__SC|, \verb|FGSS183__SC|, \verb|FGSS186__SC|, \verb|FGSS188E_SC|, \verb|GERM002__SC|, \verb|GERM044__SC|, \verb|GREK033__SC|, \verb|GREK044__SC|, \verb|HIST048__SC|, \verb|HIST143__SC|, \verb|HIST149__SC|, \verb|HIST180__SC|, \verb|HUM_195J_SC|, \verb|ITAL002__SC|, \verb|ITAL044__SC|, \verb|ITAL131__SC|, \verb|LATN044__SC|, \verb|MS__160__SC|, \verb|MUS_072__SC|, \verb|MUS_081__JM|, \verb|MUS_102__SC|, \verb|MUS_110B_SC|, \verb|MUS_112__SC|, \verb|MUS_126__SC|, \verb|PHIL118__SC|, \verb|PHIL160__SC|, \verb|POLI116__SC|, \verb|POLI118__SC|, \verb|POLI123__SC|, \verb|POLI146__SC|, \verb|POLI148__SC|, \verb|POLI153__SC|, \verb|POLI180__SC|, \verb|PSYC052__SC|, \verb|PSYC124__SC|, \verb|PSYC131L_SC|, \verb|PSYC133__SC|, \verb|PSYC168__SC|, \verb|PSYC168L_SC|, \verb|PSYC183__SC|, \verb|SPAN127__CH|, \verb|SPAN131__SC|, \verb|SPAN140__SC|, \verb|SPAN183__SC|, \verb|WRIT109__SC|, \verb|WRIT112__SC|, \verb|WRIT140__SC|, \verb|WRIT191__SC|, \verb|AFRI020__PZ|, \verb|THEA061__PO|, \verb|ARBC002__CM|, \verb|LIT_179AIHM|, \verb|ARBC044__CM|, \verb|PSYC134__HM|, \verb|ASAM124__HM|, \verb|GEOG179J_HM|, \verb|ECON126__CM|, \verb|CHIN055__PO|, \verb|ECON172__CM|, \verb|ASAM142__PO|, \verb|LGCS117__PO|, \verb|ECON191__CM|, \verb|GOVT131__CM|, \verb|GOVT070H_CM|, \verb|GOVT101__CM|, \verb|MS__102__PZ|, \verb|GOVT117__CM|, \verb|MS__124__PZ|, \verb|GOVT128__JT|, \verb|GOVT137A_CM|, \verb|MS__125__PZ|, \verb|GOVT142__CM|, \verb|PHIL036__PO|, \verb|GOVT138C_CM|, \verb|ART_140__PO|, \verb|GOVT999__CM|, \verb|CHLT085__PZ|, \verb|CHLT170__CH|, \verb|HIST056__CM|, \verb|CGS_151__PZ|, \verb|ASAM191__PO|, \verb|PHIL191__PO|, \verb|AFRI191C_AF|, \verb|HIST134__CM|, \verb|HIST166__CM|, \verb|HIST174__CM|, \verb|HIST190__CM|, \verb|KORE002__CM|, \verb|KORE044__CM|, \verb|KRNT130__CM|, \verb|LIT_058__CM|, \verb|LIT_099A_CM|, \verb|LIT_130__CM|, \verb|PHIL035__CM|, \verb|PHIL095__CM|, \verb|PHIL135__CM|, \verb|PHIL138__CM|, \verb|PORT022__CM|, \verb|PPE_001A_CM|, \verb|PPE_001B_CM|, \verb|PPE_011A_CM|, \verb|PPE_011B_CM|, \verb|PSYC132__CM|, \verb|PSYC154__CM|, \verb|PSYC188__CM|, \verb|RLST010__CM|, \verb|RLST125__CM|, \verb|SPAN002__CM|, \verb|SPAN130__CM|, \verb|ART_179M_HM|, \verb|ECON154__HM|, \verb|HSA_010__HM|, \verb|LIT_141__HM|, \verb|STS_179N_HM|, \verb|AFRI134__PO|, \verb|AMST118__PO|, \verb|ANTH053__PO|, \verb|ANTH105__PO|, \verb|ANTH107__PO|, \verb|ANTH189P_PO|, \verb|ARHI001B_PO|, \verb|ARHI131__PO|, \verb|ARHI174__PO|, \verb|ARHI186W_PO|, \verb|ARHI191__PO|, \verb|ART_021__PO|, \verb|ART_115__PO|, \verb|ART_139__PO|, \verb|ART_188__PO|, \verb|ASAM143__PO|, \verb|ASIA081__PO|, \verb|CHIN001B_PO|, \verb|CHIN051B_PO|, \verb|CHIN051H_PO|, \verb|CHIN111B_PO|, \verb|CHNT168__PO|, \verb|CHNT180__PO|, \verb|CHST028__CH|, \verb|CHST132__CH|, \verb|CLAS001__PO|, \verb|CLAS116A_PO|, \verb|DANC138__PO|, \verb|DANC139__PO|, \verb|DANC150A_PO|, \verb|DANC160__PO|, \verb|EA__185__PO|, \verb|ECON089A_PO|, \verb|ECON121__PO|, \verb|ECON122__PO|, \verb|ECON125__PO|, \verb|ECON128__PO|, \verb|ECON130__PO|, \verb|ECON155__PO|, \verb|ECON161__PO|, \verb|ECON190__PO|, \verb|ENGL020__PO|, \verb|ENGL023__PO|, \verb|ENGL064A_PO|, \verb|ENGL064B_PO|, \verb|ENGL064C_PO|, \verb|ENGL086__PO|, \verb|ENGL099__PO|, \verb|ENGL125C_AF|, \verb|FREN002__PO|, \verb|FREN150C_PO|, \verb|GERM002__PO|, \verb|HIST011__PO|, \verb|HIST021__PO|, \verb|HIST032__CH|, \verb|HIST036__PO|, \verb|HIST041__AF|, \verb|HIST062__PO|, \verb|HIST071__PO|, \verb|HIST101A_PO|, \verb|HIST101D_PO|, \verb|HIST101Q_PO|, \verb|HIST118__PO|, \verb|HIST130__CH|, \verb|HIST140__AF|, \verb|HIST148__PO|, \verb|HIST167__PO|, \verb|HIST168__PO|, \verb|HIST180__PO|, \verb|IR__118__PO|, \verb|JAPN001B_PO|, \verb|JAPN012B_PO|, \verb|JAPN014B_PO|, \verb|JAPN051B_PO|, \verb|JAPN111B_PO|, \verb|JPNT171__PO|, \verb|LGCS101__PO|, \verb|LGCS114__PO|, \verb|LGCS119__PO|, \verb|LGCS120__PO|, \verb|MS__051__PO|, \verb|MS__131__PO|, \verb|MS__148D_PO|, \verb|MS__149G_PO|, \verb|MS__149T_PO|, \verb|MS__153__PO|, \verb|MS__175__PO|, \verb|MUS_032__PO|, \verb|MUS_042B_PO|, \verb|MUS_060__PO|, \verb|MUS_062__PO|, \verb|MUS_065__PO|, \verb|MUS_068__PO|, \verb|MUS_082__PO|, \verb|MUS_082L_PO|, \verb|MUS_096B_PO|, \verb|MUS_121__PO|, \verb|PHIL030__PO|, \verb|PHIL035__PO|, \verb|PHIL037__PO|, \verb|PHIL039__PO|, \verb|PHIL042__PO|, \verb|PHIL057__JT|, \verb|PHIL080__PO|, \verb|PHIL104__PO|, \verb|POLI001B_PO|, \verb|POLI002__PO|, \verb|POLI030__PO|, \verb|POLI033B_PO|, \verb|POLI112__PO|, \verb|POLI133__JT|, \verb|POLI158__PO|, \verb|POLI168__PO|, \verb|POLI171__PO|, \verb|POLI177__PO|, \verb|POLI190C_PO|, \verb|PPE_100__PO|, \verb|PSYC106__PO|, \verb|PSYC125__AF|, \verb|PSYC131__PO|, \verb|PSYC153__AA|, \verb|PSYC162__PO|, \verb|PSYC180H_PO|, \verb|PSYC180M_PO|, \verb|PSYC180N_CH|, \verb|RLST107__PO|, \verb|RLST110__PO|, \verb|RLST146__PO|, \verb|RLST158__PO|, \verb|RLST164__PO|, \verb|RLST180__PO|, \verb|RLST184__PO|, \verb|RLST189L_PO|, \verb|RLST191__PO|, \verb|RUSS002__PO|, \verb|RUSS044__PO|, \verb|RUST112__PO|, \verb|RUST185__PO|, \verb|RUST191__PO|, \verb|SOC_090__PO|, \verb|SOC_103__JT|, \verb|SOC_146__PO|, \verb|SOC_147__PO|, \verb|SOC_150__AA|, \verb|SOC_150__CH|, \verb|SOC_154__PO|, \verb|SOC_191__PO|, \verb|SPAN002__PO|, \verb|SPAN106__PO|, \verb|SPAN108__PO|, \verb|SPAN125B_PO|, \verb|SPAN153__PO|, \verb|SPAN185__PO|, \verb|THEA001D_PO|, \verb|THEA010__PO|, \verb|THEA017__PO|, \verb|THEA022__PO|, \verb|THEA026__PO|, \verb|THEA031__PO|, \verb|THEA053HGPO|, \verb|THEA081__PO|, \verb|THEA085__PO|, \verb|THEA100K_PO|, \verb|THEA188__PO|, \verb|XGOV192__CM|, \verb|PSYC127__PZ|, \verb|PHIL084__PZ|, \verb|ANTH113__SC|, \verb|RLST061__SC|, \verb|PORT044__PZ|, \verb|SOC_157__PZ|, \verb|PHIL033__CM|, \verb|PHIL037__CM|, \verb|PHIL101C_CM|, \verb|PHIL121__CM|, \verb|PHIL192__CM|, \verb|CLAS121__JT|, \verb|CLAS150BEPZ|, \verb|ART_113__PZ|, \verb|RLST083__CM|, \verb|RLST109__CM|, \verb|RLST129__CM|, \verb|RLST135__CM|, \verb|ORST114__PZ|, \verb|ARCN120__SC|, \verb|FREN181__SC|, \verb|SOC_119__PZ|, \verb|CGS_060__PZ|, \verb|CGS_075__PZ|, \verb|FREN002__PZ|, \verb|SPAN137__PZ|, \verb|SPAN199__PZ|, \verb|LIT_061__CM|, \verb|LIT_062__CM|, \verb|LIT_068__CM|, \verb|ENGL130__PZ|, \verb|LIT_111__CM|, \verb|LIT_131__CM|, \verb|LIT_152__CM|, \verb|GOVT102__CM|, \verb|GOVT123__CM|, \verb|GOVT147__CM|, \verb|GOVT156E_CM|, \verb|GOVT164__CM|, \verb|GOVT177__CM|, \verb|GOVT178__CM|, \verb|GOVT129A_CM|, \verb|ENGL060__PZ|, \verb|LIT_099B_CM|, \verb|LIT_129IOCM|, \verb|LIT_134C_CM|, \verb|ART_199__PZ|, \verb|MS__088__PZ|, \verb|MS__087__PZ|, \verb|ENGL110__PZ|, \verb|ENGL048__PZ|, \verb|SOC_001X_PZ|, \verb|WRIT101__PZ|, \verb|HIST096__CM|, \verb|HIST101__CM|, \verb|HIST116IOCM|, \verb|HIST126__CM|, \verb|HIST135__CM|, \verb|HIST154__CM|, \verb|PSYC143__CM|, \verb|HIST197__CM|, \verb|HIST120E_CM|, \verb|LIT_161__AF|, \verb|ECON165__CM|, \verb|ECON177__CM|, \verb|ECON185__CM|, \verb|ARBC130__CM|, \verb|FREN135__CM|, \verb|LIT_165__AF|, \verb|AMST140__SC|, \verb|AFRI124__AF|, \verb|CLAS023__SC|, \verb|FREN153__SC|, \verb|HIST117__SC|, \verb|ITAL022__SC|, \verb|ITAL187C_SC|, \verb|CHST182__CH|, \verb|CHST185C_CH|, \verb|MS__139__SC|, \verb|ENGL013__PZ|, \verb|ENGL141__PZ|, \verb|ENGL047__PZ|, \verb|MS__187__SC|, \verb|MS__053__SC|, \verb|MS__125__SC|, \verb|PHIL126__SC|, \verb|PHIL136__SC|, \verb|POLI106__SC|, \verb|POLI109__SC|, \verb|POLI142__SC|, \verb|PSYC175__SC|, \verb|AFRI195D_AF|, \verb|FREN002L_CM|, \verb|POST126__PZ|, \verb|POST162__PZ|, \verb|POST131__PZ|, \verb|HIST151__SC|, \verb|POST141__PZ|, \verb|POST144__PZ|, \verb|PSYC113__SC|, \verb|PSYC114__SC|, \verb|POST107__CH|, \verb|POST060__PZ|, \verb|POST199__PZ|, \verb|PSYC121__SC|, \verb|PSYC164__SC|, \verb|SPAN141__SC|, \verb|ANTH015__PO|, \verb|ANTH077__PO|, \verb|ARHI130__PO|, \verb|ART_024__PO|, \verb|ART_122__PO|, \verb|ART_123__PO|, \verb|ART_192__PO|, \verb|ASIA083__PO|, \verb|ASIA191__PO|, \verb|ASIA192__PO|, \verb|CHIN123__PO|, \verb|CHST125__CH|, \verb|CLAS191__PO|, \verb|CLAS192__PO|, \verb|ECON115__PO|, \verb|ENGL044__PO|, \verb|ENGL072__PO|, \verb|ENGL084__PO|, \verb|ENGL087H_PO|, \verb|ENGL169__PO|, \verb|ENGL170P_PO|, \verb|ENGL170X_PO|, \verb|ENGL195B_PO|, \verb|FREN107__PO|, \verb|FREN185__PO|, \verb|GERM154F_PO|, \verb|GWS_070__PO|, \verb|GWS_081__PO|, \verb|GWS_173__PO|, \verb|GWS_183__PO|, \verb|GWS_191__PO|, \verb|HIST047__PO|, \verb|HIST101ADPO|, \verb|HIST101H_PO|, \verb|HIST101HAPO|, \verb|HIST184__PO|, \verb|HIST191__PO|, \verb|HIST192__PO|, \verb|IR__191__PO|, \verb|JAPN125__PO|, \verb|LAST191__PO|, \verb|LGCS145__PO|, \verb|LGCS188__PO|, \verb|MUS_077__PO|, \verb|MUS_177__PO|, \verb|POLI162__PO|, \verb|PPA_191__PO|, \verb|PPE_195__PO|, \verb|PSYC163__PO|, \verb|RLST189C_PO|, \verb|RUSS181__PO|, \verb|SOC_189Z_PO|, \verb|SPAN104__PO|, \verb|SPAN189A_PO|, \verb|THEA083__PO|, \verb|THEA100C_PO|, \verb|THEA132__PO|, \verb|THEA171__PO|, \verb|ART_129__SC|, \verb|ID__076__JT|, \verb|ART_127__PZ|, \verb|EA__134__PZ|, \verb|EA__197__PZ|, \verb|HIST046__PZ|, \verb|HIST065__PZ|, \verb|HIST066__PZ|, \verb|HIST150__PZ|, \verb|HIST158__PZ|, \verb|HIST199__PZ|, \verb|MLLC021__PZ|, \verb|MLLC033__PZ|, \verb|MS__050__PZ|, \verb|MS__090__PZ|, \verb|ART_179V_HM|, \verb|ECON179F_HM|, \verb|GEOG179I_HM|, \verb|HIST081__HM|, \verb|LIT_179ABHM|, \verb|MUS_003__HM|, \verb|MUS_046__HM|, \verb|MUS_179E_HM|, \verb|PHIL179F_HM|, \verb|LIT_179AHHM|, \verb|CHLT060__CH|, \verb|ECON196__CM|, \verb|MUS_131__SC|, \verb|LIT_099C_CM|, \verb|ASAM179P_AA|, \verb|ASAM179Q_AA|, \verb|ITAL141__SC|, \verb|CGS_120__PZ|, \verb|ASAM088__PZ|, \verb|ASAM115__CM|, \verb|ASAM126__HM|, \verb|ASAM128__PZ|, \verb|ASAM135B_PZ|, \verb|ASAM175__PZ|, \verb|GFS_036__PZ|, \verb|PSYC051__CM|, \verb|PSYC126__PZ|, \verb|PSYC131__CM|, \verb|PSYC140A_PZ|, \verb|PSYC160__CM|, \verb|PSYC179U_HM|, \verb|PSYC190__PO|, \verb|FREN001L_CM|, \verb|FREN033__CM|, \verb|FREN033L_CM|, \verb|FREN044__CM|, \verb|FREN044L_CM|, \verb|FREN044L_SC|, \verb|FREN100__PZ|, \verb|FREN101B_PO|, \verb|FREN110__PO|, \verb|FREN127__SC|, \verb|FREN129__PO|, \verb|FREN179__SC|, \verb|FREN188__PO|, \verb|MUS_067__HM|, \verb|MUS_099__HM|, \verb|MUS_104__SC|, \verb|MUS_179G_HM|, \verb|MUS_179I_HM|, \verb|ORST050__PZ|, \verb|ORST148__PZ|, \verb|ORST198J_PZ|, \verb|GERM112__SC|, \verb|WRIT137__SC|, \verb|AFRI014__AF|, \verb|AFRI128__AF|, \verb|AFRI190__AF|, \verb|CLAS064__PO|, \verb|CLAS088__PZ|, \verb|CLAS112__PO|, \verb|THEA056C_PO|, \verb|THEA100S_PO|, \verb|THEA187__PO|, \verb|THEA191H_PO|, \verb|THEA192H_PO|, \verb|CHLT103__CH|, \verb|CHLT138__CH|, \verb|CHLT181__CH|, \verb|RUSS180__PO|, \verb|GREK022__SC|, \verb|IR__101__PO|, \verb|LGCS010__PZ|, \verb|LGCS106__PO|, \verb|LGCS166__PZ|, \verb|RLST064__CM|, \verb|RLST082__CM|, \verb|RLST086__SC|, \verb|RLST113__HM|, \verb|RLST128__PO|, \verb|RLST130__CM|, \verb|RLST152__PO|, \verb|RLST183__HM|, \verb|SOC_099__PZ|, \verb|SOC_138__PO|, \verb|SOC_179E_HM|, \verb|SOC_179F_HM|, \verb|MS__044__PO|, \verb|MS__049__PZ|, \verb|MS__049__SC|, \verb|MS__073__PO|, \verb|MS__082__SC|, \verb|MS__089B_PO|, \verb|MS__093__PZ|, \verb|MS__120__PO|, \verb|MS__123__JT|, \verb|MS__148A_PO|, \verb|MS__148C_PO|, \verb|MS__170__HM|, \verb|MS__179I_HM|, \verb|ANTH009__PZ|, \verb|ANTH070__PZ|, \verb|ANTH075__PZ|, \verb|ANTH080__PO|, \verb|ANTH087__SC|, \verb|ANTH126__SC|, \verb|ANTH143__SC|, \verb|HIST009__SC|, \verb|HIST059__CM|, \verb|HIST070B_SC|, \verb|HIST101N_PO|, \verb|HIST101S_CH|, \verb|HIST101U_PO|, \verb|HIST101Y_PO|, \verb|HIST106__CM|, \verb|HIST108__CM|, \verb|HIST109__CM|, \verb|HIST112__CM|, \verb|HIST113__PO|, \verb|HIST113__SC|, \verb|HIST120__CM|, \verb|HIST121__CM|, \verb|HIST122__CM|, \verb|HIST129__CM|, \verb|HIST133A_CM|, \verb|HIST136C_CM|, \verb|HIST147__CM|, \verb|HIST147__PO|, \verb|HIST155__SC|, \verb|HIST164__SC|, \verb|HIST165__PO|, \verb|HIST169__PO|, \verb|HIST173__CM|, \verb|HIST174__PO|, \verb|HIST189F_PO|, \verb|HIST189G_PO|, \verb|ECON111__SC|, \verb|ECON131__PO|, \verb|ECON152__PO|, \verb|ECON165__PO|, \verb|ECON168__CM|, \verb|ECON172__PZ|, \verb|ECON179G_HM|, \verb|ARHI144B_PO|, \verb|ARHI183__SC|, \verb|ARHI186B_PZ|, \verb|ARHI186L_PO|, \verb|FGSS192__SC|, \verb|GOVT107__CM|, \verb|GOVT116__CM|, \verb|GOVT121__CM|, \verb|GOVT129__CM|, \verb|GOVT136C_CM|, \verb|GOVT141__CM|, \verb|GOVT151__CM|, \verb|GOVT157S_CM|, \verb|GOVT191__CM|, \verb|GOVT192__CM|, \verb|EA__144__PZ|, \verb|EA__171__PO|, \verb|EA__189F_PO|, \verb|POST020__PZ|, \verb|POST114__HM|, \verb|POST179D_HM|, \verb|POST182__PZ|, \verb|ENGL015__PO|, \verb|ENGL016__PZ|, \verb|ENGL018__PO|, \verb|ENGL031__PZ|, \verb|ENGL036__PZ|, \verb|ENGL037__PZ|, \verb|ENGL046__PO|, \verb|ENGL055A_PO|, \verb|ENGL061__PZ|, \verb|ENGL080__PO|, \verb|ENGL101__PZ|, \verb|ENGL113S_SC|, \verb|ENGL125__SC|, \verb|ENGL148__SC|, \verb|ENGL151__PO|, \verb|ENGL151__PZ|, \verb|ENGL160S_SC|, \verb|ENGL161__PO|, \verb|ENGL162__PZ|, \verb|ENGL163__PO|, \verb|ENGL170J_PO|, \verb|ENGL170R_PO|, \verb|ENGL194__SC|, \verb|HMSC133__SC|, \verb|LIT_081__CM|, \verb|LIT_102__CM|, \verb|LIT_134__CM|, \verb|LIT_135__CM|, \verb|LIT_144__HM|, \verb|LIT_179Y_HM|, \verb|LIT_183__CM|, \verb|CHST074__CH|, \verb|CHST081__PO|, \verb|CHST184__CH|, \verb|AMST103__PZ|, \verb|GEOG145__HM|, \verb|POLI010__PO|, \verb|POLI061__PO|, \verb|POLI070__PO|, \verb|POLI120__SC|, \verb|POLI134__PO|, \verb|POLI184__SC|, \verb|CHIN150__PO|, \verb|ID__199S3PO|, \verb|ID__199S4PO|, \verb|GRMT115__SC|, \verb|GRMT131__PO|, \verb|ARCN144__SC|, \verb|STS_010__PZ|, \verb|STS_191__SC|, \verb|ARCT179C_HM|, \verb|ART_020__PO|, \verb|ART_020__PZ|, \verb|ART_025A_PO|, \verb|ART_133__PZ|, \verb|ART_151__SC|, \verb|ART_181__SC|, \verb|DANC102__SC|, \verb|DANC153__PO|, \verb|DANC153P_PO|, \verb|DANC170__PO|, \verb|DANC182__PO|, \verb|SPAN103__SC|, \verb|SPAN125A_PO|, \verb|SPAN143__SC|, \verb|SPAN148__PZ|, \verb|SPAN156__PZ|, \verb|SPAN159__PO|, \verb|SPAN192__PO|, \verb|GWS_162__PO|, \verb|GWS_179B_HM|, \verb|PHIL033__PZ|, \verb|PHIL039__CM|, \verb|PHIL058__PZ|, \verb|PHIL122__HM|, \verb|PHIL123__CM|, \verb|PHIL133__SC|, \verb|PHIL142__CM|, \verb|PHIL151__SC|, \verb|PHIL176__CM|, \verb|PHIL182__CM|, \verb|PHIL186K_PO|, \verb|LATN033__PO|, \verb|LAMS191__PO|, \verb|RUST079__PO|, \verb|RUST110__PO|, \verb|ITAL150__SC|, \verb|ITAL191__SC|, \verb|CGS_085__PZ|, \verb|CGS_101__PZ|, \verb|CGS_199__PZ|, \verb|CHNT178__PO|, \verb|ASAM086__HM|, \verb|ASAM094__PZ|, \verb|ASAM105B_PZ|, \verb|PSYC104__PZ|, \verb|PSYC104P_PZ|, \verb|PSYC131A_PO|, \verb|PSYC147__CH|, \verb|PSYC151__SC|, \verb|PSYC179M_HM|, \verb|PSYC180C_PO|, \verb|PSYC180R_PO|, \verb|FREN001__PZ|, \verb|FREN103__PO|, \verb|FREN113__SC|, \verb|FREN126__SC|, \verb|FREN180__PO|, \verb|FREN187A_SC|, \verb|FREN191__SC|, \verb|MUS_074__PO|, \verb|MUS_092__PO|, \verb|MUS_094__PO|, \verb|MUS_119__SC|, \verb|MUS_179H_HM|, \verb|MUS_191__SC|, \verb|GLAS180IOPZ|, \verb|GERM102__PO|, \verb|GERM191__SC|, \verb|WRIT022__PZ|, \verb|WRIT125__PZ|, \verb|WRIT146__SC|, \verb|WRIT175__SC|, \verb|WRIT190__PZ|, \verb|WRIT197Q_SC|, \verb|AFRI010ASAF|, \verb|AFRI010BSAF|, \verb|AFRI138__PO|, \verb|AFRI149__AF|, \verb|AFRI159__AF|, \verb|AFRI180__AF|, \verb|CLAS114__PO|, \verb|CLAS133__SC|, \verb|CLAS162__PZ|, \verb|CLAS175__PZ|, \verb|CLAS191__SC|, \verb|MENA191__SC|, \verb|THEA084__PO|, \verb|THEA115O_PO|, \verb|THEA142__PO|, \verb|THEA172__PO|, \verb|CHLT013__CH|, \verb|CHLT151__CH|, \verb|RUSS182__PO|, \verb|RUSS191__SC|, \verb|PORT044__CM|, \verb|LGCS104__PO|, \verb|LGCS110__PZ|, \verb|LGCS118__PO|, \verb|LGCS128__PO|, \verb|RLST020__PO|, \verb|RLST040__PO|, \verb|RLST049__PO|, \verb|RLST055__CM|, \verb|RLST060__SC|, \verb|RLST098__SC|, \verb|RLST130A_CM|, \verb|RLST139__PO|, \verb|RLST144__PO|, \verb|RLST160__PO|, \verb|RLST169__CM|, \verb|RLST179__SC|, \verb|RLST180__CM|, \verb|RLST187__PO|, \verb|RLST191__SC|, \verb|SOC_030__CH|, \verb|SOC_118__PZ|, \verb|SOC_123__PO|, \verb|SOC_130__PO|, \verb|SOC_179D_HM|, \verb|SOC_189E_PO|, \verb|SOC_199C_PZ|, \verb|ARBC141__CM|, \verb|MS__060__PZ|, \verb|MS__070__PZ|, \verb|MS__092__PO|, \verb|MS__102__PO|, \verb|MS__128__PZ|, \verb|MS__130__SC|, \verb|MS__135__PO|, \verb|MS__140__PO|, \verb|MS__173__HM|, \verb|MS__190S_JT|, \verb|MS__199__SC|, \verb|ANTH099__PZ|, \verb|ANTH100__PO|, \verb|ANTH106__PZ|, \verb|ANTH115__SC|, \verb|ANTH124__SC|, \verb|ANTH137__PZ|, \verb|ANTH185__SC|, \verb|ANTH191__SC|, \verb|HIST025__CH|, \verb|HIST025__PZ|, \verb|HIST037__PO|, \verb|HIST043__PO|, \verb|HIST044__PZ|, \verb|HIST061__PO|, \verb|HIST064__PO|, \verb|HIST072__PO|, \verb|HIST090__CM|, \verb|HIST099__CM|, \verb|HIST101B_PO|, \verb|HIST101T_CH|, \verb|HIST105__PO|, \verb|HIST108__SC|, \verb|HIST113__CM|, \verb|HIST114__CM|, \verb|HIST114C_CM|, \verb|HIST116__CM|, \verb|HIST123__PO|, \verb|HIST125__CM|, \verb|HIST128__CM|, \verb|HIST133B_CM|, \verb|HIST140C_CM|, \verb|HIST146__PO|, \verb|HIST146__SC|, \verb|HIST149A_SC|, \verb|HIST150D_CM|, \verb|HIST151__HM|, \verb|HIST154__SC|, \verb|HIST157__CM|, \verb|HIST159__SC|, \verb|HIST160__PO|, \verb|HIST170__PO|, \verb|HIST175__CM|, \verb|HIST178__CM|, \verb|HIST178__PZ|, \verb|HIST179L_HM|, \verb|HIST182__CM|, \verb|HIST192__SC|, \verb|ECON120__PO|, \verb|ECON129__PO|, \verb|ECON134E_CM|, \verb|ECON140__CM|, \verb|ECON166__PO|, \verb|ECON191H_SC|, \verb|ECON199__PZ|, \verb|ARHI124__PO|, \verb|ARHI127__PO|, \verb|ARHI135__PZ|, \verb|ARHI139__PO|, \verb|ARHI141A_PO|, \verb|ARHI151__SC|, \verb|ARHI158__SC|, \verb|ARHI164__SC|, \verb|ARHI179E_HM|, \verb|ARHI186G_PO|, \verb|ARHI191__SC|, \verb|FGSS191__SC|, \verb|FGSS198__SC|, \verb|GOVT094__CM|, \verb|GOVT103__CM|, \verb|GOVT129E_CM|, \verb|GOVT142C_CM|, \verb|GOVT146__CM|, \verb|GOVT152C_CM|, \verb|GOVT154E_CM|, \verb|GOVT179C_CM|, \verb|GOVT182__CM|, \verb|EA__086__SC|, \verb|EA__106__PZ|, \verb|EA__130__PZ|, \verb|EDUC301G_CG|, \verb|EDUC424__CG|, \verb|POST066__PZ|, \verb|POST070__PZ|, \verb|POST128__PZ|, \verb|POST152__PZ|, \verb|POST153__PZ|, \verb|POST174__PZ|, \verb|POST179C_HM|, \verb|ENGL010B_PZ|, \verb|ENGL043__PO|, \verb|ENGL045__PO|, \verb|ENGL048__PO|, \verb|ENGL049__PZ|, \verb|ENGL095__PO|, \verb|ENGL111__SC|, \verb|ENGL120__PZ|, \verb|ENGL146__PO|, \verb|ENGL149S_SC|, \verb|ENGL153__PO|, \verb|ENGL157__PO|, \verb|ENGL164S_SC|, \verb|ENGL170A_PO|, \verb|ENGL170H_PO|, \verb|ENGL177__SC|, \verb|ENGL192__SC|, \verb|HMSC192__SC|, \verb|LIT_067__CM|, \verb|LIT_110__HM|, \verb|LIT_119__CM|, \verb|LIT_125__CM|, \verb|LIT_132__CM|, \verb|LIT_163__AF|, \verb|LIT_185__CM|, \verb|LIT_199__CM|, \verb|CHST067__CH|, \verb|CHST102__PO|, \verb|CHST124__PO|, \verb|CHST191__CH|, \verb|CHST192__CH|, \verb|AMST103__PO|, \verb|AMST180__PZ|, \verb|AMST191__SC|, \verb|GEOG179H_HM|, \verb|JAPN127__PO|, \verb|POLI031__PO|, \verb|POLI060__PO|, \verb|POLI115__PO|, \verb|POLI115__SC|, \verb|POLI121__SC|, \verb|POLI124__PO|, \verb|POLI126__SC|, \verb|POLI130__SC|, \verb|POLI135__SC|, \verb|POLI135L_SC|, \verb|POLI140__SC|, \verb|POLI155__SC|, \verb|POLI165__PO|, \verb|POLI167__JT|, \verb|POLI189F_PO|, \verb|POLI189G_PO|, \verb|CHIN124__PO|, \verb|ID__199P2PO|, \verb|ID__199P4PO|, \verb|PPA_191__SC|, \verb|GRMT014__PO|, \verb|GRMT114__SC|, \verb|GRMT164__PO|, \verb|ARCN101__SC|, \verb|ARCN191__SC|, \verb|STS_010__PO|, \verb|ARCT179B_HM|, \verb|ART_016__PZ|, \verb|ART_017X_PZ|, \verb|ART_022__PO|, \verb|ART_026__PO|, \verb|ART_027__PO|, \verb|ART_075__PZ|, \verb|ART_102__SC|, \verb|ART_105A_PO|, \verb|ART_112__PZ|, \verb|ART_119__PZ|, \verb|ART_125__PZ|, \verb|ART_132__SC|, \verb|ART_149__SC|, \verb|ART_150__SC|, \verb|ART_181G_SC|, \verb|DANC081A_SC|, \verb|DANC081B_SC|, \verb|DANC108A_SC|, \verb|DANC108B_SC|, \verb|DANC117B_SC|, \verb|DANC142__PO|, \verb|DANC142P_PO|, \verb|DANC166__PO|, \verb|DANC166P_PO|, \verb|DANC180__SC|, \verb|SPAN107__PZ|, \verb|SPAN125B_CM|, \verb|SPAN132__PZ|, \verb|SPAN134__SC|, \verb|SPAN145__SC|, \verb|SPAN146__PO|, \verb|SPAN150__PZ|, \verb|SPAN157__PZ|, \verb|SPAN167__CM|, \verb|SPAN171__PZ|, \verb|SPAN172__PZ|, \verb|GWS_036__PO|, \verb|GWS_140__PO|, \verb|GWS_179A_HM|, \verb|LAS_180__PO|, \verb|LAST191__SC|, \verb|PHIL043__PO|, \verb|PHIL044__PO|, \verb|PHIL046__PO|, \verb|PHIL052__PZ|, \verb|PHIL055__PZ|, \verb|PHIL070__PO|, \verb|PHIL075__PZ|, \verb|PHIL101D_CM|, \verb|PHIL104__PZ|, \verb|PHIL106__CM|, \verb|PHIL106__PO|, \verb|PHIL121__HM|, \verb|PHIL126__CM|, \verb|PHIL156__PZ|, \verb|PHIL177__CM|, \verb|PHIL179G_HM|, \verb|PHIL181__CM|, \verb|PHIL185C_PZ|, \verb|PHIL185L_PO|, \verb|PHIL185N_PZ|, \verb|PHIL185P_PO|, \verb|PHIL191__SC|, \verb|ASAM190A_PO2|, \verb|PSYC190A_PO2|, \verb|PSYC190R_PO2|, \verb|PSYC191__SC2|, \verb|PSYC197A_CM2|, \verb|PSYC197B_CM2|, \verb|PSYC198__CM2|, \verb|FREN191__PO2|, \verb|FREN192__PO2|, \verb|FREN193__PO2|, \verb|MUS_190__PO2|, \verb|GLAS194A_PZ2|, \verb|GLAS194B_PZ2|, \verb|ORST198M_PZ2|, \verb|GERM191__PO2|, \verb|GERM193__PO2|, \verb|AFRI191__AF2|, \verb|CLAS190__PO2|, \verb|THEA190__PO2|, \verb|RUSS191__PO2|, \verb|IR__190__PO2|, \verb|FLAN191__SC2|, \verb|LGCS190__PO2|, \verb|RLST190__PO2|, \verb|SOC_190__PO2|, \verb|SOC_199A_PZ2|, \verb|MS__190F_JT2|, \verb|ANTH190__PO2|, \verb|ANTH190__SC2|, \verb|HIST190__PO2|, \verb|HIST191__SC2|, \verb|HIST196__CM2|, \verb|HIST197__PZ2|, \verb|ECON191__SC2|, \verb|ECON198__PZ2|, \verb|ARHI190__PO2|, \verb|HUM_196__PO2|, \verb|EA__191__PO2|, \verb|POST195__PZ2|, \verb|ENGL190__SC2|, \verb|ENGL191__PO2|, \verb|ENGL191__SC2|, \verb|ENGL194__PZ2|, \verb|ENGL195__PO2|, \verb|ENGL199T_SC2|, \verb|HMSC191__SC2|, \verb|CHST190__CH2|, \verb|AMST190__PO2|, \verb|POLI190D_PO2|, \verb|POLI191__PO2|, \verb|POLI191__SC2|, \verb|POLI193__PO2|, \verb|POLI195__PO2|, \verb|CHIN192A_PO2|, \verb|ID__199P1PO2|, \verb|ID__199P3PO2|, \verb|ID__199S1PO2|, \verb|ID__199S2PO2|, \verb|PPA_190__PO2|, \verb|PPA_195__PO2|, \verb|RLIT191__PO2|, \verb|STS_190__PO2|, \verb|ART_190__PO2|, \verb|ART_192__SC2|, \verb|DANC190__SC2|, \verb|DANC191__SC2|, \verb|DANC192__PO2|, \verb|DANC193__SC2|, \verb|ASIA190__PO2|, \verb|SPAN191__PO2|, \verb|SPAN191__SC2|, \verb|GWS_190__PO2|, \verb|LAST190__PO2|, \verb|PHIL190__CM2|, \verb|PHIL190__PO2|, \verb|PHIL190__SC2|, \verb|PHIL191__PZ2|, \verb|PHIL198__CM2|, \verb|MS__190__JT2|, \verb|CGS_190__PZ2|, \verb|AMST191__PO2|, \verb|STS_191__PO2|, \verb|SOC_199B_PZ2|, \verb|ART_193__SC2|, \verb|ENGL193__SC2|, \verb|HUM_195J_SC2|, \verb|WRIT191__SC2|, \verb|ECON191__CM2|, \verb|ASAM191__PO2|, \verb|PHIL191__PO2|, \verb|AFRI191C_AF2|, \verb|HIST190__CM2|, \verb|ARHI191__PO2|, \verb|ECON190__PO2|, \verb|POLI190C_PO2|, \verb|RLST191__PO2|, \verb|RUST191__PO2|, \verb|SOC_191__PO2|, \verb|XGOV192__CM2|, \verb|PHIL192__CM2|, \verb|SPAN199__PZ2|, \verb|ART_199__PZ2|, \verb|HIST197__CM2|, \verb|AFRI195D_AF2|, \verb|POST199__PZ2|, \verb|ART_192__PO2|, \verb|ASIA191__PO2|, \verb|ASIA192__PO2|, \verb|CLAS191__PO2|, \verb|CLAS192__PO2|, \verb|ENGL195B_PO2|, \verb|GWS_191__PO2|, \verb|HIST191__PO2|, \verb|HIST192__PO2|, \verb|IR__191__PO2|, \verb|LAST191__PO2|, \verb|PPA_191__PO2|, \verb|PPE_195__PO2|, \verb|EA__197__PZ2|, \verb|HIST199__PZ2|, \verb|ECON196__CM2|, \verb|PSYC190__PO2|, \verb|ORST198J_PZ2|, \verb|AFRI190__AF2|, \verb|THEA191H_PO2|, \verb|THEA192H_PO2|, \verb|FGSS192__SC2|, \verb|GOVT191__CM2|, \verb|GOVT192__CM2|, \verb|ENGL194__SC2|, \verb|ID__199S3PO2|, \verb|ID__199S4PO2|, \verb|STS_191__SC2|, \verb|SPAN192__PO2|, \verb|LAMS191__PO2|, \verb|ITAL191__SC2|, \verb|CGS_199__PZ2|, \verb|FREN191__SC2|, \verb|MUS_191__SC2|, \verb|GERM191__SC2|, \verb|WRIT190__PZ2|, \verb|WRIT197Q_SC2|, \verb|CLAS191__SC2|, \verb|MENA191__SC2|, \verb|RUSS191__SC2|, \verb|RLST191__SC2|, \verb|SOC_199C_PZ2|, \verb|MS__190S_JT2|, \verb|MS__199__SC2|, \verb|ANTH191__SC2|, \verb|HIST192__SC2|, \verb|ECON191H_SC2|, \verb|ECON199__PZ2|, \verb|ARHI191__SC2|, \verb|FGSS191__SC2|, \verb|FGSS198__SC2|, \verb|ENGL192__SC2|, \verb|HMSC192__SC2|, \verb|LIT_199__CM2|, \verb|CHST191__CH2|, \verb|CHST192__CH2|, \verb|AMST191__SC2|, \verb|ID__199P2PO2|, \verb|ID__199P4PO2|, \verb|PPA_191__SC2|, \verb|ARCN191__SC2|, \verb|LAST191__SC2|, \verb|PHIL191__SC2|, \verb|HSA_000__5C|, \verb|HSA_001__5C|, \verb|HSA_002__5C|, \verb|HSA_003__5C|, \verb|HSA_004__5C|, \verb|HSA_005__5C|, \verb|HSA_006__5C|, \verb|HSA_007__5C|, \verb|HSA_008__5C|, \verb|HSA_009__5C|, \verb|HSA_000__HM|, \verb|HSA_001__HM|, \verb|HSA_002__HM|, \verb|HSA_003__HM|, \verb|HSA_100__HM|\\\hline
    \verb|HSAs_WritingIntensive| & \verb|PSYC139__HM|, \verb|RLST114__HM|, \verb|LIT_035__HM|, \verb|LIT_107__HM|, \verb|LIT_156__HM|, \verb|ART_179W_HM|, \verb|LIT_179AIHM|, \verb|PSYC134__HM|, \verb|LIT_141__HM|, \verb|LIT_179AHHM|, \verb|RLST113__HM|, \verb|RLST183__HM|, \verb|PHIL122__HM|, \verb|HIST151__HM|, \verb|GEOG179H_HM|, \verb|HSA_100__HM|\\\hline
    \verb|MATH055__HM| & \verb|MATH073__HM|\\\hline
    \verb|MATH062__HM_co| & \verb|MATH073__HM|\\\hline
    \verb|MATH181__HM_co| & \verb|MATH131__HM|\\\hline
    \verb|Maj_Bio| & \verb|BIOL054__HM|, \verb|BIOL154__HM|, \verb|BIOL101__HM|, \verb|BIOL108__HM|, \verb|BIOL109__HM|, \verb|BIOL113__HM|, \verb|CHEM056__HM|, \verb|CHEM058__HM|, \verb|CHEM105__HM|\\\hline
    \verb|Maj_BioChem| & \verb|CHEM051__HM|, \verb|CHEM056__HM|, \verb|CHEM105__HM|, \verb|BIOL054__HM|, \verb|BIOL154__HM|, \verb|BIOL111__HM|, \verb|BIOL113__HM|, \verb|CHEM058__HM|\\\hline
    \verb|Maj_BioChem_BioChem| & \verb|BIOL182__HM|, \verb|CHEM182__HM|\\\hline
    \verb|Maj_BioChem_BioChemLab| & \verb|BIOL184__HM|, \verb|CHEM184__HM|\\\hline
    \verb|Maj_BioChem_BioColloqium| & \verb|BIOL191__HM|, \verb|BIOL191__HM2|\\\hline
    \verb|Maj_BioChem_BioElective| & \verb|BIOL197__HM|, \verb|BIOL107__PO|, \verb|BIOL109__HM|, \verb|BIOL111__HM|, \verb|BIOL112__KS|, \verb|BIOL113__HM|, \verb|BIOL119__KS|, \verb|BIOL120__KS|, \verb|BIOL121__HM|, \verb|BIOL121__PO|, \verb|BIOL131L_KS|, \verb|BIOL140__PO|, \verb|BIOL141L_KS|, \verb|BIOL143__KS|, \verb|BIOL145__KS|, \verb|BIOL146L_KS|, \verb|BIOL147__KS|, \verb|BIOL154__KS|, \verb|BIOL157L_KS|, \verb|BIOL160__PO|, \verb|BIOL161__HM|, \verb|BIOL163__PO|, \verb|BIOL170L_KS|, \verb|BIOL175__KS|, \verb|BIOL177__KS|, \verb|BIOL183__KS|, \verb|BIOL184L_KS|, \verb|BIOL187JLKS|, \verb|BIOL188__HM|, \verb|BIOL188L_KS|, \verb|BIOL189__HM|, \verb|BIOL189L_KS|, \verb|BIOL111__KS|, \verb|BIOL131__KS|, \verb|BIOL135L_KS|, \verb|BIOL138__KS|, \verb|BIOL138L_KS|, \verb|BIOL156L_KS|, \verb|BIOL164__KS|, \verb|BIOL168L_KS|, \verb|BIOL176__KS|, \verb|BIOL187M_KS|, \verb|BIOL142L_KS|, \verb|BIOL101__HM|, \verb|BIOL108__HM|, \verb|BIOL154__HM|, \verb|BIOL182__HM|, \verb|BIOL184__HM|, \verb|MCBI118A_HM|, \verb|MCBI118B_HM|, \verb|BIOL169__PO|, \verb|BIOL169L_PO|, \verb|BIOL170__PO|, \verb|BIOL173__PO|, \verb|BIOL174__PO|, \verb|BIOL112__PO|, \verb|BIOL152__PO|, \verb|BIOL110__HM|, \verb|BIOL160__HM|, \verb|BIOL119__HM|, \verb|BIOL163L_PO|, \verb|BIOL182L_KS|, \verb|BIOL185G_HM|, \verb|BIOL187P_KS|, \verb|MCBI117__HM|, \verb|BIOL103__HM|, \verb|BIOL104A_PO|, \verb|BIOL108__PO|, \verb|BIOL113L_KS|, \verb|BIOL116__PO|, \verb|BIOL122__HM|, \verb|BIOL133__PO|, \verb|BIOL151L_KS|, \verb|BIOL162__PO|, \verb|BIOL173L_KS|\\\hline
    \verb|Maj_BioChem_Capstone_BioResearch| & \verb|BIOL195__HM|, \verb|BIOL195__HM2|\\\hline
    \verb|Maj_BioChem_Capstone_BioThesis| & \verb|BIOL193__HM|, \verb|BIOL193__HM2|\\\hline
    \verb|Maj_BioChem_Capstone_ChemThesis| & \verb|CHEM151__HM|, \verb|CHEM152__HM|\\\hline
    \verb|Maj_BioChem_Capstone_EngrClinic| & \verb|ENGR112__HM|, \verb|ENGR113__HM|\\\hline
    \verb|Maj_BioChem_ChemColloqium| & \verb|CHEM199__HM|, \verb|CHEM199__HM2|\\\hline
    \verb|Maj_BioChem_ChemLab| & \verb|CHEM053__HM|, \verb|CHEM109__HM|, \verb|CHEM110__HM|, \verb|CHEM111__HM|\\\hline
    \verb|Maj_BioChem_ChemWithoutCells| & \verb|CHEM103__HM|, \verb|CHEM104__HM|\\\hline
    \verb|Maj_BioChem_OrganicBio| & \verb|BIOL101__HM|, \verb|BIOL108__HM|, \verb|BIOL109__HM|\\\hline
    \verb|Maj_BioChem_Topics| & \verb|BIOL189__HM|\\\hline
    \verb|Maj_BioClim| & \verb|BIOL054__HM|, \verb|BIOL154__HM|, \verb|BIOL108__HM|, \verb|BIOL110__HM|, \verb|CLES101__HM|, \verb|CLES131__HM|\\\hline
    \verb|Maj_BioClim_BioClimateCapstone_CSClinic| & \verb|CSCI183__HM|, \verb|CSCI184__HM|\\\hline
    \verb|Maj_BioClim_BioClimateCapstone_EngClinic| & \verb|ENGR112__HM|, \verb|ENGR113__HM|\\\hline
    \verb|Maj_BioClim_BioClimateCapstone_MathClinic| & \verb|MATH193__HM|, \verb|MATH193__HM2|\\\hline
    \verb|Maj_BioClim_BioClimateCapstone_PhysClinic| & \verb|PHYS193__HM|, \verb|PHYS194__HM|\\\hline
    \verb|Maj_BioClim_BioClimateCapstone_Thesis| & \verb|BIOL193__HM|, \verb|BIOL193__HM2|\\\hline
    \verb|Maj_BioClim_BioClimateElective| & \verb|MATH082__HM|, \verb|MATH102__PZ|, \verb|MATH102__PO|, \verb|MATH102__SC|, \verb|PHYS051__HM|\\\hline
    \verb|Maj_BioClim_BioColloqium| & \verb|BIOL191__HM|, \verb|BIOL191__HM2|\\\hline
    \verb|Maj_BioClim_BioComputational| & \verb|CSCI144__HM|, \verb|CHEM080__HM|, \verb|MATH164__HM|, \verb|PHYS064__HM|\\\hline
    \verb|Maj_BioClim_BioElec| & \verb|BIOL197__HM|, \verb|BIOL107__PO|, \verb|BIOL109__HM|, \verb|BIOL111__HM|, \verb|BIOL112__KS|, \verb|BIOL113__HM|, \verb|BIOL119__KS|, \verb|BIOL120__KS|, \verb|BIOL121__HM|, \verb|BIOL121__PO|, \verb|BIOL131L_KS|, \verb|BIOL140__PO|, \verb|BIOL141L_KS|, \verb|BIOL143__KS|, \verb|BIOL145__KS|, \verb|BIOL146L_KS|, \verb|BIOL147__KS|, \verb|BIOL154__KS|, \verb|BIOL157L_KS|, \verb|BIOL160__PO|, \verb|BIOL161__HM|, \verb|BIOL163__PO|, \verb|BIOL170L_KS|, \verb|BIOL175__KS|, \verb|BIOL177__KS|, \verb|BIOL183__KS|, \verb|BIOL184L_KS|, \verb|BIOL187JLKS|, \verb|BIOL188__HM|, \verb|BIOL188L_KS|, \verb|BIOL189__HM|, \verb|BIOL189L_KS|, \verb|BIOL111__KS|, \verb|BIOL131__KS|, \verb|BIOL135L_KS|, \verb|BIOL138__KS|, \verb|BIOL138L_KS|, \verb|BIOL156L_KS|, \verb|BIOL164__KS|, \verb|BIOL168L_KS|, \verb|BIOL176__KS|, \verb|BIOL187M_KS|, \verb|BIOL142L_KS|, \verb|BIOL101__HM|, \verb|BIOL108__HM|, \verb|BIOL154__HM|, \verb|BIOL182__HM|, \verb|BIOL184__HM|, \verb|MCBI118A_HM|, \verb|MCBI118B_HM|, \verb|BIOL169__PO|, \verb|BIOL169L_PO|, \verb|BIOL170__PO|, \verb|BIOL173__PO|, \verb|BIOL174__PO|, \verb|BIOL112__PO|, \verb|BIOL152__PO|, \verb|BIOL110__HM|, \verb|BIOL160__HM|, \verb|BIOL119__HM|, \verb|BIOL163L_PO|, \verb|BIOL182L_KS|, \verb|BIOL185G_HM|, \verb|BIOL187P_KS|, \verb|MCBI117__HM|, \verb|BIOL103__HM|, \verb|BIOL104A_PO|, \verb|BIOL108__PO|, \verb|BIOL113L_KS|, \verb|BIOL116__PO|, \verb|BIOL122__HM|, \verb|BIOL133__PO|, \verb|BIOL151L_KS|, \verb|BIOL162__PO|, \verb|BIOL173L_KS|\\\hline
    \verb|Maj_BioClim_BioGroup| & \verb|BIOL101__HM|, \verb|BIOL109__HM|, \verb|BIOL113__HM|\\\hline
    \verb|Maj_BioClim_BioSeminar| & \verb|BIOL121__HM|, \verb|BIOL189__HM|\\\hline
    \verb|Maj_BioClim_ClimateColloqium| & \verb|CLES199__HM|, \verb|CLES199__HM2|\\\hline
    \verb|Maj_BioClim_ClimateContexts| & \verb|CLES130__HM|, \verb|CLES131__HM|, \verb|CHEM170__HM|, \verb|ENGR138__HM|\\\hline
    \verb|Maj_BioClim_ClimateContexts_Justice| & \verb|CLES121__HM|, \verb|CLES122__HM|\\\hline
    \verb|Maj_BioClim_ClimateInterventions_ClimateInterventionCredits| & \verb|CLES122__HM|, \verb|CLES130__HM|\\\hline
    \verb|Maj_BioClim_ClimateInterventions_Justice| & \verb|CLES121__HM|, \verb|CLES122__HM|\\\hline
    \verb|Maj_BioClim_ClimateStats| & \verb|MATH056__HM|, \verb|MATH062__HM|, \verb|MATH052__CM|, \verb|MATH052__PZ|, \verb|MATH058__PO|, \verb|MATH151__PO|, \verb|MATH152__PO|, \verb|PHYS117__HM|, \verb|BIOL154__HM|\\\hline
    \verb|Maj_BioClim_Thermodynamics| & \verb|CHEM051__HM|, \verb|ENGR082__HM|, \verb|PHYS117__HM|\\\hline
    \verb|Maj_Bio_BioColloquium| & \verb|BIOL191__HM|, \verb|BIOL191__HM2|\\\hline
    \verb|Maj_Bio_BioElecs| & \verb|BIOL197__HM|, \verb|BIOL107__PO|, \verb|BIOL109__HM|, \verb|BIOL111__HM|, \verb|BIOL112__KS|, \verb|BIOL113__HM|, \verb|BIOL119__KS|, \verb|BIOL120__KS|, \verb|BIOL121__HM|, \verb|BIOL121__PO|, \verb|BIOL131L_KS|, \verb|BIOL140__PO|, \verb|BIOL141L_KS|, \verb|BIOL143__KS|, \verb|BIOL145__KS|, \verb|BIOL146L_KS|, \verb|BIOL147__KS|, \verb|BIOL154__KS|, \verb|BIOL157L_KS|, \verb|BIOL160__PO|, \verb|BIOL161__HM|, \verb|BIOL163__PO|, \verb|BIOL170L_KS|, \verb|BIOL175__KS|, \verb|BIOL177__KS|, \verb|BIOL183__KS|, \verb|BIOL184L_KS|, \verb|BIOL187JLKS|, \verb|BIOL188__HM|, \verb|BIOL188L_KS|, \verb|BIOL189__HM|, \verb|BIOL189L_KS|, \verb|BIOL111__KS|, \verb|BIOL131__KS|, \verb|BIOL135L_KS|, \verb|BIOL138__KS|, \verb|BIOL138L_KS|, \verb|BIOL156L_KS|, \verb|BIOL164__KS|, \verb|BIOL168L_KS|, \verb|BIOL176__KS|, \verb|BIOL187M_KS|, \verb|BIOL142L_KS|, \verb|BIOL101__HM|, \verb|BIOL108__HM|, \verb|BIOL154__HM|, \verb|BIOL182__HM|, \verb|BIOL184__HM|, \verb|MCBI118A_HM|, \verb|MCBI118B_HM|, \verb|BIOL169__PO|, \verb|BIOL169L_PO|, \verb|BIOL170__PO|, \verb|BIOL173__PO|, \verb|BIOL174__PO|, \verb|BIOL112__PO|, \verb|BIOL152__PO|, \verb|BIOL110__HM|, \verb|BIOL160__HM|, \verb|BIOL119__HM|, \verb|BIOL163L_PO|, \verb|BIOL182L_KS|, \verb|BIOL185G_HM|, \verb|BIOL187P_KS|, \verb|MCBI117__HM|, \verb|BIOL103__HM|, \verb|BIOL104A_PO|, \verb|BIOL108__PO|, \verb|BIOL113L_KS|, \verb|BIOL116__PO|, \verb|BIOL122__HM|, \verb|BIOL133__PO|, \verb|BIOL151L_KS|, \verb|BIOL162__PO|, \verb|BIOL173L_KS|\\\hline
    \verb|Maj_Bio_BioLab_NoBioChem| & \verb|BIOL103__HM|, \verb|BIOL110__HM|, \verb|BIOL111__HM|\\\hline
    \verb|Maj_Bio_BioLab_WithBioChem| & \verb|BIOL184__HM|\\\hline
    \verb|Maj_Bio_BioLab_WithBioChem_Other| & \verb|BIOL103__HM|, \verb|BIOL110__HM|, \verb|BIOL111__HM|\\\hline
    \verb|Maj_Bio_BioSeminar| & \verb|BIOL121__HM|, \verb|BIOL189__HM|\\\hline
    \verb|Maj_Bio_ThesisOrClinic_BioThesis| & \verb|BIOL193__HM|, \verb|BIOL193__HM2|\\\hline
    \verb|Maj_Bio_ThesisOrClinic_CSClinic| & \verb|CSCI183__HM|, \verb|CSCI184__HM|\\\hline
    \verb|Maj_Bio_ThesisOrClinic_EngrClinic| & \verb|ENGR111__HM|, \verb|ENGR112__HM|, \verb|ENGR113__HM|\\\hline
    \verb|Maj_Bio_ThesisOrClinic_IntensiveBioThesis| & \verb|BIOL195__HM|, \verb|BIOL195__HM2|\\\hline
    \verb|Maj_Bio_ThesisOrClinic_MathClinic| & \verb|MATH193__HM|, \verb|MATH193__HM2|\\\hline
    \verb|Maj_Bio_ThesisOrClinic_PhysClinic| & \verb|PHYS193__HM|, \verb|PHYS194__HM|\\\hline
    \verb|Maj_CS| & \verb|MATH055__HM|, \verb|CSCI081__HM|, \verb|CSCI070__HM|, \verb|CSCI105__HM|, \verb|CSCI123__HM|, \verb|CSCI140__HM|, \verb|CSCI183__HM|, \verb|CSCI184__HM|\\\hline
    \verb|Maj_CSClimate| & \verb|CSCI070__HM|, \verb|MATH055__HM|, \verb|CSCI123__HM|, \verb|CLES101__HM|, \verb|CSCI183__HM|, \verb|CSCI184__HM|\\\hline
    \verb|Maj_CSClimate_Areas| & \verb|PHYS051__HM|, \verb|MATH082__HM|\\\hline
    \verb|Maj_CSClimate_Areas_Computational| & \verb|CHEM080__HM|, \verb|CSCI144__HM|, \verb|MATH164__HM|, \verb|PHYS064__HM|\\\hline
    \verb|Maj_CSClimate_Areas_ProbStats| & \verb|BIOL154__HM|, \verb|MATH056__HM|, \verb|MATH062__HM|, \verb|PHYS117__HM|\\\hline
    \verb|Maj_CSClimate_CSColloquium| & \verb|CSCI195__HM|, \verb|CSCI195__HM2|\\\hline
    \verb|Maj_CSClimate_CSElective| & \verb|MATH104__HM|, \verb|MATH106__HM|, \verb|MATH157__HM|, \verb|MATH164__HM|, \verb|MATH165__HM|, \verb|MATH187__HM|, \verb|ENGR085A_HM|, \verb|ENGR151__HM|, \verb|ENGR155__HM|, \verb|MCBI118A_HM|, \verb|MCBI118B_HM|, \verb|BIOL188__HM|, \verb|PHYS084__HM|, \verb|PHYS170__HM|, \verb|PSYC183__SC|, \verb|MCBI117__HM|, \verb|CSCI159__HM|, \verb|CSCI158__HM|, \verb|CSCI153__HM|, \verb|CSCI152__HM|, \verb|CSCI151__HM|, \verb|CSCI144__HM|, \verb|CSCI140__HM|, \verb|CSCI134__HM|, \verb|CSCI133__HM|, \verb|CSCI131__HM|, \verb|CSCI120__HM|, \verb|CSCI105__HM|\\\hline
    \verb|Maj_CSClimate_ClimateColloquium| & \verb|CLES199__HM|, \verb|CLES199__HM2|\\\hline
    \verb|Maj_CSClimate_ClimateContexts| & \verb|ECON146__HM|, \verb|ECON179G_HM|, \verb|GEOG179I_HM|, \verb|POST140__HM|, \verb|POST114__HM|, \verb|PSYC179O_HM|, \verb|RLST179G_HM|, \verb|SOC_179D_HM|\\\hline
    \verb|Maj_CSClimate_ClimateImpacts| & \verb|CLES130__HM|, \verb|CLES131__HM|, \verb|CHEM170__HM|, \verb|ENGR138__HM|\\\hline
    \verb|Maj_CSClimate_ClimateImpacts_ClimateJustice| & \verb|CLES121__HM|, \verb|CLES122__HM|\\\hline
    \verb|Maj_CSClimate_ClimateInterventions| & \verb|CLES120__HM|, \verb|CLES130__HM|\\\hline
    \verb|Maj_CSClimate_ClimateInterventions_ClimateJustice| & \verb|CLES121__HM|, \verb|CLES122__HM|\\\hline
    \verb|Maj_CSClimate_HumanCentered| & \verb|CSCI120__HM|, \verb|ENGR180__HM|, \verb|EA__185__PO|\\\hline
    \verb|Maj_CSClimate_Principles| & \verb|CSCI060__HM|\\\hline
    \verb|Maj_CSClimate_SysAlgs| & \verb|CSCI105__HM|, \verb|CSCI140__HM|\\\hline
    \verb|Maj_CSClimate_Thermodynamics| & \verb|CHEM051__HM|, \verb|ENGR082__HM|, \verb|PHYS117__HM|\\\hline
    \verb|Maj_CS_CSColloquium| & \verb|CSCI195__HM|, \verb|CSCI195__HM2|\\\hline
    \verb|Maj_CS_CSCore| & \verb|CSCI060__HM|\\\hline
    \verb|Maj_CS_CSElec| & \verb|MATH104__HM|, \verb|MATH106__HM|, \verb|MATH157__HM|, \verb|MATH164__HM|, \verb|MATH165__HM|, \verb|MATH187__HM|, \verb|ENGR085A_HM|, \verb|ENGR151__HM|, \verb|ENGR155__HM|, \verb|MCBI118A_HM|, \verb|MCBI118B_HM|, \verb|BIOL188__HM|, \verb|PHYS084__HM|, \verb|PHYS170__HM|, \verb|PSYC183__SC|, \verb|MCBI117__HM|, \verb|CSCI159__HM|, \verb|CSCI158__HM|, \verb|CSCI153__HM|, \verb|CSCI152__HM|, \verb|CSCI151__HM|, \verb|CSCI144__HM|, \verb|CSCI134__HM|, \verb|CSCI133__HM|, \verb|CSCI131__HM|, \verb|CSCI120__HM|\\\hline
    \verb|Maj_Chem| & \verb|CHEM051__HM|, \verb|CHEM056__HM|, \verb|CHEM103__HM|, \verb|CHEM104__HM|, \verb|CHEM182__HM|, \verb|PHYS051__HM|\\\hline
    \verb|Maj_ChemClimate| & \verb|CHEM051__HM|, \verb|CHEM056__HM|, \verb|CHEM103__HM|, \verb|CHEM109__HM|, \verb|PHYS051__HM|, \verb|MATH082__HM|, \verb|CLES101__HM|, \verb|CHEM170__HM|, \verb|CHEM151__HM|, \verb|CHEM152__HM|\\\hline
    \verb|Maj_ChemClimate_ChemClass| & \verb|CHEM048__HM|, \verb|CHEM080__HM|, \verb|CHEM105__HM|, \verb|CHEM111__HM|, \verb|CHEM112__HM|, \verb|CHEM114__HM|, \verb|CHEM116__HM|, \verb|CHEM122__HM|, \verb|CHEM170__HM|, \verb|CHEM171__HM|, \verb|CHEM194__HM|\\\hline
    \verb|Maj_ChemClimate_ChemClimateComputational| & \verb|CHEM080__HM|, \verb|PHYS064__HM|, \verb|CSCI144__HM|, \verb|MATH164__HM|\\\hline
    \verb|Maj_ChemClimate_ChemClimateProbStats| & \verb|MATH056__HM|, \verb|MATH062__HM|, \verb|PHYS117__HM|, \verb|BIOL154__HM|\\\hline
    \verb|Maj_ChemClimate_ChemSeminar| & \verb|CHEM199__HM|, \verb|CHEM199__HM2|\\\hline
    \verb|Maj_ChemClimate_ChemUpperDiv| & \verb|CHEM105__HM|, \verb|CHEM111__HM|, \verb|CHEM112__HM|, \verb|CHEM114__HM|, \verb|CHEM116__HM|, \verb|CHEM122__HM|, \verb|CHEM170__HM|, \verb|CHEM171__HM|, \verb|CHEM194__HM|\\\hline
    \verb|Maj_ChemClimate_ClimateColloqium| & \verb|CLES199__HM|, \verb|CLES199__HM2|\\\hline
    \verb|Maj_ChemClimate_ClimateContexts| & \verb|CLES130__HM|, \verb|CLES131__HM|, \verb|CHEM170__HM|, \verb|ENGR138__HM|\\\hline
    \verb|Maj_ChemClimate_ClimateContexts_Justice| & \verb|CLES121__HM|, \verb|CLES122__HM|\\\hline
    \verb|Maj_ChemClimate_ClimateInterventions_ClimateInterventionCredits| & \verb|CLES122__HM|, \verb|CLES130__HM|\\\hline
    \verb|Maj_ChemClimate_ClimateInterventions_Justice| & \verb|CLES121__HM|, \verb|CLES122__HM|\\\hline
    \verb|Maj_ChemClimate_Or1| & \verb|CHEM053__HM|, \verb|CHEM058__HM|, \verb|CHEM110__HM|, \verb|CHEM112__HM|, \verb|CHEM184__HM|\\\hline
    \verb|Maj_Chem_CHEMElec| & \verb|CHEM048__HM|, \verb|CHEM080__HM|, \verb|CHEM105__HM|, \verb|CHEM111__HM|, \verb|CHEM112__HM|, \verb|CHEM114__HM|, \verb|CHEM116__HM|, \verb|CHEM122__HM|, \verb|CHEM170__HM|, \verb|CHEM171__HM|, \verb|CHEM194__HM|, \verb|CLES101__HM|\\\hline
    \verb|Maj_Chem_ChemLab| & \verb|CHEM053__HM|, \verb|CHEM058__HM|, \verb|CHEM109__HM|, \verb|CHEM110__HM|, \verb|CHEM184__HM|\\\hline
    \verb|Maj_Chem_ChemSeminar| & \verb|CHEM199__HM|, \verb|CHEM199__HM2|\\\hline
    \verb|Maj_Chem_Math| & \verb|BIOL154__HM|, \verb|MATH056__HM|, \verb|MATH062__HM|, \verb|MATH082__HM|\\\hline
    \verb|Maj_Chem_ThesisOrClinic_Clinic| & \verb|ENGR112__HM|, \verb|ENGR113__HM|\\\hline
    \verb|Maj_Chem_ThesisOrClinic_Thesis| & \verb|CHEM151__HM|, \verb|CHEM152__HM|\\\hline
    \verb|Maj_CsPhys| & \verb|CSCI070__HM|, \verb|CSCI140__HM|, \verb|MATH055__HM|, \verb|MATH082__HM|, \verb|PHYS051__HM|, \verb|PHYS052__HM|, \verb|PHYS054__HM|, \verb|PHYS064__HM|, \verb|PHYS111__HM|, \verb|PHYS117__HM|\\\hline
    \verb|Maj_CsPhys_CSColloqium| & \verb|CSCI195__HM|\\\hline
    \verb|Maj_CsPhys_CSCoreElective| & \verb|CSCI081__HM|, \verb|CSCI105__HM|\\\hline
    \verb|Maj_CsPhys_CSElective| & \verb|CSCI081__HM|, \verb|CSCI159__HM|, \verb|CSCI158__HM|, \verb|CSCI153__HM|, \verb|CSCI152__HM|, \verb|CSCI151__HM|, \verb|CSCI144__HM|, \verb|CSCI134__HM|, \verb|CSCI133__HM|, \verb|CSCI131__HM|, \verb|CSCI123__HM|, \verb|CSCI120__HM|, \verb|CSCI105__HM|\\\hline
    \verb|Maj_CsPhys_Capstone_CSMTClinic| & \verb|CSMT183__HM|, \verb|CSMT184__HM|\\\hline
    \verb|Maj_CsPhys_Capstone_PhysClinic| & \verb|PHYS193__HM|, \verb|PHYS194__HM|\\\hline
    \verb|Maj_CsPhys_Capstone_PhysThesis| & \verb|PHYS199__HM|, \verb|PHYS199__HM2|\\\hline
    \verb|Maj_CsPhys_IntroCS| & \verb|CSCI060__HM|\\\hline
    \verb|Maj_CsPhys_Lab| & \verb|PHYS133__HM|, \verb|PHYS134__HM|\\\hline
    \verb|Maj_CsPhys_PhysCSElective| & \verb|CSCI081__HM|, \verb|CSCI159__HM|, \verb|CSCI158__HM|, \verb|CSCI153__HM|, \verb|CSCI152__HM|, \verb|CSCI151__HM|, \verb|CSCI144__HM|, \verb|CSCI134__HM|, \verb|CSCI133__HM|, \verb|CSCI131__HM|, \verb|CSCI123__HM|, \verb|CSCI120__HM|, \verb|CSCI105__HM|\\\hline
    \verb|Maj_CsPhys_PhysColloqium| & \verb|PHYS195__HM|, \verb|PHYS195__HM2|\\\hline
    \verb|Maj_CsPhys_Quantum| & \verb|PHYS084__HM|, \verb|PHYS116__HM|\\\hline
    \verb|Maj_Engr| & \verb|ENGR004__HM|, \verb|ENGR079__HM|, \verb|ENGR082__HM|, \verb|ENGR083__HM|, \verb|ENGR084__HM|, \verb|ENGR085__HM|, \verb|ENGR086__HM|, \verb|ENGR072__HM|, \verb|MATH056__HM|, \verb|ENGR101__HM|, \verb|ENGR102__HM|, \verb|ENGR080__HM|, \verb|ENGR111__HM|, \verb|ENGR112__HM|, \verb|ENGR113__HM|, \verb|ENGR122__HM|, \verb|ENGR124__HM|\\\hline
    \verb|Maj_Engr_EngrElec| & \verb|ENGR131__HM|, \verb|ENGR132__HM|, \verb|ENGR133__HM|, \verb|ENGR134__HM|, \verb|ENGR138__HM|, \verb|ENGR151__HM|, \verb|ENGR154__HM|, \verb|ENGR155__HM|, \verb|ENGR157__HM|, \verb|ENGR164__HM|, \verb|ENGR171__HM|, \verb|ENGR175__HM|, \verb|ENGR177__HM|, \verb|ENGR178__HM|, \verb|ENGR180__HM|, \verb|ENGR181__HM|, \verb|ENGR183__HM|, \verb|ENGR185A_HM|, \verb|ENGR185B_HM|, \verb|ENGR186__HM|, \verb|ENGR187__HM|, \verb|ENGR205__HM|, \verb|ENGR207__HM|, \verb|ENGR208__HM|, \verb|ENGR278__HM|\\\hline
    \verb|Maj_Math| & \verb|MATH055__HM|, \verb|MATH062__HM|, \verb|MATH082__HM|, \verb|MATH131__HM|, \verb|MATH171__HM|, \verb|MATH180__HM|, \verb|MATH198__HM|\\\hline
    \verb|Maj_MathCS| & \verb|MATH055__HM|, \verb|CSCI081__HM|, \verb|CSCI070__HM|, \verb|MATH062__HM|, \verb|MATH082__HM|, \verb|MATH131__HM|, \verb|MATH171__HM|, \verb|CSCI195__HM|, \verb|CSCI195__HM2|, \verb|MATH198__HM|, \verb|CSMT183__HM|, \verb|CSMT184__HM|\\\hline
    \verb|Maj_MathCS_Algs| & \verb|CSCI140__HM|\\\hline
    \verb|Maj_MathCS_CSElec| & \verb|ENGR085A_HM|, \verb|ENGR151__HM|, \verb|ENGR155__HM|, \verb|MCBI118B_HM|, \verb|BIOL188__HM|, \verb|PHYS084__HM|, \verb|PHYS170__HM|, \verb|PSYC183__SC|, \verb|CSCI159__HM|, \verb|CSCI158__HM|, \verb|CSCI153__HM|, \verb|CSCI152__HM|, \verb|CSCI151__HM|, \verb|CSCI144__HM|, \verb|CSCI134__HM|, \verb|CSCI133__HM|, \verb|CSCI131__HM|, \verb|CSCI123__HM|, \verb|CSCI120__HM|, \verb|CSCI105__HM|\\\hline
    \verb|Maj_MathCS_IntroCS| & \verb|CSCI060__HM|\\\hline
    \verb|Maj_MathCS_MathElec| & \verb|MATH104__HM|, \verb|MATH106__HM|, \verb|MATH119__HM|, \verb|MATH132__HM|, \verb|MATH136__HM|, \verb|MATH147__HM|, \verb|MATH156__HM|, \verb|MATH157__HM|, \verb|MATH158__HM|, \verb|MATH164__HM|, \verb|MATH165__HM|, \verb|MATH174__HM|, \verb|MATH175__HM|, \verb|MATH176__HM|, \verb|MATH187__HM|, \verb|MATH181__HM|, \verb|MATH179__HM|, \verb|MCBI117__HM|\\\hline
    \verb|Maj_MathCompBio| & \verb|BIOL054__HM|, \verb|BIOL154__HM|, \verb|MATH055__HM|, \verb|MATH082__HM|, \verb|MCBI118A_HM|, \verb|MCBI118B_HM|, \verb|MATH198__HM|, \verb|MCBI199__HM|\\\hline
    \verb|Maj_MathCompBio_AdvancedBio| & \verb|BIOL119__HM|, \verb|MATH119__HM|, \verb|BIOL188__HM|\\\hline
    \verb|Maj_MathCompBio_BioColloquium| & \verb|BIOL191__HM|, \verb|BIOL191__HM2|\\\hline
    \verb|Maj_MathCompBio_BioFoundations| & \verb|BIOL101__HM|, \verb|BIOL108__HM|, \verb|BIOL109__HM|, \verb|BIOL113__HM|\\\hline
    \verb|Maj_MathCompBio_BioLab| & \verb|BIOL103__HM|, \verb|BIOL110__HM|, \verb|BIOL111__HM|\\\hline
    \verb|Maj_MathCompBio_BioSeminar| & \verb|BIOL121__HM|, \verb|BIOL189__HM|\\\hline
    \verb|Maj_MathCompBio_CSElec| & \verb|CSCI159__HM|, \verb|CSCI158__HM|, \verb|CSCI153__HM|, \verb|CSCI152__HM|, \verb|CSCI151__HM|, \verb|CSCI144__HM|, \verb|CSCI140__HM|, \verb|CSCI134__HM|, \verb|CSCI133__HM|, \verb|CSCI131__HM|, \verb|CSCI123__HM|, \verb|CSCI120__HM|, \verb|CSCI105__HM|\\\hline
    \verb|Maj_MathCompBio_MathElec| & \verb|MATH104__HM|, \verb|MATH106__HM|, \verb|MATH119__HM|, \verb|MATH132__HM|, \verb|MATH136__HM|, \verb|MATH147__HM|, \verb|MATH156__HM|, \verb|MATH157__HM|, \verb|MATH158__HM|, \verb|MATH164__HM|, \verb|MATH165__HM|, \verb|MATH174__HM|, \verb|MATH175__HM|, \verb|MATH176__HM|, \verb|MATH187__HM|, \verb|MATH181__HM|, \verb|MATH179__HM|\\\hline
    \verb|Maj_MathCompBio_MathOrCSElec| & \verb|CSCI159__HM|, \verb|CSCI158__HM|, \verb|CSCI153__HM|, \verb|CSCI152__HM|, \verb|CSCI151__HM|, \verb|CSCI144__HM|, \verb|CSCI140__HM|, \verb|CSCI134__HM|, \verb|CSCI133__HM|, \verb|CSCI131__HM|, \verb|CSCI123__HM|, \verb|CSCI120__HM|, \verb|CSCI105__HM|, \verb|MATH104__HM|, \verb|MATH106__HM|, \verb|MATH119__HM|, \verb|MATH132__HM|, \verb|MATH136__HM|, \verb|MATH147__HM|, \verb|MATH156__HM|, \verb|MATH157__HM|, \verb|MATH158__HM|, \verb|MATH164__HM|, \verb|MATH165__HM|, \verb|MATH174__HM|, \verb|MATH175__HM|, \verb|MATH176__HM|, \verb|MATH187__HM|, \verb|MATH181__HM|, \verb|MATH179__HM|\\\hline
    \verb|Maj_MathCompBio_ThesisOrClinic_BioThesis| & \verb|BIOL193__HM|, \verb|BIOL193__HM2|\\\hline
    \verb|Maj_MathCompBio_ThesisOrClinic_CSClinic| & \verb|CSCI183__HM|, \verb|CSCI184__HM|\\\hline
    \verb|Maj_MathCompBio_ThesisOrClinic_IntensiveBioThesis| & \verb|BIOL195__HM|, \verb|BIOL195__HM2|\\\hline
    \verb|Maj_MathCompBio_ThesisOrClinic_MathClinic| & \verb|MATH193__HM|, \verb|MATH193__HM2|\\\hline
    \verb|Maj_MathCompBio_ThesisOrClinic_MathThesis| & \verb|MATH197__HM|, \verb|MATH197__HM2|\\\hline
    \verb|Maj_MathPhys| & \verb|MATH055__HM|, \verb|MATH082__HM|, \verb|MATH131__HM|, \verb|MATH157__HM|, \verb|MATH171__HM|, \verb|MATH180__HM|, \verb|MATH198__HM|, \verb|PHYS051__HM|, \verb|PHYS052__HM|, \verb|PHYS054__HM|, \verb|PHYS111__HM|, \verb|PHYS116__HM|, \verb|PHYS117__HM|, \verb|PHYS134__HM|\\\hline
    \verb|Maj_MathPhys_Fields| & \verb|PHYS151__HM|, \verb|PHYS154__HM|, \verb|PHYS156__HM|\\\hline
    \verb|Maj_MathPhys_PhysColloquium| & \verb|PHYS195__HM|, \verb|PHYS195__HM2|\\\hline
    \verb|Maj_MathPhys_ScientificComp| & \verb|MATH164__HM|, \verb|MATH165__HM|, \verb|PHYS170__HM|\\\hline
    \verb|Maj_MathPhys_ThesisOrClinic_MathClinic| & \verb|MATH193__HM|, \verb|MATH193__HM2|\\\hline
    \verb|Maj_MathPhys_ThesisOrClinic_MathThesis| & \verb|MATH197__HM|, \verb|MATH197__HM2|\\\hline
    \verb|Maj_MathPhys_ThesisOrClinic_PhysClinic| & \verb|PHYS193__HM|, \verb|PHYS194__HM|\\\hline
    \verb|Maj_MathPhys_ThesisOrClinic_PhysThesis| & \verb|PHYS199__HM|, \verb|PHYS199__HM2|\\\hline
    \verb|Maj_Math_Computational| & \verb|CSCI081__HM|, \verb|MATH164__HM|, \verb|MATH165__HM|, \verb|CSCI140__HM|\\\hline
    \verb|Maj_Math_MathElec| & \verb|MCBI118B_HM|, \verb|MCBI118A_HM|, \verb|MCBI117__HM|, \verb|MATH104__HM|, \verb|MATH106__HM|, \verb|MATH119__HM|, \verb|MATH132__HM|, \verb|MATH136__HM|, \verb|MATH147__HM|, \verb|MATH156__HM|, \verb|MATH157__HM|, \verb|MATH158__HM|, \verb|MATH164__HM|, \verb|MATH165__HM|, \verb|MATH174__HM|, \verb|MATH175__HM|, \verb|MATH176__HM|, \verb|MATH187__HM|, \verb|MATH181__HM|, \verb|MATH179__HM|\\\hline
    \verb|Maj_Math_ThesisOrClinic_BioThesis| & \verb|MATH197__HM|, \verb|MATH197__HM2|\\\hline
    \verb|Maj_Math_ThesisOrClinic_MathThesis| & \verb|MATH193__HM|, \verb|MATH193__HM2|\\\hline
    \verb|Maj_MolBio| & \verb|BIOL054__HM|, \verb|BIOL154__HM|, \verb|BIOL101__HM|, \verb|BIOL108__HM|, \verb|BIOL109__HM|, \verb|BIOL111__HM|, \verb|BIOL113__HM|, \verb|BIOL182__HM|, \verb|CHEM056__HM|, \verb|CHEM058__HM|, \verb|CHEM105__HM|, \verb|CHEM111__HM|\\\hline
    \verb|Maj_MolBio_BioElecs| & \verb|BIOL197__HM|, \verb|BIOL107__PO|, \verb|BIOL109__HM|, \verb|BIOL111__HM|, \verb|BIOL112__KS|, \verb|BIOL113__HM|, \verb|BIOL119__KS|, \verb|BIOL120__KS|, \verb|BIOL121__HM|, \verb|BIOL121__PO|, \verb|BIOL131L_KS|, \verb|BIOL140__PO|, \verb|BIOL141L_KS|, \verb|BIOL143__KS|, \verb|BIOL145__KS|, \verb|BIOL146L_KS|, \verb|BIOL147__KS|, \verb|BIOL154__KS|, \verb|BIOL157L_KS|, \verb|BIOL160__PO|, \verb|BIOL161__HM|, \verb|BIOL163__PO|, \verb|BIOL170L_KS|, \verb|BIOL175__KS|, \verb|BIOL177__KS|, \verb|BIOL183__KS|, \verb|BIOL184L_KS|, \verb|BIOL187JLKS|, \verb|BIOL188__HM|, \verb|BIOL188L_KS|, \verb|BIOL189__HM|, \verb|BIOL189L_KS|, \verb|BIOL111__KS|, \verb|BIOL131__KS|, \verb|BIOL135L_KS|, \verb|BIOL138__KS|, \verb|BIOL138L_KS|, \verb|BIOL156L_KS|, \verb|BIOL164__KS|, \verb|BIOL168L_KS|, \verb|BIOL176__KS|, \verb|BIOL187M_KS|, \verb|BIOL142L_KS|, \verb|BIOL101__HM|, \verb|BIOL108__HM|, \verb|BIOL154__HM|, \verb|BIOL182__HM|, \verb|BIOL184__HM|, \verb|MCBI118A_HM|, \verb|MCBI118B_HM|, \verb|BIOL169__PO|, \verb|BIOL169L_PO|, \verb|BIOL170__PO|, \verb|BIOL173__PO|, \verb|BIOL174__PO|, \verb|BIOL112__PO|, \verb|BIOL152__PO|, \verb|BIOL110__HM|, \verb|BIOL160__HM|, \verb|BIOL119__HM|, \verb|BIOL163L_PO|, \verb|BIOL182L_KS|, \verb|BIOL185G_HM|, \verb|BIOL187P_KS|, \verb|MCBI117__HM|, \verb|BIOL103__HM|, \verb|BIOL104A_PO|, \verb|BIOL108__PO|, \verb|BIOL113L_KS|, \verb|BIOL116__PO|, \verb|BIOL122__HM|, \verb|BIOL133__PO|, \verb|BIOL151L_KS|, \verb|BIOL162__PO|, \verb|BIOL173L_KS|\\\hline
    \verb|Maj_MolBio_BioLab| & \verb|BIOL103__HM|, \verb|BIOL110__HM|, \verb|BIOL184__HM|\\\hline
    \verb|Maj_MolBio_Capstone_BioThesis| & \verb|BIOL193__HM|\\\hline
    \verb|Maj_MolBio_Capstone_CSClinic| & \verb|CSCI183__HM|, \verb|CSCI184__HM|\\\hline
    \verb|Maj_MolBio_Capstone_ChemThesis| & \verb|CHEM151__HM|, \verb|CHEM152__HM|\\\hline
    \verb|Maj_MolBio_Capstone_EngrClinic| & \verb|ENGR111__HM|, \verb|ENGR112__HM|, \verb|ENGR113__HM|\\\hline
    \verb|Maj_MolBio_Capstone_IntensiveResearch| & \verb|BIOL195__HM|\\\hline
    \verb|Maj_MolBio_Capstone_MathClinic| & \verb|MATH193__HM|\\\hline
    \verb|Maj_MolBio_Capstone_PhysClinic| & \verb|PHYS193__HM|, \verb|PHYS194__HM|\\\hline
    \verb|Maj_MolBio_Colloquium| & \verb|BIOL191__HM|, \verb|CHEM199__HM|\\\hline
    \verb|Maj_MolBio_MolBio| & \verb|BIOL122__HM|, \verb|BIOL160__HM|\\\hline
    \verb|Maj_MolBio_SeminarCourse| & \verb|BIOL121__HM|, \verb|BIOL189__HM|\\\hline
    \verb|Maj_Phys| & \verb|MATH082__HM|, \verb|PHYS051__HM|, \verb|PHYS052__HM|, \verb|PHYS054__HM|, \verb|PHYS064__HM|, \verb|PHYS111__HM|, \verb|PHYS116__HM|, \verb|PHYS133__HM|, \verb|PHYS134__HM|, \verb|PHYS151__HM|\\\hline
    \verb|Maj_Phys_PhysColloquium| & \verb|PHYS195__HM|, \verb|PHYS195__HM2|\\\hline
    \verb|Maj_Phys_Thesis_Clinic| & \verb|PHYS193__HM|, \verb|PHYS194__HM|\\\hline
    \verb|Maj_Phys_Thesis_PhysThesis| & \verb|PHYS199__HM|\\\hline
    \verb|PHYS051__HM_co| & \verb|MATH082__HM|, \verb|MATH056__HM|\\\hline
    \verb|PHYS054__HM| & \verb|PHYS050__HM|, \verb|PHYS052__HM|\\\hline
    \verb|PHYS162__HM_co| & \verb|PHYS116__HM|\\\hline
    \verb|PHYS164__HM| & \verb|PHYS116__HM|\\\hline
    \caption{Requirement group contents: course prerequisites.}
    \label{tab:reqs-coreqs}
\end{longtable}

\end{landscape}
}

\section{Sample Interest and Results}\label{app:sample}

\begin{table}[H]
    \centering
    \begin{tabular}{|p{0.2\textwidth}|c||p{0.2\textwidth}|c||p{0.2\textwidth}|c|}
    \hline
    Course & Interest & Course & Interest & Course & Interest\\\hline
    \verb|MATH055__HM| & 1&
    \verb|MATH062__HM| & 3&
    \verb|MATH082__HM| & 1\\
    \verb|MATH131__HM| & 1&
    \verb|MATH171__HM| & 1&
    \verb|MATH180__HM| & 1\\
    \verb|MATH165__HM| & 2&
    \verb|MATH156__HM| & 4&
    \verb|MATH187__HM| & 5\\
    \verb|MATH158__HM| & 2&
    \verb|MATH152__PO| & 2&&\\
    \hline
\end{tabular}
    \caption{Math Major Sample Interests. Interests in the generic courses are held at 0, and interests in all other courses are set to $-10$.}
    \label{tab:math-in}
\end{table}

\begin{table}[H]
    \centering
    \begin{tabular}{|p{0.2\textwidth}|p{0.2\textwidth}|p{0.2\textwidth}|p{0.2\textwidth}|}
    \hline
    Fr. Fall & Fr. Spring & So. Fall & So. Spring\\\hline
    \verb|WRIT001__HM| \verb|BIOL023__HM| \verb|BIOL046__HM| \verb|MATH019__HM| \verb|PHYS023__HM| \verb|CSCI005__HM|&
    \verb|CHEM024__HM| \verb|CHEM042__HM| \verb|HSA_010__HM| \verb|MATH073__HM| \verb|PHYS024__HM|&
    \verb|MATH055__HM| \verb|MATH082__HM| \verb|MATH152__PO| \verb|MATH187__HM| \verb|ENGR079__HM|&
    \verb|MATH131__HM| \verb|MATH171__HM| \verb|PHYS050__HM| \verb|MATH062__HM| \verb|CORE099__HM| \verb|HSA_001__HM|\\
    \hline
    Jr. Fall & Jr. Spring & Sr. Fall & Sr. Spring\\\hline
    \verb|MATH158__HM| \verb|MATH165__HM| \verb|MATH180__HM| \verb|HSA_002__5C| \verb|HSA_003__HM|&
    \verb|MATH198__HM| \verb|STEM002__5C| \verb|HSA_001__5C| \verb|HSA_004__5C| \verb|HSA_000__HM| \verb|HSA_100__HM|&
    \verb|MATH197__HM| \verb|STEM004__5C| \verb|STEM006__5C| \verb|MATH903__HM| \verb|HSA_003__5C|&
    \verb|MATH197__HM2| \verb|STEM000__5C| \verb|STEM005__5C| \verb|HSA_000__5C| \verb|HSA_005__5C|\\
    \hline
\end{tabular}
    \caption{Math Major Sample Schedule}
    \label{tab:math-out}
\end{table}

\begin{table}[H]
    \centering
    \begin{tabular}{|p{0.2\textwidth}|c||p{0.2\textwidth}|c||p{0.2\textwidth}|c|}
    \hline
    Course & Interest & Course & Interest & Course & Interest\\\hline
    \verb|ENGR004__HM| & 3&
    \verb|ENGR085__HM| & 3&
    \verb|ENGR084__HM| & 2\\
    \verb|ENGR080__HM| & 1&
    \verb|ENGR082__HM| & 1&
    \verb|ENGR083__HM| & 1\\
    \verb|ENGR086__HM| & 1&
    \verb|ENGR154__HM| & 5&
    \verb|ENGR155__HM| & 5\\
    \verb|ENGR187__HM| & 5&&&&\\
    \hline
\end{tabular}
    \caption{Engineering Major Sample Interests. Interests in the generic courses are held at 0, and interests in all other courses are set to -10.}
    \label{tab:engr-in}
\end{table}

\begin{table}[H]
    \centering
    \begin{tabular}{|p{0.2\textwidth}|p{0.2\textwidth}|p{0.2\textwidth}|p{0.2\textwidth}|}
    \hline
    Fr. Fall & Fr. Spring & So. Fall & So. Spring\\\hline
    \verb|WRIT001__HM| \verb|BIOL023__HM| \verb|BIOL046__HM| \verb|MATH019__HM| \verb|PHYS023__HM| \verb|CSCI005__HM|&
    \verb|CHEM024__HM| \verb|CHEM042__HM| \verb|HSA_010__HM| \verb|MATH073__HM| \verb|PHYS024__HM|&
    \verb|ENGR004__HM| \verb|ENGR079__HM| \verb|ENGR085__HM| \verb|ENGR187__HM| \verb|PHYS050__HM|&
    \verb|MATH056__HM| \verb|ENGR072__HM| \verb|ENGR080__HM| \verb|ENGR154__HM| \verb|CORE099__HM|\\
    \hline
    Jr. Fall & Jr. Spring & Sr. Fall & Sr. Spring\\\hline
    \verb|ENGR082__HM| \verb|ENGR083__HM| \verb|ENGR084__HM| \verb|ENGR086__HM| \verb|ENGR122__HM| \verb|ENGR155__HM|&
    \verb|ENGR111__HM| \verb|ENGR124__HM| \verb|HSA_001__5C| \verb|HSA_003__5C| \verb|HSA_002__HM| \verb|HSA_100__HM|&
    \verb|ENGR101__HM| \verb|ENGR112__HM| \verb|HSA_000__5C| \verb|HSA_002__5C| \verb|HSA_006__5C|&
    \verb|ENGR102__HM| \verb|ENGR113__HM| \verb|HSA_004__5C| \verb|HSA_000__HM| \verb|HSA_003__HM|\\
    \hline
\end{tabular}
    \caption{Engineering Major Sample Schedule}
    \label{tab:engr-out}
\end{table}

\begin{table}[H]
    \centering
    \begin{tabular}{|p{0.2\textwidth}|c||p{0.2\textwidth}|c||p{0.2\textwidth}|c|}
    \hline
    Course & Interest & Course & Interest & Course & Interest\\\hline
    \verb|CSCI060__HM| & 1&
    \verb|MATH055__HM| & 1&
    \verb|CSCI070__HM| & 1\\
    \verb|CSCI081__HM| & 1&
    \verb|CSCI105__HM| & 2&
    \verb|CSCI123__HM| & 1\\
    \verb|CSCI140__HM| & 1&
    \verb|CSCI131__HM| & 4&
    \verb|CSCI133__HM| & 4\\
    \verb|CSCI134__HM| & 4&
    \verb|CSCI158__HM| & 3&
    \verb|CSCI153__HM| & 3\\
    \verb|CSCI152__HM| & 3&
    \verb|CSCI145__HM| & 2&&\\
    \hline
\end{tabular}
    \caption{Computer Science Major Sample Interests. Interests in the generic courses are held at 0, and interests in all other courses are set to $-10$.}
    \label{tab:cs-in}
\end{table}

\begin{table}[H]
    \centering
    \begin{tabular}{|p{0.2\textwidth}|p{0.2\textwidth}|p{0.2\textwidth}|p{0.2\textwidth}|}
    \hline
    Fr. Fall & Fr. Spring & So. Fall & So. Spring\\\hline
    \verb|WRIT001__HM| \verb|BIOL023__HM| \verb|BIOL046__HM| \verb|MATH019__HM| \verb|PHYS023__HM| \verb|CSCI005__HM|&
    \verb|CHEM024__HM| \verb|CHEM042__HM| \verb|HSA_010__HM| \verb|MATH073__HM| \verb|PHYS024__HM|&
    \verb|MATH055__HM| \verb|ENGR079__HM| \verb|CSCI060__HM| \verb|HSA_000__5C| \verb|HSA_002__5C|&
    \verb|PHYS050__HM| \verb|CSCI070__HM| \verb|CSCI081__HM| \verb|CORE099__HM| \verb|HSA_001__5C| \verb|HSA_002__HM|\\
    \hline
    Jr. Fall & Jr. Spring & Sr. Fall & Sr. Spring\\\hline
    \verb|CSCI105__HM| \verb|CSCI123__HM| \verb|CSCI140__HM| \verb|CSCI153__HM| \verb|CSCI195__HM| \verb|STEM011__5C|&
    \verb|CSCI131__HM| \verb|CSCI195__HM2| \verb|STEM002__5C| \verb|STEM005__5C| \verb|CSCI902__HM| \verb|HSA_000__HM|&
    \verb|CSCI133__HM| \verb|CSCI134__HM| \verb|CSCI183__HM| \verb|CSCI195__HM3| \verb|HSA_001__HM| \verb|HSA_100__HM|&
    \verb|CSCI184__HM| \verb|CSCI195__HM4| \verb|HSA_003__5C| \verb|HSA_004__5C| \verb|HSA_005__5C| \verb|HSA_006__5C|\\
    \hline
\end{tabular}
    \caption{Computer Science Major Sample Schedule}
    \label{tab:cs-out}
\end{table}

\section{Sensitivity Analysis Model Outputs}\label{app:sensitivity}

\begin{table}[H]
    \centering
    \input{fragility-0.tex}
    \caption{Model output with $\gamma=0$ and all other parameters as specified in Tables \ref{tab:parameters}, \ref{tab:credits}, \ref{tab:offerings}, \ref{tab:reqs-groups}, \ref{tab:reqs-prereqs}, \ref{tab:reqs-coreqs}, and \ref{tab:math-in} above.}
    \label{tab:fragility-0-out}
\end{table}

\begin{table}[H]
    \centering
    \begin{tabular}{|p{0.2\textwidth}|p{0.2\textwidth}|p{0.2\textwidth}|p{0.2\textwidth}|}
    \hline
    Fr. Fall & Fr. Spring & So. Fall & So. Spring\\\hline
    \verb|WRIT001__HM| \verb|MATH019__HM| \verb|CHEM024__HM| \verb|CHEM042__HM| \verb|PHYS023__HM| \verb|CSCI005__HM|&
    \verb|BIOL023__HM| \verb|BIOL046__HM| \verb|HSA_010__HM| \verb|MATH062__HM| \verb|MATH073__HM| \verb|PHYS024__HM|&
    \verb|MATH055__HM| \verb|MATH082__HM| \verb|MATH152__PO| \verb|MATH187__HM| \verb|ENGR079__HM|&
    \verb|MATH131__HM| \verb|MATH171__HM| \verb|PHYS050__HM| \verb|CORE099__HM| \verb|HSA_004__5C| \verb|HSA_001__HM|\\
    \hline
    Jr. Fall & Jr. Spring & Sr. Fall & Sr. Spring\\\hline
    \verb|MATH158__HM| \verb|MATH165__HM| \verb|MATH180__HM| \verb|MATH902__HM| \verb|HSA_002__5C|&
    \verb|MATH198__HM| \verb|HSA_001__5C| \verb|HSA_000__HM| \verb|HSA_003__HM| \verb|HSA_100__HM|&
    \verb|MATH193__HM| \verb|STEM000__5C| \verb|HSA_000__5C| \verb|HSA_003__5C| \verb|HSA_002__HM|&
    \verb|MATH193__HM2| \verb|STEM001__5C| \verb|MATH900__HM| \verb|MATH903__HM| \verb|MATH904__HM|\\
    \hline
\end{tabular}
    \caption{Model output with $\gamma=0.5$ and all other parameters as specified in Tables \ref{tab:parameters}, \ref{tab:credits}, \ref{tab:offerings}, \ref{tab:reqs-groups}, \ref{tab:reqs-prereqs}, \ref{tab:reqs-coreqs}, and \ref{tab:math-in} above.}
    \label{tab:fragility-05-out}
\end{table}

\begin{table}[H]
    \centering
    \begin{tabular}{|p{0.2\textwidth}|p{0.2\textwidth}|p{0.2\textwidth}|p{0.2\textwidth}|}
    \hline
    Fr. Fall & Fr. Spring & So. Fall & So. Spring\\\hline
    \verb|WRIT001__HM| \verb|BIOL023__HM| \verb|BIOL046__HM| \verb|MATH019__HM| \verb|MATH152__PO| \verb|PHYS023__HM| \verb|CSCI005__HM|&
    \verb|CHEM024__HM| \verb|CHEM042__HM| \verb|HSA_010__HM| \verb|MATH062__HM| \verb|MATH073__HM| \verb|PHYS024__HM|&
    \verb|MATH055__HM| \verb|MATH082__HM| \verb|MATH158__HM| \verb|MATH187__HM| \verb|ENGR079__HM|&
    \verb|MATH131__HM| \verb|MATH171__HM| \verb|PHYS050__HM| \verb|CORE099__HM| \verb|HSA_003__HM| \verb|HSA_100__HM|\\
    \hline
    Jr. Fall & Jr. Spring & Sr. Fall & Sr. Spring\\\hline
    \verb|MATH165__HM| \verb|MATH180__HM| \verb|STEM010__5C| \verb|HSA_003__5C| \verb|HSA_000__HM| \verb|HSA_002__HM|&
    \verb|MATH198__HM| \verb|STEM001__5C| \verb|STEM002__5C| \verb|HSA_002__5C| \verb|HSA_001__HM|&
    \verb|MATH193__HM| \verb|STEM009__5C| \verb|MATH904__HM| \verb|HSA_001__5C| \verb|HSA_004__5C| \verb|HSA_005__5C|&
    \verb|MATH193__HM2| \verb|STEM000__5C| \verb|STEM004__5C| \verb|STEM006__5C| \verb|MATH902__HM|\\
    \hline
\end{tabular}
    \caption{Model output with $O=0$ and all other parameters as specified in Tables \ref{tab:parameters}, \ref{tab:credits}, \ref{tab:offerings}, \ref{tab:reqs-groups}, \ref{tab:reqs-prereqs}, \ref{tab:reqs-coreqs}, and \ref{tab:math-in} above.}
    \label{tab:credits-0-out}
\end{table}

\begin{table}[H]
    \centering
    \input{credit-10.tex}
    \caption{Model output with $O=10$ and all other parameters as specified in Tables \ref{tab:parameters}, \ref{tab:credits}, \ref{tab:offerings}, \ref{tab:reqs-groups}, \ref{tab:reqs-prereqs}, \ref{tab:reqs-coreqs}, and \ref{tab:math-in} above.}
    \label{tab:credits-10-out}
\end{table}

\begin{table}[H]
    \centering
    \begin{tabular}{|p{0.2\textwidth}|p{0.2\textwidth}|p{0.2\textwidth}|p{0.2\textwidth}|}
    \hline
    Fr. Fall & Fr. Spring & So. Fall & So. Spring\\\hline
    \verb|NEUR189L_KS| \verb|ASTR101L_PO| \verb|WRIT001__HM| \verb|BIOL023__HM| \verb|BIOL046__HM| \verb|MATH019__HM| \verb|PHYS023__HM| \verb|CSCI005__HM| \verb|NEUR122L_PO| \verb|FREN001L_CM| \verb|FREN044L_SC|&
    \verb|CHEM024__HM| \verb|CHEM042__HM| \verb|HSA_010__HM| \verb|MATH073__HM| \verb|PHYS024__HM|&
    \verb|FREN045L_SC| \verb|MUS_175__JM| \verb|MATH055__HM| \verb|MATH082__HM| \verb|MATH187__HM| \verb|ENGR079__HM| \verb|GOVT041B_CM| \verb|EA__101L_PO| \verb|POLI195__PO| \verb|PHYS041L_PO|&
    \verb|FREN033L_SC| \verb|GERM011__PO| \verb|MATH131__HM| \verb|MATH171__HM| \verb|PHYS050__HM| \verb|BIOL138L_KS| \verb|MSL_101B_CM| \verb|MATH062__HM| \verb|FREN002L_CM| \verb|FREN033L_CM| \verb|DANC182__PO| \verb|CORE099__HM|\\
    \hline
    Jr. Fall & Jr. Spring & Sr. Fall & Sr. Spring\\\hline
    \verb|MATH152__PO| \verb|MATH158__HM| \verb|MATH165__HM| \verb|MATH180__HM| \verb|ASAM126__HM|&
    \verb|FREN193__PO| \verb|MUS_180H_SC| \verb|THEA052H_PO| \verb|PORT035__PZ| \verb|MATH157__HM| \verb|MATH198__HM| \verb|GOVT071__CM| \verb|DANC124P_PO| \verb|DANC152P_PO| \verb|DANC172__PO| \verb|DANC174__PO| \verb|CHEM023BLPO| \verb|CSCI046L_CM| \verb|GOVT999__CM| \verb|ASTR051L_PO| \verb|ENGL087H_PO| \verb|MLLC021__PZ|&
    \verb|ASAM179R_HM| \verb|FREN011__PO| \verb|MUS_015__PO| \verb|MUS_015VOPO| \verb|MATH193__HM| \verb|LIT_107__HM| \verb|PHYS196__KS| \verb|ART_179W_HM|&
    \verb|ENGR085A_HM| \verb|MATH158__PO| \verb|ASAM187__AA| \verb|ASTR123__PO| \verb|MATH156__HM| \verb|MATH193__HM2|\\
    \hline
\end{tabular}
    \caption{Model output with $I_c$ defaulting to 0 and all other parameters as specified in Tables \ref{tab:parameters}, \ref{tab:credits}, \ref{tab:offerings}, \ref{tab:reqs-groups}, \ref{tab:reqs-prereqs}, \ref{tab:reqs-coreqs}, and \ref{tab:math-in} above.}
    \label{tab:interest-0-out}
\end{table}

\end{document}